%%%%%%%%%%%%%%%%%%%%%%%%%%%%%%%Packages%%%%%%%%%%%%%%%%%%%%%%%%%%%%%%%%
\documentclass[12pt,openright]{book}
\usepackage{iftex}
\iftutex
  % Moteurs Unicode (XeLaTeX/LuaLaTeX)
  \usepackage{fontspec}
  \usepackage{polyglossia}
  \setmainlanguage{french}
  % Définition d'une commande \nombre équivalente
\newcommand{\nombre}[1]{\numprint{#1}}
\else
  % Moteur classique (pdfLaTeX)
  \usepackage[T1]{fontenc}
  \usepackage[utf8]{inputenc}
  \usepackage[french]{babel}
\fi
\usepackage[a4paper,inner=1.2cm,outer=0.8cm,top=2cm,bottom=2cm]{geometry}
\usepackage{datatool}
\DTLsetseparator{;} % Avant tout : on a besoin de définir le séparateur pour la basededonnées
\usepackage{numprint}
\npthousandsep{\,} % espace fine insécable entre les milliers (usage français)
%\usepackage[noenc]{tipa}
%\usepackage{tipa}
\usepackage{tipx}
%\usepackage[geometry,weather,misc,clock]{ifsym}
\usepackage{pifont}
\usepackage{eurosym}
%\usepackage{amsmath}
\usepackage{wasysym}
%\usepackage{amssymb,amsfonts,textcomp}
%\usepackage[T3,T1]{fontenc}
\usepackage[dvipdfm]{graphicx} % Pour les images
\usepackage{metalogo}
\usepackage{xifthen} % Programmation dans LaTeX
\usepackage{multicol}
\usepackage{array}
\usepackage{multirow}
\usepackage{makecell}
\usepackage{tabularx}
\usepackage{colortbl}
\usepackage[most]{tcolorbox}
\usepackage{enumitem}
\usepackage{booktabs}
\usepackage[table,xcdraw]{xcolor}
\usepackage{caption}
\DeclareCaptionFormat{sanslabel}{#3}
\usepackage{xspace}
\usepackage{xstring}
\usepackage{xfp}%Gestion des pourcentages
\usepackage{tikz}
\usetikzlibrary{babel}
\usepackage{pgf-pie}
\usepackage{pgfplots}
\pgfplotsset{compat=1.18}
\usepackage{changepage}
\usepackage[defaultlines=2]{nowidow}
\usepackage{xpatch} % Nécessaire pour modifier l'alignement des ndbp/Manier avec précaution
\usepackage{xparse}
\usepackage{etoolbox} % Ajouts de commandes pour résoudre des problèmes d'espacement notamment
\usepackage{setspace} % Pour changer l'espacement entre les lignes
\usepackage{shorttoc}
\usepackage{titlesec}
\usepackage{titletoc}
\usepackage{tocbibind} % rajoute la bib, l'ind, la ToC, la lof et la lot automatiquement à ToC.
\usepackage{totpages} % Pour obtenir le numéro de la dernière page
\usepackage{perpage} % Permet de recommencer les ndbp à chaque page
\MakePerPage{footnote} % Appliquer le changement des ndbp
\usepackage{fancyhdr} % Pour maîtriser les en-têtes
\usepackage{xkeyval}
\usepackage{longtable}
\usepackage{microtype} %Résolution de certains problèmes de dépassement d'hbox
\usepackage{needspace}
\usepackage{morewrites}
\usepackage{pdflscape}
\usepackage{pdfpages} % Insérer des pdf dans le fichier
\usepackage[xindy]{imakeidx} % Pour l'index
\title{Banque indexée des sujets d'Humanités, Terminale}
\author{Mikaël Quesseveur}
\date{2025-09-22}

\renewcommand{\thesubsection}{\arabic{section}.\arabic{subsection}.}

\titleformat{\subsection}
{\normalfont\large\bfseries}{\thesubsection}{1em}{}

\titleformat{\subsubsection}
{\center\normalfont\normalsize\bfseries}{\thesubsubsection}{1em}{}

%\titleformat*{\subsubsection}{\center\bfseries}

%%%Configuration de dots pour qu'il soit plus élégant et supprimer l'espace inélégante dans le cas où les points de suspension sont entre crochets%%%
\let\olddots\dots
\renewcommand\dots{\olddots{}\xspace}
\xspaceaddexceptions{]}

%%%Reconfiguration des colonnes de tabularx pour améliorer l'espacement%%%
\renewcommand{\tabularxcolumn}[1]{>{\centering\arraybackslash}m{#1}}

%%%%%%%%%%%%%%%%%%%%%%%%%%%%%%%%%%%%%%%%%%%%%%%%%%%%%%%%%%%%%%%%%%%%%%
%%%%%%%%%%%%%%%%%%%%%Commandes pour l'index%%%%%%%%%%%%%%%%%%%%%%%%%%%
%%%%%%%%%%%%%%%%%%%%%%%%%%%%%%%%%%%%%%%%%%%%%%%%%%%%%%%%%%%%%%%%%%%%%%
\newcommand{\questiontitle}{\bfseries\large}
\newcommand{\questiontype}{\itshape\normalsize}
\newcommand{\questionsize}[1]{{\normalfont\footnotesize#1}}
%%%Changement de l'indentation dans les index%%%
\makeatletter
\def\@idxitem{\par\hangindent 12pt}
\makeatother

%%%Gestion de l'espacement dans les environnements d'index
\AtBeginEnvironment{theindex}{\raggedbottom}  % nécessite etoolbox

%%%Les liens vers les banques de sujets%%%

\newcommand{\sujetshlpsanscorr}{\href{https://drive.google.com/file/d/127SoEaNlMn662R3lhdlSIw74tFSGJGNQ/view?usp=drive_link}{Banque de sujets HLP sans éléments d'évaluation}}
\newcommand{\sujetshlpcorr}{\href{https://drive.google.com/file/d/1PJyaKCNypSU1-v5-fRtyrf5YMxSI8puw/view?usp=drive_link}{Banque de sujets HLP avec éléments d'évaluation}}
\newcommand{\sujetscsv}{\href{https://docs.google.com/spreadsheets/d/1cosKjvc7rgAi1zdjNtZ4UWY1j2-Tbmz-K3ha3xRiIc4/edit?usp=sharing}{Tableau (\textit{Sheets, en ligne)} récapitulant les sujets et leur indexation}}

%%%%%%%%%%%%%%%%%%%%%%%%%%%%%%%%%%%%%%%%%%%%%%%%%%%%%%%%%%%%%%%%%%%%%%%
%%%%%%%%%%%%%%%%%%%%%%%%% Compteurs %%%%%%%%%%%%%%%%%%%%%%%%%%%%%%%%%%%
%%%%%%%%%%%%%%%%%%%%%%%%%%%%%%%%%%%%%%%%%%%%%%%%%%%%%%%%%%%%%%%%%%%%%%%
%%%%%%%%%%%%%%%%%%%%%%%%%%%%%%%%%%%%%%%%%%%%%%%%%%%%%%%%%%%%%%%%%%%%%%%
%%%%%%%%%%%%%%%%%%%%%%%%%%%%%%Compteurs%%%%%%%%%%%%%%%%%%%%%%%%%%%%%%%%
%%%%%%%%%%%%%%%%%%%%%%%%%%%%%%%%%%%%%%%%%%%%%%%%%%%%%%%%%%%%%%%%%%%%%%%
%%%Compter les sujets par thème%%%
\newcounter{nosujet}
\newcounter{sujetHQ}
\newcounter{sujetRS}
%%%Sujets simples%%%
\newcounter{sujetEXS}
\newcounter{sujetETA}
\newcounter{sujetMM}
\newcounter{sujetHV}
\newcounter{sujetCC}
\newcounter{sujetHL}
%%%Sujets mixtes%%%
\newcounter{sujetEXSTA}
\newcounter{sujetEXSMM}
\newcounter{sujetETAMM}
\newcounter{sujetOdS}
\newcounter{sujetHVL}
\newcounter{sujetHVCC}
\newcounter{sujetHLCC}
\newcounter{sujetOdH}
%%%Sujets mixtes semestres%%%
\newcounter{sujetETACC}
\newcounter{sujetETAHV}
\newcounter{sujetETAHL}
\newcounter{sujetEXSCC}
\newcounter{sujetEXSHV}
\newcounter{sujetEXSHL}
\newcounter{sujetMMCC}
\newcounter{sujetMMHV}
\newcounter{sujetMMHL}

%%%Sujets mixtes années%%
\newcounter{sujetpremterm}%%%Compteur général de sujet
\newcounter{sujetETAPdP}
\newcounter{sujetETADFI}
\newcounter{sujetETADMRC}
\newcounter{sujetETAHA}
\newcounter{sujetpremtermETA}
\newcounter{sujetEXSPdP}
\newcounter{sujetEXSDFI}
\newcounter{sujetEXSDMRC}
\newcounter{sujetEXSHA}
\newcounter{sujetpremtermEXS}
\newcounter{sujetMMPdP}
\newcounter{sujetMMDFI}
\newcounter{sujetMMDMRC}
\newcounter{sujetMMHA}
\newcounter{sujetpremtermMM}
\newcounter{sujetCCPdP}
\newcounter{sujetCCDFI}
\newcounter{sujetCCDMRC}
\newcounter{sujetCCHA}
\newcounter{sujetpremtermCC}
\newcounter{sujetHVPdP}
\newcounter{sujetHVDFI}
\newcounter{sujetHVDMRC}
\newcounter{sujetHVHA}
\newcounter{sujetpremtermHV}
\newcounter{sujetHLPdP}
\newcounter{sujetHLDFI}
\newcounter{sujetHLDMRC}
\newcounter{sujetHLHA}
\newcounter{sujetpremtermHL}
\newcounter{sujetRSPdP}
\newcounter{sujetRSRdM}
\newcounter{sujetRSDFI}
\newcounter{sujetRSDMRC}
\newcounter{sujetRSHA}
\newcounter{sujetHQPdP}
\newcounter{sujetHQRdM}
\newcounter{sujetHQDFI}
\newcounter{sujetHQDMRC}
\newcounter{sujetHQHA}

%%%Somme des compteurs%%%
\newcounter{sujetall}
\newcounter{sujetallETA}
\newcounter{sujetallEXS}
\newcounter{sujetallMM}
\newcounter{sujetallRS}
\newcounter{sujetallCC}
\newcounter{sujetallHV}
\newcounter{sujetallHL}
\newcounter{sujetallHQ}

%%%Compter les sujets par thème et types de questions%%%
%%%%%%%%%%%%%%%%%%%%%%%%%%%%%%%%%%%%%%%%%%%%%
%%%%%%%%%Question d'interprétation%%%%%%%%%%%
%%%%%%%%%%%%%%%%%%%%%%%%%%%%%%%%%%%%%%%%%%%%%

%%%%%%%%%%%%%%%Général%%%%%%%%%%%%%%%%%%%%%%%
\newcounter{nosujetint}
\newcounter{sujetHQint}
\newcounter{sujetRSint}
%%%Sujets simples%%%
\newcounter{sujetintEXS}
\newcounter{sujetintETA}
\newcounter{sujetintMM}
\newcounter{sujetintHV}
\newcounter{sujetintCC}
\newcounter{sujetintHL}
%%%Sujets mixtes%%%
\newcounter{sujetintEXSTA}
\newcounter{sujetintEXSMM}
\newcounter{sujetintETAMM}
\newcounter{sujetintOdS}
\newcounter{sujetintHVL}
\newcounter{sujetintHVCC}
\newcounter{sujetintHLCC}
\newcounter{sujetintOdH}
%%%Sujets mixtes semestres%%%
\newcounter{sujetintETACC}
\newcounter{sujetintETAHV}
\newcounter{sujetintETAHL}
\newcounter{sujetintEXSCC}
\newcounter{sujetintEXSHV}
\newcounter{sujetintEXSHL}
\newcounter{sujetintMMCC}
\newcounter{sujetintMMHV}
\newcounter{sujetintMMHL}
%%%Sujets mixtes années%%
\newcounter{sujetintpremterm}
\newcounter{sujetintpremtermRS}
\newcounter{sujetintpremtermHQ}
\newcounter{sujetintETAPdP}
\newcounter{sujetintETADFI}
\newcounter{sujetintETADMRC}
\newcounter{sujetintETAHA}
\newcounter{sujetintpremtermETA}
\newcounter{sujetintEXSPdP}
\newcounter{sujetintEXSDFI}
\newcounter{sujetintEXSDMRC}
\newcounter{sujetintEXSHA}
\newcounter{sujetintpremtermEXS}
\newcounter{sujetintMMPdP}
\newcounter{sujetintMMDFI}
\newcounter{sujetintMMDMRC}
\newcounter{sujetintMMHA}
\newcounter{sujetintpremtermMM}
\newcounter{sujetintCCPdP}
\newcounter{sujetintCCDFI}
\newcounter{sujetintCCDMRC}
\newcounter{sujetintCCHA}
\newcounter{sujetintpremtermCC}
\newcounter{sujetintHVPdP}
\newcounter{sujetintHVDFI}
\newcounter{sujetintHVDMRC}
\newcounter{sujetintHVHA}
\newcounter{sujetintpremtermHV}
\newcounter{sujetintHLPdP}
\newcounter{sujetintHLDFI}
\newcounter{sujetintHLDMRC}
\newcounter{sujetintHLHA}
\newcounter{sujetintpremtermHL}
\newcounter{sujetintRSPdP}
\newcounter{sujetintRSRdM}
\newcounter{sujetintRSDFI}
\newcounter{sujetintRSDMRC}
\newcounter{sujetintRSHA}
\newcounter{sujetintHQPdP}
\newcounter{sujetintHQRdM}
\newcounter{sujetintHQDFI}
\newcounter{sujetintHQDMRC}
\newcounter{sujetintHQHA}
%%%Somme des compteurs%%%
\newcounter{sujetallintETA}
\newcounter{sujetallintEXS}
\newcounter{sujetallintMM}
\newcounter{sujetallintRS}
\newcounter{sujetallintCC}
\newcounter{sujetallintHV}
\newcounter{sujetallintHL}
\newcounter{sujetallintHQ}

%%%%%%%%Interprétation philosophique%%%%%%%%%
\newcounter{nosujetintphilosophique}
\newcounter{sujetHQintphilosophique}
\newcounter{sujetRSintphilosophique}
%%%Sujets simples%%%
\newcounter{sujetintphilosophiqueEXS}
\newcounter{sujetintphilosophiqueETA}
\newcounter{sujetintphilosophiqueMM}
\newcounter{sujetintphilosophiqueHV}
\newcounter{sujetintphilosophiqueCC}
\newcounter{sujetintphilosophiqueHL}
%%%Sujets mixtes%%%
\newcounter{sujetintphilosophiqueEXSTA}
\newcounter{sujetintphilosophiqueEXSMM}
\newcounter{sujetintphilosophiqueETAMM}
\newcounter{sujetintphilosophiqueOdS}
\newcounter{sujetintphilosophiqueHVL}
\newcounter{sujetintphilosophiqueHVCC}
\newcounter{sujetintphilosophiqueHLCC}
\newcounter{sujetintphilosophiqueOdH}
%%%Sujets mixtes semestres%%%
\newcounter{sujetintphilosophiqueETACC}
\newcounter{sujetintphilosophiqueETAHV}
\newcounter{sujetintphilosophiqueETAHL}
\newcounter{sujetintphilosophiqueEXSCC}
\newcounter{sujetintphilosophiqueEXSHV}
\newcounter{sujetintphilosophiqueEXSHL}
\newcounter{sujetintphilosophiqueMMCC}
\newcounter{sujetintphilosophiqueMMHV}
\newcounter{sujetintphilosophiqueMMHL}
%%%Sujets mixtes années%%
\newcounter{sujetintphilosophiquepremterm}
\newcounter{sujetintphilosophiquepremtermRS}
\newcounter{sujetintphilosophiquepremtermHQ}
\newcounter{sujetintphilosophiqueETAPdP}
\newcounter{sujetintphilosophiqueETADFI}
\newcounter{sujetintphilosophiqueETADMRC}
\newcounter{sujetintphilosophiqueETAHA}
\newcounter{sujetintphilosophiquepremtermETA}
\newcounter{sujetintphilosophiqueEXSPdP}
\newcounter{sujetintphilosophiqueEXSDFI}
\newcounter{sujetintphilosophiqueEXSDMRC}
\newcounter{sujetintphilosophiqueEXSHA}
\newcounter{sujetintphilosophiquepremtermEXS}
\newcounter{sujetintphilosophiqueMMPdP}
\newcounter{sujetintphilosophiqueMMDFI}
\newcounter{sujetintphilosophiqueMMDMRC}
\newcounter{sujetintphilosophiqueMMHA}
\newcounter{sujetintphilosophiquepremtermMM}
\newcounter{sujetintphilosophiqueCCPdP}
\newcounter{sujetintphilosophiqueCCDFI}
\newcounter{sujetintphilosophiqueCCDMRC}
\newcounter{sujetintphilosophiqueCCHA}
\newcounter{sujetintphilosophiquepremtermCC}
\newcounter{sujetintphilosophiqueHVPdP}
\newcounter{sujetintphilosophiqueHVDFI}
\newcounter{sujetintphilosophiqueHVDMRC}
\newcounter{sujetintphilosophiqueHVHA}
\newcounter{sujetintphilosophiquepremtermHV}
\newcounter{sujetintphilosophiqueHLPdP}
\newcounter{sujetintphilosophiqueHLDFI}
\newcounter{sujetintphilosophiqueHLDMRC}
\newcounter{sujetintphilosophiqueHLHA}
\newcounter{sujetintphilosophiquepremtermHL}
\newcounter{sujetintphilosophiqueRSPdP}
\newcounter{sujetintphilosophiqueRSRdM}
\newcounter{sujetintphilosophiqueRSDFI}
\newcounter{sujetintphilosophiqueRSDMRC}
\newcounter{sujetintphilosophiqueRSHA}
\newcounter{sujetintphilosophiqueHQPdP}
\newcounter{sujetintphilosophiqueHQRdM}
\newcounter{sujetintphilosophiqueHQDFI}
\newcounter{sujetintphilosophiqueHQDMRC}
\newcounter{sujetintphilosophiqueHQHA}
%%%Somme des compteurs%%%
\newcounter{sujetallintphilosophiqueETA}
\newcounter{sujetallintphilosophiqueEXS}
\newcounter{sujetallintphilosophiqueMM}
\newcounter{sujetallintphilosophiqueRS}
\newcounter{sujetallintphilosophiqueCC}
\newcounter{sujetallintphilosophiqueHV}
\newcounter{sujetallintphilosophiqueHL}
\newcounter{sujetallintphilosophiqueHQ}

%%%%%%%%%Interprétation littéraire%%%%%%%%%%%
\newcounter{nosujetintlitteraire}
\newcounter{sujetHQintlitteraire}
\newcounter{sujetRSintlitteraire}
%%%Sujets simples%%%
\newcounter{sujetintlitteraireEXS}
\newcounter{sujetintlitteraireETA}
\newcounter{sujetintlitteraireMM}
\newcounter{sujetintlitteraireHV}
\newcounter{sujetintlitteraireCC}
\newcounter{sujetintlitteraireHL}
%%%Sujets mixtes%%%
\newcounter{sujetintlitteraireEXSTA}
\newcounter{sujetintlitteraireEXSMM}
\newcounter{sujetintlitteraireETAMM}
\newcounter{sujetintlitteraireOdS}
\newcounter{sujetintlitteraireHVL}
\newcounter{sujetintlitteraireHVCC}
\newcounter{sujetintlitteraireHLCC}
\newcounter{sujetintlitteraireOdH}
%%%Sujets mixtes semestres%%%
\newcounter{sujetintlitteraireETACC}
\newcounter{sujetintlitteraireETAHV}
\newcounter{sujetintlitteraireETAHL}
\newcounter{sujetintlitteraireEXSCC}
\newcounter{sujetintlitteraireEXSHV}
\newcounter{sujetintlitteraireEXSHL}
\newcounter{sujetintlitteraireMMCC}
\newcounter{sujetintlitteraireMMHV}
\newcounter{sujetintlitteraireMMHL}
%%%Sujets mixtes années%%
\newcounter{sujetintlitterairepremterm}
\newcounter{sujetintlitterairepremtermRS}
\newcounter{sujetintlitterairepremtermHQ}
\newcounter{sujetintlitteraireETAPdP}
\newcounter{sujetintlitteraireETADFI}
\newcounter{sujetintlitteraireETADMRC}
\newcounter{sujetintlitteraireETAHA}
\newcounter{sujetintlitterairepremtermETA}
\newcounter{sujetintlitteraireEXSPdP}
\newcounter{sujetintlitteraireEXSDFI}
\newcounter{sujetintlitteraireEXSDMRC}
\newcounter{sujetintlitteraireEXSHA}
\newcounter{sujetintlitterairepremtermEXS}
\newcounter{sujetintlitteraireMMPdP}
\newcounter{sujetintlitteraireMMDFI}
\newcounter{sujetintlitteraireMMDMRC}
\newcounter{sujetintlitteraireMMHA}
\newcounter{sujetintlitterairepremtermMM}
\newcounter{sujetintlitteraireCCPdP}
\newcounter{sujetintlitteraireCCDFI}
\newcounter{sujetintlitteraireCCDMRC}
\newcounter{sujetintlitteraireCCHA}
\newcounter{sujetintlitterairepremtermCC}
\newcounter{sujetintlitteraireHVPdP}
\newcounter{sujetintlitteraireHVDFI}
\newcounter{sujetintlitteraireHVDMRC}
\newcounter{sujetintlitteraireHVHA}
\newcounter{sujetintlitterairepremtermHV}
\newcounter{sujetintlitteraireHLPdP}
\newcounter{sujetintlitteraireHLDFI}
\newcounter{sujetintlitteraireHLDMRC}
\newcounter{sujetintlitteraireHLHA}
\newcounter{sujetintlitterairepremtermHL}
\newcounter{sujetintlitteraireRSPdP}
\newcounter{sujetintlitteraireRSRdM}
\newcounter{sujetintlitteraireRSDFI}
\newcounter{sujetintlitteraireRSDMRC}
\newcounter{sujetintlitteraireRSHA}
\newcounter{sujetintlitteraireHQPdP}
\newcounter{sujetintlitteraireHQRdM}
\newcounter{sujetintlitteraireHQDFI}
\newcounter{sujetintlitteraireHQDMRC}
\newcounter{sujetintlitteraireHQHA}
%%%Somme des compteurs%%%
\newcounter{sujetallintlitteraireETA}
\newcounter{sujetallintlitteraireEXS}
\newcounter{sujetallintlitteraireMM}
\newcounter{sujetallintlitteraireRS}
\newcounter{sujetallintlitteraireCC}
\newcounter{sujetallintlitteraireHV}
\newcounter{sujetallintlitteraireHL}
\newcounter{sujetallintlitteraireHQ}

%%%%%%%%%%%%%%%%%%%%%%%%%%%%%%%%%%%%%%%%%%%%%
%%%%%%%%%%%%Question de réflexion%%%%%%%%%%%%
%%%%%%%%%%%%%%%%%%%%%%%%%%%%%%%%%%%%%%%%%%%%%
\newcounter{nosujetref}
%%%Sujets simples%%%
\newcounter{sujetrefEXS}
\newcounter{sujetrefETA}
\newcounter{sujetrefMM}
\newcounter{sujetrefHV}
\newcounter{sujetrefCC}
\newcounter{sujetrefHL}
%%%Sujets mixtes%%%
\newcounter{sujetrefEXSTA}
\newcounter{sujetrefEXSMM}
\newcounter{sujetrefETAMM}
\newcounter{sujetrefOdS}
\newcounter{sujetrefHVL}
\newcounter{sujetrefHVCC}
\newcounter{sujetrefHLCC}
\newcounter{sujetrefOdH}
%%%Sujets mixtes semestres%%%
\newcounter{sujetrefETACC}
\newcounter{sujetrefETAHV}
\newcounter{sujetrefETAHL}
\newcounter{sujetrefEXSCC}
\newcounter{sujetrefEXSHV}
\newcounter{sujetrefEXSHL}
\newcounter{sujetrefMMCC}
\newcounter{sujetrefMMHV}
\newcounter{sujetrefMMHL}
%%%Sujets mixtes années%%
\newcounter{sujetrefpremterm}%%%Compteur général
\newcounter{sujetrefETAPdP}
\newcounter{sujetrefETADFI}
\newcounter{sujetrefETADMRC}
\newcounter{sujetrefETAHA}
\newcounter{sujetrefpremtermETA}
\newcounter{sujetrefEXSPdP}
\newcounter{sujetrefEXSDFI}
\newcounter{sujetrefEXSDMRC}
\newcounter{sujetrefEXSHA}
\newcounter{sujetrefpremtermEXS}
\newcounter{sujetrefMMPdP}
\newcounter{sujetrefMMDFI}
\newcounter{sujetrefMMDMRC}
\newcounter{sujetrefMMHA}
\newcounter{sujetrefpremtermMM}
\newcounter{sujetrefCCPdP}
\newcounter{sujetrefCCDFI}
\newcounter{sujetrefCCDMRC}
\newcounter{sujetrefCCHA}
\newcounter{sujetrefpremtermCC}
\newcounter{sujetrefHVPdP}
\newcounter{sujetrefHVDFI}
\newcounter{sujetrefHVDMRC}
\newcounter{sujetrefHVHA}
\newcounter{sujetrefpremtermHV}
\newcounter{sujetrefHLPdP}
\newcounter{sujetrefHLDFI}
\newcounter{sujetrefHLDMRC}
\newcounter{sujetrefHLHA}
\newcounter{sujetrefpremtermHL}
\newcounter{sujetrefRSPdP}
\newcounter{sujetrefRSRdM}
\newcounter{sujetrefRSDFI}
\newcounter{sujetrefRSDMRC}
\newcounter{sujetrefRSHA}
\newcounter{sujetrefHQPdP}
\newcounter{sujetrefHQRdM}
\newcounter{sujetrefHQDFI}
\newcounter{sujetrefHQDMRC}
\newcounter{sujetrefHQHA}
%%%Somme des compteurs%%%
\newcounter{sujetallrefETA}
\newcounter{sujetallrefEXS}
\newcounter{sujetallrefMM}
\newcounter{sujetallrefRS}
\newcounter{sujetallrefCC}
\newcounter{sujetallrefHV}
\newcounter{sujetallrefHL}
\newcounter{sujetallrefHQ}

%%%%%%%%Réflexion philosophique%%%%%%%%%
\newcounter{nosujetrefphilosophique}
\newcounter{sujetHQrefphilosophique}
\newcounter{sujetRSrefphilosophique}
%%%Sujets simples%%%
\newcounter{sujetrefphilosophiqueEXS}
\newcounter{sujetrefphilosophiqueETA}
\newcounter{sujetrefphilosophiqueMM}
\newcounter{sujetrefphilosophiqueHV}
\newcounter{sujetrefphilosophiqueCC}
\newcounter{sujetrefphilosophiqueHL}
%%%Sujets mixtes%%%
\newcounter{sujetrefphilosophiqueEXSTA}
\newcounter{sujetrefphilosophiqueEXSMM}
\newcounter{sujetrefphilosophiqueETAMM}
\newcounter{sujetrefphilosophiqueOdS}
\newcounter{sujetrefphilosophiqueHVL}
\newcounter{sujetrefphilosophiqueHVCC}
\newcounter{sujetrefphilosophiqueHLCC}
\newcounter{sujetrefphilosophiqueOdH}
%%%Sujets mixtes semestres%%%
\newcounter{sujetrefphilosophiqueETACC}
\newcounter{sujetrefphilosophiqueETAHV}
\newcounter{sujetrefphilosophiqueETAHL}
\newcounter{sujetrefphilosophiqueEXSCC}
\newcounter{sujetrefphilosophiqueEXSHV}
\newcounter{sujetrefphilosophiqueEXSHL}
\newcounter{sujetrefphilosophiqueMMCC}
\newcounter{sujetrefphilosophiqueMMHV}
\newcounter{sujetrefphilosophiqueMMHL}
%%%Sujets mixtes années%%
\newcounter{sujetrefphilosophiquepremterm}
\newcounter{sujetrefphilosophiquepremtermRS}
\newcounter{sujetrefphilosophiquepremtermHQ}
\newcounter{sujetrefphilosophiqueETAPdP}
\newcounter{sujetrefphilosophiqueETADFI}
\newcounter{sujetrefphilosophiqueETADMRC}
\newcounter{sujetrefphilosophiqueETAHA}\newcounter{sujetrefphilosophiquepremtermETA}
\newcounter{sujetrefphilosophiqueEXSPdP}
\newcounter{sujetrefphilosophiqueEXSDFI}
\newcounter{sujetrefphilosophiqueEXSDMRC}
\newcounter{sujetrefphilosophiqueEXSHA}
\newcounter{sujetrefphilosophiquepremtermEXS}
\newcounter{sujetrefphilosophiqueMMPdP}
\newcounter{sujetrefphilosophiqueMMDFI}
\newcounter{sujetrefphilosophiqueMMDMRC}
\newcounter{sujetrefphilosophiqueMMHA}
\newcounter{sujetrefphilosophiquepremtermMM}
\newcounter{sujetrefphilosophiqueCCPdP}
\newcounter{sujetrefphilosophiqueCCDFI}
\newcounter{sujetrefphilosophiqueCCDMRC}
\newcounter{sujetrefphilosophiqueCCHA}
\newcounter{sujetrefphilosophiquepremtermCC}
\newcounter{sujetrefphilosophiqueHVPdP}
\newcounter{sujetrefphilosophiqueHVDFI}
\newcounter{sujetrefphilosophiqueHVDMRC}
\newcounter{sujetrefphilosophiqueHVHA}
\newcounter{sujetrefphilosophiquepremtermHV}
\newcounter{sujetrefphilosophiqueHLPdP}
\newcounter{sujetrefphilosophiqueHLDFI}
\newcounter{sujetrefphilosophiqueHLDMRC}
\newcounter{sujetrefphilosophiqueHLHA}
\newcounter{sujetrefphilosophiquepremtermHL}
\newcounter{sujetrefphilosophiqueRSPdP}
\newcounter{sujetrefphilosophiqueRSRdM}
\newcounter{sujetrefphilosophiqueRSDFI}
\newcounter{sujetrefphilosophiqueRSDMRC}
\newcounter{sujetrefphilosophiqueRSHA}
\newcounter{sujetrefphilosophiqueHQPdP}
\newcounter{sujetrefphilosophiqueHQRdM}
\newcounter{sujetrefphilosophiqueHQDFI}
\newcounter{sujetrefphilosophiqueHQDMRC}
\newcounter{sujetrefphilosophiqueHQHA}
%%%Somme des compteurs%%%
\newcounter{sujetallrefphilosophiqueETA}
\newcounter{sujetallrefphilosophiqueEXS}
\newcounter{sujetallrefphilosophiqueMM}
\newcounter{sujetallrefphilosophiqueRS}
\newcounter{sujetallrefphilosophiqueCC}
\newcounter{sujetallrefphilosophiqueHV}
\newcounter{sujetallrefphilosophiqueHL}
\newcounter{sujetallrefphilosophiqueHQ}

%%%%%%%%%Réflexion littéraire%%%%%%%%%%%
\newcounter{nosujetreflitteraire}
\newcounter{sujetHQreflitteraire}
\newcounter{sujetRSreflitteraire}
%%%Sujets simples%%%
\newcounter{sujetreflitteraireEXS}
\newcounter{sujetreflitteraireETA}
\newcounter{sujetreflitteraireMM}
\newcounter{sujetreflitteraireHV}
\newcounter{sujetreflitteraireCC}
\newcounter{sujetreflitteraireHL}
%%%Sujets mixtes%%%
\newcounter{sujetreflitteraireEXSTA}
\newcounter{sujetreflitteraireEXSMM}
\newcounter{sujetreflitteraireETAMM}
\newcounter{sujetreflitteraireOdS}
\newcounter{sujetreflitteraireHVL}
\newcounter{sujetreflitteraireHVCC}
\newcounter{sujetreflitteraireHLCC}
\newcounter{sujetreflitteraireOdH}
%%%Sujets mixtes semestres%%%
\newcounter{sujetreflitteraireETACC}
\newcounter{sujetreflitteraireETAHV}
\newcounter{sujetreflitteraireETAHL}
\newcounter{sujetreflitteraireEXSCC}
\newcounter{sujetreflitteraireEXSHV}
\newcounter{sujetreflitteraireEXSHL}
\newcounter{sujetreflitteraireMMCC}
\newcounter{sujetreflitteraireMMHV}
\newcounter{sujetreflitteraireMMHL}
%%%Sujets mixtes années%%
\newcounter{sujetreflitterairepremterm}
\newcounter{sujetreflitterairepremtermRS}
\newcounter{sujetreflitterairepremtermHQ}
\newcounter{sujetreflitteraireETAPdP}
\newcounter{sujetreflitteraireETADFI}
\newcounter{sujetreflitteraireETADMRC}
\newcounter{sujetreflitteraireETAHA}
\newcounter{sujetreflitterairepremtermETA}
\newcounter{sujetreflitteraireEXSPdP}
\newcounter{sujetreflitteraireEXSDFI}
\newcounter{sujetreflitteraireEXSDMRC}
\newcounter{sujetreflitteraireEXSHA}
\newcounter{sujetreflitterairepremtermEXS}
\newcounter{sujetreflitteraireMMPdP}
\newcounter{sujetreflitteraireMMDFI}
\newcounter{sujetreflitteraireMMDMRC}
\newcounter{sujetreflitteraireMMHA}
\newcounter{sujetreflitterairepremtermMM}
\newcounter{sujetreflitteraireCCPdP}
\newcounter{sujetreflitteraireCCDFI}
\newcounter{sujetreflitteraireCCDMRC}
\newcounter{sujetreflitteraireCCHA}
\newcounter{sujetreflitterairepremtermCC}
\newcounter{sujetreflitteraireHVPdP}
\newcounter{sujetreflitteraireHVDFI}
\newcounter{sujetreflitteraireHVDMRC}
\newcounter{sujetreflitteraireHVHA}
\newcounter{sujetreflitterairepremtermHV}
\newcounter{sujetreflitteraireHLPdP}
\newcounter{sujetreflitteraireHLDFI}
\newcounter{sujetreflitteraireHLDMRC}
\newcounter{sujetreflitteraireHLHA}
\newcounter{sujetreflitterairepremtermHL}
\newcounter{sujetreflitteraireRSPdP}
\newcounter{sujetreflitteraireRSRdM}
\newcounter{sujetreflitteraireRSDFI}
\newcounter{sujetreflitteraireRSDMRC}
\newcounter{sujetreflitteraireRSHA}
\newcounter{sujetreflitteraireHQPdP}
\newcounter{sujetreflitteraireHQRdM}
\newcounter{sujetreflitteraireHQDFI}
\newcounter{sujetreflitteraireHQDMRC}
\newcounter{sujetreflitteraireHQHA}
%%%Somme des compteurs%%%
\newcounter{sujetallreflitteraireETA}
\newcounter{sujetallreflitteraireEXS}
\newcounter{sujetallreflitteraireMM}
\newcounter{sujetallreflitteraireCC}
\newcounter{sujetallreflitteraireRS}
\newcounter{sujetallreflitteraireHV}
\newcounter{sujetallreflitteraireHL}
\newcounter{sujetallreflitteraireHQ}
%%%Compteurs de statistiques globales sur les genres de question%%%
\newcounter{sujetexclRS}
\newcounter{sujetmixtRS}
\newcounter{sujetexclHQ}
\newcounter{sujetmixtHQ}
\newcounter{sujetmixtHQRS}
%%%Compteurs de statistiques sur les transversaux (divers semestres, première et terminale)
\newcounter{sujettransintlitteraire}
\newcounter{sujettransreflitteraire}
\newcounter{sujettransintphilosophique}
\newcounter{sujettransrefphilosophique}
%%% Compteurs du nombre de sujets par type
\newcounter{sujetintlitteraire}
\newcounter{sujetreflitteraire}
\newcounter{sujetintphilosophique}
\newcounter{sujetrefphilosophique}
%%% Compteurs par discipline
%%Litt
\newcounter{sujetlitteraireETA}
\newcounter{sujetlitteraireEXS}
\newcounter{sujetlitteraireMM}
\newcounter{sujetlitteraireCC}
\newcounter{sujetlitteraireHV}
\newcounter{sujetlitteraireHL}
\newcounter{sujetlitteraireEXSTA}
\newcounter{sujetlitteraireETAMM}
\newcounter{sujetlitteraireEXSMM}
\newcounter{sujetlitteraireHVCC}
\newcounter{sujetlitteraireHVL}
\newcounter{sujetlitteraireaut}
%%Phil
\newcounter{sujetphilosophiqueETA}
\newcounter{sujetphilosophiqueEXS}
\newcounter{sujetphilosophiqueMM}
\newcounter{sujetphilosophiqueCC}
\newcounter{sujetphilosophiqueHV}
\newcounter{sujetphilosophiqueHL}
\newcounter{sujetphilosophiqueEXSTA}
\newcounter{sujetphilosophiqueETAMM}
\newcounter{sujetphilosophiqueEXSMM}
\newcounter{sujetphilosophiqueHVCC}
\newcounter{sujetphilosophiqueHVL}
\newcounter{sujetphilosophiqueaut}
%%% Autres (Ods et Odh) comptés ensemble
\newcounter{sujetallaut}
%%%Pour les graphiques%%%
\newcounter{sujetlitterairetotal}
\newcounter{sujetphilosophiquetotal}
\newcounter{sujettotal}

%%%Compteurs par centres d'examen%%%
%%%%%%%%%%%%%%%%%%%%%%%%%%%%%%%%%%%%%%%%%%%%%%%%%%%%%%%%%%%%%%%%%%%%%%%
%%%%%%%%%%%%%%%%%%%%%%%%%%%%%%Compteurs%%%%%%%%%%%%%%%%%%%%%%%%%%%%%%%%
%%%%%%%%%%%%%%%%%%%%%%%%%%%%%%%%%%%%%%%%%%%%%%%%%%%%%%%%%%%%%%%%%%%%%%%
%%%Compter les sujets par thème%%%
%\newcounter{nosujetamnord}
\newcounter{sujetamnordall}%Total de questions pour le centre
\newcounter{sujetamnordtotal}%Total des questions en comptant en double les questions pour les graphiques
\newcounter{sujetamnordHQ}
\newcounter{sujetamnordRS}
%%%sujetamnords simples%%%
\newcounter{sujetamnordEXS}
\newcounter{sujetamnordETA}
\newcounter{sujetamnordMM}
\newcounter{sujetamnordHV}
\newcounter{sujetamnordCC}
\newcounter{sujetamnordHL}
%%%sujetamnords mixtes%%%
\newcounter{sujetamnordEXSTA}
\newcounter{sujetamnordEXSMM}
\newcounter{sujetamnordETAMM}
\newcounter{sujetamnordOdS}
\newcounter{sujetamnordHVL}
\newcounter{sujetamnordHVCC}
\newcounter{sujetamnordHLCC}
\newcounter{sujetamnordOdH}
%%%sujetamnords mixtes semestres%%%
\newcounter{sujetamnordETACC}
\newcounter{sujetamnordETAHV}
\newcounter{sujetamnordETAHL}
\newcounter{sujetamnordEXSCC}
\newcounter{sujetamnordEXSHV}
\newcounter{sujetamnordEXSHL}
\newcounter{sujetamnordMMCC}
\newcounter{sujetamnordMMHV}
\newcounter{sujetamnordMMHL}

%%%sujetamnords mixtes années%%
\newcounter{sujetamnordpremterm}%%%Compteur général de sujetamnord
\newcounter{sujetamnordETAPdP}
\newcounter{sujetamnordETADFI}
\newcounter{sujetamnordETADMRC}
\newcounter{sujetamnordETAHA}
\newcounter{sujetamnordpremtermETA}
\newcounter{sujetamnordEXSPdP}
\newcounter{sujetamnordEXSDFI}
\newcounter{sujetamnordEXSDMRC}
\newcounter{sujetamnordEXSHA}
\newcounter{sujetamnordpremtermEXS}
\newcounter{sujetamnordMMPdP}
\newcounter{sujetamnordMMDFI}
\newcounter{sujetamnordMMDMRC}
\newcounter{sujetamnordMMHA}
\newcounter{sujetamnordpremtermMM}
\newcounter{sujetamnordCCPdP}
\newcounter{sujetamnordCCDFI}
\newcounter{sujetamnordCCDMRC}
\newcounter{sujetamnordCCHA}
\newcounter{sujetamnordpremtermCC}
\newcounter{sujetamnordHVPdP}
\newcounter{sujetamnordHVDFI}
\newcounter{sujetamnordHVDMRC}
\newcounter{sujetamnordHVHA}
\newcounter{sujetamnordpremtermHV}
\newcounter{sujetamnordHLPdP}
\newcounter{sujetamnordHLDFI}
\newcounter{sujetamnordHLDMRC}
\newcounter{sujetamnordHLHA}
\newcounter{sujetamnordpremtermHL}
\newcounter{sujetamnordRSPdP}
\newcounter{sujetamnordRSRdM}
\newcounter{sujetamnordRSDFI}
\newcounter{sujetamnordRSDMRC}
\newcounter{sujetamnordRSHA}
\newcounter{sujetamnordHQPdP}
\newcounter{sujetamnordHQRdM}
\newcounter{sujetamnordHQDFI}
\newcounter{sujetamnordHQDMRC}
\newcounter{sujetamnordHQHA}

%%%Somme des compteurs%%%
\newcounter{sujetamnordallETA}
\newcounter{sujetamnordallEXS}
\newcounter{sujetamnordallMM}
\newcounter{sujetamnordallRS}
\newcounter{sujetamnordallCC}
\newcounter{sujetamnordallHV}
\newcounter{sujetamnordallHL}
\newcounter{sujetamnordallHQ}

%%%Compter les sujetamnords par thème et types de questions%%%
%%%%%%%%%%%%%%%%%%%%%%%%%%%%%%%%%%%%%%%%%%%%%
%%%%%%%%%Question d'interprétation%%%%%%%%%%%
%%%%%%%%%%%%%%%%%%%%%%%%%%%%%%%%%%%%%%%%%%%%%

%%%%%%%%%%%%%%%Général%%%%%%%%%%%%%%%%%%%%%%%
\newcounter{nosujetamnordint}
\newcounter{sujetamnordHQint}
\newcounter{sujetamnordRSint}
%%%sujetamnords simples%%%
\newcounter{sujetamnordintEXS}
\newcounter{sujetamnordintETA}
\newcounter{sujetamnordintMM}
\newcounter{sujetamnordintHV}
\newcounter{sujetamnordintCC}
\newcounter{sujetamnordintHL}
%%%sujetamnords mixtes%%%
\newcounter{sujetamnordintEXSTA}
\newcounter{sujetamnordintEXSMM}
\newcounter{sujetamnordintETAMM}
\newcounter{sujetamnordintOdS}
\newcounter{sujetamnordintHVL}
\newcounter{sujetamnordintHVCC}
\newcounter{sujetamnordintHLCC}
\newcounter{sujetamnordintOdH}
%%%sujetamnords mixtes semestres%%%
\newcounter{sujetamnordintETACC}
\newcounter{sujetamnordintETAHV}
\newcounter{sujetamnordintETAHL}
\newcounter{sujetamnordintEXSCC}
\newcounter{sujetamnordintEXSHV}
\newcounter{sujetamnordintEXSHL}
\newcounter{sujetamnordintMMCC}
\newcounter{sujetamnordintMMHV}
\newcounter{sujetamnordintMMHL}
%%%sujetamnords mixtes années%%
\newcounter{sujetamnordintpremterm}
\newcounter{sujetamnordintpremtermRS}
\newcounter{sujetamnordintpremtermHQ}
\newcounter{sujetamnordintETAPdP}
\newcounter{sujetamnordintETADFI}
\newcounter{sujetamnordintETADMRC}
\newcounter{sujetamnordintETAHA}
\newcounter{sujetamnordintpremtermETA}
\newcounter{sujetamnordintEXSPdP}
\newcounter{sujetamnordintEXSDFI}
\newcounter{sujetamnordintEXSDMRC}
\newcounter{sujetamnordintEXSHA}
\newcounter{sujetamnordintpremtermEXS}
\newcounter{sujetamnordintMMPdP}
\newcounter{sujetamnordintMMDFI}
\newcounter{sujetamnordintMMDMRC}
\newcounter{sujetamnordintMMHA}
\newcounter{sujetamnordintpremtermMM}
\newcounter{sujetamnordintCCPdP}
\newcounter{sujetamnordintCCDFI}
\newcounter{sujetamnordintCCDMRC}
\newcounter{sujetamnordintCCHA}
\newcounter{sujetamnordintpremtermCC}
\newcounter{sujetamnordintHVPdP}
\newcounter{sujetamnordintHVDFI}
\newcounter{sujetamnordintHVDMRC}
\newcounter{sujetamnordintHVHA}
\newcounter{sujetamnordintpremtermHV}
\newcounter{sujetamnordintHLPdP}
\newcounter{sujetamnordintHLDFI}
\newcounter{sujetamnordintHLDMRC}
\newcounter{sujetamnordintHLHA}
\newcounter{sujetamnordintpremtermHL}
\newcounter{sujetamnordintRSPdP}
\newcounter{sujetamnordintRSRdM}
\newcounter{sujetamnordintRSDFI}
\newcounter{sujetamnordintRSDMRC}
\newcounter{sujetamnordintRSHA}
\newcounter{sujetamnordintHQPdP}
\newcounter{sujetamnordintHQRdM}
\newcounter{sujetamnordintHQDFI}
\newcounter{sujetamnordintHQDMRC}
\newcounter{sujetamnordintHQHA}
%%%Somme des compteurs%%%
\newcounter{sujetamnordallintETA}
\newcounter{sujetamnordallintEXS}
\newcounter{sujetamnordallintMM}
\newcounter{sujetamnordallintRS}
\newcounter{sujetamnordallintCC}
\newcounter{sujetamnordallintHV}
\newcounter{sujetamnordallintHL}
\newcounter{sujetamnordallintHQ}

%%%%%%%%Interprétation philosophique%%%%%%%%%
\newcounter{nosujetamnordintphilosophique}
\newcounter{sujetamnordHQintphilosophique}
\newcounter{sujetamnordRSintphilosophique}
%%%sujetamnords simples%%%
\newcounter{sujetamnordintphilosophiqueEXS}
\newcounter{sujetamnordintphilosophiqueETA}
\newcounter{sujetamnordintphilosophiqueMM}
\newcounter{sujetamnordintphilosophiqueHV}
\newcounter{sujetamnordintphilosophiqueCC}
\newcounter{sujetamnordintphilosophiqueHL}
%%%sujetamnords mixtes%%%
\newcounter{sujetamnordintphilosophiqueEXSTA}
\newcounter{sujetamnordintphilosophiqueEXSMM}
\newcounter{sujetamnordintphilosophiqueETAMM}
\newcounter{sujetamnordintphilosophiqueOdS}
\newcounter{sujetamnordintphilosophiqueHVL}
\newcounter{sujetamnordintphilosophiqueHVCC}
\newcounter{sujetamnordintphilosophiqueHLCC}
\newcounter{sujetamnordintphilosophiqueOdH}
%%%sujetamnords mixtes semestres%%%
\newcounter{sujetamnordintphilosophiqueETACC}
\newcounter{sujetamnordintphilosophiqueETAHV}
\newcounter{sujetamnordintphilosophiqueETAHL}
\newcounter{sujetamnordintphilosophiqueEXSCC}
\newcounter{sujetamnordintphilosophiqueEXSHV}
\newcounter{sujetamnordintphilosophiqueEXSHL}
\newcounter{sujetamnordintphilosophiqueMMCC}
\newcounter{sujetamnordintphilosophiqueMMHV}
\newcounter{sujetamnordintphilosophiqueMMHL}
%%%sujetamnords mixtes années%%
\newcounter{sujetamnordintphilosophiquepremterm}
\newcounter{sujetamnordintphilosophiquepremtermRS}
\newcounter{sujetamnordintphilosophiquepremtermHQ}
\newcounter{sujetamnordintphilosophiqueETAPdP}
\newcounter{sujetamnordintphilosophiqueETADFI}
\newcounter{sujetamnordintphilosophiqueETADMRC}
\newcounter{sujetamnordintphilosophiqueETAHA}
\newcounter{sujetamnordintphilosophiquepremtermETA}
\newcounter{sujetamnordintphilosophiqueEXSPdP}
\newcounter{sujetamnordintphilosophiqueEXSDFI}
\newcounter{sujetamnordintphilosophiqueEXSDMRC}
\newcounter{sujetamnordintphilosophiqueEXSHA}
\newcounter{sujetamnordintphilosophiquepremtermEXS}
\newcounter{sujetamnordintphilosophiqueMMPdP}
\newcounter{sujetamnordintphilosophiqueMMDFI}
\newcounter{sujetamnordintphilosophiqueMMDMRC}
\newcounter{sujetamnordintphilosophiqueMMHA}
\newcounter{sujetamnordintphilosophiquepremtermMM}
\newcounter{sujetamnordintphilosophiqueCCPdP}
\newcounter{sujetamnordintphilosophiqueCCDFI}
\newcounter{sujetamnordintphilosophiqueCCDMRC}
\newcounter{sujetamnordintphilosophiqueCCHA}
\newcounter{sujetamnordintphilosophiquepremtermCC}
\newcounter{sujetamnordintphilosophiqueHVPdP}
\newcounter{sujetamnordintphilosophiqueHVDFI}
\newcounter{sujetamnordintphilosophiqueHVDMRC}
\newcounter{sujetamnordintphilosophiqueHVHA}
\newcounter{sujetamnordintphilosophiquepremtermHV}
\newcounter{sujetamnordintphilosophiqueHLPdP}
\newcounter{sujetamnordintphilosophiqueHLDFI}
\newcounter{sujetamnordintphilosophiqueHLDMRC}
\newcounter{sujetamnordintphilosophiqueHLHA}
\newcounter{sujetamnordintphilosophiquepremtermHL}
\newcounter{sujetamnordintphilosophiqueRSPdP}
\newcounter{sujetamnordintphilosophiqueRSRdM}
\newcounter{sujetamnordintphilosophiqueRSDFI}
\newcounter{sujetamnordintphilosophiqueRSDMRC}
\newcounter{sujetamnordintphilosophiqueRSHA}
\newcounter{sujetamnordintphilosophiqueHQPdP}
\newcounter{sujetamnordintphilosophiqueHQRdM}
\newcounter{sujetamnordintphilosophiqueHQDFI}
\newcounter{sujetamnordintphilosophiqueHQDMRC}
\newcounter{sujetamnordintphilosophiqueHQHA}
%%%Somme des compteurs%%%
\newcounter{sujetamnordallintphilosophiqueETA}
\newcounter{sujetamnordallintphilosophiqueEXS}
\newcounter{sujetamnordallintphilosophiqueMM}
\newcounter{sujetamnordallintphilosophiqueRS}
\newcounter{sujetamnordallintphilosophiqueCC}
\newcounter{sujetamnordallintphilosophiqueHV}
\newcounter{sujetamnordallintphilosophiqueHL}
\newcounter{sujetamnordallintphilosophiqueHQ}

%%%%%%%%%Interprétation littéraire%%%%%%%%%%%
\newcounter{nosujetamnordintlitteraire}
\newcounter{sujetamnordHQintlitteraire}
\newcounter{sujetamnordRSintlitteraire}
%%%sujetamnords simples%%%
\newcounter{sujetamnordintlitteraireEXS}
\newcounter{sujetamnordintlitteraireETA}
\newcounter{sujetamnordintlitteraireMM}
\newcounter{sujetamnordintlitteraireHV}
\newcounter{sujetamnordintlitteraireCC}
\newcounter{sujetamnordintlitteraireHL}
%%%sujetamnords mixtes%%%
\newcounter{sujetamnordintlitteraireEXSTA}
\newcounter{sujetamnordintlitteraireEXSMM}
\newcounter{sujetamnordintlitteraireETAMM}
\newcounter{sujetamnordintlitteraireOdS}
\newcounter{sujetamnordintlitteraireHVL}
\newcounter{sujetamnordintlitteraireHVCC}
\newcounter{sujetamnordintlitteraireHLCC}
\newcounter{sujetamnordintlitteraireOdH}
%%%sujetamnords mixtes semestres%%%
\newcounter{sujetamnordintlitteraireETACC}
\newcounter{sujetamnordintlitteraireETAHV}
\newcounter{sujetamnordintlitteraireETAHL}
\newcounter{sujetamnordintlitteraireEXSCC}
\newcounter{sujetamnordintlitteraireEXSHV}
\newcounter{sujetamnordintlitteraireEXSHL}
\newcounter{sujetamnordintlitteraireMMCC}
\newcounter{sujetamnordintlitteraireMMHV}
\newcounter{sujetamnordintlitteraireMMHL}
%%%sujetamnords mixtes années%%
\newcounter{sujetamnordintlitterairepremterm}
\newcounter{sujetamnordintlitterairepremtermRS}
\newcounter{sujetamnordintlitterairepremtermHQ}
\newcounter{sujetamnordintlitteraireETAPdP}
\newcounter{sujetamnordintlitteraireETADFI}
\newcounter{sujetamnordintlitteraireETADMRC}
\newcounter{sujetamnordintlitteraireETAHA}
\newcounter{sujetamnordintlitterairepremtermETA}
\newcounter{sujetamnordintlitteraireEXSPdP}
\newcounter{sujetamnordintlitteraireEXSDFI}
\newcounter{sujetamnordintlitteraireEXSDMRC}
\newcounter{sujetamnordintlitteraireEXSHA}
\newcounter{sujetamnordintlitterairepremtermEXS}
\newcounter{sujetamnordintlitteraireMMPdP}
\newcounter{sujetamnordintlitteraireMMDFI}
\newcounter{sujetamnordintlitteraireMMDMRC}
\newcounter{sujetamnordintlitteraireMMHA}
\newcounter{sujetamnordintlitterairepremtermMM}
\newcounter{sujetamnordintlitteraireCCPdP}
\newcounter{sujetamnordintlitteraireCCDFI}
\newcounter{sujetamnordintlitteraireCCDMRC}
\newcounter{sujetamnordintlitteraireCCHA}
\newcounter{sujetamnordintlitterairepremtermCC}
\newcounter{sujetamnordintlitteraireHVPdP}
\newcounter{sujetamnordintlitteraireHVDFI}
\newcounter{sujetamnordintlitteraireHVDMRC}
\newcounter{sujetamnordintlitteraireHVHA}
\newcounter{sujetamnordintlitterairepremtermHV}
\newcounter{sujetamnordintlitteraireHLPdP}
\newcounter{sujetamnordintlitteraireHLDFI}
\newcounter{sujetamnordintlitteraireHLDMRC}
\newcounter{sujetamnordintlitteraireHLHA}
\newcounter{sujetamnordintlitterairepremtermHL}
\newcounter{sujetamnordintlitteraireRSPdP}
\newcounter{sujetamnordintlitteraireRSRdM}
\newcounter{sujetamnordintlitteraireRSDFI}
\newcounter{sujetamnordintlitteraireRSDMRC}
\newcounter{sujetamnordintlitteraireRSHA}
\newcounter{sujetamnordintlitteraireHQPdP}
\newcounter{sujetamnordintlitteraireHQRdM}
\newcounter{sujetamnordintlitteraireHQDFI}
\newcounter{sujetamnordintlitteraireHQDMRC}
\newcounter{sujetamnordintlitteraireHQHA}
%%%Somme des compteurs%%%
\newcounter{sujetamnordallintlitteraireETA}
\newcounter{sujetamnordallintlitteraireEXS}
\newcounter{sujetamnordallintlitteraireMM}
\newcounter{sujetamnordallintlitteraireRS}
\newcounter{sujetamnordallintlitteraireCC}
\newcounter{sujetamnordallintlitteraireHV}
\newcounter{sujetamnordallintlitteraireHL}
\newcounter{sujetamnordallintlitteraireHQ}

%%%%%%%%%%%%%%%%%%%%%%%%%%%%%%%%%%%%%%%%%%%%%
%%%%%%%%%%%%Question de réflexion%%%%%%%%%%%%
%%%%%%%%%%%%%%%%%%%%%%%%%%%%%%%%%%%%%%%%%%%%%
\newcounter{nosujetamnordref}
%%%sujetamnords simples%%%
\newcounter{sujetamnordrefEXS}
\newcounter{sujetamnordrefETA}
\newcounter{sujetamnordrefMM}
\newcounter{sujetamnordrefHV}
\newcounter{sujetamnordrefCC}
\newcounter{sujetamnordrefHL}
%%%sujetamnords mixtes%%%
\newcounter{sujetamnordrefEXSTA}
\newcounter{sujetamnordrefEXSMM}
\newcounter{sujetamnordrefETAMM}
\newcounter{sujetamnordrefOdS}
\newcounter{sujetamnordrefHVL}
\newcounter{sujetamnordrefHVCC}
\newcounter{sujetamnordrefHLCC}
\newcounter{sujetamnordrefOdH}
%%%sujetamnords mixtes semestres%%%
\newcounter{sujetamnordrefETACC}
\newcounter{sujetamnordrefETAHV}
\newcounter{sujetamnordrefETAHL}
\newcounter{sujetamnordrefEXSCC}
\newcounter{sujetamnordrefEXSHV}
\newcounter{sujetamnordrefEXSHL}
\newcounter{sujetamnordrefMMCC}
\newcounter{sujetamnordrefMMHV}
\newcounter{sujetamnordrefMMHL}
%%%sujetamnords mixtes années%%
\newcounter{sujetamnordrefpremterm}%%%Compteur général
\newcounter{sujetamnordrefETAPdP}
\newcounter{sujetamnordrefETADFI}
\newcounter{sujetamnordrefETADMRC}
\newcounter{sujetamnordrefETAHA}
\newcounter{sujetamnordrefpremtermETA}
\newcounter{sujetamnordrefEXSPdP}
\newcounter{sujetamnordrefEXSDFI}
\newcounter{sujetamnordrefEXSDMRC}
\newcounter{sujetamnordrefEXSHA}
\newcounter{sujetamnordrefpremtermEXS}
\newcounter{sujetamnordrefMMPdP}
\newcounter{sujetamnordrefMMDFI}
\newcounter{sujetamnordrefMMDMRC}
\newcounter{sujetamnordrefMMHA}
\newcounter{sujetamnordrefpremtermMM}
\newcounter{sujetamnordrefCCPdP}
\newcounter{sujetamnordrefCCDFI}
\newcounter{sujetamnordrefCCDMRC}
\newcounter{sujetamnordrefCCHA}
\newcounter{sujetamnordrefpremtermCC}
\newcounter{sujetamnordrefHVPdP}
\newcounter{sujetamnordrefHVDFI}
\newcounter{sujetamnordrefHVDMRC}
\newcounter{sujetamnordrefHVHA}
\newcounter{sujetamnordrefpremtermHV}
\newcounter{sujetamnordrefHLPdP}
\newcounter{sujetamnordrefHLDFI}
\newcounter{sujetamnordrefHLDMRC}
\newcounter{sujetamnordrefHLHA}
\newcounter{sujetamnordrefpremtermHL}
\newcounter{sujetamnordrefRSPdP}
\newcounter{sujetamnordrefRSRdM}
\newcounter{sujetamnordrefRSDFI}
\newcounter{sujetamnordrefRSDMRC}
\newcounter{sujetamnordrefRSHA}
\newcounter{sujetamnordrefHQPdP}
\newcounter{sujetamnordrefHQRdM}
\newcounter{sujetamnordrefHQDFI}
\newcounter{sujetamnordrefHQDMRC}
\newcounter{sujetamnordrefHQHA}
%%%Somme des compteurs%%%
\newcounter{sujetamnordallrefETA}
\newcounter{sujetamnordallrefEXS}
\newcounter{sujetamnordallrefMM}
\newcounter{sujetamnordallrefRS}
\newcounter{sujetamnordallrefCC}
\newcounter{sujetamnordallrefHV}
\newcounter{sujetamnordallrefHL}
\newcounter{sujetamnordallrefHQ}

%%%%%%%%Réflexion philosophique%%%%%%%%%
\newcounter{nosujetamnordrefphilosophique}
\newcounter{sujetamnordHQrefphilosophique}
\newcounter{sujetamnordRSrefphilosophique}
%%%sujetamnords simples%%%
\newcounter{sujetamnordrefphilosophiqueEXS}
\newcounter{sujetamnordrefphilosophiqueETA}
\newcounter{sujetamnordrefphilosophiqueMM}
\newcounter{sujetamnordrefphilosophiqueHV}
\newcounter{sujetamnordrefphilosophiqueCC}
\newcounter{sujetamnordrefphilosophiqueHL}
%%%sujetamnords mixtes%%%
\newcounter{sujetamnordrefphilosophiqueEXSTA}
\newcounter{sujetamnordrefphilosophiqueEXSMM}
\newcounter{sujetamnordrefphilosophiqueETAMM}
\newcounter{sujetamnordrefphilosophiqueOdS}
\newcounter{sujetamnordrefphilosophiqueHVL}
\newcounter{sujetamnordrefphilosophiqueHVCC}
\newcounter{sujetamnordrefphilosophiqueHLCC}
\newcounter{sujetamnordrefphilosophiqueOdH}
%%%sujetamnords mixtes semestres%%%
\newcounter{sujetamnordrefphilosophiqueETACC}
\newcounter{sujetamnordrefphilosophiqueETAHV}
\newcounter{sujetamnordrefphilosophiqueETAHL}
\newcounter{sujetamnordrefphilosophiqueEXSCC}
\newcounter{sujetamnordrefphilosophiqueEXSHV}
\newcounter{sujetamnordrefphilosophiqueEXSHL}
\newcounter{sujetamnordrefphilosophiqueMMCC}
\newcounter{sujetamnordrefphilosophiqueMMHV}
\newcounter{sujetamnordrefphilosophiqueMMHL}
%%%sujetamnords mixtes années%%
\newcounter{sujetamnordrefphilosophiquepremterm}
\newcounter{sujetamnordrefphilosophiquepremtermRS}
\newcounter{sujetamnordrefphilosophiquepremtermHQ}
\newcounter{sujetamnordrefphilosophiqueETAPdP}
\newcounter{sujetamnordrefphilosophiqueETADFI}
\newcounter{sujetamnordrefphilosophiqueETADMRC}
\newcounter{sujetamnordrefphilosophiqueETAHA}\newcounter{sujetamnordrefphilosophiquepremtermETA}
\newcounter{sujetamnordrefphilosophiqueEXSPdP}
\newcounter{sujetamnordrefphilosophiqueEXSDFI}
\newcounter{sujetamnordrefphilosophiqueEXSDMRC}
\newcounter{sujetamnordrefphilosophiqueEXSHA}
\newcounter{sujetamnordrefphilosophiquepremtermEXS}
\newcounter{sujetamnordrefphilosophiqueMMPdP}
\newcounter{sujetamnordrefphilosophiqueMMDFI}
\newcounter{sujetamnordrefphilosophiqueMMDMRC}
\newcounter{sujetamnordrefphilosophiqueMMHA}
\newcounter{sujetamnordrefphilosophiquepremtermMM}
\newcounter{sujetamnordrefphilosophiqueCCPdP}
\newcounter{sujetamnordrefphilosophiqueCCDFI}
\newcounter{sujetamnordrefphilosophiqueCCDMRC}
\newcounter{sujetamnordrefphilosophiqueCCHA}
\newcounter{sujetamnordrefphilosophiquepremtermCC}
\newcounter{sujetamnordrefphilosophiqueHVPdP}
\newcounter{sujetamnordrefphilosophiqueHVDFI}
\newcounter{sujetamnordrefphilosophiqueHVDMRC}
\newcounter{sujetamnordrefphilosophiqueHVHA}
\newcounter{sujetamnordrefphilosophiquepremtermHV}
\newcounter{sujetamnordrefphilosophiqueHLPdP}
\newcounter{sujetamnordrefphilosophiqueHLDFI}
\newcounter{sujetamnordrefphilosophiqueHLDMRC}
\newcounter{sujetamnordrefphilosophiqueHLHA}
\newcounter{sujetamnordrefphilosophiquepremtermHL}
\newcounter{sujetamnordrefphilosophiqueRSPdP}
\newcounter{sujetamnordrefphilosophiqueRSRdM}
\newcounter{sujetamnordrefphilosophiqueRSDFI}
\newcounter{sujetamnordrefphilosophiqueRSDMRC}
\newcounter{sujetamnordrefphilosophiqueRSHA}
\newcounter{sujetamnordrefphilosophiqueHQPdP}
\newcounter{sujetamnordrefphilosophiqueHQRdM}
\newcounter{sujetamnordrefphilosophiqueHQDFI}
\newcounter{sujetamnordrefphilosophiqueHQDMRC}
\newcounter{sujetamnordrefphilosophiqueHQHA}
%%%Somme des compteurs%%%
\newcounter{sujetamnordallrefphilosophiqueETA}
\newcounter{sujetamnordallrefphilosophiqueEXS}
\newcounter{sujetamnordallrefphilosophiqueMM}
\newcounter{sujetamnordallrefphilosophiqueRS}
\newcounter{sujetamnordallrefphilosophiqueCC}
\newcounter{sujetamnordallrefphilosophiqueHV}
\newcounter{sujetamnordallrefphilosophiqueHL}
\newcounter{sujetamnordallrefphilosophiqueHQ}

%%%%%%%%%Réflexion littéraire%%%%%%%%%%%
\newcounter{nosujetamnordreflitteraire}
\newcounter{sujetamnordHQreflitteraire}
\newcounter{sujetamnordRSreflitteraire}
%%%sujetamnords simples%%%
\newcounter{sujetamnordreflitteraireEXS}
\newcounter{sujetamnordreflitteraireETA}
\newcounter{sujetamnordreflitteraireMM}
\newcounter{sujetamnordreflitteraireHV}
\newcounter{sujetamnordreflitteraireCC}
\newcounter{sujetamnordreflitteraireHL}
%%%sujetamnords mixtes%%%
\newcounter{sujetamnordreflitteraireEXSTA}
\newcounter{sujetamnordreflitteraireEXSMM}
\newcounter{sujetamnordreflitteraireETAMM}
\newcounter{sujetamnordreflitteraireOdS}
\newcounter{sujetamnordreflitteraireHVL}
\newcounter{sujetamnordreflitteraireHVCC}
\newcounter{sujetamnordreflitteraireHLCC}
\newcounter{sujetamnordreflitteraireOdH}
%%%sujetamnords mixtes semestres%%%
\newcounter{sujetamnordreflitteraireETACC}
\newcounter{sujetamnordreflitteraireETAHV}
\newcounter{sujetamnordreflitteraireETAHL}
\newcounter{sujetamnordreflitteraireEXSCC}
\newcounter{sujetamnordreflitteraireEXSHV}
\newcounter{sujetamnordreflitteraireEXSHL}
\newcounter{sujetamnordreflitteraireMMCC}
\newcounter{sujetamnordreflitteraireMMHV}
\newcounter{sujetamnordreflitteraireMMHL}
%%%sujetamnords mixtes années%%
\newcounter{sujetamnordreflitterairepremterm}
\newcounter{sujetamnordreflitterairepremtermRS}
\newcounter{sujetamnordreflitterairepremtermHQ}
\newcounter{sujetamnordreflitteraireETAPdP}
\newcounter{sujetamnordreflitteraireETADFI}
\newcounter{sujetamnordreflitteraireETADMRC}
\newcounter{sujetamnordreflitteraireETAHA}
\newcounter{sujetamnordreflitterairepremtermETA}
\newcounter{sujetamnordreflitteraireEXSPdP}
\newcounter{sujetamnordreflitteraireEXSDFI}
\newcounter{sujetamnordreflitteraireEXSDMRC}
\newcounter{sujetamnordreflitteraireEXSHA}
\newcounter{sujetamnordreflitterairepremtermEXS}
\newcounter{sujetamnordreflitteraireMMPdP}
\newcounter{sujetamnordreflitteraireMMDFI}
\newcounter{sujetamnordreflitteraireMMDMRC}
\newcounter{sujetamnordreflitteraireMMHA}
\newcounter{sujetamnordreflitterairepremtermMM}
\newcounter{sujetamnordreflitteraireCCPdP}
\newcounter{sujetamnordreflitteraireCCDFI}
\newcounter{sujetamnordreflitteraireCCDMRC}
\newcounter{sujetamnordreflitteraireCCHA}
\newcounter{sujetamnordreflitterairepremtermCC}
\newcounter{sujetamnordreflitteraireHVPdP}
\newcounter{sujetamnordreflitteraireHVDFI}
\newcounter{sujetamnordreflitteraireHVDMRC}
\newcounter{sujetamnordreflitteraireHVHA}
\newcounter{sujetamnordreflitterairepremtermHV}
\newcounter{sujetamnordreflitteraireHLPdP}
\newcounter{sujetamnordreflitteraireHLDFI}
\newcounter{sujetamnordreflitteraireHLDMRC}
\newcounter{sujetamnordreflitteraireHLHA}
\newcounter{sujetamnordreflitterairepremtermHL}
\newcounter{sujetamnordreflitteraireRSPdP}
\newcounter{sujetamnordreflitteraireRSRdM}
\newcounter{sujetamnordreflitteraireRSDFI}
\newcounter{sujetamnordreflitteraireRSDMRC}
\newcounter{sujetamnordreflitteraireRSHA}
\newcounter{sujetamnordreflitteraireHQPdP}
\newcounter{sujetamnordreflitteraireHQRdM}
\newcounter{sujetamnordreflitteraireHQDFI}
\newcounter{sujetamnordreflitteraireHQDMRC}
\newcounter{sujetamnordreflitteraireHQHA}
%%%Somme des compteurs%%%
\newcounter{sujetamnordallreflitteraireETA}
\newcounter{sujetamnordallreflitteraireEXS}
\newcounter{sujetamnordallreflitteraireMM}
\newcounter{sujetamnordallreflitteraireCC}
\newcounter{sujetamnordallreflitteraireRS}
\newcounter{sujetamnordallreflitteraireHV}
\newcounter{sujetamnordallreflitteraireHL}
\newcounter{sujetamnordallreflitteraireHQ}
%%%Compteurs de statistiques globales sur les genres de question%%%
\newcounter{sujetamnordexclRS}
\newcounter{sujetamnordmixtRS}
\newcounter{sujetamnordexclHQ}
\newcounter{sujetamnordmixtHQ}
\newcounter{sujetamnordmixtHQRS}
%%%Compteurs de statistiques sur les transversaux (divers semestres, première et terminale)
\newcounter{sujettransamnordintlitteraire}
\newcounter{sujettransamnordreflitteraire}
\newcounter{sujettransamnordintphilosophique}
\newcounter{sujettransamnordrefphilosophique}
%%% Compteurs du nombre de sujets par type
\newcounter{sujetamnordintlitteraire}
\newcounter{sujetamnordreflitteraire}
\newcounter{sujetamnordintphilosophique}
\newcounter{sujetamnordrefphilosophique}
%%% Compteurs par discipline
%%Litt
\newcounter{sujetamnordlitteraireETA}
\newcounter{sujetamnordlitteraireEXS}
\newcounter{sujetamnordlitteraireMM}
\newcounter{sujetamnordlitteraireCC}
\newcounter{sujetamnordlitteraireHV}
\newcounter{sujetamnordlitteraireHL}
\newcounter{sujetamnordlitteraireEXSTA}
\newcounter{sujetamnordlitteraireETAMM}
\newcounter{sujetamnordlitteraireEXSMM}
\newcounter{sujetamnordlitteraireHVCC}
\newcounter{sujetamnordlitteraireHVL}
\newcounter{sujetamnordlitteraireaut}
%%Phil
\newcounter{sujetamnordphilosophiqueETA}
\newcounter{sujetamnordphilosophiqueEXS}
\newcounter{sujetamnordphilosophiqueMM}
\newcounter{sujetamnordphilosophiqueCC}
\newcounter{sujetamnordphilosophiqueHV}
\newcounter{sujetamnordphilosophiqueHL}
\newcounter{sujetamnordphilosophiqueEXSTA}
\newcounter{sujetamnordphilosophiqueETAMM}
\newcounter{sujetamnordphilosophiqueEXSMM}
\newcounter{sujetamnordphilosophiqueHVCC}
\newcounter{sujetamnordphilosophiqueHVL}
\newcounter{sujetamnordphilosophiqueaut}
%%Totaux pour les graphiques
\newcounter{sujetamnordlitterairetotal}
\newcounter{sujetamnordphilosophiquetotal}
%%% Autres (Ods et Odh) comptés ensemble
\newcounter{sujetamnordallaut}
%%%%%%%%%%%%%%%%%%%%%%%%%%%%%%%%%%%%%%%%%%%%%%%%%%%%%%%%%%%%%%%%%%%%%%%
%%%%%%%%%%%%%%%%%%%%%%%%%%%%%%Compteurs%%%%%%%%%%%%%%%%%%%%%%%%%%%%%%%%
%%%%%%%%%%%%%%%%%%%%%%%%%%%%%%%%%%%%%%%%%%%%%%%%%%%%%%%%%%%%%%%%%%%%%%%
%%%Compter les sujets par thème%%%
%\newcounter{nosujetamsud}
\newcounter{sujetamsudall}%Total de questions pour le centre
\newcounter{sujetamsudtotal}%Total des questions en comptant en double les questions pour les graphiques
\newcounter{sujetamsudHQ}
\newcounter{sujetamsudRS}
%%%sujetamsuds simples%%%
\newcounter{sujetamsudEXS}
\newcounter{sujetamsudETA}
\newcounter{sujetamsudMM}
\newcounter{sujetamsudHV}
\newcounter{sujetamsudCC}
\newcounter{sujetamsudHL}
%%%sujetamsuds mixtes%%%
\newcounter{sujetamsudEXSTA}
\newcounter{sujetamsudEXSMM}
\newcounter{sujetamsudETAMM}
\newcounter{sujetamsudOdS}
\newcounter{sujetamsudHVL}
\newcounter{sujetamsudHVCC}
\newcounter{sujetamsudHLCC}
\newcounter{sujetamsudOdH}
%%%sujetamsuds mixtes semestres%%%
\newcounter{sujetamsudETACC}
\newcounter{sujetamsudETAHV}
\newcounter{sujetamsudETAHL}
\newcounter{sujetamsudEXSCC}
\newcounter{sujetamsudEXSHV}
\newcounter{sujetamsudEXSHL}
\newcounter{sujetamsudMMCC}
\newcounter{sujetamsudMMHV}
\newcounter{sujetamsudMMHL}

%%%sujetamsuds mixtes années%%
\newcounter{sujetamsudpremterm}%%%Compteur général de sujetamsud
\newcounter{sujetamsudETAPdP}
\newcounter{sujetamsudETADFI}
\newcounter{sujetamsudETADMRC}
\newcounter{sujetamsudETAHA}
\newcounter{sujetamsudpremtermETA}
\newcounter{sujetamsudEXSPdP}
\newcounter{sujetamsudEXSDFI}
\newcounter{sujetamsudEXSDMRC}
\newcounter{sujetamsudEXSHA}
\newcounter{sujetamsudpremtermEXS}
\newcounter{sujetamsudMMPdP}
\newcounter{sujetamsudMMDFI}
\newcounter{sujetamsudMMDMRC}
\newcounter{sujetamsudMMHA}
\newcounter{sujetamsudpremtermMM}
\newcounter{sujetamsudCCPdP}
\newcounter{sujetamsudCCDFI}
\newcounter{sujetamsudCCDMRC}
\newcounter{sujetamsudCCHA}
\newcounter{sujetamsudpremtermCC}
\newcounter{sujetamsudHVPdP}
\newcounter{sujetamsudHVDFI}
\newcounter{sujetamsudHVDMRC}
\newcounter{sujetamsudHVHA}
\newcounter{sujetamsudpremtermHV}
\newcounter{sujetamsudHLPdP}
\newcounter{sujetamsudHLDFI}
\newcounter{sujetamsudHLDMRC}
\newcounter{sujetamsudHLHA}
\newcounter{sujetamsudpremtermHL}
\newcounter{sujetamsudRSPdP}
\newcounter{sujetamsudRSRdM}
\newcounter{sujetamsudRSDFI}
\newcounter{sujetamsudRSDMRC}
\newcounter{sujetamsudRSHA}
\newcounter{sujetamsudHQPdP}
\newcounter{sujetamsudHQRdM}
\newcounter{sujetamsudHQDFI}
\newcounter{sujetamsudHQDMRC}
\newcounter{sujetamsudHQHA}

%%%Somme des compteurs%%%
\newcounter{sujetamsudallETA}
\newcounter{sujetamsudallEXS}
\newcounter{sujetamsudallMM}
\newcounter{sujetamsudallRS}
\newcounter{sujetamsudallCC}
\newcounter{sujetamsudallHV}
\newcounter{sujetamsudallHL}
\newcounter{sujetamsudallHQ}

%%%Compter les sujetamsuds par thème et types de questions%%%
%%%%%%%%%%%%%%%%%%%%%%%%%%%%%%%%%%%%%%%%%%%%%
%%%%%%%%%Question d'interprétation%%%%%%%%%%%
%%%%%%%%%%%%%%%%%%%%%%%%%%%%%%%%%%%%%%%%%%%%%

%%%%%%%%%%%%%%%Général%%%%%%%%%%%%%%%%%%%%%%%
\newcounter{nosujetamsudint}
\newcounter{sujetamsudHQint}
\newcounter{sujetamsudRSint}
%%%sujetamsuds simples%%%
\newcounter{sujetamsudintEXS}
\newcounter{sujetamsudintETA}
\newcounter{sujetamsudintMM}
\newcounter{sujetamsudintHV}
\newcounter{sujetamsudintCC}
\newcounter{sujetamsudintHL}
%%%sujetamsuds mixtes%%%
\newcounter{sujetamsudintEXSTA}
\newcounter{sujetamsudintEXSMM}
\newcounter{sujetamsudintETAMM}
\newcounter{sujetamsudintOdS}
\newcounter{sujetamsudintHVL}
\newcounter{sujetamsudintHVCC}
\newcounter{sujetamsudintHLCC}
\newcounter{sujetamsudintOdH}
%%%sujetamsuds mixtes semestres%%%
\newcounter{sujetamsudintETACC}
\newcounter{sujetamsudintETAHV}
\newcounter{sujetamsudintETAHL}
\newcounter{sujetamsudintEXSCC}
\newcounter{sujetamsudintEXSHV}
\newcounter{sujetamsudintEXSHL}
\newcounter{sujetamsudintMMCC}
\newcounter{sujetamsudintMMHV}
\newcounter{sujetamsudintMMHL}
%%%sujetamsuds mixtes années%%
\newcounter{sujetamsudintpremterm}
\newcounter{sujetamsudintpremtermRS}
\newcounter{sujetamsudintpremtermHQ}
\newcounter{sujetamsudintETAPdP}
\newcounter{sujetamsudintETADFI}
\newcounter{sujetamsudintETADMRC}
\newcounter{sujetamsudintETAHA}
\newcounter{sujetamsudintpremtermETA}
\newcounter{sujetamsudintEXSPdP}
\newcounter{sujetamsudintEXSDFI}
\newcounter{sujetamsudintEXSDMRC}
\newcounter{sujetamsudintEXSHA}
\newcounter{sujetamsudintpremtermEXS}
\newcounter{sujetamsudintMMPdP}
\newcounter{sujetamsudintMMDFI}
\newcounter{sujetamsudintMMDMRC}
\newcounter{sujetamsudintMMHA}
\newcounter{sujetamsudintpremtermMM}
\newcounter{sujetamsudintCCPdP}
\newcounter{sujetamsudintCCDFI}
\newcounter{sujetamsudintCCDMRC}
\newcounter{sujetamsudintCCHA}
\newcounter{sujetamsudintpremtermCC}
\newcounter{sujetamsudintHVPdP}
\newcounter{sujetamsudintHVDFI}
\newcounter{sujetamsudintHVDMRC}
\newcounter{sujetamsudintHVHA}
\newcounter{sujetamsudintpremtermHV}
\newcounter{sujetamsudintHLPdP}
\newcounter{sujetamsudintHLDFI}
\newcounter{sujetamsudintHLDMRC}
\newcounter{sujetamsudintHLHA}
\newcounter{sujetamsudintpremtermHL}
\newcounter{sujetamsudintRSPdP}
\newcounter{sujetamsudintRSRdM}
\newcounter{sujetamsudintRSDFI}
\newcounter{sujetamsudintRSDMRC}
\newcounter{sujetamsudintRSHA}
\newcounter{sujetamsudintHQPdP}
\newcounter{sujetamsudintHQRdM}
\newcounter{sujetamsudintHQDFI}
\newcounter{sujetamsudintHQDMRC}
\newcounter{sujetamsudintHQHA}
%%%Somme des compteurs%%%
\newcounter{sujetamsudallintETA}
\newcounter{sujetamsudallintEXS}
\newcounter{sujetamsudallintMM}
\newcounter{sujetamsudallintRS}
\newcounter{sujetamsudallintCC}
\newcounter{sujetamsudallintHV}
\newcounter{sujetamsudallintHL}
\newcounter{sujetamsudallintHQ}

%%%%%%%%Interprétation philosophique%%%%%%%%%
\newcounter{nosujetamsudintphilosophique}
\newcounter{sujetamsudHQintphilosophique}
\newcounter{sujetamsudRSintphilosophique}
%%%sujetamsuds simples%%%
\newcounter{sujetamsudintphilosophiqueEXS}
\newcounter{sujetamsudintphilosophiqueETA}
\newcounter{sujetamsudintphilosophiqueMM}
\newcounter{sujetamsudintphilosophiqueHV}
\newcounter{sujetamsudintphilosophiqueCC}
\newcounter{sujetamsudintphilosophiqueHL}
%%%sujetamsuds mixtes%%%
\newcounter{sujetamsudintphilosophiqueEXSTA}
\newcounter{sujetamsudintphilosophiqueEXSMM}
\newcounter{sujetamsudintphilosophiqueETAMM}
\newcounter{sujetamsudintphilosophiqueOdS}
\newcounter{sujetamsudintphilosophiqueHVL}
\newcounter{sujetamsudintphilosophiqueHVCC}
\newcounter{sujetamsudintphilosophiqueHLCC}
\newcounter{sujetamsudintphilosophiqueOdH}
%%%sujetamsuds mixtes semestres%%%
\newcounter{sujetamsudintphilosophiqueETACC}
\newcounter{sujetamsudintphilosophiqueETAHV}
\newcounter{sujetamsudintphilosophiqueETAHL}
\newcounter{sujetamsudintphilosophiqueEXSCC}
\newcounter{sujetamsudintphilosophiqueEXSHV}
\newcounter{sujetamsudintphilosophiqueEXSHL}
\newcounter{sujetamsudintphilosophiqueMMCC}
\newcounter{sujetamsudintphilosophiqueMMHV}
\newcounter{sujetamsudintphilosophiqueMMHL}
%%%sujetamsuds mixtes années%%
\newcounter{sujetamsudintphilosophiquepremterm}
\newcounter{sujetamsudintphilosophiquepremtermRS}
\newcounter{sujetamsudintphilosophiquepremtermHQ}
\newcounter{sujetamsudintphilosophiqueETAPdP}
\newcounter{sujetamsudintphilosophiqueETADFI}
\newcounter{sujetamsudintphilosophiqueETADMRC}
\newcounter{sujetamsudintphilosophiqueETAHA}
\newcounter{sujetamsudintphilosophiquepremtermETA}
\newcounter{sujetamsudintphilosophiqueEXSPdP}
\newcounter{sujetamsudintphilosophiqueEXSDFI}
\newcounter{sujetamsudintphilosophiqueEXSDMRC}
\newcounter{sujetamsudintphilosophiqueEXSHA}
\newcounter{sujetamsudintphilosophiquepremtermEXS}
\newcounter{sujetamsudintphilosophiqueMMPdP}
\newcounter{sujetamsudintphilosophiqueMMDFI}
\newcounter{sujetamsudintphilosophiqueMMDMRC}
\newcounter{sujetamsudintphilosophiqueMMHA}
\newcounter{sujetamsudintphilosophiquepremtermMM}
\newcounter{sujetamsudintphilosophiqueCCPdP}
\newcounter{sujetamsudintphilosophiqueCCDFI}
\newcounter{sujetamsudintphilosophiqueCCDMRC}
\newcounter{sujetamsudintphilosophiqueCCHA}
\newcounter{sujetamsudintphilosophiquepremtermCC}
\newcounter{sujetamsudintphilosophiqueHVPdP}
\newcounter{sujetamsudintphilosophiqueHVDFI}
\newcounter{sujetamsudintphilosophiqueHVDMRC}
\newcounter{sujetamsudintphilosophiqueHVHA}
\newcounter{sujetamsudintphilosophiquepremtermHV}
\newcounter{sujetamsudintphilosophiqueHLPdP}
\newcounter{sujetamsudintphilosophiqueHLDFI}
\newcounter{sujetamsudintphilosophiqueHLDMRC}
\newcounter{sujetamsudintphilosophiqueHLHA}
\newcounter{sujetamsudintphilosophiquepremtermHL}
\newcounter{sujetamsudintphilosophiqueRSPdP}
\newcounter{sujetamsudintphilosophiqueRSRdM}
\newcounter{sujetamsudintphilosophiqueRSDFI}
\newcounter{sujetamsudintphilosophiqueRSDMRC}
\newcounter{sujetamsudintphilosophiqueRSHA}
\newcounter{sujetamsudintphilosophiqueHQPdP}
\newcounter{sujetamsudintphilosophiqueHQRdM}
\newcounter{sujetamsudintphilosophiqueHQDFI}
\newcounter{sujetamsudintphilosophiqueHQDMRC}
\newcounter{sujetamsudintphilosophiqueHQHA}
%%%Somme des compteurs%%%
\newcounter{sujetamsudallintphilosophiqueETA}
\newcounter{sujetamsudallintphilosophiqueEXS}
\newcounter{sujetamsudallintphilosophiqueMM}
\newcounter{sujetamsudallintphilosophiqueRS}
\newcounter{sujetamsudallintphilosophiqueCC}
\newcounter{sujetamsudallintphilosophiqueHV}
\newcounter{sujetamsudallintphilosophiqueHL}
\newcounter{sujetamsudallintphilosophiqueHQ}

%%%%%%%%%Interprétation littéraire%%%%%%%%%%%
\newcounter{nosujetamsudintlitteraire}
\newcounter{sujetamsudHQintlitteraire}
\newcounter{sujetamsudRSintlitteraire}
%%%sujetamsuds simples%%%
\newcounter{sujetamsudintlitteraireEXS}
\newcounter{sujetamsudintlitteraireETA}
\newcounter{sujetamsudintlitteraireMM}
\newcounter{sujetamsudintlitteraireHV}
\newcounter{sujetamsudintlitteraireCC}
\newcounter{sujetamsudintlitteraireHL}
%%%sujetamsuds mixtes%%%
\newcounter{sujetamsudintlitteraireEXSTA}
\newcounter{sujetamsudintlitteraireEXSMM}
\newcounter{sujetamsudintlitteraireETAMM}
\newcounter{sujetamsudintlitteraireOdS}
\newcounter{sujetamsudintlitteraireHVL}
\newcounter{sujetamsudintlitteraireHVCC}
\newcounter{sujetamsudintlitteraireHLCC}
\newcounter{sujetamsudintlitteraireOdH}
%%%sujetamsuds mixtes semestres%%%
\newcounter{sujetamsudintlitteraireETACC}
\newcounter{sujetamsudintlitteraireETAHV}
\newcounter{sujetamsudintlitteraireETAHL}
\newcounter{sujetamsudintlitteraireEXSCC}
\newcounter{sujetamsudintlitteraireEXSHV}
\newcounter{sujetamsudintlitteraireEXSHL}
\newcounter{sujetamsudintlitteraireMMCC}
\newcounter{sujetamsudintlitteraireMMHV}
\newcounter{sujetamsudintlitteraireMMHL}
%%%sujetamsuds mixtes années%%
\newcounter{sujetamsudintlitterairepremterm}
\newcounter{sujetamsudintlitterairepremtermRS}
\newcounter{sujetamsudintlitterairepremtermHQ}
\newcounter{sujetamsudintlitteraireETAPdP}
\newcounter{sujetamsudintlitteraireETADFI}
\newcounter{sujetamsudintlitteraireETADMRC}
\newcounter{sujetamsudintlitteraireETAHA}
\newcounter{sujetamsudintlitterairepremtermETA}
\newcounter{sujetamsudintlitteraireEXSPdP}
\newcounter{sujetamsudintlitteraireEXSDFI}
\newcounter{sujetamsudintlitteraireEXSDMRC}
\newcounter{sujetamsudintlitteraireEXSHA}
\newcounter{sujetamsudintlitterairepremtermEXS}
\newcounter{sujetamsudintlitteraireMMPdP}
\newcounter{sujetamsudintlitteraireMMDFI}
\newcounter{sujetamsudintlitteraireMMDMRC}
\newcounter{sujetamsudintlitteraireMMHA}
\newcounter{sujetamsudintlitterairepremtermMM}
\newcounter{sujetamsudintlitteraireCCPdP}
\newcounter{sujetamsudintlitteraireCCDFI}
\newcounter{sujetamsudintlitteraireCCDMRC}
\newcounter{sujetamsudintlitteraireCCHA}
\newcounter{sujetamsudintlitterairepremtermCC}
\newcounter{sujetamsudintlitteraireHVPdP}
\newcounter{sujetamsudintlitteraireHVDFI}
\newcounter{sujetamsudintlitteraireHVDMRC}
\newcounter{sujetamsudintlitteraireHVHA}
\newcounter{sujetamsudintlitterairepremtermHV}
\newcounter{sujetamsudintlitteraireHLPdP}
\newcounter{sujetamsudintlitteraireHLDFI}
\newcounter{sujetamsudintlitteraireHLDMRC}
\newcounter{sujetamsudintlitteraireHLHA}
\newcounter{sujetamsudintlitterairepremtermHL}
\newcounter{sujetamsudintlitteraireRSPdP}
\newcounter{sujetamsudintlitteraireRSRdM}
\newcounter{sujetamsudintlitteraireRSDFI}
\newcounter{sujetamsudintlitteraireRSDMRC}
\newcounter{sujetamsudintlitteraireRSHA}
\newcounter{sujetamsudintlitteraireHQPdP}
\newcounter{sujetamsudintlitteraireHQRdM}
\newcounter{sujetamsudintlitteraireHQDFI}
\newcounter{sujetamsudintlitteraireHQDMRC}
\newcounter{sujetamsudintlitteraireHQHA}
%%%Somme des compteurs%%%
\newcounter{sujetamsudallintlitteraireETA}
\newcounter{sujetamsudallintlitteraireEXS}
\newcounter{sujetamsudallintlitteraireMM}
\newcounter{sujetamsudallintlitteraireRS}
\newcounter{sujetamsudallintlitteraireCC}
\newcounter{sujetamsudallintlitteraireHV}
\newcounter{sujetamsudallintlitteraireHL}
\newcounter{sujetamsudallintlitteraireHQ}

%%%%%%%%%%%%%%%%%%%%%%%%%%%%%%%%%%%%%%%%%%%%%
%%%%%%%%%%%%Question de réflexion%%%%%%%%%%%%
%%%%%%%%%%%%%%%%%%%%%%%%%%%%%%%%%%%%%%%%%%%%%
\newcounter{nosujetamsudref}
%%%sujetamsuds simples%%%
\newcounter{sujetamsudrefEXS}
\newcounter{sujetamsudrefETA}
\newcounter{sujetamsudrefMM}
\newcounter{sujetamsudrefHV}
\newcounter{sujetamsudrefCC}
\newcounter{sujetamsudrefHL}
%%%sujetamsuds mixtes%%%
\newcounter{sujetamsudrefEXSTA}
\newcounter{sujetamsudrefEXSMM}
\newcounter{sujetamsudrefETAMM}
\newcounter{sujetamsudrefOdS}
\newcounter{sujetamsudrefHVL}
\newcounter{sujetamsudrefHVCC}
\newcounter{sujetamsudrefHLCC}
\newcounter{sujetamsudrefOdH}
%%%sujetamsuds mixtes semestres%%%
\newcounter{sujetamsudrefETACC}
\newcounter{sujetamsudrefETAHV}
\newcounter{sujetamsudrefETAHL}
\newcounter{sujetamsudrefEXSCC}
\newcounter{sujetamsudrefEXSHV}
\newcounter{sujetamsudrefEXSHL}
\newcounter{sujetamsudrefMMCC}
\newcounter{sujetamsudrefMMHV}
\newcounter{sujetamsudrefMMHL}
%%%sujetamsuds mixtes années%%
\newcounter{sujetamsudrefpremterm}%%%Compteur général
\newcounter{sujetamsudrefETAPdP}
\newcounter{sujetamsudrefETADFI}
\newcounter{sujetamsudrefETADMRC}
\newcounter{sujetamsudrefETAHA}
\newcounter{sujetamsudrefpremtermETA}
\newcounter{sujetamsudrefEXSPdP}
\newcounter{sujetamsudrefEXSDFI}
\newcounter{sujetamsudrefEXSDMRC}
\newcounter{sujetamsudrefEXSHA}
\newcounter{sujetamsudrefpremtermEXS}
\newcounter{sujetamsudrefMMPdP}
\newcounter{sujetamsudrefMMDFI}
\newcounter{sujetamsudrefMMDMRC}
\newcounter{sujetamsudrefMMHA}
\newcounter{sujetamsudrefpremtermMM}
\newcounter{sujetamsudrefCCPdP}
\newcounter{sujetamsudrefCCDFI}
\newcounter{sujetamsudrefCCDMRC}
\newcounter{sujetamsudrefCCHA}
\newcounter{sujetamsudrefpremtermCC}
\newcounter{sujetamsudrefHVPdP}
\newcounter{sujetamsudrefHVDFI}
\newcounter{sujetamsudrefHVDMRC}
\newcounter{sujetamsudrefHVHA}
\newcounter{sujetamsudrefpremtermHV}
\newcounter{sujetamsudrefHLPdP}
\newcounter{sujetamsudrefHLDFI}
\newcounter{sujetamsudrefHLDMRC}
\newcounter{sujetamsudrefHLHA}
\newcounter{sujetamsudrefpremtermHL}
\newcounter{sujetamsudrefRSPdP}
\newcounter{sujetamsudrefRSRdM}
\newcounter{sujetamsudrefRSDFI}
\newcounter{sujetamsudrefRSDMRC}
\newcounter{sujetamsudrefRSHA}
\newcounter{sujetamsudrefHQPdP}
\newcounter{sujetamsudrefHQRdM}
\newcounter{sujetamsudrefHQDFI}
\newcounter{sujetamsudrefHQDMRC}
\newcounter{sujetamsudrefHQHA}
%%%Somme des compteurs%%%
\newcounter{sujetamsudallrefETA}
\newcounter{sujetamsudallrefEXS}
\newcounter{sujetamsudallrefMM}
\newcounter{sujetamsudallrefRS}
\newcounter{sujetamsudallrefCC}
\newcounter{sujetamsudallrefHV}
\newcounter{sujetamsudallrefHL}
\newcounter{sujetamsudallrefHQ}

%%%%%%%%Réflexion philosophique%%%%%%%%%
\newcounter{nosujetamsudrefphilosophique}
\newcounter{sujetamsudHQrefphilosophique}
\newcounter{sujetamsudRSrefphilosophique}
%%%sujetamsuds simples%%%
\newcounter{sujetamsudrefphilosophiqueEXS}
\newcounter{sujetamsudrefphilosophiqueETA}
\newcounter{sujetamsudrefphilosophiqueMM}
\newcounter{sujetamsudrefphilosophiqueHV}
\newcounter{sujetamsudrefphilosophiqueCC}
\newcounter{sujetamsudrefphilosophiqueHL}
%%%sujetamsuds mixtes%%%
\newcounter{sujetamsudrefphilosophiqueEXSTA}
\newcounter{sujetamsudrefphilosophiqueEXSMM}
\newcounter{sujetamsudrefphilosophiqueETAMM}
\newcounter{sujetamsudrefphilosophiqueOdS}
\newcounter{sujetamsudrefphilosophiqueHVL}
\newcounter{sujetamsudrefphilosophiqueHVCC}
\newcounter{sujetamsudrefphilosophiqueHLCC}
\newcounter{sujetamsudrefphilosophiqueOdH}
%%%sujetamsuds mixtes semestres%%%
\newcounter{sujetamsudrefphilosophiqueETACC}
\newcounter{sujetamsudrefphilosophiqueETAHV}
\newcounter{sujetamsudrefphilosophiqueETAHL}
\newcounter{sujetamsudrefphilosophiqueEXSCC}
\newcounter{sujetamsudrefphilosophiqueEXSHV}
\newcounter{sujetamsudrefphilosophiqueEXSHL}
\newcounter{sujetamsudrefphilosophiqueMMCC}
\newcounter{sujetamsudrefphilosophiqueMMHV}
\newcounter{sujetamsudrefphilosophiqueMMHL}
%%%sujetamsuds mixtes années%%
\newcounter{sujetamsudrefphilosophiquepremterm}
\newcounter{sujetamsudrefphilosophiquepremtermRS}
\newcounter{sujetamsudrefphilosophiquepremtermHQ}
\newcounter{sujetamsudrefphilosophiqueETAPdP}
\newcounter{sujetamsudrefphilosophiqueETADFI}
\newcounter{sujetamsudrefphilosophiqueETADMRC}
\newcounter{sujetamsudrefphilosophiqueETAHA}\newcounter{sujetamsudrefphilosophiquepremtermETA}
\newcounter{sujetamsudrefphilosophiqueEXSPdP}
\newcounter{sujetamsudrefphilosophiqueEXSDFI}
\newcounter{sujetamsudrefphilosophiqueEXSDMRC}
\newcounter{sujetamsudrefphilosophiqueEXSHA}
\newcounter{sujetamsudrefphilosophiquepremtermEXS}
\newcounter{sujetamsudrefphilosophiqueMMPdP}
\newcounter{sujetamsudrefphilosophiqueMMDFI}
\newcounter{sujetamsudrefphilosophiqueMMDMRC}
\newcounter{sujetamsudrefphilosophiqueMMHA}
\newcounter{sujetamsudrefphilosophiquepremtermMM}
\newcounter{sujetamsudrefphilosophiqueCCPdP}
\newcounter{sujetamsudrefphilosophiqueCCDFI}
\newcounter{sujetamsudrefphilosophiqueCCDMRC}
\newcounter{sujetamsudrefphilosophiqueCCHA}
\newcounter{sujetamsudrefphilosophiquepremtermCC}
\newcounter{sujetamsudrefphilosophiqueHVPdP}
\newcounter{sujetamsudrefphilosophiqueHVDFI}
\newcounter{sujetamsudrefphilosophiqueHVDMRC}
\newcounter{sujetamsudrefphilosophiqueHVHA}
\newcounter{sujetamsudrefphilosophiquepremtermHV}
\newcounter{sujetamsudrefphilosophiqueHLPdP}
\newcounter{sujetamsudrefphilosophiqueHLDFI}
\newcounter{sujetamsudrefphilosophiqueHLDMRC}
\newcounter{sujetamsudrefphilosophiqueHLHA}
\newcounter{sujetamsudrefphilosophiquepremtermHL}
\newcounter{sujetamsudrefphilosophiqueRSPdP}
\newcounter{sujetamsudrefphilosophiqueRSRdM}
\newcounter{sujetamsudrefphilosophiqueRSDFI}
\newcounter{sujetamsudrefphilosophiqueRSDMRC}
\newcounter{sujetamsudrefphilosophiqueRSHA}
\newcounter{sujetamsudrefphilosophiqueHQPdP}
\newcounter{sujetamsudrefphilosophiqueHQRdM}
\newcounter{sujetamsudrefphilosophiqueHQDFI}
\newcounter{sujetamsudrefphilosophiqueHQDMRC}
\newcounter{sujetamsudrefphilosophiqueHQHA}
%%%Somme des compteurs%%%
\newcounter{sujetamsudallrefphilosophiqueETA}
\newcounter{sujetamsudallrefphilosophiqueEXS}
\newcounter{sujetamsudallrefphilosophiqueMM}
\newcounter{sujetamsudallrefphilosophiqueRS}
\newcounter{sujetamsudallrefphilosophiqueCC}
\newcounter{sujetamsudallrefphilosophiqueHV}
\newcounter{sujetamsudallrefphilosophiqueHL}
\newcounter{sujetamsudallrefphilosophiqueHQ}

%%%%%%%%%Réflexion littéraire%%%%%%%%%%%
\newcounter{nosujetamsudreflitteraire}
\newcounter{sujetamsudHQreflitteraire}
\newcounter{sujetamsudRSreflitteraire}
%%%sujetamsuds simples%%%
\newcounter{sujetamsudreflitteraireEXS}
\newcounter{sujetamsudreflitteraireETA}
\newcounter{sujetamsudreflitteraireMM}
\newcounter{sujetamsudreflitteraireHV}
\newcounter{sujetamsudreflitteraireCC}
\newcounter{sujetamsudreflitteraireHL}
%%%sujetamsuds mixtes%%%
\newcounter{sujetamsudreflitteraireEXSTA}
\newcounter{sujetamsudreflitteraireEXSMM}
\newcounter{sujetamsudreflitteraireETAMM}
\newcounter{sujetamsudreflitteraireOdS}
\newcounter{sujetamsudreflitteraireHVL}
\newcounter{sujetamsudreflitteraireHVCC}
\newcounter{sujetamsudreflitteraireHLCC}
\newcounter{sujetamsudreflitteraireOdH}
%%%sujetamsuds mixtes semestres%%%
\newcounter{sujetamsudreflitteraireETACC}
\newcounter{sujetamsudreflitteraireETAHV}
\newcounter{sujetamsudreflitteraireETAHL}
\newcounter{sujetamsudreflitteraireEXSCC}
\newcounter{sujetamsudreflitteraireEXSHV}
\newcounter{sujetamsudreflitteraireEXSHL}
\newcounter{sujetamsudreflitteraireMMCC}
\newcounter{sujetamsudreflitteraireMMHV}
\newcounter{sujetamsudreflitteraireMMHL}
%%%sujetamsuds mixtes années%%
\newcounter{sujetamsudreflitterairepremterm}
\newcounter{sujetamsudreflitterairepremtermRS}
\newcounter{sujetamsudreflitterairepremtermHQ}
\newcounter{sujetamsudreflitteraireETAPdP}
\newcounter{sujetamsudreflitteraireETADFI}
\newcounter{sujetamsudreflitteraireETADMRC}
\newcounter{sujetamsudreflitteraireETAHA}
\newcounter{sujetamsudreflitterairepremtermETA}
\newcounter{sujetamsudreflitteraireEXSPdP}
\newcounter{sujetamsudreflitteraireEXSDFI}
\newcounter{sujetamsudreflitteraireEXSDMRC}
\newcounter{sujetamsudreflitteraireEXSHA}
\newcounter{sujetamsudreflitterairepremtermEXS}
\newcounter{sujetamsudreflitteraireMMPdP}
\newcounter{sujetamsudreflitteraireMMDFI}
\newcounter{sujetamsudreflitteraireMMDMRC}
\newcounter{sujetamsudreflitteraireMMHA}
\newcounter{sujetamsudreflitterairepremtermMM}
\newcounter{sujetamsudreflitteraireCCPdP}
\newcounter{sujetamsudreflitteraireCCDFI}
\newcounter{sujetamsudreflitteraireCCDMRC}
\newcounter{sujetamsudreflitteraireCCHA}
\newcounter{sujetamsudreflitterairepremtermCC}
\newcounter{sujetamsudreflitteraireHVPdP}
\newcounter{sujetamsudreflitteraireHVDFI}
\newcounter{sujetamsudreflitteraireHVDMRC}
\newcounter{sujetamsudreflitteraireHVHA}
\newcounter{sujetamsudreflitterairepremtermHV}
\newcounter{sujetamsudreflitteraireHLPdP}
\newcounter{sujetamsudreflitteraireHLDFI}
\newcounter{sujetamsudreflitteraireHLDMRC}
\newcounter{sujetamsudreflitteraireHLHA}
\newcounter{sujetamsudreflitterairepremtermHL}
\newcounter{sujetamsudreflitteraireRSPdP}
\newcounter{sujetamsudreflitteraireRSRdM}
\newcounter{sujetamsudreflitteraireRSDFI}
\newcounter{sujetamsudreflitteraireRSDMRC}
\newcounter{sujetamsudreflitteraireRSHA}
\newcounter{sujetamsudreflitteraireHQPdP}
\newcounter{sujetamsudreflitteraireHQRdM}
\newcounter{sujetamsudreflitteraireHQDFI}
\newcounter{sujetamsudreflitteraireHQDMRC}
\newcounter{sujetamsudreflitteraireHQHA}
%%%Somme des compteurs%%%
\newcounter{sujetamsudallreflitteraireETA}
\newcounter{sujetamsudallreflitteraireEXS}
\newcounter{sujetamsudallreflitteraireMM}
\newcounter{sujetamsudallreflitteraireCC}
\newcounter{sujetamsudallreflitteraireRS}
\newcounter{sujetamsudallreflitteraireHV}
\newcounter{sujetamsudallreflitteraireHL}
\newcounter{sujetamsudallreflitteraireHQ}
%%%Compteurs de statistiques globales sur les genres de question%%%
\newcounter{sujetamsudexclRS}
\newcounter{sujetamsudmixtRS}
\newcounter{sujetamsudexclHQ}
\newcounter{sujetamsudmixtHQ}
\newcounter{sujetamsudmixtHQRS}
%%%Compteurs de statistiques sur les transversaux (divers semestres, première et terminale)
\newcounter{sujettransamsudintlitteraire}
\newcounter{sujettransamsudreflitteraire}
\newcounter{sujettransamsudintphilosophique}
\newcounter{sujettransamsudrefphilosophique}
%%% Compteurs du nombre de sujets par type
\newcounter{sujetamsudintlitteraire}
\newcounter{sujetamsudreflitteraire}
\newcounter{sujetamsudintphilosophique}
\newcounter{sujetamsudrefphilosophique}
%%% Compteurs par discipline
%%Litt
\newcounter{sujetamsudlitteraireETA}
\newcounter{sujetamsudlitteraireEXS}
\newcounter{sujetamsudlitteraireMM}
\newcounter{sujetamsudlitteraireCC}
\newcounter{sujetamsudlitteraireHV}
\newcounter{sujetamsudlitteraireHL}
\newcounter{sujetamsudlitteraireEXSTA}
\newcounter{sujetamsudlitteraireETAMM}
\newcounter{sujetamsudlitteraireEXSMM}
\newcounter{sujetamsudlitteraireHVCC}
\newcounter{sujetamsudlitteraireHVL}
\newcounter{sujetamsudlitteraireaut}
%%Phil
\newcounter{sujetamsudphilosophiqueETA}
\newcounter{sujetamsudphilosophiqueEXS}
\newcounter{sujetamsudphilosophiqueMM}
\newcounter{sujetamsudphilosophiqueCC}
\newcounter{sujetamsudphilosophiqueHV}
\newcounter{sujetamsudphilosophiqueHL}
\newcounter{sujetamsudphilosophiqueEXSTA}
\newcounter{sujetamsudphilosophiqueETAMM}
\newcounter{sujetamsudphilosophiqueEXSMM}
\newcounter{sujetamsudphilosophiqueHVCC}
\newcounter{sujetamsudphilosophiqueHVL}
\newcounter{sujetamsudphilosophiqueaut}
%%Totaux pour les graphiques
\newcounter{sujetamsudlitterairetotal}
\newcounter{sujetamsudphilosophiquetotal}
%%% Autres (Ods et Odh) comptés ensemble
\newcounter{sujetamsudallaut}
%%%%%%%%%%%%%%%%%%%%%%%%%%%%%%%%%%%%%%%%%%%%%%%%%%%%%%%%%%%%%%%%%%%%%%%
%%%%%%%%%%%%%%%%%%%%%%%%%%%%%%Compteurs%%%%%%%%%%%%%%%%%%%%%%%%%%%%%%%%
%%%%%%%%%%%%%%%%%%%%%%%%%%%%%%%%%%%%%%%%%%%%%%%%%%%%%%%%%%%%%%%%%%%%%%%
%%%Compter les sujets par thème%%%
%\newcounter{nosujetang}
\newcounter{sujetangall}%Total de questions pour le centre
\newcounter{sujetangtotal}%Total des questions en comptant en double les questions pour les graphiques
\newcounter{sujetangHQ}
\newcounter{sujetangRS}
%%%sujetangs simples%%%
\newcounter{sujetangEXS}
\newcounter{sujetangETA}
\newcounter{sujetangMM}
\newcounter{sujetangHV}
\newcounter{sujetangCC}
\newcounter{sujetangHL}
%%%sujetangs mixtes%%%
\newcounter{sujetangEXSTA}
\newcounter{sujetangEXSMM}
\newcounter{sujetangETAMM}
\newcounter{sujetangOdS}
\newcounter{sujetangHVL}
\newcounter{sujetangHVCC}
\newcounter{sujetangHLCC}
\newcounter{sujetangOdH}
%%%sujetangs mixtes semestres%%%
\newcounter{sujetangETACC}
\newcounter{sujetangETAHV}
\newcounter{sujetangETAHL}
\newcounter{sujetangEXSCC}
\newcounter{sujetangEXSHV}
\newcounter{sujetangEXSHL}
\newcounter{sujetangMMCC}
\newcounter{sujetangMMHV}
\newcounter{sujetangMMHL}

%%%sujetangs mixtes années%%
\newcounter{sujetangpremterm}%%%Compteur général de sujetang
\newcounter{sujetangETAPdP}
\newcounter{sujetangETADFI}
\newcounter{sujetangETADMRC}
\newcounter{sujetangETAHA}
\newcounter{sujetangpremtermETA}
\newcounter{sujetangEXSPdP}
\newcounter{sujetangEXSDFI}
\newcounter{sujetangEXSDMRC}
\newcounter{sujetangEXSHA}
\newcounter{sujetangpremtermEXS}
\newcounter{sujetangMMPdP}
\newcounter{sujetangMMDFI}
\newcounter{sujetangMMDMRC}
\newcounter{sujetangMMHA}
\newcounter{sujetangpremtermMM}
\newcounter{sujetangCCPdP}
\newcounter{sujetangCCDFI}
\newcounter{sujetangCCDMRC}
\newcounter{sujetangCCHA}
\newcounter{sujetangpremtermCC}
\newcounter{sujetangHVPdP}
\newcounter{sujetangHVDFI}
\newcounter{sujetangHVDMRC}
\newcounter{sujetangHVHA}
\newcounter{sujetangpremtermHV}
\newcounter{sujetangHLPdP}
\newcounter{sujetangHLDFI}
\newcounter{sujetangHLDMRC}
\newcounter{sujetangHLHA}
\newcounter{sujetangpremtermHL}
\newcounter{sujetangRSPdP}
\newcounter{sujetangRSRdM}
\newcounter{sujetangRSDFI}
\newcounter{sujetangRSDMRC}
\newcounter{sujetangRSHA}
\newcounter{sujetangHQPdP}
\newcounter{sujetangHQRdM}
\newcounter{sujetangHQDFI}
\newcounter{sujetangHQDMRC}
\newcounter{sujetangHQHA}

%%%Somme des compteurs%%%
\newcounter{sujetangallETA}
\newcounter{sujetangallEXS}
\newcounter{sujetangallMM}
\newcounter{sujetangallRS}
\newcounter{sujetangallCC}
\newcounter{sujetangallHV}
\newcounter{sujetangallHL}
\newcounter{sujetangallHQ}

%%%Compter les sujetangs par thème et types de questions%%%
%%%%%%%%%%%%%%%%%%%%%%%%%%%%%%%%%%%%%%%%%%%%%
%%%%%%%%%Question d'interprétation%%%%%%%%%%%
%%%%%%%%%%%%%%%%%%%%%%%%%%%%%%%%%%%%%%%%%%%%%

%%%%%%%%%%%%%%%Général%%%%%%%%%%%%%%%%%%%%%%%
\newcounter{nosujetangint}
\newcounter{sujetangHQint}
\newcounter{sujetangRSint}
%%%sujetangs simples%%%
\newcounter{sujetangintEXS}
\newcounter{sujetangintETA}
\newcounter{sujetangintMM}
\newcounter{sujetangintHV}
\newcounter{sujetangintCC}
\newcounter{sujetangintHL}
%%%sujetangs mixtes%%%
\newcounter{sujetangintEXSTA}
\newcounter{sujetangintEXSMM}
\newcounter{sujetangintETAMM}
\newcounter{sujetangintOdS}
\newcounter{sujetangintHVL}
\newcounter{sujetangintHVCC}
\newcounter{sujetangintHLCC}
\newcounter{sujetangintOdH}
%%%sujetangs mixtes semestres%%%
\newcounter{sujetangintETACC}
\newcounter{sujetangintETAHV}
\newcounter{sujetangintETAHL}
\newcounter{sujetangintEXSCC}
\newcounter{sujetangintEXSHV}
\newcounter{sujetangintEXSHL}
\newcounter{sujetangintMMCC}
\newcounter{sujetangintMMHV}
\newcounter{sujetangintMMHL}
%%%sujetangs mixtes années%%
\newcounter{sujetangintpremterm}
\newcounter{sujetangintpremtermRS}
\newcounter{sujetangintpremtermHQ}
\newcounter{sujetangintETAPdP}
\newcounter{sujetangintETADFI}
\newcounter{sujetangintETADMRC}
\newcounter{sujetangintETAHA}
\newcounter{sujetangintpremtermETA}
\newcounter{sujetangintEXSPdP}
\newcounter{sujetangintEXSDFI}
\newcounter{sujetangintEXSDMRC}
\newcounter{sujetangintEXSHA}
\newcounter{sujetangintpremtermEXS}
\newcounter{sujetangintMMPdP}
\newcounter{sujetangintMMDFI}
\newcounter{sujetangintMMDMRC}
\newcounter{sujetangintMMHA}
\newcounter{sujetangintpremtermMM}
\newcounter{sujetangintCCPdP}
\newcounter{sujetangintCCDFI}
\newcounter{sujetangintCCDMRC}
\newcounter{sujetangintCCHA}
\newcounter{sujetangintpremtermCC}
\newcounter{sujetangintHVPdP}
\newcounter{sujetangintHVDFI}
\newcounter{sujetangintHVDMRC}
\newcounter{sujetangintHVHA}
\newcounter{sujetangintpremtermHV}
\newcounter{sujetangintHLPdP}
\newcounter{sujetangintHLDFI}
\newcounter{sujetangintHLDMRC}
\newcounter{sujetangintHLHA}
\newcounter{sujetangintpremtermHL}
\newcounter{sujetangintRSPdP}
\newcounter{sujetangintRSRdM}
\newcounter{sujetangintRSDFI}
\newcounter{sujetangintRSDMRC}
\newcounter{sujetangintRSHA}
\newcounter{sujetangintHQPdP}
\newcounter{sujetangintHQRdM}
\newcounter{sujetangintHQDFI}
\newcounter{sujetangintHQDMRC}
\newcounter{sujetangintHQHA}
%%%Somme des compteurs%%%
\newcounter{sujetangallintETA}
\newcounter{sujetangallintEXS}
\newcounter{sujetangallintMM}
\newcounter{sujetangallintRS}
\newcounter{sujetangallintCC}
\newcounter{sujetangallintHV}
\newcounter{sujetangallintHL}
\newcounter{sujetangallintHQ}

%%%%%%%%Interprétation philosophique%%%%%%%%%
\newcounter{nosujetangintphilosophique}
\newcounter{sujetangHQintphilosophique}
\newcounter{sujetangRSintphilosophique}
%%%sujetangs simples%%%
\newcounter{sujetangintphilosophiqueEXS}
\newcounter{sujetangintphilosophiqueETA}
\newcounter{sujetangintphilosophiqueMM}
\newcounter{sujetangintphilosophiqueHV}
\newcounter{sujetangintphilosophiqueCC}
\newcounter{sujetangintphilosophiqueHL}
%%%sujetangs mixtes%%%
\newcounter{sujetangintphilosophiqueEXSTA}
\newcounter{sujetangintphilosophiqueEXSMM}
\newcounter{sujetangintphilosophiqueETAMM}
\newcounter{sujetangintphilosophiqueOdS}
\newcounter{sujetangintphilosophiqueHVL}
\newcounter{sujetangintphilosophiqueHVCC}
\newcounter{sujetangintphilosophiqueHLCC}
\newcounter{sujetangintphilosophiqueOdH}
%%%sujetangs mixtes semestres%%%
\newcounter{sujetangintphilosophiqueETACC}
\newcounter{sujetangintphilosophiqueETAHV}
\newcounter{sujetangintphilosophiqueETAHL}
\newcounter{sujetangintphilosophiqueEXSCC}
\newcounter{sujetangintphilosophiqueEXSHV}
\newcounter{sujetangintphilosophiqueEXSHL}
\newcounter{sujetangintphilosophiqueMMCC}
\newcounter{sujetangintphilosophiqueMMHV}
\newcounter{sujetangintphilosophiqueMMHL}
%%%sujetangs mixtes années%%
\newcounter{sujetangintphilosophiquepremterm}
\newcounter{sujetangintphilosophiquepremtermRS}
\newcounter{sujetangintphilosophiquepremtermHQ}
\newcounter{sujetangintphilosophiqueETAPdP}
\newcounter{sujetangintphilosophiqueETADFI}
\newcounter{sujetangintphilosophiqueETADMRC}
\newcounter{sujetangintphilosophiqueETAHA}
\newcounter{sujetangintphilosophiquepremtermETA}
\newcounter{sujetangintphilosophiqueEXSPdP}
\newcounter{sujetangintphilosophiqueEXSDFI}
\newcounter{sujetangintphilosophiqueEXSDMRC}
\newcounter{sujetangintphilosophiqueEXSHA}
\newcounter{sujetangintphilosophiquepremtermEXS}
\newcounter{sujetangintphilosophiqueMMPdP}
\newcounter{sujetangintphilosophiqueMMDFI}
\newcounter{sujetangintphilosophiqueMMDMRC}
\newcounter{sujetangintphilosophiqueMMHA}
\newcounter{sujetangintphilosophiquepremtermMM}
\newcounter{sujetangintphilosophiqueCCPdP}
\newcounter{sujetangintphilosophiqueCCDFI}
\newcounter{sujetangintphilosophiqueCCDMRC}
\newcounter{sujetangintphilosophiqueCCHA}
\newcounter{sujetangintphilosophiquepremtermCC}
\newcounter{sujetangintphilosophiqueHVPdP}
\newcounter{sujetangintphilosophiqueHVDFI}
\newcounter{sujetangintphilosophiqueHVDMRC}
\newcounter{sujetangintphilosophiqueHVHA}
\newcounter{sujetangintphilosophiquepremtermHV}
\newcounter{sujetangintphilosophiqueHLPdP}
\newcounter{sujetangintphilosophiqueHLDFI}
\newcounter{sujetangintphilosophiqueHLDMRC}
\newcounter{sujetangintphilosophiqueHLHA}
\newcounter{sujetangintphilosophiquepremtermHL}
\newcounter{sujetangintphilosophiqueRSPdP}
\newcounter{sujetangintphilosophiqueRSRdM}
\newcounter{sujetangintphilosophiqueRSDFI}
\newcounter{sujetangintphilosophiqueRSDMRC}
\newcounter{sujetangintphilosophiqueRSHA}
\newcounter{sujetangintphilosophiqueHQPdP}
\newcounter{sujetangintphilosophiqueHQRdM}
\newcounter{sujetangintphilosophiqueHQDFI}
\newcounter{sujetangintphilosophiqueHQDMRC}
\newcounter{sujetangintphilosophiqueHQHA}
%%%Somme des compteurs%%%
\newcounter{sujetangallintphilosophiqueETA}
\newcounter{sujetangallintphilosophiqueEXS}
\newcounter{sujetangallintphilosophiqueMM}
\newcounter{sujetangallintphilosophiqueRS}
\newcounter{sujetangallintphilosophiqueCC}
\newcounter{sujetangallintphilosophiqueHV}
\newcounter{sujetangallintphilosophiqueHL}
\newcounter{sujetangallintphilosophiqueHQ}

%%%%%%%%%Interprétation littéraire%%%%%%%%%%%
\newcounter{nosujetangintlitteraire}
\newcounter{sujetangHQintlitteraire}
\newcounter{sujetangRSintlitteraire}
%%%sujetangs simples%%%
\newcounter{sujetangintlitteraireEXS}
\newcounter{sujetangintlitteraireETA}
\newcounter{sujetangintlitteraireMM}
\newcounter{sujetangintlitteraireHV}
\newcounter{sujetangintlitteraireCC}
\newcounter{sujetangintlitteraireHL}
%%%sujetangs mixtes%%%
\newcounter{sujetangintlitteraireEXSTA}
\newcounter{sujetangintlitteraireEXSMM}
\newcounter{sujetangintlitteraireETAMM}
\newcounter{sujetangintlitteraireOdS}
\newcounter{sujetangintlitteraireHVL}
\newcounter{sujetangintlitteraireHVCC}
\newcounter{sujetangintlitteraireHLCC}
\newcounter{sujetangintlitteraireOdH}
%%%sujetangs mixtes semestres%%%
\newcounter{sujetangintlitteraireETACC}
\newcounter{sujetangintlitteraireETAHV}
\newcounter{sujetangintlitteraireETAHL}
\newcounter{sujetangintlitteraireEXSCC}
\newcounter{sujetangintlitteraireEXSHV}
\newcounter{sujetangintlitteraireEXSHL}
\newcounter{sujetangintlitteraireMMCC}
\newcounter{sujetangintlitteraireMMHV}
\newcounter{sujetangintlitteraireMMHL}
%%%sujetangs mixtes années%%
\newcounter{sujetangintlitterairepremterm}
\newcounter{sujetangintlitterairepremtermRS}
\newcounter{sujetangintlitterairepremtermHQ}
\newcounter{sujetangintlitteraireETAPdP}
\newcounter{sujetangintlitteraireETADFI}
\newcounter{sujetangintlitteraireETADMRC}
\newcounter{sujetangintlitteraireETAHA}
\newcounter{sujetangintlitterairepremtermETA}
\newcounter{sujetangintlitteraireEXSPdP}
\newcounter{sujetangintlitteraireEXSDFI}
\newcounter{sujetangintlitteraireEXSDMRC}
\newcounter{sujetangintlitteraireEXSHA}
\newcounter{sujetangintlitterairepremtermEXS}
\newcounter{sujetangintlitteraireMMPdP}
\newcounter{sujetangintlitteraireMMDFI}
\newcounter{sujetangintlitteraireMMDMRC}
\newcounter{sujetangintlitteraireMMHA}
\newcounter{sujetangintlitterairepremtermMM}
\newcounter{sujetangintlitteraireCCPdP}
\newcounter{sujetangintlitteraireCCDFI}
\newcounter{sujetangintlitteraireCCDMRC}
\newcounter{sujetangintlitteraireCCHA}
\newcounter{sujetangintlitterairepremtermCC}
\newcounter{sujetangintlitteraireHVPdP}
\newcounter{sujetangintlitteraireHVDFI}
\newcounter{sujetangintlitteraireHVDMRC}
\newcounter{sujetangintlitteraireHVHA}
\newcounter{sujetangintlitterairepremtermHV}
\newcounter{sujetangintlitteraireHLPdP}
\newcounter{sujetangintlitteraireHLDFI}
\newcounter{sujetangintlitteraireHLDMRC}
\newcounter{sujetangintlitteraireHLHA}
\newcounter{sujetangintlitterairepremtermHL}
\newcounter{sujetangintlitteraireRSPdP}
\newcounter{sujetangintlitteraireRSRdM}
\newcounter{sujetangintlitteraireRSDFI}
\newcounter{sujetangintlitteraireRSDMRC}
\newcounter{sujetangintlitteraireRSHA}
\newcounter{sujetangintlitteraireHQPdP}
\newcounter{sujetangintlitteraireHQRdM}
\newcounter{sujetangintlitteraireHQDFI}
\newcounter{sujetangintlitteraireHQDMRC}
\newcounter{sujetangintlitteraireHQHA}
%%%Somme des compteurs%%%
\newcounter{sujetangallintlitteraireETA}
\newcounter{sujetangallintlitteraireEXS}
\newcounter{sujetangallintlitteraireMM}
\newcounter{sujetangallintlitteraireRS}
\newcounter{sujetangallintlitteraireCC}
\newcounter{sujetangallintlitteraireHV}
\newcounter{sujetangallintlitteraireHL}
\newcounter{sujetangallintlitteraireHQ}

%%%%%%%%%%%%%%%%%%%%%%%%%%%%%%%%%%%%%%%%%%%%%
%%%%%%%%%%%%Question de réflexion%%%%%%%%%%%%
%%%%%%%%%%%%%%%%%%%%%%%%%%%%%%%%%%%%%%%%%%%%%
\newcounter{nosujetangref}
%%%sujetangs simples%%%
\newcounter{sujetangrefEXS}
\newcounter{sujetangrefETA}
\newcounter{sujetangrefMM}
\newcounter{sujetangrefHV}
\newcounter{sujetangrefCC}
\newcounter{sujetangrefHL}
%%%sujetangs mixtes%%%
\newcounter{sujetangrefEXSTA}
\newcounter{sujetangrefEXSMM}
\newcounter{sujetangrefETAMM}
\newcounter{sujetangrefOdS}
\newcounter{sujetangrefHVL}
\newcounter{sujetangrefHVCC}
\newcounter{sujetangrefHLCC}
\newcounter{sujetangrefOdH}
%%%sujetangs mixtes semestres%%%
\newcounter{sujetangrefETACC}
\newcounter{sujetangrefETAHV}
\newcounter{sujetangrefETAHL}
\newcounter{sujetangrefEXSCC}
\newcounter{sujetangrefEXSHV}
\newcounter{sujetangrefEXSHL}
\newcounter{sujetangrefMMCC}
\newcounter{sujetangrefMMHV}
\newcounter{sujetangrefMMHL}
%%%sujetangs mixtes années%%
\newcounter{sujetangrefpremterm}%%%Compteur général
\newcounter{sujetangrefETAPdP}
\newcounter{sujetangrefETADFI}
\newcounter{sujetangrefETADMRC}
\newcounter{sujetangrefETAHA}
\newcounter{sujetangrefpremtermETA}
\newcounter{sujetangrefEXSPdP}
\newcounter{sujetangrefEXSDFI}
\newcounter{sujetangrefEXSDMRC}
\newcounter{sujetangrefEXSHA}
\newcounter{sujetangrefpremtermEXS}
\newcounter{sujetangrefMMPdP}
\newcounter{sujetangrefMMDFI}
\newcounter{sujetangrefMMDMRC}
\newcounter{sujetangrefMMHA}
\newcounter{sujetangrefpremtermMM}
\newcounter{sujetangrefCCPdP}
\newcounter{sujetangrefCCDFI}
\newcounter{sujetangrefCCDMRC}
\newcounter{sujetangrefCCHA}
\newcounter{sujetangrefpremtermCC}
\newcounter{sujetangrefHVPdP}
\newcounter{sujetangrefHVDFI}
\newcounter{sujetangrefHVDMRC}
\newcounter{sujetangrefHVHA}
\newcounter{sujetangrefpremtermHV}
\newcounter{sujetangrefHLPdP}
\newcounter{sujetangrefHLDFI}
\newcounter{sujetangrefHLDMRC}
\newcounter{sujetangrefHLHA}
\newcounter{sujetangrefpremtermHL}
\newcounter{sujetangrefRSPdP}
\newcounter{sujetangrefRSRdM}
\newcounter{sujetangrefRSDFI}
\newcounter{sujetangrefRSDMRC}
\newcounter{sujetangrefRSHA}
\newcounter{sujetangrefHQPdP}
\newcounter{sujetangrefHQRdM}
\newcounter{sujetangrefHQDFI}
\newcounter{sujetangrefHQDMRC}
\newcounter{sujetangrefHQHA}
%%%Somme des compteurs%%%
\newcounter{sujetangallrefETA}
\newcounter{sujetangallrefEXS}
\newcounter{sujetangallrefMM}
\newcounter{sujetangallrefRS}
\newcounter{sujetangallrefCC}
\newcounter{sujetangallrefHV}
\newcounter{sujetangallrefHL}
\newcounter{sujetangallrefHQ}

%%%%%%%%Réflexion philosophique%%%%%%%%%
\newcounter{nosujetangrefphilosophique}
\newcounter{sujetangHQrefphilosophique}
\newcounter{sujetangRSrefphilosophique}
%%%sujetangs simples%%%
\newcounter{sujetangrefphilosophiqueEXS}
\newcounter{sujetangrefphilosophiqueETA}
\newcounter{sujetangrefphilosophiqueMM}
\newcounter{sujetangrefphilosophiqueHV}
\newcounter{sujetangrefphilosophiqueCC}
\newcounter{sujetangrefphilosophiqueHL}
%%%sujetangs mixtes%%%
\newcounter{sujetangrefphilosophiqueEXSTA}
\newcounter{sujetangrefphilosophiqueEXSMM}
\newcounter{sujetangrefphilosophiqueETAMM}
\newcounter{sujetangrefphilosophiqueOdS}
\newcounter{sujetangrefphilosophiqueHVL}
\newcounter{sujetangrefphilosophiqueHVCC}
\newcounter{sujetangrefphilosophiqueHLCC}
\newcounter{sujetangrefphilosophiqueOdH}
%%%sujetangs mixtes semestres%%%
\newcounter{sujetangrefphilosophiqueETACC}
\newcounter{sujetangrefphilosophiqueETAHV}
\newcounter{sujetangrefphilosophiqueETAHL}
\newcounter{sujetangrefphilosophiqueEXSCC}
\newcounter{sujetangrefphilosophiqueEXSHV}
\newcounter{sujetangrefphilosophiqueEXSHL}
\newcounter{sujetangrefphilosophiqueMMCC}
\newcounter{sujetangrefphilosophiqueMMHV}
\newcounter{sujetangrefphilosophiqueMMHL}
%%%sujetangs mixtes années%%
\newcounter{sujetangrefphilosophiquepremterm}
\newcounter{sujetangrefphilosophiquepremtermRS}
\newcounter{sujetangrefphilosophiquepremtermHQ}
\newcounter{sujetangrefphilosophiqueETAPdP}
\newcounter{sujetangrefphilosophiqueETADFI}
\newcounter{sujetangrefphilosophiqueETADMRC}
\newcounter{sujetangrefphilosophiqueETAHA}\newcounter{sujetangrefphilosophiquepremtermETA}
\newcounter{sujetangrefphilosophiqueEXSPdP}
\newcounter{sujetangrefphilosophiqueEXSDFI}
\newcounter{sujetangrefphilosophiqueEXSDMRC}
\newcounter{sujetangrefphilosophiqueEXSHA}
\newcounter{sujetangrefphilosophiquepremtermEXS}
\newcounter{sujetangrefphilosophiqueMMPdP}
\newcounter{sujetangrefphilosophiqueMMDFI}
\newcounter{sujetangrefphilosophiqueMMDMRC}
\newcounter{sujetangrefphilosophiqueMMHA}
\newcounter{sujetangrefphilosophiquepremtermMM}
\newcounter{sujetangrefphilosophiqueCCPdP}
\newcounter{sujetangrefphilosophiqueCCDFI}
\newcounter{sujetangrefphilosophiqueCCDMRC}
\newcounter{sujetangrefphilosophiqueCCHA}
\newcounter{sujetangrefphilosophiquepremtermCC}
\newcounter{sujetangrefphilosophiqueHVPdP}
\newcounter{sujetangrefphilosophiqueHVDFI}
\newcounter{sujetangrefphilosophiqueHVDMRC}
\newcounter{sujetangrefphilosophiqueHVHA}
\newcounter{sujetangrefphilosophiquepremtermHV}
\newcounter{sujetangrefphilosophiqueHLPdP}
\newcounter{sujetangrefphilosophiqueHLDFI}
\newcounter{sujetangrefphilosophiqueHLDMRC}
\newcounter{sujetangrefphilosophiqueHLHA}
\newcounter{sujetangrefphilosophiquepremtermHL}
\newcounter{sujetangrefphilosophiqueRSPdP}
\newcounter{sujetangrefphilosophiqueRSRdM}
\newcounter{sujetangrefphilosophiqueRSDFI}
\newcounter{sujetangrefphilosophiqueRSDMRC}
\newcounter{sujetangrefphilosophiqueRSHA}
\newcounter{sujetangrefphilosophiqueHQPdP}
\newcounter{sujetangrefphilosophiqueHQRdM}
\newcounter{sujetangrefphilosophiqueHQDFI}
\newcounter{sujetangrefphilosophiqueHQDMRC}
\newcounter{sujetangrefphilosophiqueHQHA}
%%%Somme des compteurs%%%
\newcounter{sujetangallrefphilosophiqueETA}
\newcounter{sujetangallrefphilosophiqueEXS}
\newcounter{sujetangallrefphilosophiqueMM}
\newcounter{sujetangallrefphilosophiqueRS}
\newcounter{sujetangallrefphilosophiqueCC}
\newcounter{sujetangallrefphilosophiqueHV}
\newcounter{sujetangallrefphilosophiqueHL}
\newcounter{sujetangallrefphilosophiqueHQ}

%%%%%%%%%Réflexion littéraire%%%%%%%%%%%
\newcounter{nosujetangreflitteraire}
\newcounter{sujetangHQreflitteraire}
\newcounter{sujetangRSreflitteraire}
%%%sujetangs simples%%%
\newcounter{sujetangreflitteraireEXS}
\newcounter{sujetangreflitteraireETA}
\newcounter{sujetangreflitteraireMM}
\newcounter{sujetangreflitteraireHV}
\newcounter{sujetangreflitteraireCC}
\newcounter{sujetangreflitteraireHL}
%%%sujetangs mixtes%%%
\newcounter{sujetangreflitteraireEXSTA}
\newcounter{sujetangreflitteraireEXSMM}
\newcounter{sujetangreflitteraireETAMM}
\newcounter{sujetangreflitteraireOdS}
\newcounter{sujetangreflitteraireHVL}
\newcounter{sujetangreflitteraireHVCC}
\newcounter{sujetangreflitteraireHLCC}
\newcounter{sujetangreflitteraireOdH}
%%%sujetangs mixtes semestres%%%
\newcounter{sujetangreflitteraireETACC}
\newcounter{sujetangreflitteraireETAHV}
\newcounter{sujetangreflitteraireETAHL}
\newcounter{sujetangreflitteraireEXSCC}
\newcounter{sujetangreflitteraireEXSHV}
\newcounter{sujetangreflitteraireEXSHL}
\newcounter{sujetangreflitteraireMMCC}
\newcounter{sujetangreflitteraireMMHV}
\newcounter{sujetangreflitteraireMMHL}
%%%sujetangs mixtes années%%
\newcounter{sujetangreflitterairepremterm}
\newcounter{sujetangreflitterairepremtermRS}
\newcounter{sujetangreflitterairepremtermHQ}
\newcounter{sujetangreflitteraireETAPdP}
\newcounter{sujetangreflitteraireETADFI}
\newcounter{sujetangreflitteraireETADMRC}
\newcounter{sujetangreflitteraireETAHA}
\newcounter{sujetangreflitterairepremtermETA}
\newcounter{sujetangreflitteraireEXSPdP}
\newcounter{sujetangreflitteraireEXSDFI}
\newcounter{sujetangreflitteraireEXSDMRC}
\newcounter{sujetangreflitteraireEXSHA}
\newcounter{sujetangreflitterairepremtermEXS}
\newcounter{sujetangreflitteraireMMPdP}
\newcounter{sujetangreflitteraireMMDFI}
\newcounter{sujetangreflitteraireMMDMRC}
\newcounter{sujetangreflitteraireMMHA}
\newcounter{sujetangreflitterairepremtermMM}
\newcounter{sujetangreflitteraireCCPdP}
\newcounter{sujetangreflitteraireCCDFI}
\newcounter{sujetangreflitteraireCCDMRC}
\newcounter{sujetangreflitteraireCCHA}
\newcounter{sujetangreflitterairepremtermCC}
\newcounter{sujetangreflitteraireHVPdP}
\newcounter{sujetangreflitteraireHVDFI}
\newcounter{sujetangreflitteraireHVDMRC}
\newcounter{sujetangreflitteraireHVHA}
\newcounter{sujetangreflitterairepremtermHV}
\newcounter{sujetangreflitteraireHLPdP}
\newcounter{sujetangreflitteraireHLDFI}
\newcounter{sujetangreflitteraireHLDMRC}
\newcounter{sujetangreflitteraireHLHA}
\newcounter{sujetangreflitterairepremtermHL}
\newcounter{sujetangreflitteraireRSPdP}
\newcounter{sujetangreflitteraireRSRdM}
\newcounter{sujetangreflitteraireRSDFI}
\newcounter{sujetangreflitteraireRSDMRC}
\newcounter{sujetangreflitteraireRSHA}
\newcounter{sujetangreflitteraireHQPdP}
\newcounter{sujetangreflitteraireHQRdM}
\newcounter{sujetangreflitteraireHQDFI}
\newcounter{sujetangreflitteraireHQDMRC}
\newcounter{sujetangreflitteraireHQHA}
%%%Somme des compteurs%%%
\newcounter{sujetangallreflitteraireETA}
\newcounter{sujetangallreflitteraireEXS}
\newcounter{sujetangallreflitteraireMM}
\newcounter{sujetangallreflitteraireCC}
\newcounter{sujetangallreflitteraireRS}
\newcounter{sujetangallreflitteraireHV}
\newcounter{sujetangallreflitteraireHL}
\newcounter{sujetangallreflitteraireHQ}
%%%Compteurs de statistiques globales sur les genres de question%%%
\newcounter{sujetangexclRS}
\newcounter{sujetangmixtRS}
\newcounter{sujetangexclHQ}
\newcounter{sujetangmixtHQ}
\newcounter{sujetangmixtHQRS}
%%%Compteurs de statistiques sur les transversaux (divers semestres, première et terminale)
\newcounter{sujettransangintlitteraire}
\newcounter{sujettransangreflitteraire}
\newcounter{sujettransangintphilosophique}
\newcounter{sujettransangrefphilosophique}
%%% Compteurs du nombre de sujets par type
\newcounter{sujetangintlitteraire}
\newcounter{sujetangreflitteraire}
\newcounter{sujetangintphilosophique}
\newcounter{sujetangrefphilosophique}
%%% Compteurs par discipline
%%Litt
\newcounter{sujetanglitteraireETA}
\newcounter{sujetanglitteraireEXS}
\newcounter{sujetanglitteraireMM}
\newcounter{sujetanglitteraireCC}
\newcounter{sujetanglitteraireHV}
\newcounter{sujetanglitteraireHL}
\newcounter{sujetanglitteraireEXSTA}
\newcounter{sujetanglitteraireETAMM}
\newcounter{sujetanglitteraireEXSMM}
\newcounter{sujetanglitteraireHVCC}
\newcounter{sujetanglitteraireHVL}
\newcounter{sujetanglitteraireaut}
%%Phil
\newcounter{sujetangphilosophiqueETA}
\newcounter{sujetangphilosophiqueEXS}
\newcounter{sujetangphilosophiqueMM}
\newcounter{sujetangphilosophiqueCC}
\newcounter{sujetangphilosophiqueHV}
\newcounter{sujetangphilosophiqueHL}
\newcounter{sujetangphilosophiqueEXSTA}
\newcounter{sujetangphilosophiqueETAMM}
\newcounter{sujetangphilosophiqueEXSMM}
\newcounter{sujetangphilosophiqueHVCC}
\newcounter{sujetangphilosophiqueHVL}
\newcounter{sujetangphilosophiqueaut}
%%Totaux pour les graphiques
\newcounter{sujetanglitterairetotal}
\newcounter{sujetangphilosophiquetotal}
%%% Autres (Ods et Odh) comptés ensemble
\newcounter{sujetangallaut}
%%%%%%%%%%%%%%%%%%%%%%%%%%%%%%%%%%%%%%%%%%%%%%%%%%%%%%%%%%%%%%%%%%%%%%%
%%%%%%%%%%%%%%%%%%%%%%%%%%%%%%Compteurs%%%%%%%%%%%%%%%%%%%%%%%%%%%%%%%%
%%%%%%%%%%%%%%%%%%%%%%%%%%%%%%%%%%%%%%%%%%%%%%%%%%%%%%%%%%%%%%%%%%%%%%%
%%%Compter les sujets par thème%%%
%\newcounter{nosujetasi}
\newcounter{sujetasiall}%Total de questions pour le centre
\newcounter{sujetasitotal}%Total des questions en comptant en double les questions pour les graphiques
\newcounter{sujetasiHQ}
\newcounter{sujetasiRS}
%%%sujetasis simples%%%
\newcounter{sujetasiEXS}
\newcounter{sujetasiETA}
\newcounter{sujetasiMM}
\newcounter{sujetasiHV}
\newcounter{sujetasiCC}
\newcounter{sujetasiHL}
%%%sujetasis mixtes%%%
\newcounter{sujetasiEXSTA}
\newcounter{sujetasiEXSMM}
\newcounter{sujetasiETAMM}
\newcounter{sujetasiOdS}
\newcounter{sujetasiHVL}
\newcounter{sujetasiHVCC}
\newcounter{sujetasiHLCC}
\newcounter{sujetasiOdH}
%%%sujetasis mixtes semestres%%%
\newcounter{sujetasiETACC}
\newcounter{sujetasiETAHV}
\newcounter{sujetasiETAHL}
\newcounter{sujetasiEXSCC}
\newcounter{sujetasiEXSHV}
\newcounter{sujetasiEXSHL}
\newcounter{sujetasiMMCC}
\newcounter{sujetasiMMHV}
\newcounter{sujetasiMMHL}

%%%sujetasis mixtes années%%
\newcounter{sujetasipremterm}%%%Compteur général de sujetasi
\newcounter{sujetasiETAPdP}
\newcounter{sujetasiETADFI}
\newcounter{sujetasiETADMRC}
\newcounter{sujetasiETAHA}
\newcounter{sujetasipremtermETA}
\newcounter{sujetasiEXSPdP}
\newcounter{sujetasiEXSDFI}
\newcounter{sujetasiEXSDMRC}
\newcounter{sujetasiEXSHA}
\newcounter{sujetasipremtermEXS}
\newcounter{sujetasiMMPdP}
\newcounter{sujetasiMMDFI}
\newcounter{sujetasiMMDMRC}
\newcounter{sujetasiMMHA}
\newcounter{sujetasipremtermMM}
\newcounter{sujetasiCCPdP}
\newcounter{sujetasiCCDFI}
\newcounter{sujetasiCCDMRC}
\newcounter{sujetasiCCHA}
\newcounter{sujetasipremtermCC}
\newcounter{sujetasiHVPdP}
\newcounter{sujetasiHVDFI}
\newcounter{sujetasiHVDMRC}
\newcounter{sujetasiHVHA}
\newcounter{sujetasipremtermHV}
\newcounter{sujetasiHLPdP}
\newcounter{sujetasiHLDFI}
\newcounter{sujetasiHLDMRC}
\newcounter{sujetasiHLHA}
\newcounter{sujetasipremtermHL}
\newcounter{sujetasiRSPdP}
\newcounter{sujetasiRSRdM}
\newcounter{sujetasiRSDFI}
\newcounter{sujetasiRSDMRC}
\newcounter{sujetasiRSHA}
\newcounter{sujetasiHQPdP}
\newcounter{sujetasiHQRdM}
\newcounter{sujetasiHQDFI}
\newcounter{sujetasiHQDMRC}
\newcounter{sujetasiHQHA}

%%%Somme des compteurs%%%
\newcounter{sujetasiallETA}
\newcounter{sujetasiallEXS}
\newcounter{sujetasiallMM}
\newcounter{sujetasiallRS}
\newcounter{sujetasiallCC}
\newcounter{sujetasiallHV}
\newcounter{sujetasiallHL}
\newcounter{sujetasiallHQ}

%%%Compter les sujetasis par thème et types de questions%%%
%%%%%%%%%%%%%%%%%%%%%%%%%%%%%%%%%%%%%%%%%%%%%
%%%%%%%%%Question d'interprétation%%%%%%%%%%%
%%%%%%%%%%%%%%%%%%%%%%%%%%%%%%%%%%%%%%%%%%%%%

%%%%%%%%%%%%%%%Général%%%%%%%%%%%%%%%%%%%%%%%
\newcounter{nosujetasiint}
\newcounter{sujetasiHQint}
\newcounter{sujetasiRSint}
%%%sujetasis simples%%%
\newcounter{sujetasiintEXS}
\newcounter{sujetasiintETA}
\newcounter{sujetasiintMM}
\newcounter{sujetasiintHV}
\newcounter{sujetasiintCC}
\newcounter{sujetasiintHL}
%%%sujetasis mixtes%%%
\newcounter{sujetasiintEXSTA}
\newcounter{sujetasiintEXSMM}
\newcounter{sujetasiintETAMM}
\newcounter{sujetasiintOdS}
\newcounter{sujetasiintHVL}
\newcounter{sujetasiintHVCC}
\newcounter{sujetasiintHLCC}
\newcounter{sujetasiintOdH}
%%%sujetasis mixtes semestres%%%
\newcounter{sujetasiintETACC}
\newcounter{sujetasiintETAHV}
\newcounter{sujetasiintETAHL}
\newcounter{sujetasiintEXSCC}
\newcounter{sujetasiintEXSHV}
\newcounter{sujetasiintEXSHL}
\newcounter{sujetasiintMMCC}
\newcounter{sujetasiintMMHV}
\newcounter{sujetasiintMMHL}
%%%sujetasis mixtes années%%
\newcounter{sujetasiintpremterm}
\newcounter{sujetasiintpremtermRS}
\newcounter{sujetasiintpremtermHQ}
\newcounter{sujetasiintETAPdP}
\newcounter{sujetasiintETADFI}
\newcounter{sujetasiintETADMRC}
\newcounter{sujetasiintETAHA}
\newcounter{sujetasiintpremtermETA}
\newcounter{sujetasiintEXSPdP}
\newcounter{sujetasiintEXSDFI}
\newcounter{sujetasiintEXSDMRC}
\newcounter{sujetasiintEXSHA}
\newcounter{sujetasiintpremtermEXS}
\newcounter{sujetasiintMMPdP}
\newcounter{sujetasiintMMDFI}
\newcounter{sujetasiintMMDMRC}
\newcounter{sujetasiintMMHA}
\newcounter{sujetasiintpremtermMM}
\newcounter{sujetasiintCCPdP}
\newcounter{sujetasiintCCDFI}
\newcounter{sujetasiintCCDMRC}
\newcounter{sujetasiintCCHA}
\newcounter{sujetasiintpremtermCC}
\newcounter{sujetasiintHVPdP}
\newcounter{sujetasiintHVDFI}
\newcounter{sujetasiintHVDMRC}
\newcounter{sujetasiintHVHA}
\newcounter{sujetasiintpremtermHV}
\newcounter{sujetasiintHLPdP}
\newcounter{sujetasiintHLDFI}
\newcounter{sujetasiintHLDMRC}
\newcounter{sujetasiintHLHA}
\newcounter{sujetasiintpremtermHL}
\newcounter{sujetasiintRSPdP}
\newcounter{sujetasiintRSRdM}
\newcounter{sujetasiintRSDFI}
\newcounter{sujetasiintRSDMRC}
\newcounter{sujetasiintRSHA}
\newcounter{sujetasiintHQPdP}
\newcounter{sujetasiintHQRdM}
\newcounter{sujetasiintHQDFI}
\newcounter{sujetasiintHQDMRC}
\newcounter{sujetasiintHQHA}
%%%Somme des compteurs%%%
\newcounter{sujetasiallintETA}
\newcounter{sujetasiallintEXS}
\newcounter{sujetasiallintMM}
\newcounter{sujetasiallintRS}
\newcounter{sujetasiallintCC}
\newcounter{sujetasiallintHV}
\newcounter{sujetasiallintHL}
\newcounter{sujetasiallintHQ}

%%%%%%%%Interprétation philosophique%%%%%%%%%
\newcounter{nosujetasiintphilosophique}
\newcounter{sujetasiHQintphilosophique}
\newcounter{sujetasiRSintphilosophique}
%%%sujetasis simples%%%
\newcounter{sujetasiintphilosophiqueEXS}
\newcounter{sujetasiintphilosophiqueETA}
\newcounter{sujetasiintphilosophiqueMM}
\newcounter{sujetasiintphilosophiqueHV}
\newcounter{sujetasiintphilosophiqueCC}
\newcounter{sujetasiintphilosophiqueHL}
%%%sujetasis mixtes%%%
\newcounter{sujetasiintphilosophiqueEXSTA}
\newcounter{sujetasiintphilosophiqueEXSMM}
\newcounter{sujetasiintphilosophiqueETAMM}
\newcounter{sujetasiintphilosophiqueOdS}
\newcounter{sujetasiintphilosophiqueHVL}
\newcounter{sujetasiintphilosophiqueHVCC}
\newcounter{sujetasiintphilosophiqueHLCC}
\newcounter{sujetasiintphilosophiqueOdH}
%%%sujetasis mixtes semestres%%%
\newcounter{sujetasiintphilosophiqueETACC}
\newcounter{sujetasiintphilosophiqueETAHV}
\newcounter{sujetasiintphilosophiqueETAHL}
\newcounter{sujetasiintphilosophiqueEXSCC}
\newcounter{sujetasiintphilosophiqueEXSHV}
\newcounter{sujetasiintphilosophiqueEXSHL}
\newcounter{sujetasiintphilosophiqueMMCC}
\newcounter{sujetasiintphilosophiqueMMHV}
\newcounter{sujetasiintphilosophiqueMMHL}
%%%sujetasis mixtes années%%
\newcounter{sujetasiintphilosophiquepremterm}
\newcounter{sujetasiintphilosophiquepremtermRS}
\newcounter{sujetasiintphilosophiquepremtermHQ}
\newcounter{sujetasiintphilosophiqueETAPdP}
\newcounter{sujetasiintphilosophiqueETADFI}
\newcounter{sujetasiintphilosophiqueETADMRC}
\newcounter{sujetasiintphilosophiqueETAHA}
\newcounter{sujetasiintphilosophiquepremtermETA}
\newcounter{sujetasiintphilosophiqueEXSPdP}
\newcounter{sujetasiintphilosophiqueEXSDFI}
\newcounter{sujetasiintphilosophiqueEXSDMRC}
\newcounter{sujetasiintphilosophiqueEXSHA}
\newcounter{sujetasiintphilosophiquepremtermEXS}
\newcounter{sujetasiintphilosophiqueMMPdP}
\newcounter{sujetasiintphilosophiqueMMDFI}
\newcounter{sujetasiintphilosophiqueMMDMRC}
\newcounter{sujetasiintphilosophiqueMMHA}
\newcounter{sujetasiintphilosophiquepremtermMM}
\newcounter{sujetasiintphilosophiqueCCPdP}
\newcounter{sujetasiintphilosophiqueCCDFI}
\newcounter{sujetasiintphilosophiqueCCDMRC}
\newcounter{sujetasiintphilosophiqueCCHA}
\newcounter{sujetasiintphilosophiquepremtermCC}
\newcounter{sujetasiintphilosophiqueHVPdP}
\newcounter{sujetasiintphilosophiqueHVDFI}
\newcounter{sujetasiintphilosophiqueHVDMRC}
\newcounter{sujetasiintphilosophiqueHVHA}
\newcounter{sujetasiintphilosophiquepremtermHV}
\newcounter{sujetasiintphilosophiqueHLPdP}
\newcounter{sujetasiintphilosophiqueHLDFI}
\newcounter{sujetasiintphilosophiqueHLDMRC}
\newcounter{sujetasiintphilosophiqueHLHA}
\newcounter{sujetasiintphilosophiquepremtermHL}
\newcounter{sujetasiintphilosophiqueRSPdP}
\newcounter{sujetasiintphilosophiqueRSRdM}
\newcounter{sujetasiintphilosophiqueRSDFI}
\newcounter{sujetasiintphilosophiqueRSDMRC}
\newcounter{sujetasiintphilosophiqueRSHA}
\newcounter{sujetasiintphilosophiqueHQPdP}
\newcounter{sujetasiintphilosophiqueHQRdM}
\newcounter{sujetasiintphilosophiqueHQDFI}
\newcounter{sujetasiintphilosophiqueHQDMRC}
\newcounter{sujetasiintphilosophiqueHQHA}
%%%Somme des compteurs%%%
\newcounter{sujetasiallintphilosophiqueETA}
\newcounter{sujetasiallintphilosophiqueEXS}
\newcounter{sujetasiallintphilosophiqueMM}
\newcounter{sujetasiallintphilosophiqueRS}
\newcounter{sujetasiallintphilosophiqueCC}
\newcounter{sujetasiallintphilosophiqueHV}
\newcounter{sujetasiallintphilosophiqueHL}
\newcounter{sujetasiallintphilosophiqueHQ}

%%%%%%%%%Interprétation littéraire%%%%%%%%%%%
\newcounter{nosujetasiintlitteraire}
\newcounter{sujetasiHQintlitteraire}
\newcounter{sujetasiRSintlitteraire}
%%%sujetasis simples%%%
\newcounter{sujetasiintlitteraireEXS}
\newcounter{sujetasiintlitteraireETA}
\newcounter{sujetasiintlitteraireMM}
\newcounter{sujetasiintlitteraireHV}
\newcounter{sujetasiintlitteraireCC}
\newcounter{sujetasiintlitteraireHL}
%%%sujetasis mixtes%%%
\newcounter{sujetasiintlitteraireEXSTA}
\newcounter{sujetasiintlitteraireEXSMM}
\newcounter{sujetasiintlitteraireETAMM}
\newcounter{sujetasiintlitteraireOdS}
\newcounter{sujetasiintlitteraireHVL}
\newcounter{sujetasiintlitteraireHVCC}
\newcounter{sujetasiintlitteraireHLCC}
\newcounter{sujetasiintlitteraireOdH}
%%%sujetasis mixtes semestres%%%
\newcounter{sujetasiintlitteraireETACC}
\newcounter{sujetasiintlitteraireETAHV}
\newcounter{sujetasiintlitteraireETAHL}
\newcounter{sujetasiintlitteraireEXSCC}
\newcounter{sujetasiintlitteraireEXSHV}
\newcounter{sujetasiintlitteraireEXSHL}
\newcounter{sujetasiintlitteraireMMCC}
\newcounter{sujetasiintlitteraireMMHV}
\newcounter{sujetasiintlitteraireMMHL}
%%%sujetasis mixtes années%%
\newcounter{sujetasiintlitterairepremterm}
\newcounter{sujetasiintlitterairepremtermRS}
\newcounter{sujetasiintlitterairepremtermHQ}
\newcounter{sujetasiintlitteraireETAPdP}
\newcounter{sujetasiintlitteraireETADFI}
\newcounter{sujetasiintlitteraireETADMRC}
\newcounter{sujetasiintlitteraireETAHA}
\newcounter{sujetasiintlitterairepremtermETA}
\newcounter{sujetasiintlitteraireEXSPdP}
\newcounter{sujetasiintlitteraireEXSDFI}
\newcounter{sujetasiintlitteraireEXSDMRC}
\newcounter{sujetasiintlitteraireEXSHA}
\newcounter{sujetasiintlitterairepremtermEXS}
\newcounter{sujetasiintlitteraireMMPdP}
\newcounter{sujetasiintlitteraireMMDFI}
\newcounter{sujetasiintlitteraireMMDMRC}
\newcounter{sujetasiintlitteraireMMHA}
\newcounter{sujetasiintlitterairepremtermMM}
\newcounter{sujetasiintlitteraireCCPdP}
\newcounter{sujetasiintlitteraireCCDFI}
\newcounter{sujetasiintlitteraireCCDMRC}
\newcounter{sujetasiintlitteraireCCHA}
\newcounter{sujetasiintlitterairepremtermCC}
\newcounter{sujetasiintlitteraireHVPdP}
\newcounter{sujetasiintlitteraireHVDFI}
\newcounter{sujetasiintlitteraireHVDMRC}
\newcounter{sujetasiintlitteraireHVHA}
\newcounter{sujetasiintlitterairepremtermHV}
\newcounter{sujetasiintlitteraireHLPdP}
\newcounter{sujetasiintlitteraireHLDFI}
\newcounter{sujetasiintlitteraireHLDMRC}
\newcounter{sujetasiintlitteraireHLHA}
\newcounter{sujetasiintlitterairepremtermHL}
\newcounter{sujetasiintlitteraireRSPdP}
\newcounter{sujetasiintlitteraireRSRdM}
\newcounter{sujetasiintlitteraireRSDFI}
\newcounter{sujetasiintlitteraireRSDMRC}
\newcounter{sujetasiintlitteraireRSHA}
\newcounter{sujetasiintlitteraireHQPdP}
\newcounter{sujetasiintlitteraireHQRdM}
\newcounter{sujetasiintlitteraireHQDFI}
\newcounter{sujetasiintlitteraireHQDMRC}
\newcounter{sujetasiintlitteraireHQHA}
%%%Somme des compteurs%%%
\newcounter{sujetasiallintlitteraireETA}
\newcounter{sujetasiallintlitteraireEXS}
\newcounter{sujetasiallintlitteraireMM}
\newcounter{sujetasiallintlitteraireRS}
\newcounter{sujetasiallintlitteraireCC}
\newcounter{sujetasiallintlitteraireHV}
\newcounter{sujetasiallintlitteraireHL}
\newcounter{sujetasiallintlitteraireHQ}

%%%%%%%%%%%%%%%%%%%%%%%%%%%%%%%%%%%%%%%%%%%%%
%%%%%%%%%%%%Question de réflexion%%%%%%%%%%%%
%%%%%%%%%%%%%%%%%%%%%%%%%%%%%%%%%%%%%%%%%%%%%
\newcounter{nosujetasiref}
%%%sujetasis simples%%%
\newcounter{sujetasirefEXS}
\newcounter{sujetasirefETA}
\newcounter{sujetasirefMM}
\newcounter{sujetasirefHV}
\newcounter{sujetasirefCC}
\newcounter{sujetasirefHL}
%%%sujetasis mixtes%%%
\newcounter{sujetasirefEXSTA}
\newcounter{sujetasirefEXSMM}
\newcounter{sujetasirefETAMM}
\newcounter{sujetasirefOdS}
\newcounter{sujetasirefHVL}
\newcounter{sujetasirefHVCC}
\newcounter{sujetasirefHLCC}
\newcounter{sujetasirefOdH}
%%%sujetasis mixtes semestres%%%
\newcounter{sujetasirefETACC}
\newcounter{sujetasirefETAHV}
\newcounter{sujetasirefETAHL}
\newcounter{sujetasirefEXSCC}
\newcounter{sujetasirefEXSHV}
\newcounter{sujetasirefEXSHL}
\newcounter{sujetasirefMMCC}
\newcounter{sujetasirefMMHV}
\newcounter{sujetasirefMMHL}
%%%sujetasis mixtes années%%
\newcounter{sujetasirefpremterm}%%%Compteur général
\newcounter{sujetasirefETAPdP}
\newcounter{sujetasirefETADFI}
\newcounter{sujetasirefETADMRC}
\newcounter{sujetasirefETAHA}
\newcounter{sujetasirefpremtermETA}
\newcounter{sujetasirefEXSPdP}
\newcounter{sujetasirefEXSDFI}
\newcounter{sujetasirefEXSDMRC}
\newcounter{sujetasirefEXSHA}
\newcounter{sujetasirefpremtermEXS}
\newcounter{sujetasirefMMPdP}
\newcounter{sujetasirefMMDFI}
\newcounter{sujetasirefMMDMRC}
\newcounter{sujetasirefMMHA}
\newcounter{sujetasirefpremtermMM}
\newcounter{sujetasirefCCPdP}
\newcounter{sujetasirefCCDFI}
\newcounter{sujetasirefCCDMRC}
\newcounter{sujetasirefCCHA}
\newcounter{sujetasirefpremtermCC}
\newcounter{sujetasirefHVPdP}
\newcounter{sujetasirefHVDFI}
\newcounter{sujetasirefHVDMRC}
\newcounter{sujetasirefHVHA}
\newcounter{sujetasirefpremtermHV}
\newcounter{sujetasirefHLPdP}
\newcounter{sujetasirefHLDFI}
\newcounter{sujetasirefHLDMRC}
\newcounter{sujetasirefHLHA}
\newcounter{sujetasirefpremtermHL}
\newcounter{sujetasirefRSPdP}
\newcounter{sujetasirefRSRdM}
\newcounter{sujetasirefRSDFI}
\newcounter{sujetasirefRSDMRC}
\newcounter{sujetasirefRSHA}
\newcounter{sujetasirefHQPdP}
\newcounter{sujetasirefHQRdM}
\newcounter{sujetasirefHQDFI}
\newcounter{sujetasirefHQDMRC}
\newcounter{sujetasirefHQHA}
%%%Somme des compteurs%%%
\newcounter{sujetasiallrefETA}
\newcounter{sujetasiallrefEXS}
\newcounter{sujetasiallrefMM}
\newcounter{sujetasiallrefRS}
\newcounter{sujetasiallrefCC}
\newcounter{sujetasiallrefHV}
\newcounter{sujetasiallrefHL}
\newcounter{sujetasiallrefHQ}

%%%%%%%%Réflexion philosophique%%%%%%%%%
\newcounter{nosujetasirefphilosophique}
\newcounter{sujetasiHQrefphilosophique}
\newcounter{sujetasiRSrefphilosophique}
%%%sujetasis simples%%%
\newcounter{sujetasirefphilosophiqueEXS}
\newcounter{sujetasirefphilosophiqueETA}
\newcounter{sujetasirefphilosophiqueMM}
\newcounter{sujetasirefphilosophiqueHV}
\newcounter{sujetasirefphilosophiqueCC}
\newcounter{sujetasirefphilosophiqueHL}
%%%sujetasis mixtes%%%
\newcounter{sujetasirefphilosophiqueEXSTA}
\newcounter{sujetasirefphilosophiqueEXSMM}
\newcounter{sujetasirefphilosophiqueETAMM}
\newcounter{sujetasirefphilosophiqueOdS}
\newcounter{sujetasirefphilosophiqueHVL}
\newcounter{sujetasirefphilosophiqueHVCC}
\newcounter{sujetasirefphilosophiqueHLCC}
\newcounter{sujetasirefphilosophiqueOdH}
%%%sujetasis mixtes semestres%%%
\newcounter{sujetasirefphilosophiqueETACC}
\newcounter{sujetasirefphilosophiqueETAHV}
\newcounter{sujetasirefphilosophiqueETAHL}
\newcounter{sujetasirefphilosophiqueEXSCC}
\newcounter{sujetasirefphilosophiqueEXSHV}
\newcounter{sujetasirefphilosophiqueEXSHL}
\newcounter{sujetasirefphilosophiqueMMCC}
\newcounter{sujetasirefphilosophiqueMMHV}
\newcounter{sujetasirefphilosophiqueMMHL}
%%%sujetasis mixtes années%%
\newcounter{sujetasirefphilosophiquepremterm}
\newcounter{sujetasirefphilosophiquepremtermRS}
\newcounter{sujetasirefphilosophiquepremtermHQ}
\newcounter{sujetasirefphilosophiqueETAPdP}
\newcounter{sujetasirefphilosophiqueETADFI}
\newcounter{sujetasirefphilosophiqueETADMRC}
\newcounter{sujetasirefphilosophiqueETAHA}\newcounter{sujetasirefphilosophiquepremtermETA}
\newcounter{sujetasirefphilosophiqueEXSPdP}
\newcounter{sujetasirefphilosophiqueEXSDFI}
\newcounter{sujetasirefphilosophiqueEXSDMRC}
\newcounter{sujetasirefphilosophiqueEXSHA}
\newcounter{sujetasirefphilosophiquepremtermEXS}
\newcounter{sujetasirefphilosophiqueMMPdP}
\newcounter{sujetasirefphilosophiqueMMDFI}
\newcounter{sujetasirefphilosophiqueMMDMRC}
\newcounter{sujetasirefphilosophiqueMMHA}
\newcounter{sujetasirefphilosophiquepremtermMM}
\newcounter{sujetasirefphilosophiqueCCPdP}
\newcounter{sujetasirefphilosophiqueCCDFI}
\newcounter{sujetasirefphilosophiqueCCDMRC}
\newcounter{sujetasirefphilosophiqueCCHA}
\newcounter{sujetasirefphilosophiquepremtermCC}
\newcounter{sujetasirefphilosophiqueHVPdP}
\newcounter{sujetasirefphilosophiqueHVDFI}
\newcounter{sujetasirefphilosophiqueHVDMRC}
\newcounter{sujetasirefphilosophiqueHVHA}
\newcounter{sujetasirefphilosophiquepremtermHV}
\newcounter{sujetasirefphilosophiqueHLPdP}
\newcounter{sujetasirefphilosophiqueHLDFI}
\newcounter{sujetasirefphilosophiqueHLDMRC}
\newcounter{sujetasirefphilosophiqueHLHA}
\newcounter{sujetasirefphilosophiquepremtermHL}
\newcounter{sujetasirefphilosophiqueRSPdP}
\newcounter{sujetasirefphilosophiqueRSRdM}
\newcounter{sujetasirefphilosophiqueRSDFI}
\newcounter{sujetasirefphilosophiqueRSDMRC}
\newcounter{sujetasirefphilosophiqueRSHA}
\newcounter{sujetasirefphilosophiqueHQPdP}
\newcounter{sujetasirefphilosophiqueHQRdM}
\newcounter{sujetasirefphilosophiqueHQDFI}
\newcounter{sujetasirefphilosophiqueHQDMRC}
\newcounter{sujetasirefphilosophiqueHQHA}
%%%Somme des compteurs%%%
\newcounter{sujetasiallrefphilosophiqueETA}
\newcounter{sujetasiallrefphilosophiqueEXS}
\newcounter{sujetasiallrefphilosophiqueMM}
\newcounter{sujetasiallrefphilosophiqueRS}
\newcounter{sujetasiallrefphilosophiqueCC}
\newcounter{sujetasiallrefphilosophiqueHV}
\newcounter{sujetasiallrefphilosophiqueHL}
\newcounter{sujetasiallrefphilosophiqueHQ}

%%%%%%%%%Réflexion littéraire%%%%%%%%%%%
\newcounter{nosujetasireflitteraire}
\newcounter{sujetasiHQreflitteraire}
\newcounter{sujetasiRSreflitteraire}
%%%sujetasis simples%%%
\newcounter{sujetasireflitteraireEXS}
\newcounter{sujetasireflitteraireETA}
\newcounter{sujetasireflitteraireMM}
\newcounter{sujetasireflitteraireHV}
\newcounter{sujetasireflitteraireCC}
\newcounter{sujetasireflitteraireHL}
%%%sujetasis mixtes%%%
\newcounter{sujetasireflitteraireEXSTA}
\newcounter{sujetasireflitteraireEXSMM}
\newcounter{sujetasireflitteraireETAMM}
\newcounter{sujetasireflitteraireOdS}
\newcounter{sujetasireflitteraireHVL}
\newcounter{sujetasireflitteraireHVCC}
\newcounter{sujetasireflitteraireHLCC}
\newcounter{sujetasireflitteraireOdH}
%%%sujetasis mixtes semestres%%%
\newcounter{sujetasireflitteraireETACC}
\newcounter{sujetasireflitteraireETAHV}
\newcounter{sujetasireflitteraireETAHL}
\newcounter{sujetasireflitteraireEXSCC}
\newcounter{sujetasireflitteraireEXSHV}
\newcounter{sujetasireflitteraireEXSHL}
\newcounter{sujetasireflitteraireMMCC}
\newcounter{sujetasireflitteraireMMHV}
\newcounter{sujetasireflitteraireMMHL}
%%%sujetasis mixtes années%%
\newcounter{sujetasireflitterairepremterm}
\newcounter{sujetasireflitterairepremtermRS}
\newcounter{sujetasireflitterairepremtermHQ}
\newcounter{sujetasireflitteraireETAPdP}
\newcounter{sujetasireflitteraireETADFI}
\newcounter{sujetasireflitteraireETADMRC}
\newcounter{sujetasireflitteraireETAHA}
\newcounter{sujetasireflitterairepremtermETA}
\newcounter{sujetasireflitteraireEXSPdP}
\newcounter{sujetasireflitteraireEXSDFI}
\newcounter{sujetasireflitteraireEXSDMRC}
\newcounter{sujetasireflitteraireEXSHA}
\newcounter{sujetasireflitterairepremtermEXS}
\newcounter{sujetasireflitteraireMMPdP}
\newcounter{sujetasireflitteraireMMDFI}
\newcounter{sujetasireflitteraireMMDMRC}
\newcounter{sujetasireflitteraireMMHA}
\newcounter{sujetasireflitterairepremtermMM}
\newcounter{sujetasireflitteraireCCPdP}
\newcounter{sujetasireflitteraireCCDFI}
\newcounter{sujetasireflitteraireCCDMRC}
\newcounter{sujetasireflitteraireCCHA}
\newcounter{sujetasireflitterairepremtermCC}
\newcounter{sujetasireflitteraireHVPdP}
\newcounter{sujetasireflitteraireHVDFI}
\newcounter{sujetasireflitteraireHVDMRC}
\newcounter{sujetasireflitteraireHVHA}
\newcounter{sujetasireflitterairepremtermHV}
\newcounter{sujetasireflitteraireHLPdP}
\newcounter{sujetasireflitteraireHLDFI}
\newcounter{sujetasireflitteraireHLDMRC}
\newcounter{sujetasireflitteraireHLHA}
\newcounter{sujetasireflitterairepremtermHL}
\newcounter{sujetasireflitteraireRSPdP}
\newcounter{sujetasireflitteraireRSRdM}
\newcounter{sujetasireflitteraireRSDFI}
\newcounter{sujetasireflitteraireRSDMRC}
\newcounter{sujetasireflitteraireRSHA}
\newcounter{sujetasireflitteraireHQPdP}
\newcounter{sujetasireflitteraireHQRdM}
\newcounter{sujetasireflitteraireHQDFI}
\newcounter{sujetasireflitteraireHQDMRC}
\newcounter{sujetasireflitteraireHQHA}
%%%Somme des compteurs%%%
\newcounter{sujetasiallreflitteraireETA}
\newcounter{sujetasiallreflitteraireEXS}
\newcounter{sujetasiallreflitteraireMM}
\newcounter{sujetasiallreflitteraireCC}
\newcounter{sujetasiallreflitteraireRS}
\newcounter{sujetasiallreflitteraireHV}
\newcounter{sujetasiallreflitteraireHL}
\newcounter{sujetasiallreflitteraireHQ}
%%%Compteurs de statistiques globales sur les genres de question%%%
\newcounter{sujetasiexclRS}
\newcounter{sujetasimixtRS}
\newcounter{sujetasiexclHQ}
\newcounter{sujetasimixtHQ}
\newcounter{sujetasimixtHQRS}
%%%Compteurs de statistiques sur les transversaux (divers semestres, première et terminale)
\newcounter{sujettransasiintlitteraire}
\newcounter{sujettransasireflitteraire}
\newcounter{sujettransasiintphilosophique}
\newcounter{sujettransasirefphilosophique}
%%% Compteurs du nombre de sujets par type
\newcounter{sujetasiintlitteraire}
\newcounter{sujetasireflitteraire}
\newcounter{sujetasiintphilosophique}
\newcounter{sujetasirefphilosophique}
%%% Compteurs par discipline
%%Litt
\newcounter{sujetasilitteraireETA}
\newcounter{sujetasilitteraireEXS}
\newcounter{sujetasilitteraireMM}
\newcounter{sujetasilitteraireCC}
\newcounter{sujetasilitteraireHV}
\newcounter{sujetasilitteraireHL}
\newcounter{sujetasilitteraireEXSTA}
\newcounter{sujetasilitteraireETAMM}
\newcounter{sujetasilitteraireEXSMM}
\newcounter{sujetasilitteraireHVCC}
\newcounter{sujetasilitteraireHVL}
\newcounter{sujetasilitteraireaut}
%%Phil
\newcounter{sujetasiphilosophiqueETA}
\newcounter{sujetasiphilosophiqueEXS}
\newcounter{sujetasiphilosophiqueMM}
\newcounter{sujetasiphilosophiqueCC}
\newcounter{sujetasiphilosophiqueHV}
\newcounter{sujetasiphilosophiqueHL}
\newcounter{sujetasiphilosophiqueEXSTA}
\newcounter{sujetasiphilosophiqueETAMM}
\newcounter{sujetasiphilosophiqueEXSMM}
\newcounter{sujetasiphilosophiqueHVCC}
\newcounter{sujetasiphilosophiqueHVL}
\newcounter{sujetasiphilosophiqueaut}
%%Totaux pour les graphiques
\newcounter{sujetasilitterairetotal}
\newcounter{sujetasiphilosophiquetotal}
%%% Autres (Ods et Odh) comptés ensemble
\newcounter{sujetasiallaut}
%%%%%%%%%%%%%%%%%%%%%%%%%%%%%%%%%%%%%%%%%%%%%%%%%%%%%%%%%%%%%%%%%%%%%%%
%%%%%%%%%%%%%%%%%%%%%%%%%%%%%%Compteurs%%%%%%%%%%%%%%%%%%%%%%%%%%%%%%%%
%%%%%%%%%%%%%%%%%%%%%%%%%%%%%%%%%%%%%%%%%%%%%%%%%%%%%%%%%%%%%%%%%%%%%%%
%%%Compter les sujets par thème%%%
%\newcounter{nosujetcea}
\newcounter{sujetceaall}%Total de questions pour le centre
\newcounter{sujetceatotal}%Total des questions en comptant en double les questions pour les graphiques
\newcounter{sujetceaHQ}
\newcounter{sujetceaRS}
%%%sujetceas simples%%%
\newcounter{sujetceaEXS}
\newcounter{sujetceaETA}
\newcounter{sujetceaMM}
\newcounter{sujetceaHV}
\newcounter{sujetceaCC}
\newcounter{sujetceaHL}
%%%sujetceas mixtes%%%
\newcounter{sujetceaEXSTA}
\newcounter{sujetceaEXSMM}
\newcounter{sujetceaETAMM}
\newcounter{sujetceaOdS}
\newcounter{sujetceaHVL}
\newcounter{sujetceaHVCC}
\newcounter{sujetceaHLCC}
\newcounter{sujetceaOdH}
%%%sujetceas mixtes semestres%%%
\newcounter{sujetceaETACC}
\newcounter{sujetceaETAHV}
\newcounter{sujetceaETAHL}
\newcounter{sujetceaEXSCC}
\newcounter{sujetceaEXSHV}
\newcounter{sujetceaEXSHL}
\newcounter{sujetceaMMCC}
\newcounter{sujetceaMMHV}
\newcounter{sujetceaMMHL}

%%%sujetceas mixtes années%%
\newcounter{sujetceapremterm}%%%Compteur général de sujetcea
\newcounter{sujetceaETAPdP}
\newcounter{sujetceaETADFI}
\newcounter{sujetceaETADMRC}
\newcounter{sujetceaETAHA}
\newcounter{sujetceapremtermETA}
\newcounter{sujetceaEXSPdP}
\newcounter{sujetceaEXSDFI}
\newcounter{sujetceaEXSDMRC}
\newcounter{sujetceaEXSHA}
\newcounter{sujetceapremtermEXS}
\newcounter{sujetceaMMPdP}
\newcounter{sujetceaMMDFI}
\newcounter{sujetceaMMDMRC}
\newcounter{sujetceaMMHA}
\newcounter{sujetceapremtermMM}
\newcounter{sujetceaCCPdP}
\newcounter{sujetceaCCDFI}
\newcounter{sujetceaCCDMRC}
\newcounter{sujetceaCCHA}
\newcounter{sujetceapremtermCC}
\newcounter{sujetceaHVPdP}
\newcounter{sujetceaHVDFI}
\newcounter{sujetceaHVDMRC}
\newcounter{sujetceaHVHA}
\newcounter{sujetceapremtermHV}
\newcounter{sujetceaHLPdP}
\newcounter{sujetceaHLDFI}
\newcounter{sujetceaHLDMRC}
\newcounter{sujetceaHLHA}
\newcounter{sujetceapremtermHL}
\newcounter{sujetceaRSPdP}
\newcounter{sujetceaRSRdM}
\newcounter{sujetceaRSDFI}
\newcounter{sujetceaRSDMRC}
\newcounter{sujetceaRSHA}
\newcounter{sujetceaHQPdP}
\newcounter{sujetceaHQRdM}
\newcounter{sujetceaHQDFI}
\newcounter{sujetceaHQDMRC}
\newcounter{sujetceaHQHA}

%%%Somme des compteurs%%%
\newcounter{sujetceaallETA}
\newcounter{sujetceaallEXS}
\newcounter{sujetceaallMM}
\newcounter{sujetceaallRS}
\newcounter{sujetceaallCC}
\newcounter{sujetceaallHV}
\newcounter{sujetceaallHL}
\newcounter{sujetceaallHQ}

%%%Compter les sujetceas par thème et types de questions%%%
%%%%%%%%%%%%%%%%%%%%%%%%%%%%%%%%%%%%%%%%%%%%%
%%%%%%%%%Question d'interprétation%%%%%%%%%%%
%%%%%%%%%%%%%%%%%%%%%%%%%%%%%%%%%%%%%%%%%%%%%

%%%%%%%%%%%%%%%Général%%%%%%%%%%%%%%%%%%%%%%%
\newcounter{nosujetceaint}
\newcounter{sujetceaHQint}
\newcounter{sujetceaRSint}
%%%sujetceas simples%%%
\newcounter{sujetceaintEXS}
\newcounter{sujetceaintETA}
\newcounter{sujetceaintMM}
\newcounter{sujetceaintHV}
\newcounter{sujetceaintCC}
\newcounter{sujetceaintHL}
%%%sujetceas mixtes%%%
\newcounter{sujetceaintEXSTA}
\newcounter{sujetceaintEXSMM}
\newcounter{sujetceaintETAMM}
\newcounter{sujetceaintOdS}
\newcounter{sujetceaintHVL}
\newcounter{sujetceaintHVCC}
\newcounter{sujetceaintHLCC}
\newcounter{sujetceaintOdH}
%%%sujetceas mixtes semestres%%%
\newcounter{sujetceaintETACC}
\newcounter{sujetceaintETAHV}
\newcounter{sujetceaintETAHL}
\newcounter{sujetceaintEXSCC}
\newcounter{sujetceaintEXSHV}
\newcounter{sujetceaintEXSHL}
\newcounter{sujetceaintMMCC}
\newcounter{sujetceaintMMHV}
\newcounter{sujetceaintMMHL}
%%%sujetceas mixtes années%%
\newcounter{sujetceaintpremterm}
\newcounter{sujetceaintpremtermRS}
\newcounter{sujetceaintpremtermHQ}
\newcounter{sujetceaintETAPdP}
\newcounter{sujetceaintETADFI}
\newcounter{sujetceaintETADMRC}
\newcounter{sujetceaintETAHA}
\newcounter{sujetceaintpremtermETA}
\newcounter{sujetceaintEXSPdP}
\newcounter{sujetceaintEXSDFI}
\newcounter{sujetceaintEXSDMRC}
\newcounter{sujetceaintEXSHA}
\newcounter{sujetceaintpremtermEXS}
\newcounter{sujetceaintMMPdP}
\newcounter{sujetceaintMMDFI}
\newcounter{sujetceaintMMDMRC}
\newcounter{sujetceaintMMHA}
\newcounter{sujetceaintpremtermMM}
\newcounter{sujetceaintCCPdP}
\newcounter{sujetceaintCCDFI}
\newcounter{sujetceaintCCDMRC}
\newcounter{sujetceaintCCHA}
\newcounter{sujetceaintpremtermCC}
\newcounter{sujetceaintHVPdP}
\newcounter{sujetceaintHVDFI}
\newcounter{sujetceaintHVDMRC}
\newcounter{sujetceaintHVHA}
\newcounter{sujetceaintpremtermHV}
\newcounter{sujetceaintHLPdP}
\newcounter{sujetceaintHLDFI}
\newcounter{sujetceaintHLDMRC}
\newcounter{sujetceaintHLHA}
\newcounter{sujetceaintpremtermHL}
\newcounter{sujetceaintRSPdP}
\newcounter{sujetceaintRSRdM}
\newcounter{sujetceaintRSDFI}
\newcounter{sujetceaintRSDMRC}
\newcounter{sujetceaintRSHA}
\newcounter{sujetceaintHQPdP}
\newcounter{sujetceaintHQRdM}
\newcounter{sujetceaintHQDFI}
\newcounter{sujetceaintHQDMRC}
\newcounter{sujetceaintHQHA}
%%%Somme des compteurs%%%
\newcounter{sujetceaallintETA}
\newcounter{sujetceaallintEXS}
\newcounter{sujetceaallintMM}
\newcounter{sujetceaallintRS}
\newcounter{sujetceaallintCC}
\newcounter{sujetceaallintHV}
\newcounter{sujetceaallintHL}
\newcounter{sujetceaallintHQ}

%%%%%%%%Interprétation philosophique%%%%%%%%%
\newcounter{nosujetceaintphilosophique}
\newcounter{sujetceaHQintphilosophique}
\newcounter{sujetceaRSintphilosophique}
%%%sujetceas simples%%%
\newcounter{sujetceaintphilosophiqueEXS}
\newcounter{sujetceaintphilosophiqueETA}
\newcounter{sujetceaintphilosophiqueMM}
\newcounter{sujetceaintphilosophiqueHV}
\newcounter{sujetceaintphilosophiqueCC}
\newcounter{sujetceaintphilosophiqueHL}
%%%sujetceas mixtes%%%
\newcounter{sujetceaintphilosophiqueEXSTA}
\newcounter{sujetceaintphilosophiqueEXSMM}
\newcounter{sujetceaintphilosophiqueETAMM}
\newcounter{sujetceaintphilosophiqueOdS}
\newcounter{sujetceaintphilosophiqueHVL}
\newcounter{sujetceaintphilosophiqueHVCC}
\newcounter{sujetceaintphilosophiqueHLCC}
\newcounter{sujetceaintphilosophiqueOdH}
%%%sujetceas mixtes semestres%%%
\newcounter{sujetceaintphilosophiqueETACC}
\newcounter{sujetceaintphilosophiqueETAHV}
\newcounter{sujetceaintphilosophiqueETAHL}
\newcounter{sujetceaintphilosophiqueEXSCC}
\newcounter{sujetceaintphilosophiqueEXSHV}
\newcounter{sujetceaintphilosophiqueEXSHL}
\newcounter{sujetceaintphilosophiqueMMCC}
\newcounter{sujetceaintphilosophiqueMMHV}
\newcounter{sujetceaintphilosophiqueMMHL}
%%%sujetceas mixtes années%%
\newcounter{sujetceaintphilosophiquepremterm}
\newcounter{sujetceaintphilosophiquepremtermRS}
\newcounter{sujetceaintphilosophiquepremtermHQ}
\newcounter{sujetceaintphilosophiqueETAPdP}
\newcounter{sujetceaintphilosophiqueETADFI}
\newcounter{sujetceaintphilosophiqueETADMRC}
\newcounter{sujetceaintphilosophiqueETAHA}
\newcounter{sujetceaintphilosophiquepremtermETA}
\newcounter{sujetceaintphilosophiqueEXSPdP}
\newcounter{sujetceaintphilosophiqueEXSDFI}
\newcounter{sujetceaintphilosophiqueEXSDMRC}
\newcounter{sujetceaintphilosophiqueEXSHA}
\newcounter{sujetceaintphilosophiquepremtermEXS}
\newcounter{sujetceaintphilosophiqueMMPdP}
\newcounter{sujetceaintphilosophiqueMMDFI}
\newcounter{sujetceaintphilosophiqueMMDMRC}
\newcounter{sujetceaintphilosophiqueMMHA}
\newcounter{sujetceaintphilosophiquepremtermMM}
\newcounter{sujetceaintphilosophiqueCCPdP}
\newcounter{sujetceaintphilosophiqueCCDFI}
\newcounter{sujetceaintphilosophiqueCCDMRC}
\newcounter{sujetceaintphilosophiqueCCHA}
\newcounter{sujetceaintphilosophiquepremtermCC}
\newcounter{sujetceaintphilosophiqueHVPdP}
\newcounter{sujetceaintphilosophiqueHVDFI}
\newcounter{sujetceaintphilosophiqueHVDMRC}
\newcounter{sujetceaintphilosophiqueHVHA}
\newcounter{sujetceaintphilosophiquepremtermHV}
\newcounter{sujetceaintphilosophiqueHLPdP}
\newcounter{sujetceaintphilosophiqueHLDFI}
\newcounter{sujetceaintphilosophiqueHLDMRC}
\newcounter{sujetceaintphilosophiqueHLHA}
\newcounter{sujetceaintphilosophiquepremtermHL}
\newcounter{sujetceaintphilosophiqueRSPdP}
\newcounter{sujetceaintphilosophiqueRSRdM}
\newcounter{sujetceaintphilosophiqueRSDFI}
\newcounter{sujetceaintphilosophiqueRSDMRC}
\newcounter{sujetceaintphilosophiqueRSHA}
\newcounter{sujetceaintphilosophiqueHQPdP}
\newcounter{sujetceaintphilosophiqueHQRdM}
\newcounter{sujetceaintphilosophiqueHQDFI}
\newcounter{sujetceaintphilosophiqueHQDMRC}
\newcounter{sujetceaintphilosophiqueHQHA}
%%%Somme des compteurs%%%
\newcounter{sujetceaallintphilosophiqueETA}
\newcounter{sujetceaallintphilosophiqueEXS}
\newcounter{sujetceaallintphilosophiqueMM}
\newcounter{sujetceaallintphilosophiqueRS}
\newcounter{sujetceaallintphilosophiqueCC}
\newcounter{sujetceaallintphilosophiqueHV}
\newcounter{sujetceaallintphilosophiqueHL}
\newcounter{sujetceaallintphilosophiqueHQ}

%%%%%%%%%Interprétation littéraire%%%%%%%%%%%
\newcounter{nosujetceaintlitteraire}
\newcounter{sujetceaHQintlitteraire}
\newcounter{sujetceaRSintlitteraire}
%%%sujetceas simples%%%
\newcounter{sujetceaintlitteraireEXS}
\newcounter{sujetceaintlitteraireETA}
\newcounter{sujetceaintlitteraireMM}
\newcounter{sujetceaintlitteraireHV}
\newcounter{sujetceaintlitteraireCC}
\newcounter{sujetceaintlitteraireHL}
%%%sujetceas mixtes%%%
\newcounter{sujetceaintlitteraireEXSTA}
\newcounter{sujetceaintlitteraireEXSMM}
\newcounter{sujetceaintlitteraireETAMM}
\newcounter{sujetceaintlitteraireOdS}
\newcounter{sujetceaintlitteraireHVL}
\newcounter{sujetceaintlitteraireHVCC}
\newcounter{sujetceaintlitteraireHLCC}
\newcounter{sujetceaintlitteraireOdH}
%%%sujetceas mixtes semestres%%%
\newcounter{sujetceaintlitteraireETACC}
\newcounter{sujetceaintlitteraireETAHV}
\newcounter{sujetceaintlitteraireETAHL}
\newcounter{sujetceaintlitteraireEXSCC}
\newcounter{sujetceaintlitteraireEXSHV}
\newcounter{sujetceaintlitteraireEXSHL}
\newcounter{sujetceaintlitteraireMMCC}
\newcounter{sujetceaintlitteraireMMHV}
\newcounter{sujetceaintlitteraireMMHL}
%%%sujetceas mixtes années%%
\newcounter{sujetceaintlitterairepremterm}
\newcounter{sujetceaintlitterairepremtermRS}
\newcounter{sujetceaintlitterairepremtermHQ}
\newcounter{sujetceaintlitteraireETAPdP}
\newcounter{sujetceaintlitteraireETADFI}
\newcounter{sujetceaintlitteraireETADMRC}
\newcounter{sujetceaintlitteraireETAHA}
\newcounter{sujetceaintlitterairepremtermETA}
\newcounter{sujetceaintlitteraireEXSPdP}
\newcounter{sujetceaintlitteraireEXSDFI}
\newcounter{sujetceaintlitteraireEXSDMRC}
\newcounter{sujetceaintlitteraireEXSHA}
\newcounter{sujetceaintlitterairepremtermEXS}
\newcounter{sujetceaintlitteraireMMPdP}
\newcounter{sujetceaintlitteraireMMDFI}
\newcounter{sujetceaintlitteraireMMDMRC}
\newcounter{sujetceaintlitteraireMMHA}
\newcounter{sujetceaintlitterairepremtermMM}
\newcounter{sujetceaintlitteraireCCPdP}
\newcounter{sujetceaintlitteraireCCDFI}
\newcounter{sujetceaintlitteraireCCDMRC}
\newcounter{sujetceaintlitteraireCCHA}
\newcounter{sujetceaintlitterairepremtermCC}
\newcounter{sujetceaintlitteraireHVPdP}
\newcounter{sujetceaintlitteraireHVDFI}
\newcounter{sujetceaintlitteraireHVDMRC}
\newcounter{sujetceaintlitteraireHVHA}
\newcounter{sujetceaintlitterairepremtermHV}
\newcounter{sujetceaintlitteraireHLPdP}
\newcounter{sujetceaintlitteraireHLDFI}
\newcounter{sujetceaintlitteraireHLDMRC}
\newcounter{sujetceaintlitteraireHLHA}
\newcounter{sujetceaintlitterairepremtermHL}
\newcounter{sujetceaintlitteraireRSPdP}
\newcounter{sujetceaintlitteraireRSRdM}
\newcounter{sujetceaintlitteraireRSDFI}
\newcounter{sujetceaintlitteraireRSDMRC}
\newcounter{sujetceaintlitteraireRSHA}
\newcounter{sujetceaintlitteraireHQPdP}
\newcounter{sujetceaintlitteraireHQRdM}
\newcounter{sujetceaintlitteraireHQDFI}
\newcounter{sujetceaintlitteraireHQDMRC}
\newcounter{sujetceaintlitteraireHQHA}
%%%Somme des compteurs%%%
\newcounter{sujetceaallintlitteraireETA}
\newcounter{sujetceaallintlitteraireEXS}
\newcounter{sujetceaallintlitteraireMM}
\newcounter{sujetceaallintlitteraireRS}
\newcounter{sujetceaallintlitteraireCC}
\newcounter{sujetceaallintlitteraireHV}
\newcounter{sujetceaallintlitteraireHL}
\newcounter{sujetceaallintlitteraireHQ}

%%%%%%%%%%%%%%%%%%%%%%%%%%%%%%%%%%%%%%%%%%%%%
%%%%%%%%%%%%Question de réflexion%%%%%%%%%%%%
%%%%%%%%%%%%%%%%%%%%%%%%%%%%%%%%%%%%%%%%%%%%%
\newcounter{nosujetcearef}
%%%sujetceas simples%%%
\newcounter{sujetcearefEXS}
\newcounter{sujetcearefETA}
\newcounter{sujetcearefMM}
\newcounter{sujetcearefHV}
\newcounter{sujetcearefCC}
\newcounter{sujetcearefHL}
%%%sujetceas mixtes%%%
\newcounter{sujetcearefEXSTA}
\newcounter{sujetcearefEXSMM}
\newcounter{sujetcearefETAMM}
\newcounter{sujetcearefOdS}
\newcounter{sujetcearefHVL}
\newcounter{sujetcearefHVCC}
\newcounter{sujetcearefHLCC}
\newcounter{sujetcearefOdH}
%%%sujetceas mixtes semestres%%%
\newcounter{sujetcearefETACC}
\newcounter{sujetcearefETAHV}
\newcounter{sujetcearefETAHL}
\newcounter{sujetcearefEXSCC}
\newcounter{sujetcearefEXSHV}
\newcounter{sujetcearefEXSHL}
\newcounter{sujetcearefMMCC}
\newcounter{sujetcearefMMHV}
\newcounter{sujetcearefMMHL}
%%%sujetceas mixtes années%%
\newcounter{sujetcearefpremterm}%%%Compteur général
\newcounter{sujetcearefETAPdP}
\newcounter{sujetcearefETADFI}
\newcounter{sujetcearefETADMRC}
\newcounter{sujetcearefETAHA}
\newcounter{sujetcearefpremtermETA}
\newcounter{sujetcearefEXSPdP}
\newcounter{sujetcearefEXSDFI}
\newcounter{sujetcearefEXSDMRC}
\newcounter{sujetcearefEXSHA}
\newcounter{sujetcearefpremtermEXS}
\newcounter{sujetcearefMMPdP}
\newcounter{sujetcearefMMDFI}
\newcounter{sujetcearefMMDMRC}
\newcounter{sujetcearefMMHA}
\newcounter{sujetcearefpremtermMM}
\newcounter{sujetcearefCCPdP}
\newcounter{sujetcearefCCDFI}
\newcounter{sujetcearefCCDMRC}
\newcounter{sujetcearefCCHA}
\newcounter{sujetcearefpremtermCC}
\newcounter{sujetcearefHVPdP}
\newcounter{sujetcearefHVDFI}
\newcounter{sujetcearefHVDMRC}
\newcounter{sujetcearefHVHA}
\newcounter{sujetcearefpremtermHV}
\newcounter{sujetcearefHLPdP}
\newcounter{sujetcearefHLDFI}
\newcounter{sujetcearefHLDMRC}
\newcounter{sujetcearefHLHA}
\newcounter{sujetcearefpremtermHL}
\newcounter{sujetcearefRSPdP}
\newcounter{sujetcearefRSRdM}
\newcounter{sujetcearefRSDFI}
\newcounter{sujetcearefRSDMRC}
\newcounter{sujetcearefRSHA}
\newcounter{sujetcearefHQPdP}
\newcounter{sujetcearefHQRdM}
\newcounter{sujetcearefHQDFI}
\newcounter{sujetcearefHQDMRC}
\newcounter{sujetcearefHQHA}
%%%Somme des compteurs%%%
\newcounter{sujetceaallrefETA}
\newcounter{sujetceaallrefEXS}
\newcounter{sujetceaallrefMM}
\newcounter{sujetceaallrefRS}
\newcounter{sujetceaallrefCC}
\newcounter{sujetceaallrefHV}
\newcounter{sujetceaallrefHL}
\newcounter{sujetceaallrefHQ}

%%%%%%%%Réflexion philosophique%%%%%%%%%
\newcounter{nosujetcearefphilosophique}
\newcounter{sujetceaHQrefphilosophique}
\newcounter{sujetceaRSrefphilosophique}
%%%sujetceas simples%%%
\newcounter{sujetcearefphilosophiqueEXS}
\newcounter{sujetcearefphilosophiqueETA}
\newcounter{sujetcearefphilosophiqueMM}
\newcounter{sujetcearefphilosophiqueHV}
\newcounter{sujetcearefphilosophiqueCC}
\newcounter{sujetcearefphilosophiqueHL}
%%%sujetceas mixtes%%%
\newcounter{sujetcearefphilosophiqueEXSTA}
\newcounter{sujetcearefphilosophiqueEXSMM}
\newcounter{sujetcearefphilosophiqueETAMM}
\newcounter{sujetcearefphilosophiqueOdS}
\newcounter{sujetcearefphilosophiqueHVL}
\newcounter{sujetcearefphilosophiqueHVCC}
\newcounter{sujetcearefphilosophiqueHLCC}
\newcounter{sujetcearefphilosophiqueOdH}
%%%sujetceas mixtes semestres%%%
\newcounter{sujetcearefphilosophiqueETACC}
\newcounter{sujetcearefphilosophiqueETAHV}
\newcounter{sujetcearefphilosophiqueETAHL}
\newcounter{sujetcearefphilosophiqueEXSCC}
\newcounter{sujetcearefphilosophiqueEXSHV}
\newcounter{sujetcearefphilosophiqueEXSHL}
\newcounter{sujetcearefphilosophiqueMMCC}
\newcounter{sujetcearefphilosophiqueMMHV}
\newcounter{sujetcearefphilosophiqueMMHL}
%%%sujetceas mixtes années%%
\newcounter{sujetcearefphilosophiquepremterm}
\newcounter{sujetcearefphilosophiquepremtermRS}
\newcounter{sujetcearefphilosophiquepremtermHQ}
\newcounter{sujetcearefphilosophiqueETAPdP}
\newcounter{sujetcearefphilosophiqueETADFI}
\newcounter{sujetcearefphilosophiqueETADMRC}
\newcounter{sujetcearefphilosophiqueETAHA}\newcounter{sujetcearefphilosophiquepremtermETA}
\newcounter{sujetcearefphilosophiqueEXSPdP}
\newcounter{sujetcearefphilosophiqueEXSDFI}
\newcounter{sujetcearefphilosophiqueEXSDMRC}
\newcounter{sujetcearefphilosophiqueEXSHA}
\newcounter{sujetcearefphilosophiquepremtermEXS}
\newcounter{sujetcearefphilosophiqueMMPdP}
\newcounter{sujetcearefphilosophiqueMMDFI}
\newcounter{sujetcearefphilosophiqueMMDMRC}
\newcounter{sujetcearefphilosophiqueMMHA}
\newcounter{sujetcearefphilosophiquepremtermMM}
\newcounter{sujetcearefphilosophiqueCCPdP}
\newcounter{sujetcearefphilosophiqueCCDFI}
\newcounter{sujetcearefphilosophiqueCCDMRC}
\newcounter{sujetcearefphilosophiqueCCHA}
\newcounter{sujetcearefphilosophiquepremtermCC}
\newcounter{sujetcearefphilosophiqueHVPdP}
\newcounter{sujetcearefphilosophiqueHVDFI}
\newcounter{sujetcearefphilosophiqueHVDMRC}
\newcounter{sujetcearefphilosophiqueHVHA}
\newcounter{sujetcearefphilosophiquepremtermHV}
\newcounter{sujetcearefphilosophiqueHLPdP}
\newcounter{sujetcearefphilosophiqueHLDFI}
\newcounter{sujetcearefphilosophiqueHLDMRC}
\newcounter{sujetcearefphilosophiqueHLHA}
\newcounter{sujetcearefphilosophiquepremtermHL}
\newcounter{sujetcearefphilosophiqueRSPdP}
\newcounter{sujetcearefphilosophiqueRSRdM}
\newcounter{sujetcearefphilosophiqueRSDFI}
\newcounter{sujetcearefphilosophiqueRSDMRC}
\newcounter{sujetcearefphilosophiqueRSHA}
\newcounter{sujetcearefphilosophiqueHQPdP}
\newcounter{sujetcearefphilosophiqueHQRdM}
\newcounter{sujetcearefphilosophiqueHQDFI}
\newcounter{sujetcearefphilosophiqueHQDMRC}
\newcounter{sujetcearefphilosophiqueHQHA}
%%%Somme des compteurs%%%
\newcounter{sujetceaallrefphilosophiqueETA}
\newcounter{sujetceaallrefphilosophiqueEXS}
\newcounter{sujetceaallrefphilosophiqueMM}
\newcounter{sujetceaallrefphilosophiqueRS}
\newcounter{sujetceaallrefphilosophiqueCC}
\newcounter{sujetceaallrefphilosophiqueHV}
\newcounter{sujetceaallrefphilosophiqueHL}
\newcounter{sujetceaallrefphilosophiqueHQ}

%%%%%%%%%Réflexion littéraire%%%%%%%%%%%
\newcounter{nosujetceareflitteraire}
\newcounter{sujetceaHQreflitteraire}
\newcounter{sujetceaRSreflitteraire}
%%%sujetceas simples%%%
\newcounter{sujetceareflitteraireEXS}
\newcounter{sujetceareflitteraireETA}
\newcounter{sujetceareflitteraireMM}
\newcounter{sujetceareflitteraireHV}
\newcounter{sujetceareflitteraireCC}
\newcounter{sujetceareflitteraireHL}
%%%sujetceas mixtes%%%
\newcounter{sujetceareflitteraireEXSTA}
\newcounter{sujetceareflitteraireEXSMM}
\newcounter{sujetceareflitteraireETAMM}
\newcounter{sujetceareflitteraireOdS}
\newcounter{sujetceareflitteraireHVL}
\newcounter{sujetceareflitteraireHVCC}
\newcounter{sujetceareflitteraireHLCC}
\newcounter{sujetceareflitteraireOdH}
%%%sujetceas mixtes semestres%%%
\newcounter{sujetceareflitteraireETACC}
\newcounter{sujetceareflitteraireETAHV}
\newcounter{sujetceareflitteraireETAHL}
\newcounter{sujetceareflitteraireEXSCC}
\newcounter{sujetceareflitteraireEXSHV}
\newcounter{sujetceareflitteraireEXSHL}
\newcounter{sujetceareflitteraireMMCC}
\newcounter{sujetceareflitteraireMMHV}
\newcounter{sujetceareflitteraireMMHL}
%%%sujetceas mixtes années%%
\newcounter{sujetceareflitterairepremterm}
\newcounter{sujetceareflitterairepremtermRS}
\newcounter{sujetceareflitterairepremtermHQ}
\newcounter{sujetceareflitteraireETAPdP}
\newcounter{sujetceareflitteraireETADFI}
\newcounter{sujetceareflitteraireETADMRC}
\newcounter{sujetceareflitteraireETAHA}
\newcounter{sujetceareflitterairepremtermETA}
\newcounter{sujetceareflitteraireEXSPdP}
\newcounter{sujetceareflitteraireEXSDFI}
\newcounter{sujetceareflitteraireEXSDMRC}
\newcounter{sujetceareflitteraireEXSHA}
\newcounter{sujetceareflitterairepremtermEXS}
\newcounter{sujetceareflitteraireMMPdP}
\newcounter{sujetceareflitteraireMMDFI}
\newcounter{sujetceareflitteraireMMDMRC}
\newcounter{sujetceareflitteraireMMHA}
\newcounter{sujetceareflitterairepremtermMM}
\newcounter{sujetceareflitteraireCCPdP}
\newcounter{sujetceareflitteraireCCDFI}
\newcounter{sujetceareflitteraireCCDMRC}
\newcounter{sujetceareflitteraireCCHA}
\newcounter{sujetceareflitterairepremtermCC}
\newcounter{sujetceareflitteraireHVPdP}
\newcounter{sujetceareflitteraireHVDFI}
\newcounter{sujetceareflitteraireHVDMRC}
\newcounter{sujetceareflitteraireHVHA}
\newcounter{sujetceareflitterairepremtermHV}
\newcounter{sujetceareflitteraireHLPdP}
\newcounter{sujetceareflitteraireHLDFI}
\newcounter{sujetceareflitteraireHLDMRC}
\newcounter{sujetceareflitteraireHLHA}
\newcounter{sujetceareflitterairepremtermHL}
\newcounter{sujetceareflitteraireRSPdP}
\newcounter{sujetceareflitteraireRSRdM}
\newcounter{sujetceareflitteraireRSDFI}
\newcounter{sujetceareflitteraireRSDMRC}
\newcounter{sujetceareflitteraireRSHA}
\newcounter{sujetceareflitteraireHQPdP}
\newcounter{sujetceareflitteraireHQRdM}
\newcounter{sujetceareflitteraireHQDFI}
\newcounter{sujetceareflitteraireHQDMRC}
\newcounter{sujetceareflitteraireHQHA}
%%%Somme des compteurs%%%
\newcounter{sujetceaallreflitteraireETA}
\newcounter{sujetceaallreflitteraireEXS}
\newcounter{sujetceaallreflitteraireMM}
\newcounter{sujetceaallreflitteraireCC}
\newcounter{sujetceaallreflitteraireRS}
\newcounter{sujetceaallreflitteraireHV}
\newcounter{sujetceaallreflitteraireHL}
\newcounter{sujetceaallreflitteraireHQ}
%%%Compteurs de statistiques globales sur les genres de question%%%
\newcounter{sujetceaexclRS}
\newcounter{sujetceamixtRS}
\newcounter{sujetceaexclHQ}
\newcounter{sujetceamixtHQ}
\newcounter{sujetceamixtHQRS}
%%%Compteurs de statistiques sur les transversaux (divers semestres, première et terminale)
\newcounter{sujettransceaintlitteraire}
\newcounter{sujettransceareflitteraire}
\newcounter{sujettransceaintphilosophique}
\newcounter{sujettranscearefphilosophique}
%%% Compteurs du nombre de sujets par type
\newcounter{sujetceaintlitteraire}
\newcounter{sujetceareflitteraire}
\newcounter{sujetceaintphilosophique}
\newcounter{sujetcearefphilosophique}
%%% Compteurs par discipline
%%Litt
\newcounter{sujetcealitteraireETA}
\newcounter{sujetcealitteraireEXS}
\newcounter{sujetcealitteraireMM}
\newcounter{sujetcealitteraireCC}
\newcounter{sujetcealitteraireHV}
\newcounter{sujetcealitteraireHL}
\newcounter{sujetcealitteraireEXSTA}
\newcounter{sujetcealitteraireETAMM}
\newcounter{sujetcealitteraireEXSMM}
\newcounter{sujetcealitteraireHVCC}
\newcounter{sujetcealitteraireHVL}
\newcounter{sujetcealitteraireaut}
%%Phil
\newcounter{sujetceaphilosophiqueETA}
\newcounter{sujetceaphilosophiqueEXS}
\newcounter{sujetceaphilosophiqueMM}
\newcounter{sujetceaphilosophiqueCC}
\newcounter{sujetceaphilosophiqueHV}
\newcounter{sujetceaphilosophiqueHL}
\newcounter{sujetceaphilosophiqueEXSTA}
\newcounter{sujetceaphilosophiqueETAMM}
\newcounter{sujetceaphilosophiqueEXSMM}
\newcounter{sujetceaphilosophiqueHVCC}
\newcounter{sujetceaphilosophiqueHVL}
\newcounter{sujetceaphilosophiqueaut}
%%Totaux pour les graphiques
\newcounter{sujetcealitterairetotal}
\newcounter{sujetceaphilosophiquetotal}
%%% Autres (Ods et Odh) comptés ensemble
\newcounter{sujetceaallaut}
%%%%%%%%%%%%%%%%%%%%%%%%%%%%%%%%%%%%%%%%%%%%%%%%%%%%%%%%%%%%%%%%%%%%%%%
%%%%%%%%%%%%%%%%%%%%%%%%%%%%%%Compteurs%%%%%%%%%%%%%%%%%%%%%%%%%%%%%%%%
%%%%%%%%%%%%%%%%%%%%%%%%%%%%%%%%%%%%%%%%%%%%%%%%%%%%%%%%%%%%%%%%%%%%%%%
%%%Compter les sujets par thème%%%
%\newcounter{nosujetlib}
\newcounter{sujetliball}%Total de questions pour le centre
\newcounter{sujetlibtotal}%Total des questions en comptant en double les questions pour les graphiques
\newcounter{sujetlibHQ}
\newcounter{sujetlibRS}
%%%sujetlibs simples%%%
\newcounter{sujetlibEXS}
\newcounter{sujetlibETA}
\newcounter{sujetlibMM}
\newcounter{sujetlibHV}
\newcounter{sujetlibCC}
\newcounter{sujetlibHL}
%%%sujetlibs mixtes%%%
\newcounter{sujetlibEXSTA}
\newcounter{sujetlibEXSMM}
\newcounter{sujetlibETAMM}
\newcounter{sujetlibOdS}
\newcounter{sujetlibHVL}
\newcounter{sujetlibHVCC}
\newcounter{sujetlibHLCC}
\newcounter{sujetlibOdH}
%%%sujetlibs mixtes semestres%%%
\newcounter{sujetlibETACC}
\newcounter{sujetlibETAHV}
\newcounter{sujetlibETAHL}
\newcounter{sujetlibEXSCC}
\newcounter{sujetlibEXSHV}
\newcounter{sujetlibEXSHL}
\newcounter{sujetlibMMCC}
\newcounter{sujetlibMMHV}
\newcounter{sujetlibMMHL}

%%%sujetlibs mixtes années%%
\newcounter{sujetlibpremterm}%%%Compteur général de sujetlib
\newcounter{sujetlibETAPdP}
\newcounter{sujetlibETADFI}
\newcounter{sujetlibETADMRC}
\newcounter{sujetlibETAHA}
\newcounter{sujetlibpremtermETA}
\newcounter{sujetlibEXSPdP}
\newcounter{sujetlibEXSDFI}
\newcounter{sujetlibEXSDMRC}
\newcounter{sujetlibEXSHA}
\newcounter{sujetlibpremtermEXS}
\newcounter{sujetlibMMPdP}
\newcounter{sujetlibMMDFI}
\newcounter{sujetlibMMDMRC}
\newcounter{sujetlibMMHA}
\newcounter{sujetlibpremtermMM}
\newcounter{sujetlibCCPdP}
\newcounter{sujetlibCCDFI}
\newcounter{sujetlibCCDMRC}
\newcounter{sujetlibCCHA}
\newcounter{sujetlibpremtermCC}
\newcounter{sujetlibHVPdP}
\newcounter{sujetlibHVDFI}
\newcounter{sujetlibHVDMRC}
\newcounter{sujetlibHVHA}
\newcounter{sujetlibpremtermHV}
\newcounter{sujetlibHLPdP}
\newcounter{sujetlibHLDFI}
\newcounter{sujetlibHLDMRC}
\newcounter{sujetlibHLHA}
\newcounter{sujetlibpremtermHL}
\newcounter{sujetlibRSPdP}
\newcounter{sujetlibRSRdM}
\newcounter{sujetlibRSDFI}
\newcounter{sujetlibRSDMRC}
\newcounter{sujetlibRSHA}
\newcounter{sujetlibHQPdP}
\newcounter{sujetlibHQRdM}
\newcounter{sujetlibHQDFI}
\newcounter{sujetlibHQDMRC}
\newcounter{sujetlibHQHA}

%%%Somme des compteurs%%%
\newcounter{sujetliballETA}
\newcounter{sujetliballEXS}
\newcounter{sujetliballMM}
\newcounter{sujetliballRS}
\newcounter{sujetliballCC}
\newcounter{sujetliballHV}
\newcounter{sujetliballHL}
\newcounter{sujetliballHQ}

%%%Compter les sujetlibs par thème et types de questions%%%
%%%%%%%%%%%%%%%%%%%%%%%%%%%%%%%%%%%%%%%%%%%%%
%%%%%%%%%Question d'interprétation%%%%%%%%%%%
%%%%%%%%%%%%%%%%%%%%%%%%%%%%%%%%%%%%%%%%%%%%%

%%%%%%%%%%%%%%%Général%%%%%%%%%%%%%%%%%%%%%%%
\newcounter{nosujetlibint}
\newcounter{sujetlibHQint}
\newcounter{sujetlibRSint}
%%%sujetlibs simples%%%
\newcounter{sujetlibintEXS}
\newcounter{sujetlibintETA}
\newcounter{sujetlibintMM}
\newcounter{sujetlibintHV}
\newcounter{sujetlibintCC}
\newcounter{sujetlibintHL}
%%%sujetlibs mixtes%%%
\newcounter{sujetlibintEXSTA}
\newcounter{sujetlibintEXSMM}
\newcounter{sujetlibintETAMM}
\newcounter{sujetlibintOdS}
\newcounter{sujetlibintHVL}
\newcounter{sujetlibintHVCC}
\newcounter{sujetlibintHLCC}
\newcounter{sujetlibintOdH}
%%%sujetlibs mixtes semestres%%%
\newcounter{sujetlibintETACC}
\newcounter{sujetlibintETAHV}
\newcounter{sujetlibintETAHL}
\newcounter{sujetlibintEXSCC}
\newcounter{sujetlibintEXSHV}
\newcounter{sujetlibintEXSHL}
\newcounter{sujetlibintMMCC}
\newcounter{sujetlibintMMHV}
\newcounter{sujetlibintMMHL}
%%%sujetlibs mixtes années%%
\newcounter{sujetlibintpremterm}
\newcounter{sujetlibintpremtermRS}
\newcounter{sujetlibintpremtermHQ}
\newcounter{sujetlibintETAPdP}
\newcounter{sujetlibintETADFI}
\newcounter{sujetlibintETADMRC}
\newcounter{sujetlibintETAHA}
\newcounter{sujetlibintpremtermETA}
\newcounter{sujetlibintEXSPdP}
\newcounter{sujetlibintEXSDFI}
\newcounter{sujetlibintEXSDMRC}
\newcounter{sujetlibintEXSHA}
\newcounter{sujetlibintpremtermEXS}
\newcounter{sujetlibintMMPdP}
\newcounter{sujetlibintMMDFI}
\newcounter{sujetlibintMMDMRC}
\newcounter{sujetlibintMMHA}
\newcounter{sujetlibintpremtermMM}
\newcounter{sujetlibintCCPdP}
\newcounter{sujetlibintCCDFI}
\newcounter{sujetlibintCCDMRC}
\newcounter{sujetlibintCCHA}
\newcounter{sujetlibintpremtermCC}
\newcounter{sujetlibintHVPdP}
\newcounter{sujetlibintHVDFI}
\newcounter{sujetlibintHVDMRC}
\newcounter{sujetlibintHVHA}
\newcounter{sujetlibintpremtermHV}
\newcounter{sujetlibintHLPdP}
\newcounter{sujetlibintHLDFI}
\newcounter{sujetlibintHLDMRC}
\newcounter{sujetlibintHLHA}
\newcounter{sujetlibintpremtermHL}
\newcounter{sujetlibintRSPdP}
\newcounter{sujetlibintRSRdM}
\newcounter{sujetlibintRSDFI}
\newcounter{sujetlibintRSDMRC}
\newcounter{sujetlibintRSHA}
\newcounter{sujetlibintHQPdP}
\newcounter{sujetlibintHQRdM}
\newcounter{sujetlibintHQDFI}
\newcounter{sujetlibintHQDMRC}
\newcounter{sujetlibintHQHA}
%%%Somme des compteurs%%%
\newcounter{sujetliballintETA}
\newcounter{sujetliballintEXS}
\newcounter{sujetliballintMM}
\newcounter{sujetliballintRS}
\newcounter{sujetliballintCC}
\newcounter{sujetliballintHV}
\newcounter{sujetliballintHL}
\newcounter{sujetliballintHQ}

%%%%%%%%Interprétation philosophique%%%%%%%%%
\newcounter{nosujetlibintphilosophique}
\newcounter{sujetlibHQintphilosophique}
\newcounter{sujetlibRSintphilosophique}
%%%sujetlibs simples%%%
\newcounter{sujetlibintphilosophiqueEXS}
\newcounter{sujetlibintphilosophiqueETA}
\newcounter{sujetlibintphilosophiqueMM}
\newcounter{sujetlibintphilosophiqueHV}
\newcounter{sujetlibintphilosophiqueCC}
\newcounter{sujetlibintphilosophiqueHL}
%%%sujetlibs mixtes%%%
\newcounter{sujetlibintphilosophiqueEXSTA}
\newcounter{sujetlibintphilosophiqueEXSMM}
\newcounter{sujetlibintphilosophiqueETAMM}
\newcounter{sujetlibintphilosophiqueOdS}
\newcounter{sujetlibintphilosophiqueHVL}
\newcounter{sujetlibintphilosophiqueHVCC}
\newcounter{sujetlibintphilosophiqueHLCC}
\newcounter{sujetlibintphilosophiqueOdH}
%%%sujetlibs mixtes semestres%%%
\newcounter{sujetlibintphilosophiqueETACC}
\newcounter{sujetlibintphilosophiqueETAHV}
\newcounter{sujetlibintphilosophiqueETAHL}
\newcounter{sujetlibintphilosophiqueEXSCC}
\newcounter{sujetlibintphilosophiqueEXSHV}
\newcounter{sujetlibintphilosophiqueEXSHL}
\newcounter{sujetlibintphilosophiqueMMCC}
\newcounter{sujetlibintphilosophiqueMMHV}
\newcounter{sujetlibintphilosophiqueMMHL}
%%%sujetlibs mixtes années%%
\newcounter{sujetlibintphilosophiquepremterm}
\newcounter{sujetlibintphilosophiquepremtermRS}
\newcounter{sujetlibintphilosophiquepremtermHQ}
\newcounter{sujetlibintphilosophiqueETAPdP}
\newcounter{sujetlibintphilosophiqueETADFI}
\newcounter{sujetlibintphilosophiqueETADMRC}
\newcounter{sujetlibintphilosophiqueETAHA}
\newcounter{sujetlibintphilosophiquepremtermETA}
\newcounter{sujetlibintphilosophiqueEXSPdP}
\newcounter{sujetlibintphilosophiqueEXSDFI}
\newcounter{sujetlibintphilosophiqueEXSDMRC}
\newcounter{sujetlibintphilosophiqueEXSHA}
\newcounter{sujetlibintphilosophiquepremtermEXS}
\newcounter{sujetlibintphilosophiqueMMPdP}
\newcounter{sujetlibintphilosophiqueMMDFI}
\newcounter{sujetlibintphilosophiqueMMDMRC}
\newcounter{sujetlibintphilosophiqueMMHA}
\newcounter{sujetlibintphilosophiquepremtermMM}
\newcounter{sujetlibintphilosophiqueCCPdP}
\newcounter{sujetlibintphilosophiqueCCDFI}
\newcounter{sujetlibintphilosophiqueCCDMRC}
\newcounter{sujetlibintphilosophiqueCCHA}
\newcounter{sujetlibintphilosophiquepremtermCC}
\newcounter{sujetlibintphilosophiqueHVPdP}
\newcounter{sujetlibintphilosophiqueHVDFI}
\newcounter{sujetlibintphilosophiqueHVDMRC}
\newcounter{sujetlibintphilosophiqueHVHA}
\newcounter{sujetlibintphilosophiquepremtermHV}
\newcounter{sujetlibintphilosophiqueHLPdP}
\newcounter{sujetlibintphilosophiqueHLDFI}
\newcounter{sujetlibintphilosophiqueHLDMRC}
\newcounter{sujetlibintphilosophiqueHLHA}
\newcounter{sujetlibintphilosophiquepremtermHL}
\newcounter{sujetlibintphilosophiqueRSPdP}
\newcounter{sujetlibintphilosophiqueRSRdM}
\newcounter{sujetlibintphilosophiqueRSDFI}
\newcounter{sujetlibintphilosophiqueRSDMRC}
\newcounter{sujetlibintphilosophiqueRSHA}
\newcounter{sujetlibintphilosophiqueHQPdP}
\newcounter{sujetlibintphilosophiqueHQRdM}
\newcounter{sujetlibintphilosophiqueHQDFI}
\newcounter{sujetlibintphilosophiqueHQDMRC}
\newcounter{sujetlibintphilosophiqueHQHA}
%%%Somme des compteurs%%%
\newcounter{sujetliballintphilosophiqueETA}
\newcounter{sujetliballintphilosophiqueEXS}
\newcounter{sujetliballintphilosophiqueMM}
\newcounter{sujetliballintphilosophiqueRS}
\newcounter{sujetliballintphilosophiqueCC}
\newcounter{sujetliballintphilosophiqueHV}
\newcounter{sujetliballintphilosophiqueHL}
\newcounter{sujetliballintphilosophiqueHQ}

%%%%%%%%%Interprétation littéraire%%%%%%%%%%%
\newcounter{nosujetlibintlitteraire}
\newcounter{sujetlibHQintlitteraire}
\newcounter{sujetlibRSintlitteraire}
%%%sujetlibs simples%%%
\newcounter{sujetlibintlitteraireEXS}
\newcounter{sujetlibintlitteraireETA}
\newcounter{sujetlibintlitteraireMM}
\newcounter{sujetlibintlitteraireHV}
\newcounter{sujetlibintlitteraireCC}
\newcounter{sujetlibintlitteraireHL}
%%%sujetlibs mixtes%%%
\newcounter{sujetlibintlitteraireEXSTA}
\newcounter{sujetlibintlitteraireEXSMM}
\newcounter{sujetlibintlitteraireETAMM}
\newcounter{sujetlibintlitteraireOdS}
\newcounter{sujetlibintlitteraireHVL}
\newcounter{sujetlibintlitteraireHVCC}
\newcounter{sujetlibintlitteraireHLCC}
\newcounter{sujetlibintlitteraireOdH}
%%%sujetlibs mixtes semestres%%%
\newcounter{sujetlibintlitteraireETACC}
\newcounter{sujetlibintlitteraireETAHV}
\newcounter{sujetlibintlitteraireETAHL}
\newcounter{sujetlibintlitteraireEXSCC}
\newcounter{sujetlibintlitteraireEXSHV}
\newcounter{sujetlibintlitteraireEXSHL}
\newcounter{sujetlibintlitteraireMMCC}
\newcounter{sujetlibintlitteraireMMHV}
\newcounter{sujetlibintlitteraireMMHL}
%%%sujetlibs mixtes années%%
\newcounter{sujetlibintlitterairepremterm}
\newcounter{sujetlibintlitterairepremtermRS}
\newcounter{sujetlibintlitterairepremtermHQ}
\newcounter{sujetlibintlitteraireETAPdP}
\newcounter{sujetlibintlitteraireETADFI}
\newcounter{sujetlibintlitteraireETADMRC}
\newcounter{sujetlibintlitteraireETAHA}
\newcounter{sujetlibintlitterairepremtermETA}
\newcounter{sujetlibintlitteraireEXSPdP}
\newcounter{sujetlibintlitteraireEXSDFI}
\newcounter{sujetlibintlitteraireEXSDMRC}
\newcounter{sujetlibintlitteraireEXSHA}
\newcounter{sujetlibintlitterairepremtermEXS}
\newcounter{sujetlibintlitteraireMMPdP}
\newcounter{sujetlibintlitteraireMMDFI}
\newcounter{sujetlibintlitteraireMMDMRC}
\newcounter{sujetlibintlitteraireMMHA}
\newcounter{sujetlibintlitterairepremtermMM}
\newcounter{sujetlibintlitteraireCCPdP}
\newcounter{sujetlibintlitteraireCCDFI}
\newcounter{sujetlibintlitteraireCCDMRC}
\newcounter{sujetlibintlitteraireCCHA}
\newcounter{sujetlibintlitterairepremtermCC}
\newcounter{sujetlibintlitteraireHVPdP}
\newcounter{sujetlibintlitteraireHVDFI}
\newcounter{sujetlibintlitteraireHVDMRC}
\newcounter{sujetlibintlitteraireHVHA}
\newcounter{sujetlibintlitterairepremtermHV}
\newcounter{sujetlibintlitteraireHLPdP}
\newcounter{sujetlibintlitteraireHLDFI}
\newcounter{sujetlibintlitteraireHLDMRC}
\newcounter{sujetlibintlitteraireHLHA}
\newcounter{sujetlibintlitterairepremtermHL}
\newcounter{sujetlibintlitteraireRSPdP}
\newcounter{sujetlibintlitteraireRSRdM}
\newcounter{sujetlibintlitteraireRSDFI}
\newcounter{sujetlibintlitteraireRSDMRC}
\newcounter{sujetlibintlitteraireRSHA}
\newcounter{sujetlibintlitteraireHQPdP}
\newcounter{sujetlibintlitteraireHQRdM}
\newcounter{sujetlibintlitteraireHQDFI}
\newcounter{sujetlibintlitteraireHQDMRC}
\newcounter{sujetlibintlitteraireHQHA}
%%%Somme des compteurs%%%
\newcounter{sujetliballintlitteraireETA}
\newcounter{sujetliballintlitteraireEXS}
\newcounter{sujetliballintlitteraireMM}
\newcounter{sujetliballintlitteraireRS}
\newcounter{sujetliballintlitteraireCC}
\newcounter{sujetliballintlitteraireHV}
\newcounter{sujetliballintlitteraireHL}
\newcounter{sujetliballintlitteraireHQ}

%%%%%%%%%%%%%%%%%%%%%%%%%%%%%%%%%%%%%%%%%%%%%
%%%%%%%%%%%%Question de réflexion%%%%%%%%%%%%
%%%%%%%%%%%%%%%%%%%%%%%%%%%%%%%%%%%%%%%%%%%%%
\newcounter{nosujetlibref}
%%%sujetlibs simples%%%
\newcounter{sujetlibrefEXS}
\newcounter{sujetlibrefETA}
\newcounter{sujetlibrefMM}
\newcounter{sujetlibrefHV}
\newcounter{sujetlibrefCC}
\newcounter{sujetlibrefHL}
%%%sujetlibs mixtes%%%
\newcounter{sujetlibrefEXSTA}
\newcounter{sujetlibrefEXSMM}
\newcounter{sujetlibrefETAMM}
\newcounter{sujetlibrefOdS}
\newcounter{sujetlibrefHVL}
\newcounter{sujetlibrefHVCC}
\newcounter{sujetlibrefHLCC}
\newcounter{sujetlibrefOdH}
%%%sujetlibs mixtes semestres%%%
\newcounter{sujetlibrefETACC}
\newcounter{sujetlibrefETAHV}
\newcounter{sujetlibrefETAHL}
\newcounter{sujetlibrefEXSCC}
\newcounter{sujetlibrefEXSHV}
\newcounter{sujetlibrefEXSHL}
\newcounter{sujetlibrefMMCC}
\newcounter{sujetlibrefMMHV}
\newcounter{sujetlibrefMMHL}
%%%sujetlibs mixtes années%%
\newcounter{sujetlibrefpremterm}%%%Compteur général
\newcounter{sujetlibrefETAPdP}
\newcounter{sujetlibrefETADFI}
\newcounter{sujetlibrefETADMRC}
\newcounter{sujetlibrefETAHA}
\newcounter{sujetlibrefpremtermETA}
\newcounter{sujetlibrefEXSPdP}
\newcounter{sujetlibrefEXSDFI}
\newcounter{sujetlibrefEXSDMRC}
\newcounter{sujetlibrefEXSHA}
\newcounter{sujetlibrefpremtermEXS}
\newcounter{sujetlibrefMMPdP}
\newcounter{sujetlibrefMMDFI}
\newcounter{sujetlibrefMMDMRC}
\newcounter{sujetlibrefMMHA}
\newcounter{sujetlibrefpremtermMM}
\newcounter{sujetlibrefCCPdP}
\newcounter{sujetlibrefCCDFI}
\newcounter{sujetlibrefCCDMRC}
\newcounter{sujetlibrefCCHA}
\newcounter{sujetlibrefpremtermCC}
\newcounter{sujetlibrefHVPdP}
\newcounter{sujetlibrefHVDFI}
\newcounter{sujetlibrefHVDMRC}
\newcounter{sujetlibrefHVHA}
\newcounter{sujetlibrefpremtermHV}
\newcounter{sujetlibrefHLPdP}
\newcounter{sujetlibrefHLDFI}
\newcounter{sujetlibrefHLDMRC}
\newcounter{sujetlibrefHLHA}
\newcounter{sujetlibrefpremtermHL}
\newcounter{sujetlibrefRSPdP}
\newcounter{sujetlibrefRSRdM}
\newcounter{sujetlibrefRSDFI}
\newcounter{sujetlibrefRSDMRC}
\newcounter{sujetlibrefRSHA}
\newcounter{sujetlibrefHQPdP}
\newcounter{sujetlibrefHQRdM}
\newcounter{sujetlibrefHQDFI}
\newcounter{sujetlibrefHQDMRC}
\newcounter{sujetlibrefHQHA}
%%%Somme des compteurs%%%
\newcounter{sujetliballrefETA}
\newcounter{sujetliballrefEXS}
\newcounter{sujetliballrefMM}
\newcounter{sujetliballrefRS}
\newcounter{sujetliballrefCC}
\newcounter{sujetliballrefHV}
\newcounter{sujetliballrefHL}
\newcounter{sujetliballrefHQ}

%%%%%%%%Réflexion philosophique%%%%%%%%%
\newcounter{nosujetlibrefphilosophique}
\newcounter{sujetlibHQrefphilosophique}
\newcounter{sujetlibRSrefphilosophique}
%%%sujetlibs simples%%%
\newcounter{sujetlibrefphilosophiqueEXS}
\newcounter{sujetlibrefphilosophiqueETA}
\newcounter{sujetlibrefphilosophiqueMM}
\newcounter{sujetlibrefphilosophiqueHV}
\newcounter{sujetlibrefphilosophiqueCC}
\newcounter{sujetlibrefphilosophiqueHL}
%%%sujetlibs mixtes%%%
\newcounter{sujetlibrefphilosophiqueEXSTA}
\newcounter{sujetlibrefphilosophiqueEXSMM}
\newcounter{sujetlibrefphilosophiqueETAMM}
\newcounter{sujetlibrefphilosophiqueOdS}
\newcounter{sujetlibrefphilosophiqueHVL}
\newcounter{sujetlibrefphilosophiqueHVCC}
\newcounter{sujetlibrefphilosophiqueHLCC}
\newcounter{sujetlibrefphilosophiqueOdH}
%%%sujetlibs mixtes semestres%%%
\newcounter{sujetlibrefphilosophiqueETACC}
\newcounter{sujetlibrefphilosophiqueETAHV}
\newcounter{sujetlibrefphilosophiqueETAHL}
\newcounter{sujetlibrefphilosophiqueEXSCC}
\newcounter{sujetlibrefphilosophiqueEXSHV}
\newcounter{sujetlibrefphilosophiqueEXSHL}
\newcounter{sujetlibrefphilosophiqueMMCC}
\newcounter{sujetlibrefphilosophiqueMMHV}
\newcounter{sujetlibrefphilosophiqueMMHL}
%%%sujetlibs mixtes années%%
\newcounter{sujetlibrefphilosophiquepremterm}
\newcounter{sujetlibrefphilosophiquepremtermRS}
\newcounter{sujetlibrefphilosophiquepremtermHQ}
\newcounter{sujetlibrefphilosophiqueETAPdP}
\newcounter{sujetlibrefphilosophiqueETADFI}
\newcounter{sujetlibrefphilosophiqueETADMRC}
\newcounter{sujetlibrefphilosophiqueETAHA}\newcounter{sujetlibrefphilosophiquepremtermETA}
\newcounter{sujetlibrefphilosophiqueEXSPdP}
\newcounter{sujetlibrefphilosophiqueEXSDFI}
\newcounter{sujetlibrefphilosophiqueEXSDMRC}
\newcounter{sujetlibrefphilosophiqueEXSHA}
\newcounter{sujetlibrefphilosophiquepremtermEXS}
\newcounter{sujetlibrefphilosophiqueMMPdP}
\newcounter{sujetlibrefphilosophiqueMMDFI}
\newcounter{sujetlibrefphilosophiqueMMDMRC}
\newcounter{sujetlibrefphilosophiqueMMHA}
\newcounter{sujetlibrefphilosophiquepremtermMM}
\newcounter{sujetlibrefphilosophiqueCCPdP}
\newcounter{sujetlibrefphilosophiqueCCDFI}
\newcounter{sujetlibrefphilosophiqueCCDMRC}
\newcounter{sujetlibrefphilosophiqueCCHA}
\newcounter{sujetlibrefphilosophiquepremtermCC}
\newcounter{sujetlibrefphilosophiqueHVPdP}
\newcounter{sujetlibrefphilosophiqueHVDFI}
\newcounter{sujetlibrefphilosophiqueHVDMRC}
\newcounter{sujetlibrefphilosophiqueHVHA}
\newcounter{sujetlibrefphilosophiquepremtermHV}
\newcounter{sujetlibrefphilosophiqueHLPdP}
\newcounter{sujetlibrefphilosophiqueHLDFI}
\newcounter{sujetlibrefphilosophiqueHLDMRC}
\newcounter{sujetlibrefphilosophiqueHLHA}
\newcounter{sujetlibrefphilosophiquepremtermHL}
\newcounter{sujetlibrefphilosophiqueRSPdP}
\newcounter{sujetlibrefphilosophiqueRSRdM}
\newcounter{sujetlibrefphilosophiqueRSDFI}
\newcounter{sujetlibrefphilosophiqueRSDMRC}
\newcounter{sujetlibrefphilosophiqueRSHA}
\newcounter{sujetlibrefphilosophiqueHQPdP}
\newcounter{sujetlibrefphilosophiqueHQRdM}
\newcounter{sujetlibrefphilosophiqueHQDFI}
\newcounter{sujetlibrefphilosophiqueHQDMRC}
\newcounter{sujetlibrefphilosophiqueHQHA}
%%%Somme des compteurs%%%
\newcounter{sujetliballrefphilosophiqueETA}
\newcounter{sujetliballrefphilosophiqueEXS}
\newcounter{sujetliballrefphilosophiqueMM}
\newcounter{sujetliballrefphilosophiqueRS}
\newcounter{sujetliballrefphilosophiqueCC}
\newcounter{sujetliballrefphilosophiqueHV}
\newcounter{sujetliballrefphilosophiqueHL}
\newcounter{sujetliballrefphilosophiqueHQ}

%%%%%%%%%Réflexion littéraire%%%%%%%%%%%
\newcounter{nosujetlibreflitteraire}
\newcounter{sujetlibHQreflitteraire}
\newcounter{sujetlibRSreflitteraire}
%%%sujetlibs simples%%%
\newcounter{sujetlibreflitteraireEXS}
\newcounter{sujetlibreflitteraireETA}
\newcounter{sujetlibreflitteraireMM}
\newcounter{sujetlibreflitteraireHV}
\newcounter{sujetlibreflitteraireCC}
\newcounter{sujetlibreflitteraireHL}
%%%sujetlibs mixtes%%%
\newcounter{sujetlibreflitteraireEXSTA}
\newcounter{sujetlibreflitteraireEXSMM}
\newcounter{sujetlibreflitteraireETAMM}
\newcounter{sujetlibreflitteraireOdS}
\newcounter{sujetlibreflitteraireHVL}
\newcounter{sujetlibreflitteraireHVCC}
\newcounter{sujetlibreflitteraireHLCC}
\newcounter{sujetlibreflitteraireOdH}
%%%sujetlibs mixtes semestres%%%
\newcounter{sujetlibreflitteraireETACC}
\newcounter{sujetlibreflitteraireETAHV}
\newcounter{sujetlibreflitteraireETAHL}
\newcounter{sujetlibreflitteraireEXSCC}
\newcounter{sujetlibreflitteraireEXSHV}
\newcounter{sujetlibreflitteraireEXSHL}
\newcounter{sujetlibreflitteraireMMCC}
\newcounter{sujetlibreflitteraireMMHV}
\newcounter{sujetlibreflitteraireMMHL}
%%%sujetlibs mixtes années%%
\newcounter{sujetlibreflitterairepremterm}
\newcounter{sujetlibreflitterairepremtermRS}
\newcounter{sujetlibreflitterairepremtermHQ}
\newcounter{sujetlibreflitteraireETAPdP}
\newcounter{sujetlibreflitteraireETADFI}
\newcounter{sujetlibreflitteraireETADMRC}
\newcounter{sujetlibreflitteraireETAHA}
\newcounter{sujetlibreflitterairepremtermETA}
\newcounter{sujetlibreflitteraireEXSPdP}
\newcounter{sujetlibreflitteraireEXSDFI}
\newcounter{sujetlibreflitteraireEXSDMRC}
\newcounter{sujetlibreflitteraireEXSHA}
\newcounter{sujetlibreflitterairepremtermEXS}
\newcounter{sujetlibreflitteraireMMPdP}
\newcounter{sujetlibreflitteraireMMDFI}
\newcounter{sujetlibreflitteraireMMDMRC}
\newcounter{sujetlibreflitteraireMMHA}
\newcounter{sujetlibreflitterairepremtermMM}
\newcounter{sujetlibreflitteraireCCPdP}
\newcounter{sujetlibreflitteraireCCDFI}
\newcounter{sujetlibreflitteraireCCDMRC}
\newcounter{sujetlibreflitteraireCCHA}
\newcounter{sujetlibreflitterairepremtermCC}
\newcounter{sujetlibreflitteraireHVPdP}
\newcounter{sujetlibreflitteraireHVDFI}
\newcounter{sujetlibreflitteraireHVDMRC}
\newcounter{sujetlibreflitteraireHVHA}
\newcounter{sujetlibreflitterairepremtermHV}
\newcounter{sujetlibreflitteraireHLPdP}
\newcounter{sujetlibreflitteraireHLDFI}
\newcounter{sujetlibreflitteraireHLDMRC}
\newcounter{sujetlibreflitteraireHLHA}
\newcounter{sujetlibreflitterairepremtermHL}
\newcounter{sujetlibreflitteraireRSPdP}
\newcounter{sujetlibreflitteraireRSRdM}
\newcounter{sujetlibreflitteraireRSDFI}
\newcounter{sujetlibreflitteraireRSDMRC}
\newcounter{sujetlibreflitteraireRSHA}
\newcounter{sujetlibreflitteraireHQPdP}
\newcounter{sujetlibreflitteraireHQRdM}
\newcounter{sujetlibreflitteraireHQDFI}
\newcounter{sujetlibreflitteraireHQDMRC}
\newcounter{sujetlibreflitteraireHQHA}
%%%Somme des compteurs%%%
\newcounter{sujetliballreflitteraireETA}
\newcounter{sujetliballreflitteraireEXS}
\newcounter{sujetliballreflitteraireMM}
\newcounter{sujetliballreflitteraireCC}
\newcounter{sujetliballreflitteraireRS}
\newcounter{sujetliballreflitteraireHV}
\newcounter{sujetliballreflitteraireHL}
\newcounter{sujetliballreflitteraireHQ}
%%%Compteurs de statistiques globales sur les genres de question%%%
\newcounter{sujetlibexclRS}
\newcounter{sujetlibmixtRS}
\newcounter{sujetlibexclHQ}
\newcounter{sujetlibmixtHQ}
\newcounter{sujetlibmixtHQRS}
%%%Compteurs de statistiques sur les transversaux (divers semestres, première et terminale)
\newcounter{sujettranslibintlitteraire}
\newcounter{sujettranslibreflitteraire}
\newcounter{sujettranslibintphilosophique}
\newcounter{sujettranslibrefphilosophique}
%%% Compteurs du nombre de sujets par type
\newcounter{sujetlibintlitteraire}
\newcounter{sujetlibreflitteraire}
\newcounter{sujetlibintphilosophique}
\newcounter{sujetlibrefphilosophique}
%%% Compteurs par discipline
%%Litt
\newcounter{sujetliblitteraireETA}
\newcounter{sujetliblitteraireEXS}
\newcounter{sujetliblitteraireMM}
\newcounter{sujetliblitteraireCC}
\newcounter{sujetliblitteraireHV}
\newcounter{sujetliblitteraireHL}
\newcounter{sujetliblitteraireEXSTA}
\newcounter{sujetliblitteraireETAMM}
\newcounter{sujetliblitteraireEXSMM}
\newcounter{sujetliblitteraireHVCC}
\newcounter{sujetliblitteraireHVL}
\newcounter{sujetliblitteraireaut}
%%Phil
\newcounter{sujetlibphilosophiqueETA}
\newcounter{sujetlibphilosophiqueEXS}
\newcounter{sujetlibphilosophiqueMM}
\newcounter{sujetlibphilosophiqueCC}
\newcounter{sujetlibphilosophiqueHV}
\newcounter{sujetlibphilosophiqueHL}
\newcounter{sujetlibphilosophiqueEXSTA}
\newcounter{sujetlibphilosophiqueETAMM}
\newcounter{sujetlibphilosophiqueEXSMM}
\newcounter{sujetlibphilosophiqueHVCC}
\newcounter{sujetlibphilosophiqueHVL}
\newcounter{sujetlibphilosophiqueaut}
%%Totaux pour les graphiques
\newcounter{sujetliblitterairetotal}
\newcounter{sujetlibphilosophiquetotal}
%%% Autres (Ods et Odh) comptés ensemble
\newcounter{sujetliballaut}
%%%%%%%%%%%%%%%%%%%%%%%%%%%%%%%%%%%%%%%%%%%%%%%%%%%%%%%%%%%%%%%%%%%%%%%
%%%%%%%%%%%%%%%%%%%%%%%%%%%%%%Compteurs%%%%%%%%%%%%%%%%%%%%%%%%%%%%%%%%
%%%%%%%%%%%%%%%%%%%%%%%%%%%%%%%%%%%%%%%%%%%%%%%%%%%%%%%%%%%%%%%%%%%%%%%
%%%Compter les sujets par thème%%%
%\newcounter{nosujetunk}
\newcounter{sujetunkall}%Total de questions pour le centre
\newcounter{sujetunktotal}%Total des questions en comptant en double les questions pour les graphiques
\newcounter{sujetunkHQ}
\newcounter{sujetunkRS}
%%%sujetunks simples%%%
\newcounter{sujetunkEXS}
\newcounter{sujetunkETA}
\newcounter{sujetunkMM}
\newcounter{sujetunkHV}
\newcounter{sujetunkCC}
\newcounter{sujetunkHL}
%%%sujetunks mixtes%%%
\newcounter{sujetunkEXSTA}
\newcounter{sujetunkEXSMM}
\newcounter{sujetunkETAMM}
\newcounter{sujetunkOdS}
\newcounter{sujetunkHVL}
\newcounter{sujetunkHVCC}
\newcounter{sujetunkHLCC}
\newcounter{sujetunkOdH}
%%%sujetunks mixtes semestres%%%
\newcounter{sujetunkETACC}
\newcounter{sujetunkETAHV}
\newcounter{sujetunkETAHL}
\newcounter{sujetunkEXSCC}
\newcounter{sujetunkEXSHV}
\newcounter{sujetunkEXSHL}
\newcounter{sujetunkMMCC}
\newcounter{sujetunkMMHV}
\newcounter{sujetunkMMHL}

%%%sujetunks mixtes années%%
\newcounter{sujetunkpremterm}%%%Compteur général de sujetunk
\newcounter{sujetunkETAPdP}
\newcounter{sujetunkETADFI}
\newcounter{sujetunkETADMRC}
\newcounter{sujetunkETAHA}
\newcounter{sujetunkpremtermETA}
\newcounter{sujetunkEXSPdP}
\newcounter{sujetunkEXSDFI}
\newcounter{sujetunkEXSDMRC}
\newcounter{sujetunkEXSHA}
\newcounter{sujetunkpremtermEXS}
\newcounter{sujetunkMMPdP}
\newcounter{sujetunkMMDFI}
\newcounter{sujetunkMMDMRC}
\newcounter{sujetunkMMHA}
\newcounter{sujetunkpremtermMM}
\newcounter{sujetunkCCPdP}
\newcounter{sujetunkCCDFI}
\newcounter{sujetunkCCDMRC}
\newcounter{sujetunkCCHA}
\newcounter{sujetunkpremtermCC}
\newcounter{sujetunkHVPdP}
\newcounter{sujetunkHVDFI}
\newcounter{sujetunkHVDMRC}
\newcounter{sujetunkHVHA}
\newcounter{sujetunkpremtermHV}
\newcounter{sujetunkHLPdP}
\newcounter{sujetunkHLDFI}
\newcounter{sujetunkHLDMRC}
\newcounter{sujetunkHLHA}
\newcounter{sujetunkpremtermHL}
\newcounter{sujetunkRSPdP}
\newcounter{sujetunkRSRdM}
\newcounter{sujetunkRSDFI}
\newcounter{sujetunkRSDMRC}
\newcounter{sujetunkRSHA}
\newcounter{sujetunkHQPdP}
\newcounter{sujetunkHQRdM}
\newcounter{sujetunkHQDFI}
\newcounter{sujetunkHQDMRC}
\newcounter{sujetunkHQHA}

%%%Somme des compteurs%%%
\newcounter{sujetunkallETA}
\newcounter{sujetunkallEXS}
\newcounter{sujetunkallMM}
\newcounter{sujetunkallRS}
\newcounter{sujetunkallCC}
\newcounter{sujetunkallHV}
\newcounter{sujetunkallHL}
\newcounter{sujetunkallHQ}

%%%Compter les sujetunks par thème et types de questions%%%
%%%%%%%%%%%%%%%%%%%%%%%%%%%%%%%%%%%%%%%%%%%%%
%%%%%%%%%Question d'interprétation%%%%%%%%%%%
%%%%%%%%%%%%%%%%%%%%%%%%%%%%%%%%%%%%%%%%%%%%%

%%%%%%%%%%%%%%%Général%%%%%%%%%%%%%%%%%%%%%%%
\newcounter{nosujetunkint}
\newcounter{sujetunkHQint}
\newcounter{sujetunkRSint}
%%%sujetunks simples%%%
\newcounter{sujetunkintEXS}
\newcounter{sujetunkintETA}
\newcounter{sujetunkintMM}
\newcounter{sujetunkintHV}
\newcounter{sujetunkintCC}
\newcounter{sujetunkintHL}
%%%sujetunks mixtes%%%
\newcounter{sujetunkintEXSTA}
\newcounter{sujetunkintEXSMM}
\newcounter{sujetunkintETAMM}
\newcounter{sujetunkintOdS}
\newcounter{sujetunkintHVL}
\newcounter{sujetunkintHVCC}
\newcounter{sujetunkintHLCC}
\newcounter{sujetunkintOdH}
%%%sujetunks mixtes semestres%%%
\newcounter{sujetunkintETACC}
\newcounter{sujetunkintETAHV}
\newcounter{sujetunkintETAHL}
\newcounter{sujetunkintEXSCC}
\newcounter{sujetunkintEXSHV}
\newcounter{sujetunkintEXSHL}
\newcounter{sujetunkintMMCC}
\newcounter{sujetunkintMMHV}
\newcounter{sujetunkintMMHL}
%%%sujetunks mixtes années%%
\newcounter{sujetunkintpremterm}
\newcounter{sujetunkintpremtermRS}
\newcounter{sujetunkintpremtermHQ}
\newcounter{sujetunkintETAPdP}
\newcounter{sujetunkintETADFI}
\newcounter{sujetunkintETADMRC}
\newcounter{sujetunkintETAHA}
\newcounter{sujetunkintpremtermETA}
\newcounter{sujetunkintEXSPdP}
\newcounter{sujetunkintEXSDFI}
\newcounter{sujetunkintEXSDMRC}
\newcounter{sujetunkintEXSHA}
\newcounter{sujetunkintpremtermEXS}
\newcounter{sujetunkintMMPdP}
\newcounter{sujetunkintMMDFI}
\newcounter{sujetunkintMMDMRC}
\newcounter{sujetunkintMMHA}
\newcounter{sujetunkintpremtermMM}
\newcounter{sujetunkintCCPdP}
\newcounter{sujetunkintCCDFI}
\newcounter{sujetunkintCCDMRC}
\newcounter{sujetunkintCCHA}
\newcounter{sujetunkintpremtermCC}
\newcounter{sujetunkintHVPdP}
\newcounter{sujetunkintHVDFI}
\newcounter{sujetunkintHVDMRC}
\newcounter{sujetunkintHVHA}
\newcounter{sujetunkintpremtermHV}
\newcounter{sujetunkintHLPdP}
\newcounter{sujetunkintHLDFI}
\newcounter{sujetunkintHLDMRC}
\newcounter{sujetunkintHLHA}
\newcounter{sujetunkintpremtermHL}
\newcounter{sujetunkintRSPdP}
\newcounter{sujetunkintRSRdM}
\newcounter{sujetunkintRSDFI}
\newcounter{sujetunkintRSDMRC}
\newcounter{sujetunkintRSHA}
\newcounter{sujetunkintHQPdP}
\newcounter{sujetunkintHQRdM}
\newcounter{sujetunkintHQDFI}
\newcounter{sujetunkintHQDMRC}
\newcounter{sujetunkintHQHA}
%%%Somme des compteurs%%%
\newcounter{sujetunkallintETA}
\newcounter{sujetunkallintEXS}
\newcounter{sujetunkallintMM}
\newcounter{sujetunkallintRS}
\newcounter{sujetunkallintCC}
\newcounter{sujetunkallintHV}
\newcounter{sujetunkallintHL}
\newcounter{sujetunkallintHQ}

%%%%%%%%Interprétation philosophique%%%%%%%%%
\newcounter{nosujetunkintphilosophique}
\newcounter{sujetunkHQintphilosophique}
\newcounter{sujetunkRSintphilosophique}
%%%sujetunks simples%%%
\newcounter{sujetunkintphilosophiqueEXS}
\newcounter{sujetunkintphilosophiqueETA}
\newcounter{sujetunkintphilosophiqueMM}
\newcounter{sujetunkintphilosophiqueHV}
\newcounter{sujetunkintphilosophiqueCC}
\newcounter{sujetunkintphilosophiqueHL}
%%%sujetunks mixtes%%%
\newcounter{sujetunkintphilosophiqueEXSTA}
\newcounter{sujetunkintphilosophiqueEXSMM}
\newcounter{sujetunkintphilosophiqueETAMM}
\newcounter{sujetunkintphilosophiqueOdS}
\newcounter{sujetunkintphilosophiqueHVL}
\newcounter{sujetunkintphilosophiqueHVCC}
\newcounter{sujetunkintphilosophiqueHLCC}
\newcounter{sujetunkintphilosophiqueOdH}
%%%sujetunks mixtes semestres%%%
\newcounter{sujetunkintphilosophiqueETACC}
\newcounter{sujetunkintphilosophiqueETAHV}
\newcounter{sujetunkintphilosophiqueETAHL}
\newcounter{sujetunkintphilosophiqueEXSCC}
\newcounter{sujetunkintphilosophiqueEXSHV}
\newcounter{sujetunkintphilosophiqueEXSHL}
\newcounter{sujetunkintphilosophiqueMMCC}
\newcounter{sujetunkintphilosophiqueMMHV}
\newcounter{sujetunkintphilosophiqueMMHL}
%%%sujetunks mixtes années%%
\newcounter{sujetunkintphilosophiquepremterm}
\newcounter{sujetunkintphilosophiquepremtermRS}
\newcounter{sujetunkintphilosophiquepremtermHQ}
\newcounter{sujetunkintphilosophiqueETAPdP}
\newcounter{sujetunkintphilosophiqueETADFI}
\newcounter{sujetunkintphilosophiqueETADMRC}
\newcounter{sujetunkintphilosophiqueETAHA}
\newcounter{sujetunkintphilosophiquepremtermETA}
\newcounter{sujetunkintphilosophiqueEXSPdP}
\newcounter{sujetunkintphilosophiqueEXSDFI}
\newcounter{sujetunkintphilosophiqueEXSDMRC}
\newcounter{sujetunkintphilosophiqueEXSHA}
\newcounter{sujetunkintphilosophiquepremtermEXS}
\newcounter{sujetunkintphilosophiqueMMPdP}
\newcounter{sujetunkintphilosophiqueMMDFI}
\newcounter{sujetunkintphilosophiqueMMDMRC}
\newcounter{sujetunkintphilosophiqueMMHA}
\newcounter{sujetunkintphilosophiquepremtermMM}
\newcounter{sujetunkintphilosophiqueCCPdP}
\newcounter{sujetunkintphilosophiqueCCDFI}
\newcounter{sujetunkintphilosophiqueCCDMRC}
\newcounter{sujetunkintphilosophiqueCCHA}
\newcounter{sujetunkintphilosophiquepremtermCC}
\newcounter{sujetunkintphilosophiqueHVPdP}
\newcounter{sujetunkintphilosophiqueHVDFI}
\newcounter{sujetunkintphilosophiqueHVDMRC}
\newcounter{sujetunkintphilosophiqueHVHA}
\newcounter{sujetunkintphilosophiquepremtermHV}
\newcounter{sujetunkintphilosophiqueHLPdP}
\newcounter{sujetunkintphilosophiqueHLDFI}
\newcounter{sujetunkintphilosophiqueHLDMRC}
\newcounter{sujetunkintphilosophiqueHLHA}
\newcounter{sujetunkintphilosophiquepremtermHL}
\newcounter{sujetunkintphilosophiqueRSPdP}
\newcounter{sujetunkintphilosophiqueRSRdM}
\newcounter{sujetunkintphilosophiqueRSDFI}
\newcounter{sujetunkintphilosophiqueRSDMRC}
\newcounter{sujetunkintphilosophiqueRSHA}
\newcounter{sujetunkintphilosophiqueHQPdP}
\newcounter{sujetunkintphilosophiqueHQRdM}
\newcounter{sujetunkintphilosophiqueHQDFI}
\newcounter{sujetunkintphilosophiqueHQDMRC}
\newcounter{sujetunkintphilosophiqueHQHA}
%%%Somme des compteurs%%%
\newcounter{sujetunkallintphilosophiqueETA}
\newcounter{sujetunkallintphilosophiqueEXS}
\newcounter{sujetunkallintphilosophiqueMM}
\newcounter{sujetunkallintphilosophiqueRS}
\newcounter{sujetunkallintphilosophiqueCC}
\newcounter{sujetunkallintphilosophiqueHV}
\newcounter{sujetunkallintphilosophiqueHL}
\newcounter{sujetunkallintphilosophiqueHQ}

%%%%%%%%%Interprétation littéraire%%%%%%%%%%%
\newcounter{nosujetunkintlitteraire}
\newcounter{sujetunkHQintlitteraire}
\newcounter{sujetunkRSintlitteraire}
%%%sujetunks simples%%%
\newcounter{sujetunkintlitteraireEXS}
\newcounter{sujetunkintlitteraireETA}
\newcounter{sujetunkintlitteraireMM}
\newcounter{sujetunkintlitteraireHV}
\newcounter{sujetunkintlitteraireCC}
\newcounter{sujetunkintlitteraireHL}
%%%sujetunks mixtes%%%
\newcounter{sujetunkintlitteraireEXSTA}
\newcounter{sujetunkintlitteraireEXSMM}
\newcounter{sujetunkintlitteraireETAMM}
\newcounter{sujetunkintlitteraireOdS}
\newcounter{sujetunkintlitteraireHVL}
\newcounter{sujetunkintlitteraireHVCC}
\newcounter{sujetunkintlitteraireHLCC}
\newcounter{sujetunkintlitteraireOdH}
%%%sujetunks mixtes semestres%%%
\newcounter{sujetunkintlitteraireETACC}
\newcounter{sujetunkintlitteraireETAHV}
\newcounter{sujetunkintlitteraireETAHL}
\newcounter{sujetunkintlitteraireEXSCC}
\newcounter{sujetunkintlitteraireEXSHV}
\newcounter{sujetunkintlitteraireEXSHL}
\newcounter{sujetunkintlitteraireMMCC}
\newcounter{sujetunkintlitteraireMMHV}
\newcounter{sujetunkintlitteraireMMHL}
%%%sujetunks mixtes années%%
\newcounter{sujetunkintlitterairepremterm}
\newcounter{sujetunkintlitterairepremtermRS}
\newcounter{sujetunkintlitterairepremtermHQ}
\newcounter{sujetunkintlitteraireETAPdP}
\newcounter{sujetunkintlitteraireETADFI}
\newcounter{sujetunkintlitteraireETADMRC}
\newcounter{sujetunkintlitteraireETAHA}
\newcounter{sujetunkintlitterairepremtermETA}
\newcounter{sujetunkintlitteraireEXSPdP}
\newcounter{sujetunkintlitteraireEXSDFI}
\newcounter{sujetunkintlitteraireEXSDMRC}
\newcounter{sujetunkintlitteraireEXSHA}
\newcounter{sujetunkintlitterairepremtermEXS}
\newcounter{sujetunkintlitteraireMMPdP}
\newcounter{sujetunkintlitteraireMMDFI}
\newcounter{sujetunkintlitteraireMMDMRC}
\newcounter{sujetunkintlitteraireMMHA}
\newcounter{sujetunkintlitterairepremtermMM}
\newcounter{sujetunkintlitteraireCCPdP}
\newcounter{sujetunkintlitteraireCCDFI}
\newcounter{sujetunkintlitteraireCCDMRC}
\newcounter{sujetunkintlitteraireCCHA}
\newcounter{sujetunkintlitterairepremtermCC}
\newcounter{sujetunkintlitteraireHVPdP}
\newcounter{sujetunkintlitteraireHVDFI}
\newcounter{sujetunkintlitteraireHVDMRC}
\newcounter{sujetunkintlitteraireHVHA}
\newcounter{sujetunkintlitterairepremtermHV}
\newcounter{sujetunkintlitteraireHLPdP}
\newcounter{sujetunkintlitteraireHLDFI}
\newcounter{sujetunkintlitteraireHLDMRC}
\newcounter{sujetunkintlitteraireHLHA}
\newcounter{sujetunkintlitterairepremtermHL}
\newcounter{sujetunkintlitteraireRSPdP}
\newcounter{sujetunkintlitteraireRSRdM}
\newcounter{sujetunkintlitteraireRSDFI}
\newcounter{sujetunkintlitteraireRSDMRC}
\newcounter{sujetunkintlitteraireRSHA}
\newcounter{sujetunkintlitteraireHQPdP}
\newcounter{sujetunkintlitteraireHQRdM}
\newcounter{sujetunkintlitteraireHQDFI}
\newcounter{sujetunkintlitteraireHQDMRC}
\newcounter{sujetunkintlitteraireHQHA}
%%%Somme des compteurs%%%
\newcounter{sujetunkallintlitteraireETA}
\newcounter{sujetunkallintlitteraireEXS}
\newcounter{sujetunkallintlitteraireMM}
\newcounter{sujetunkallintlitteraireRS}
\newcounter{sujetunkallintlitteraireCC}
\newcounter{sujetunkallintlitteraireHV}
\newcounter{sujetunkallintlitteraireHL}
\newcounter{sujetunkallintlitteraireHQ}

%%%%%%%%%%%%%%%%%%%%%%%%%%%%%%%%%%%%%%%%%%%%%
%%%%%%%%%%%%Question de réflexion%%%%%%%%%%%%
%%%%%%%%%%%%%%%%%%%%%%%%%%%%%%%%%%%%%%%%%%%%%
\newcounter{nosujetunkref}
%%%sujetunks simples%%%
\newcounter{sujetunkrefEXS}
\newcounter{sujetunkrefETA}
\newcounter{sujetunkrefMM}
\newcounter{sujetunkrefHV}
\newcounter{sujetunkrefCC}
\newcounter{sujetunkrefHL}
%%%sujetunks mixtes%%%
\newcounter{sujetunkrefEXSTA}
\newcounter{sujetunkrefEXSMM}
\newcounter{sujetunkrefETAMM}
\newcounter{sujetunkrefOdS}
\newcounter{sujetunkrefHVL}
\newcounter{sujetunkrefHVCC}
\newcounter{sujetunkrefHLCC}
\newcounter{sujetunkrefOdH}
%%%sujetunks mixtes semestres%%%
\newcounter{sujetunkrefETACC}
\newcounter{sujetunkrefETAHV}
\newcounter{sujetunkrefETAHL}
\newcounter{sujetunkrefEXSCC}
\newcounter{sujetunkrefEXSHV}
\newcounter{sujetunkrefEXSHL}
\newcounter{sujetunkrefMMCC}
\newcounter{sujetunkrefMMHV}
\newcounter{sujetunkrefMMHL}
%%%sujetunks mixtes années%%
\newcounter{sujetunkrefpremterm}%%%Compteur général
\newcounter{sujetunkrefETAPdP}
\newcounter{sujetunkrefETADFI}
\newcounter{sujetunkrefETADMRC}
\newcounter{sujetunkrefETAHA}
\newcounter{sujetunkrefpremtermETA}
\newcounter{sujetunkrefEXSPdP}
\newcounter{sujetunkrefEXSDFI}
\newcounter{sujetunkrefEXSDMRC}
\newcounter{sujetunkrefEXSHA}
\newcounter{sujetunkrefpremtermEXS}
\newcounter{sujetunkrefMMPdP}
\newcounter{sujetunkrefMMDFI}
\newcounter{sujetunkrefMMDMRC}
\newcounter{sujetunkrefMMHA}
\newcounter{sujetunkrefpremtermMM}
\newcounter{sujetunkrefCCPdP}
\newcounter{sujetunkrefCCDFI}
\newcounter{sujetunkrefCCDMRC}
\newcounter{sujetunkrefCCHA}
\newcounter{sujetunkrefpremtermCC}
\newcounter{sujetunkrefHVPdP}
\newcounter{sujetunkrefHVDFI}
\newcounter{sujetunkrefHVDMRC}
\newcounter{sujetunkrefHVHA}
\newcounter{sujetunkrefpremtermHV}
\newcounter{sujetunkrefHLPdP}
\newcounter{sujetunkrefHLDFI}
\newcounter{sujetunkrefHLDMRC}
\newcounter{sujetunkrefHLHA}
\newcounter{sujetunkrefpremtermHL}
\newcounter{sujetunkrefRSPdP}
\newcounter{sujetunkrefRSRdM}
\newcounter{sujetunkrefRSDFI}
\newcounter{sujetunkrefRSDMRC}
\newcounter{sujetunkrefRSHA}
\newcounter{sujetunkrefHQPdP}
\newcounter{sujetunkrefHQRdM}
\newcounter{sujetunkrefHQDFI}
\newcounter{sujetunkrefHQDMRC}
\newcounter{sujetunkrefHQHA}
%%%Somme des compteurs%%%
\newcounter{sujetunkallrefETA}
\newcounter{sujetunkallrefEXS}
\newcounter{sujetunkallrefMM}
\newcounter{sujetunkallrefRS}
\newcounter{sujetunkallrefCC}
\newcounter{sujetunkallrefHV}
\newcounter{sujetunkallrefHL}
\newcounter{sujetunkallrefHQ}

%%%%%%%%Réflexion philosophique%%%%%%%%%
\newcounter{nosujetunkrefphilosophique}
\newcounter{sujetunkHQrefphilosophique}
\newcounter{sujetunkRSrefphilosophique}
%%%sujetunks simples%%%
\newcounter{sujetunkrefphilosophiqueEXS}
\newcounter{sujetunkrefphilosophiqueETA}
\newcounter{sujetunkrefphilosophiqueMM}
\newcounter{sujetunkrefphilosophiqueHV}
\newcounter{sujetunkrefphilosophiqueCC}
\newcounter{sujetunkrefphilosophiqueHL}
%%%sujetunks mixtes%%%
\newcounter{sujetunkrefphilosophiqueEXSTA}
\newcounter{sujetunkrefphilosophiqueEXSMM}
\newcounter{sujetunkrefphilosophiqueETAMM}
\newcounter{sujetunkrefphilosophiqueOdS}
\newcounter{sujetunkrefphilosophiqueHVL}
\newcounter{sujetunkrefphilosophiqueHVCC}
\newcounter{sujetunkrefphilosophiqueHLCC}
\newcounter{sujetunkrefphilosophiqueOdH}
%%%sujetunks mixtes semestres%%%
\newcounter{sujetunkrefphilosophiqueETACC}
\newcounter{sujetunkrefphilosophiqueETAHV}
\newcounter{sujetunkrefphilosophiqueETAHL}
\newcounter{sujetunkrefphilosophiqueEXSCC}
\newcounter{sujetunkrefphilosophiqueEXSHV}
\newcounter{sujetunkrefphilosophiqueEXSHL}
\newcounter{sujetunkrefphilosophiqueMMCC}
\newcounter{sujetunkrefphilosophiqueMMHV}
\newcounter{sujetunkrefphilosophiqueMMHL}
%%%sujetunks mixtes années%%
\newcounter{sujetunkrefphilosophiquepremterm}
\newcounter{sujetunkrefphilosophiquepremtermRS}
\newcounter{sujetunkrefphilosophiquepremtermHQ}
\newcounter{sujetunkrefphilosophiqueETAPdP}
\newcounter{sujetunkrefphilosophiqueETADFI}
\newcounter{sujetunkrefphilosophiqueETADMRC}
\newcounter{sujetunkrefphilosophiqueETAHA}\newcounter{sujetunkrefphilosophiquepremtermETA}
\newcounter{sujetunkrefphilosophiqueEXSPdP}
\newcounter{sujetunkrefphilosophiqueEXSDFI}
\newcounter{sujetunkrefphilosophiqueEXSDMRC}
\newcounter{sujetunkrefphilosophiqueEXSHA}
\newcounter{sujetunkrefphilosophiquepremtermEXS}
\newcounter{sujetunkrefphilosophiqueMMPdP}
\newcounter{sujetunkrefphilosophiqueMMDFI}
\newcounter{sujetunkrefphilosophiqueMMDMRC}
\newcounter{sujetunkrefphilosophiqueMMHA}
\newcounter{sujetunkrefphilosophiquepremtermMM}
\newcounter{sujetunkrefphilosophiqueCCPdP}
\newcounter{sujetunkrefphilosophiqueCCDFI}
\newcounter{sujetunkrefphilosophiqueCCDMRC}
\newcounter{sujetunkrefphilosophiqueCCHA}
\newcounter{sujetunkrefphilosophiquepremtermCC}
\newcounter{sujetunkrefphilosophiqueHVPdP}
\newcounter{sujetunkrefphilosophiqueHVDFI}
\newcounter{sujetunkrefphilosophiqueHVDMRC}
\newcounter{sujetunkrefphilosophiqueHVHA}
\newcounter{sujetunkrefphilosophiquepremtermHV}
\newcounter{sujetunkrefphilosophiqueHLPdP}
\newcounter{sujetunkrefphilosophiqueHLDFI}
\newcounter{sujetunkrefphilosophiqueHLDMRC}
\newcounter{sujetunkrefphilosophiqueHLHA}
\newcounter{sujetunkrefphilosophiquepremtermHL}
\newcounter{sujetunkrefphilosophiqueRSPdP}
\newcounter{sujetunkrefphilosophiqueRSRdM}
\newcounter{sujetunkrefphilosophiqueRSDFI}
\newcounter{sujetunkrefphilosophiqueRSDMRC}
\newcounter{sujetunkrefphilosophiqueRSHA}
\newcounter{sujetunkrefphilosophiqueHQPdP}
\newcounter{sujetunkrefphilosophiqueHQRdM}
\newcounter{sujetunkrefphilosophiqueHQDFI}
\newcounter{sujetunkrefphilosophiqueHQDMRC}
\newcounter{sujetunkrefphilosophiqueHQHA}
%%%Somme des compteurs%%%
\newcounter{sujetunkallrefphilosophiqueETA}
\newcounter{sujetunkallrefphilosophiqueEXS}
\newcounter{sujetunkallrefphilosophiqueMM}
\newcounter{sujetunkallrefphilosophiqueRS}
\newcounter{sujetunkallrefphilosophiqueCC}
\newcounter{sujetunkallrefphilosophiqueHV}
\newcounter{sujetunkallrefphilosophiqueHL}
\newcounter{sujetunkallrefphilosophiqueHQ}

%%%%%%%%%Réflexion littéraire%%%%%%%%%%%
\newcounter{nosujetunkreflitteraire}
\newcounter{sujetunkHQreflitteraire}
\newcounter{sujetunkRSreflitteraire}
%%%sujetunks simples%%%
\newcounter{sujetunkreflitteraireEXS}
\newcounter{sujetunkreflitteraireETA}
\newcounter{sujetunkreflitteraireMM}
\newcounter{sujetunkreflitteraireHV}
\newcounter{sujetunkreflitteraireCC}
\newcounter{sujetunkreflitteraireHL}
%%%sujetunks mixtes%%%
\newcounter{sujetunkreflitteraireEXSTA}
\newcounter{sujetunkreflitteraireEXSMM}
\newcounter{sujetunkreflitteraireETAMM}
\newcounter{sujetunkreflitteraireOdS}
\newcounter{sujetunkreflitteraireHVL}
\newcounter{sujetunkreflitteraireHVCC}
\newcounter{sujetunkreflitteraireHLCC}
\newcounter{sujetunkreflitteraireOdH}
%%%sujetunks mixtes semestres%%%
\newcounter{sujetunkreflitteraireETACC}
\newcounter{sujetunkreflitteraireETAHV}
\newcounter{sujetunkreflitteraireETAHL}
\newcounter{sujetunkreflitteraireEXSCC}
\newcounter{sujetunkreflitteraireEXSHV}
\newcounter{sujetunkreflitteraireEXSHL}
\newcounter{sujetunkreflitteraireMMCC}
\newcounter{sujetunkreflitteraireMMHV}
\newcounter{sujetunkreflitteraireMMHL}
%%%sujetunks mixtes années%%
\newcounter{sujetunkreflitterairepremterm}
\newcounter{sujetunkreflitterairepremtermRS}
\newcounter{sujetunkreflitterairepremtermHQ}
\newcounter{sujetunkreflitteraireETAPdP}
\newcounter{sujetunkreflitteraireETADFI}
\newcounter{sujetunkreflitteraireETADMRC}
\newcounter{sujetunkreflitteraireETAHA}
\newcounter{sujetunkreflitterairepremtermETA}
\newcounter{sujetunkreflitteraireEXSPdP}
\newcounter{sujetunkreflitteraireEXSDFI}
\newcounter{sujetunkreflitteraireEXSDMRC}
\newcounter{sujetunkreflitteraireEXSHA}
\newcounter{sujetunkreflitterairepremtermEXS}
\newcounter{sujetunkreflitteraireMMPdP}
\newcounter{sujetunkreflitteraireMMDFI}
\newcounter{sujetunkreflitteraireMMDMRC}
\newcounter{sujetunkreflitteraireMMHA}
\newcounter{sujetunkreflitterairepremtermMM}
\newcounter{sujetunkreflitteraireCCPdP}
\newcounter{sujetunkreflitteraireCCDFI}
\newcounter{sujetunkreflitteraireCCDMRC}
\newcounter{sujetunkreflitteraireCCHA}
\newcounter{sujetunkreflitterairepremtermCC}
\newcounter{sujetunkreflitteraireHVPdP}
\newcounter{sujetunkreflitteraireHVDFI}
\newcounter{sujetunkreflitteraireHVDMRC}
\newcounter{sujetunkreflitteraireHVHA}
\newcounter{sujetunkreflitterairepremtermHV}
\newcounter{sujetunkreflitteraireHLPdP}
\newcounter{sujetunkreflitteraireHLDFI}
\newcounter{sujetunkreflitteraireHLDMRC}
\newcounter{sujetunkreflitteraireHLHA}
\newcounter{sujetunkreflitterairepremtermHL}
\newcounter{sujetunkreflitteraireRSPdP}
\newcounter{sujetunkreflitteraireRSRdM}
\newcounter{sujetunkreflitteraireRSDFI}
\newcounter{sujetunkreflitteraireRSDMRC}
\newcounter{sujetunkreflitteraireRSHA}
\newcounter{sujetunkreflitteraireHQPdP}
\newcounter{sujetunkreflitteraireHQRdM}
\newcounter{sujetunkreflitteraireHQDFI}
\newcounter{sujetunkreflitteraireHQDMRC}
\newcounter{sujetunkreflitteraireHQHA}
%%%Somme des compteurs%%%
\newcounter{sujetunkallreflitteraireETA}
\newcounter{sujetunkallreflitteraireEXS}
\newcounter{sujetunkallreflitteraireMM}
\newcounter{sujetunkallreflitteraireCC}
\newcounter{sujetunkallreflitteraireRS}
\newcounter{sujetunkallreflitteraireHV}
\newcounter{sujetunkallreflitteraireHL}
\newcounter{sujetunkallreflitteraireHQ}
%%%Compteurs de statistiques globales sur les genres de question%%%
\newcounter{sujetunkexclRS}
\newcounter{sujetunkmixtRS}
\newcounter{sujetunkexclHQ}
\newcounter{sujetunkmixtHQ}
\newcounter{sujetunkmixtHQRS}
%%%Compteurs de statistiques sur les transversaux (divers semestres, première et terminale)
\newcounter{sujettransunkintlitteraire}
\newcounter{sujettransunkreflitteraire}
\newcounter{sujettransunkintphilosophique}
\newcounter{sujettransunkrefphilosophique}
%%% Compteurs du nombre de sujets par type
\newcounter{sujetunkintlitteraire}
\newcounter{sujetunkreflitteraire}
\newcounter{sujetunkintphilosophique}
\newcounter{sujetunkrefphilosophique}
%%% Compteurs par discipline
%%Litt
\newcounter{sujetunklitteraireETA}
\newcounter{sujetunklitteraireEXS}
\newcounter{sujetunklitteraireMM}
\newcounter{sujetunklitteraireCC}
\newcounter{sujetunklitteraireHV}
\newcounter{sujetunklitteraireHL}
\newcounter{sujetunklitteraireEXSTA}
\newcounter{sujetunklitteraireETAMM}
\newcounter{sujetunklitteraireEXSMM}
\newcounter{sujetunklitteraireHVCC}
\newcounter{sujetunklitteraireHVL}
\newcounter{sujetunklitteraireaut}
%%Phil
\newcounter{sujetunkphilosophiqueETA}
\newcounter{sujetunkphilosophiqueEXS}
\newcounter{sujetunkphilosophiqueMM}
\newcounter{sujetunkphilosophiqueCC}
\newcounter{sujetunkphilosophiqueHV}
\newcounter{sujetunkphilosophiqueHL}
\newcounter{sujetunkphilosophiqueEXSTA}
\newcounter{sujetunkphilosophiqueETAMM}
\newcounter{sujetunkphilosophiqueEXSMM}
\newcounter{sujetunkphilosophiqueHVCC}
\newcounter{sujetunkphilosophiqueHVL}
\newcounter{sujetunkphilosophiqueaut}
%%Totaux pour les graphiques
\newcounter{sujetunklitterairetotal}
\newcounter{sujetunkphilosophiquetotal}
%%% Autres (Ods et Odh) comptés ensemble
\newcounter{sujetunkallaut}
%%%%%%%%%%%%%%%%%%%%%%%%%%%%%%%%%%%%%%%%%%%%%%%%%%%%%%%%%%%%%%%%%%%%%%%
%%%%%%%%%%%%%%%%%%%%%%%%%%%%%%Compteurs%%%%%%%%%%%%%%%%%%%%%%%%%%%%%%%%
%%%%%%%%%%%%%%%%%%%%%%%%%%%%%%%%%%%%%%%%%%%%%%%%%%%%%%%%%%%%%%%%%%%%%%%
%%%Compter les sujets par thème%%%
%\newcounter{nosujetmet}
\newcounter{sujetmetall}%Total de questions pour le centre
\newcounter{sujetmettotal}%Total des questions en comptant en double les questions pour les graphiques
\newcounter{sujetmetHQ}
\newcounter{sujetmetRS}
%%%sujetmets simples%%%
\newcounter{sujetmetEXS}
\newcounter{sujetmetETA}
\newcounter{sujetmetMM}
\newcounter{sujetmetHV}
\newcounter{sujetmetCC}
\newcounter{sujetmetHL}
%%%sujetmets mixtes%%%
\newcounter{sujetmetEXSTA}
\newcounter{sujetmetEXSMM}
\newcounter{sujetmetETAMM}
\newcounter{sujetmetOdS}
\newcounter{sujetmetHVL}
\newcounter{sujetmetHVCC}
\newcounter{sujetmetHLCC}
\newcounter{sujetmetOdH}
%%%sujetmets mixtes semestres%%%
\newcounter{sujetmetETACC}
\newcounter{sujetmetETAHV}
\newcounter{sujetmetETAHL}
\newcounter{sujetmetEXSCC}
\newcounter{sujetmetEXSHV}
\newcounter{sujetmetEXSHL}
\newcounter{sujetmetMMCC}
\newcounter{sujetmetMMHV}
\newcounter{sujetmetMMHL}

%%%sujetmets mixtes années%%
\newcounter{sujetmetpremterm}%%%Compteur général de sujetmet
\newcounter{sujetmetETAPdP}
\newcounter{sujetmetETADFI}
\newcounter{sujetmetETADMRC}
\newcounter{sujetmetETAHA}
\newcounter{sujetmetpremtermETA}
\newcounter{sujetmetEXSPdP}
\newcounter{sujetmetEXSDFI}
\newcounter{sujetmetEXSDMRC}
\newcounter{sujetmetEXSHA}
\newcounter{sujetmetpremtermEXS}
\newcounter{sujetmetMMPdP}
\newcounter{sujetmetMMDFI}
\newcounter{sujetmetMMDMRC}
\newcounter{sujetmetMMHA}
\newcounter{sujetmetpremtermMM}
\newcounter{sujetmetCCPdP}
\newcounter{sujetmetCCDFI}
\newcounter{sujetmetCCDMRC}
\newcounter{sujetmetCCHA}
\newcounter{sujetmetpremtermCC}
\newcounter{sujetmetHVPdP}
\newcounter{sujetmetHVDFI}
\newcounter{sujetmetHVDMRC}
\newcounter{sujetmetHVHA}
\newcounter{sujetmetpremtermHV}
\newcounter{sujetmetHLPdP}
\newcounter{sujetmetHLDFI}
\newcounter{sujetmetHLDMRC}
\newcounter{sujetmetHLHA}
\newcounter{sujetmetpremtermHL}
\newcounter{sujetmetRSPdP}
\newcounter{sujetmetRSRdM}
\newcounter{sujetmetRSDFI}
\newcounter{sujetmetRSDMRC}
\newcounter{sujetmetRSHA}
\newcounter{sujetmetHQPdP}
\newcounter{sujetmetHQRdM}
\newcounter{sujetmetHQDFI}
\newcounter{sujetmetHQDMRC}
\newcounter{sujetmetHQHA}

%%%Somme des compteurs%%%
\newcounter{sujetmetallETA}
\newcounter{sujetmetallEXS}
\newcounter{sujetmetallMM}
\newcounter{sujetmetallRS}
\newcounter{sujetmetallCC}
\newcounter{sujetmetallHV}
\newcounter{sujetmetallHL}
\newcounter{sujetmetallHQ}

%%%Compter les sujetmets par thème et types de questions%%%
%%%%%%%%%%%%%%%%%%%%%%%%%%%%%%%%%%%%%%%%%%%%%
%%%%%%%%%Question d'interprétation%%%%%%%%%%%
%%%%%%%%%%%%%%%%%%%%%%%%%%%%%%%%%%%%%%%%%%%%%

%%%%%%%%%%%%%%%Général%%%%%%%%%%%%%%%%%%%%%%%
\newcounter{nosujetmetint}
\newcounter{sujetmetHQint}
\newcounter{sujetmetRSint}
%%%sujetmets simples%%%
\newcounter{sujetmetintEXS}
\newcounter{sujetmetintETA}
\newcounter{sujetmetintMM}
\newcounter{sujetmetintHV}
\newcounter{sujetmetintCC}
\newcounter{sujetmetintHL}
%%%sujetmets mixtes%%%
\newcounter{sujetmetintEXSTA}
\newcounter{sujetmetintEXSMM}
\newcounter{sujetmetintETAMM}
\newcounter{sujetmetintOdS}
\newcounter{sujetmetintHVL}
\newcounter{sujetmetintHVCC}
\newcounter{sujetmetintHLCC}
\newcounter{sujetmetintOdH}
%%%sujetmets mixtes semestres%%%
\newcounter{sujetmetintETACC}
\newcounter{sujetmetintETAHV}
\newcounter{sujetmetintETAHL}
\newcounter{sujetmetintEXSCC}
\newcounter{sujetmetintEXSHV}
\newcounter{sujetmetintEXSHL}
\newcounter{sujetmetintMMCC}
\newcounter{sujetmetintMMHV}
\newcounter{sujetmetintMMHL}
%%%sujetmets mixtes années%%
\newcounter{sujetmetintpremterm}
\newcounter{sujetmetintpremtermRS}
\newcounter{sujetmetintpremtermHQ}
\newcounter{sujetmetintETAPdP}
\newcounter{sujetmetintETADFI}
\newcounter{sujetmetintETADMRC}
\newcounter{sujetmetintETAHA}
\newcounter{sujetmetintpremtermETA}
\newcounter{sujetmetintEXSPdP}
\newcounter{sujetmetintEXSDFI}
\newcounter{sujetmetintEXSDMRC}
\newcounter{sujetmetintEXSHA}
\newcounter{sujetmetintpremtermEXS}
\newcounter{sujetmetintMMPdP}
\newcounter{sujetmetintMMDFI}
\newcounter{sujetmetintMMDMRC}
\newcounter{sujetmetintMMHA}
\newcounter{sujetmetintpremtermMM}
\newcounter{sujetmetintCCPdP}
\newcounter{sujetmetintCCDFI}
\newcounter{sujetmetintCCDMRC}
\newcounter{sujetmetintCCHA}
\newcounter{sujetmetintpremtermCC}
\newcounter{sujetmetintHVPdP}
\newcounter{sujetmetintHVDFI}
\newcounter{sujetmetintHVDMRC}
\newcounter{sujetmetintHVHA}
\newcounter{sujetmetintpremtermHV}
\newcounter{sujetmetintHLPdP}
\newcounter{sujetmetintHLDFI}
\newcounter{sujetmetintHLDMRC}
\newcounter{sujetmetintHLHA}
\newcounter{sujetmetintpremtermHL}
\newcounter{sujetmetintRSPdP}
\newcounter{sujetmetintRSRdM}
\newcounter{sujetmetintRSDFI}
\newcounter{sujetmetintRSDMRC}
\newcounter{sujetmetintRSHA}
\newcounter{sujetmetintHQPdP}
\newcounter{sujetmetintHQRdM}
\newcounter{sujetmetintHQDFI}
\newcounter{sujetmetintHQDMRC}
\newcounter{sujetmetintHQHA}
%%%Somme des compteurs%%%
\newcounter{sujetmetallintETA}
\newcounter{sujetmetallintEXS}
\newcounter{sujetmetallintMM}
\newcounter{sujetmetallintRS}
\newcounter{sujetmetallintCC}
\newcounter{sujetmetallintHV}
\newcounter{sujetmetallintHL}
\newcounter{sujetmetallintHQ}

%%%%%%%%Interprétation philosophique%%%%%%%%%
\newcounter{nosujetmetintphilosophique}
\newcounter{sujetmetHQintphilosophique}
\newcounter{sujetmetRSintphilosophique}
%%%sujetmets simples%%%
\newcounter{sujetmetintphilosophiqueEXS}
\newcounter{sujetmetintphilosophiqueETA}
\newcounter{sujetmetintphilosophiqueMM}
\newcounter{sujetmetintphilosophiqueHV}
\newcounter{sujetmetintphilosophiqueCC}
\newcounter{sujetmetintphilosophiqueHL}
%%%sujetmets mixtes%%%
\newcounter{sujetmetintphilosophiqueEXSTA}
\newcounter{sujetmetintphilosophiqueEXSMM}
\newcounter{sujetmetintphilosophiqueETAMM}
\newcounter{sujetmetintphilosophiqueOdS}
\newcounter{sujetmetintphilosophiqueHVL}
\newcounter{sujetmetintphilosophiqueHVCC}
\newcounter{sujetmetintphilosophiqueHLCC}
\newcounter{sujetmetintphilosophiqueOdH}
%%%sujetmets mixtes semestres%%%
\newcounter{sujetmetintphilosophiqueETACC}
\newcounter{sujetmetintphilosophiqueETAHV}
\newcounter{sujetmetintphilosophiqueETAHL}
\newcounter{sujetmetintphilosophiqueEXSCC}
\newcounter{sujetmetintphilosophiqueEXSHV}
\newcounter{sujetmetintphilosophiqueEXSHL}
\newcounter{sujetmetintphilosophiqueMMCC}
\newcounter{sujetmetintphilosophiqueMMHV}
\newcounter{sujetmetintphilosophiqueMMHL}
%%%sujetmets mixtes années%%
\newcounter{sujetmetintphilosophiquepremterm}
\newcounter{sujetmetintphilosophiquepremtermRS}
\newcounter{sujetmetintphilosophiquepremtermHQ}
\newcounter{sujetmetintphilosophiqueETAPdP}
\newcounter{sujetmetintphilosophiqueETADFI}
\newcounter{sujetmetintphilosophiqueETADMRC}
\newcounter{sujetmetintphilosophiqueETAHA}
\newcounter{sujetmetintphilosophiquepremtermETA}
\newcounter{sujetmetintphilosophiqueEXSPdP}
\newcounter{sujetmetintphilosophiqueEXSDFI}
\newcounter{sujetmetintphilosophiqueEXSDMRC}
\newcounter{sujetmetintphilosophiqueEXSHA}
\newcounter{sujetmetintphilosophiquepremtermEXS}
\newcounter{sujetmetintphilosophiqueMMPdP}
\newcounter{sujetmetintphilosophiqueMMDFI}
\newcounter{sujetmetintphilosophiqueMMDMRC}
\newcounter{sujetmetintphilosophiqueMMHA}
\newcounter{sujetmetintphilosophiquepremtermMM}
\newcounter{sujetmetintphilosophiqueCCPdP}
\newcounter{sujetmetintphilosophiqueCCDFI}
\newcounter{sujetmetintphilosophiqueCCDMRC}
\newcounter{sujetmetintphilosophiqueCCHA}
\newcounter{sujetmetintphilosophiquepremtermCC}
\newcounter{sujetmetintphilosophiqueHVPdP}
\newcounter{sujetmetintphilosophiqueHVDFI}
\newcounter{sujetmetintphilosophiqueHVDMRC}
\newcounter{sujetmetintphilosophiqueHVHA}
\newcounter{sujetmetintphilosophiquepremtermHV}
\newcounter{sujetmetintphilosophiqueHLPdP}
\newcounter{sujetmetintphilosophiqueHLDFI}
\newcounter{sujetmetintphilosophiqueHLDMRC}
\newcounter{sujetmetintphilosophiqueHLHA}
\newcounter{sujetmetintphilosophiquepremtermHL}
\newcounter{sujetmetintphilosophiqueRSPdP}
\newcounter{sujetmetintphilosophiqueRSRdM}
\newcounter{sujetmetintphilosophiqueRSDFI}
\newcounter{sujetmetintphilosophiqueRSDMRC}
\newcounter{sujetmetintphilosophiqueRSHA}
\newcounter{sujetmetintphilosophiqueHQPdP}
\newcounter{sujetmetintphilosophiqueHQRdM}
\newcounter{sujetmetintphilosophiqueHQDFI}
\newcounter{sujetmetintphilosophiqueHQDMRC}
\newcounter{sujetmetintphilosophiqueHQHA}
%%%Somme des compteurs%%%
\newcounter{sujetmetallintphilosophiqueETA}
\newcounter{sujetmetallintphilosophiqueEXS}
\newcounter{sujetmetallintphilosophiqueMM}
\newcounter{sujetmetallintphilosophiqueRS}
\newcounter{sujetmetallintphilosophiqueCC}
\newcounter{sujetmetallintphilosophiqueHV}
\newcounter{sujetmetallintphilosophiqueHL}
\newcounter{sujetmetallintphilosophiqueHQ}

%%%%%%%%%Interprétation littéraire%%%%%%%%%%%
\newcounter{nosujetmetintlitteraire}
\newcounter{sujetmetHQintlitteraire}
\newcounter{sujetmetRSintlitteraire}
%%%sujetmets simples%%%
\newcounter{sujetmetintlitteraireEXS}
\newcounter{sujetmetintlitteraireETA}
\newcounter{sujetmetintlitteraireMM}
\newcounter{sujetmetintlitteraireHV}
\newcounter{sujetmetintlitteraireCC}
\newcounter{sujetmetintlitteraireHL}
%%%sujetmets mixtes%%%
\newcounter{sujetmetintlitteraireEXSTA}
\newcounter{sujetmetintlitteraireEXSMM}
\newcounter{sujetmetintlitteraireETAMM}
\newcounter{sujetmetintlitteraireOdS}
\newcounter{sujetmetintlitteraireHVL}
\newcounter{sujetmetintlitteraireHVCC}
\newcounter{sujetmetintlitteraireHLCC}
\newcounter{sujetmetintlitteraireOdH}
%%%sujetmets mixtes semestres%%%
\newcounter{sujetmetintlitteraireETACC}
\newcounter{sujetmetintlitteraireETAHV}
\newcounter{sujetmetintlitteraireETAHL}
\newcounter{sujetmetintlitteraireEXSCC}
\newcounter{sujetmetintlitteraireEXSHV}
\newcounter{sujetmetintlitteraireEXSHL}
\newcounter{sujetmetintlitteraireMMCC}
\newcounter{sujetmetintlitteraireMMHV}
\newcounter{sujetmetintlitteraireMMHL}
%%%sujetmets mixtes années%%
\newcounter{sujetmetintlitterairepremterm}
\newcounter{sujetmetintlitterairepremtermRS}
\newcounter{sujetmetintlitterairepremtermHQ}
\newcounter{sujetmetintlitteraireETAPdP}
\newcounter{sujetmetintlitteraireETADFI}
\newcounter{sujetmetintlitteraireETADMRC}
\newcounter{sujetmetintlitteraireETAHA}
\newcounter{sujetmetintlitterairepremtermETA}
\newcounter{sujetmetintlitteraireEXSPdP}
\newcounter{sujetmetintlitteraireEXSDFI}
\newcounter{sujetmetintlitteraireEXSDMRC}
\newcounter{sujetmetintlitteraireEXSHA}
\newcounter{sujetmetintlitterairepremtermEXS}
\newcounter{sujetmetintlitteraireMMPdP}
\newcounter{sujetmetintlitteraireMMDFI}
\newcounter{sujetmetintlitteraireMMDMRC}
\newcounter{sujetmetintlitteraireMMHA}
\newcounter{sujetmetintlitterairepremtermMM}
\newcounter{sujetmetintlitteraireCCPdP}
\newcounter{sujetmetintlitteraireCCDFI}
\newcounter{sujetmetintlitteraireCCDMRC}
\newcounter{sujetmetintlitteraireCCHA}
\newcounter{sujetmetintlitterairepremtermCC}
\newcounter{sujetmetintlitteraireHVPdP}
\newcounter{sujetmetintlitteraireHVDFI}
\newcounter{sujetmetintlitteraireHVDMRC}
\newcounter{sujetmetintlitteraireHVHA}
\newcounter{sujetmetintlitterairepremtermHV}
\newcounter{sujetmetintlitteraireHLPdP}
\newcounter{sujetmetintlitteraireHLDFI}
\newcounter{sujetmetintlitteraireHLDMRC}
\newcounter{sujetmetintlitteraireHLHA}
\newcounter{sujetmetintlitterairepremtermHL}
\newcounter{sujetmetintlitteraireRSPdP}
\newcounter{sujetmetintlitteraireRSRdM}
\newcounter{sujetmetintlitteraireRSDFI}
\newcounter{sujetmetintlitteraireRSDMRC}
\newcounter{sujetmetintlitteraireRSHA}
\newcounter{sujetmetintlitteraireHQPdP}
\newcounter{sujetmetintlitteraireHQRdM}
\newcounter{sujetmetintlitteraireHQDFI}
\newcounter{sujetmetintlitteraireHQDMRC}
\newcounter{sujetmetintlitteraireHQHA}
%%%Somme des compteurs%%%
\newcounter{sujetmetallintlitteraireETA}
\newcounter{sujetmetallintlitteraireEXS}
\newcounter{sujetmetallintlitteraireMM}
\newcounter{sujetmetallintlitteraireRS}
\newcounter{sujetmetallintlitteraireCC}
\newcounter{sujetmetallintlitteraireHV}
\newcounter{sujetmetallintlitteraireHL}
\newcounter{sujetmetallintlitteraireHQ}

%%%%%%%%%%%%%%%%%%%%%%%%%%%%%%%%%%%%%%%%%%%%%
%%%%%%%%%%%%Question de réflexion%%%%%%%%%%%%
%%%%%%%%%%%%%%%%%%%%%%%%%%%%%%%%%%%%%%%%%%%%%
\newcounter{nosujetmetref}
%%%sujetmets simples%%%
\newcounter{sujetmetrefEXS}
\newcounter{sujetmetrefETA}
\newcounter{sujetmetrefMM}
\newcounter{sujetmetrefHV}
\newcounter{sujetmetrefCC}
\newcounter{sujetmetrefHL}
%%%sujetmets mixtes%%%
\newcounter{sujetmetrefEXSTA}
\newcounter{sujetmetrefEXSMM}
\newcounter{sujetmetrefETAMM}
\newcounter{sujetmetrefOdS}
\newcounter{sujetmetrefHVL}
\newcounter{sujetmetrefHVCC}
\newcounter{sujetmetrefHLCC}
\newcounter{sujetmetrefOdH}
%%%sujetmets mixtes semestres%%%
\newcounter{sujetmetrefETACC}
\newcounter{sujetmetrefETAHV}
\newcounter{sujetmetrefETAHL}
\newcounter{sujetmetrefEXSCC}
\newcounter{sujetmetrefEXSHV}
\newcounter{sujetmetrefEXSHL}
\newcounter{sujetmetrefMMCC}
\newcounter{sujetmetrefMMHV}
\newcounter{sujetmetrefMMHL}
%%%sujetmets mixtes années%%
\newcounter{sujetmetrefpremterm}%%%Compteur général
\newcounter{sujetmetrefETAPdP}
\newcounter{sujetmetrefETADFI}
\newcounter{sujetmetrefETADMRC}
\newcounter{sujetmetrefETAHA}
\newcounter{sujetmetrefpremtermETA}
\newcounter{sujetmetrefEXSPdP}
\newcounter{sujetmetrefEXSDFI}
\newcounter{sujetmetrefEXSDMRC}
\newcounter{sujetmetrefEXSHA}
\newcounter{sujetmetrefpremtermEXS}
\newcounter{sujetmetrefMMPdP}
\newcounter{sujetmetrefMMDFI}
\newcounter{sujetmetrefMMDMRC}
\newcounter{sujetmetrefMMHA}
\newcounter{sujetmetrefpremtermMM}
\newcounter{sujetmetrefCCPdP}
\newcounter{sujetmetrefCCDFI}
\newcounter{sujetmetrefCCDMRC}
\newcounter{sujetmetrefCCHA}
\newcounter{sujetmetrefpremtermCC}
\newcounter{sujetmetrefHVPdP}
\newcounter{sujetmetrefHVDFI}
\newcounter{sujetmetrefHVDMRC}
\newcounter{sujetmetrefHVHA}
\newcounter{sujetmetrefpremtermHV}
\newcounter{sujetmetrefHLPdP}
\newcounter{sujetmetrefHLDFI}
\newcounter{sujetmetrefHLDMRC}
\newcounter{sujetmetrefHLHA}
\newcounter{sujetmetrefpremtermHL}
\newcounter{sujetmetrefRSPdP}
\newcounter{sujetmetrefRSRdM}
\newcounter{sujetmetrefRSDFI}
\newcounter{sujetmetrefRSDMRC}
\newcounter{sujetmetrefRSHA}
\newcounter{sujetmetrefHQPdP}
\newcounter{sujetmetrefHQRdM}
\newcounter{sujetmetrefHQDFI}
\newcounter{sujetmetrefHQDMRC}
\newcounter{sujetmetrefHQHA}
%%%Somme des compteurs%%%
\newcounter{sujetmetallrefETA}
\newcounter{sujetmetallrefEXS}
\newcounter{sujetmetallrefMM}
\newcounter{sujetmetallrefRS}
\newcounter{sujetmetallrefCC}
\newcounter{sujetmetallrefHV}
\newcounter{sujetmetallrefHL}
\newcounter{sujetmetallrefHQ}

%%%%%%%%Réflexion philosophique%%%%%%%%%
\newcounter{nosujetmetrefphilosophique}
\newcounter{sujetmetHQrefphilosophique}
\newcounter{sujetmetRSrefphilosophique}
%%%sujetmets simples%%%
\newcounter{sujetmetrefphilosophiqueEXS}
\newcounter{sujetmetrefphilosophiqueETA}
\newcounter{sujetmetrefphilosophiqueMM}
\newcounter{sujetmetrefphilosophiqueHV}
\newcounter{sujetmetrefphilosophiqueCC}
\newcounter{sujetmetrefphilosophiqueHL}
%%%sujetmets mixtes%%%
\newcounter{sujetmetrefphilosophiqueEXSTA}
\newcounter{sujetmetrefphilosophiqueEXSMM}
\newcounter{sujetmetrefphilosophiqueETAMM}
\newcounter{sujetmetrefphilosophiqueOdS}
\newcounter{sujetmetrefphilosophiqueHVL}
\newcounter{sujetmetrefphilosophiqueHVCC}
\newcounter{sujetmetrefphilosophiqueHLCC}
\newcounter{sujetmetrefphilosophiqueOdH}
%%%sujetmets mixtes semestres%%%
\newcounter{sujetmetrefphilosophiqueETACC}
\newcounter{sujetmetrefphilosophiqueETAHV}
\newcounter{sujetmetrefphilosophiqueETAHL}
\newcounter{sujetmetrefphilosophiqueEXSCC}
\newcounter{sujetmetrefphilosophiqueEXSHV}
\newcounter{sujetmetrefphilosophiqueEXSHL}
\newcounter{sujetmetrefphilosophiqueMMCC}
\newcounter{sujetmetrefphilosophiqueMMHV}
\newcounter{sujetmetrefphilosophiqueMMHL}
%%%sujetmets mixtes années%%
\newcounter{sujetmetrefphilosophiquepremterm}
\newcounter{sujetmetrefphilosophiquepremtermRS}
\newcounter{sujetmetrefphilosophiquepremtermHQ}
\newcounter{sujetmetrefphilosophiqueETAPdP}
\newcounter{sujetmetrefphilosophiqueETADFI}
\newcounter{sujetmetrefphilosophiqueETADMRC}
\newcounter{sujetmetrefphilosophiqueETAHA}\newcounter{sujetmetrefphilosophiquepremtermETA}
\newcounter{sujetmetrefphilosophiqueEXSPdP}
\newcounter{sujetmetrefphilosophiqueEXSDFI}
\newcounter{sujetmetrefphilosophiqueEXSDMRC}
\newcounter{sujetmetrefphilosophiqueEXSHA}
\newcounter{sujetmetrefphilosophiquepremtermEXS}
\newcounter{sujetmetrefphilosophiqueMMPdP}
\newcounter{sujetmetrefphilosophiqueMMDFI}
\newcounter{sujetmetrefphilosophiqueMMDMRC}
\newcounter{sujetmetrefphilosophiqueMMHA}
\newcounter{sujetmetrefphilosophiquepremtermMM}
\newcounter{sujetmetrefphilosophiqueCCPdP}
\newcounter{sujetmetrefphilosophiqueCCDFI}
\newcounter{sujetmetrefphilosophiqueCCDMRC}
\newcounter{sujetmetrefphilosophiqueCCHA}
\newcounter{sujetmetrefphilosophiquepremtermCC}
\newcounter{sujetmetrefphilosophiqueHVPdP}
\newcounter{sujetmetrefphilosophiqueHVDFI}
\newcounter{sujetmetrefphilosophiqueHVDMRC}
\newcounter{sujetmetrefphilosophiqueHVHA}
\newcounter{sujetmetrefphilosophiquepremtermHV}
\newcounter{sujetmetrefphilosophiqueHLPdP}
\newcounter{sujetmetrefphilosophiqueHLDFI}
\newcounter{sujetmetrefphilosophiqueHLDMRC}
\newcounter{sujetmetrefphilosophiqueHLHA}
\newcounter{sujetmetrefphilosophiquepremtermHL}
\newcounter{sujetmetrefphilosophiqueRSPdP}
\newcounter{sujetmetrefphilosophiqueRSRdM}
\newcounter{sujetmetrefphilosophiqueRSDFI}
\newcounter{sujetmetrefphilosophiqueRSDMRC}
\newcounter{sujetmetrefphilosophiqueRSHA}
\newcounter{sujetmetrefphilosophiqueHQPdP}
\newcounter{sujetmetrefphilosophiqueHQRdM}
\newcounter{sujetmetrefphilosophiqueHQDFI}
\newcounter{sujetmetrefphilosophiqueHQDMRC}
\newcounter{sujetmetrefphilosophiqueHQHA}
%%%Somme des compteurs%%%
\newcounter{sujetmetallrefphilosophiqueETA}
\newcounter{sujetmetallrefphilosophiqueEXS}
\newcounter{sujetmetallrefphilosophiqueMM}
\newcounter{sujetmetallrefphilosophiqueRS}
\newcounter{sujetmetallrefphilosophiqueCC}
\newcounter{sujetmetallrefphilosophiqueHV}
\newcounter{sujetmetallrefphilosophiqueHL}
\newcounter{sujetmetallrefphilosophiqueHQ}

%%%%%%%%%Réflexion littéraire%%%%%%%%%%%
\newcounter{nosujetmetreflitteraire}
\newcounter{sujetmetHQreflitteraire}
\newcounter{sujetmetRSreflitteraire}
%%%sujetmets simples%%%
\newcounter{sujetmetreflitteraireEXS}
\newcounter{sujetmetreflitteraireETA}
\newcounter{sujetmetreflitteraireMM}
\newcounter{sujetmetreflitteraireHV}
\newcounter{sujetmetreflitteraireCC}
\newcounter{sujetmetreflitteraireHL}
%%%sujetmets mixtes%%%
\newcounter{sujetmetreflitteraireEXSTA}
\newcounter{sujetmetreflitteraireEXSMM}
\newcounter{sujetmetreflitteraireETAMM}
\newcounter{sujetmetreflitteraireOdS}
\newcounter{sujetmetreflitteraireHVL}
\newcounter{sujetmetreflitteraireHVCC}
\newcounter{sujetmetreflitteraireHLCC}
\newcounter{sujetmetreflitteraireOdH}
%%%sujetmets mixtes semestres%%%
\newcounter{sujetmetreflitteraireETACC}
\newcounter{sujetmetreflitteraireETAHV}
\newcounter{sujetmetreflitteraireETAHL}
\newcounter{sujetmetreflitteraireEXSCC}
\newcounter{sujetmetreflitteraireEXSHV}
\newcounter{sujetmetreflitteraireEXSHL}
\newcounter{sujetmetreflitteraireMMCC}
\newcounter{sujetmetreflitteraireMMHV}
\newcounter{sujetmetreflitteraireMMHL}
%%%sujetmets mixtes années%%
\newcounter{sujetmetreflitterairepremterm}
\newcounter{sujetmetreflitterairepremtermRS}
\newcounter{sujetmetreflitterairepremtermHQ}
\newcounter{sujetmetreflitteraireETAPdP}
\newcounter{sujetmetreflitteraireETADFI}
\newcounter{sujetmetreflitteraireETADMRC}
\newcounter{sujetmetreflitteraireETAHA}
\newcounter{sujetmetreflitterairepremtermETA}
\newcounter{sujetmetreflitteraireEXSPdP}
\newcounter{sujetmetreflitteraireEXSDFI}
\newcounter{sujetmetreflitteraireEXSDMRC}
\newcounter{sujetmetreflitteraireEXSHA}
\newcounter{sujetmetreflitterairepremtermEXS}
\newcounter{sujetmetreflitteraireMMPdP}
\newcounter{sujetmetreflitteraireMMDFI}
\newcounter{sujetmetreflitteraireMMDMRC}
\newcounter{sujetmetreflitteraireMMHA}
\newcounter{sujetmetreflitterairepremtermMM}
\newcounter{sujetmetreflitteraireCCPdP}
\newcounter{sujetmetreflitteraireCCDFI}
\newcounter{sujetmetreflitteraireCCDMRC}
\newcounter{sujetmetreflitteraireCCHA}
\newcounter{sujetmetreflitterairepremtermCC}
\newcounter{sujetmetreflitteraireHVPdP}
\newcounter{sujetmetreflitteraireHVDFI}
\newcounter{sujetmetreflitteraireHVDMRC}
\newcounter{sujetmetreflitteraireHVHA}
\newcounter{sujetmetreflitterairepremtermHV}
\newcounter{sujetmetreflitteraireHLPdP}
\newcounter{sujetmetreflitteraireHLDFI}
\newcounter{sujetmetreflitteraireHLDMRC}
\newcounter{sujetmetreflitteraireHLHA}
\newcounter{sujetmetreflitterairepremtermHL}
\newcounter{sujetmetreflitteraireRSPdP}
\newcounter{sujetmetreflitteraireRSRdM}
\newcounter{sujetmetreflitteraireRSDFI}
\newcounter{sujetmetreflitteraireRSDMRC}
\newcounter{sujetmetreflitteraireRSHA}
\newcounter{sujetmetreflitteraireHQPdP}
\newcounter{sujetmetreflitteraireHQRdM}
\newcounter{sujetmetreflitteraireHQDFI}
\newcounter{sujetmetreflitteraireHQDMRC}
\newcounter{sujetmetreflitteraireHQHA}
%%%Somme des compteurs%%%
\newcounter{sujetmetallreflitteraireETA}
\newcounter{sujetmetallreflitteraireEXS}
\newcounter{sujetmetallreflitteraireMM}
\newcounter{sujetmetallreflitteraireCC}
\newcounter{sujetmetallreflitteraireRS}
\newcounter{sujetmetallreflitteraireHV}
\newcounter{sujetmetallreflitteraireHL}
\newcounter{sujetmetallreflitteraireHQ}
%%%Compteurs de statistiques globales sur les genres de question%%%
\newcounter{sujetmetexclRS}
\newcounter{sujetmetmixtRS}
\newcounter{sujetmetexclHQ}
\newcounter{sujetmetmixtHQ}
\newcounter{sujetmetmixtHQRS}
%%%Compteurs de statistiques sur les transversaux (divers semestres, première et terminale)
\newcounter{sujettransmetintlitteraire}
\newcounter{sujettransmetreflitteraire}
\newcounter{sujettransmetintphilosophique}
\newcounter{sujettransmetrefphilosophique}
%%% Compteurs du nombre de sujets par type
\newcounter{sujetmetintlitteraire}
\newcounter{sujetmetreflitteraire}
\newcounter{sujetmetintphilosophique}
\newcounter{sujetmetrefphilosophique}
%%% Compteurs par discipline
%%Litt
\newcounter{sujetmetlitteraireETA}
\newcounter{sujetmetlitteraireEXS}
\newcounter{sujetmetlitteraireMM}
\newcounter{sujetmetlitteraireCC}
\newcounter{sujetmetlitteraireHV}
\newcounter{sujetmetlitteraireHL}
\newcounter{sujetmetlitteraireEXSTA}
\newcounter{sujetmetlitteraireETAMM}
\newcounter{sujetmetlitteraireEXSMM}
\newcounter{sujetmetlitteraireHVCC}
\newcounter{sujetmetlitteraireHVL}
\newcounter{sujetmetlitteraireaut}
%%Phil
\newcounter{sujetmetphilosophiqueETA}
\newcounter{sujetmetphilosophiqueEXS}
\newcounter{sujetmetphilosophiqueMM}
\newcounter{sujetmetphilosophiqueCC}
\newcounter{sujetmetphilosophiqueHV}
\newcounter{sujetmetphilosophiqueHL}
\newcounter{sujetmetphilosophiqueEXSTA}
\newcounter{sujetmetphilosophiqueETAMM}
\newcounter{sujetmetphilosophiqueEXSMM}
\newcounter{sujetmetphilosophiqueHVCC}
\newcounter{sujetmetphilosophiqueHVL}
\newcounter{sujetmetphilosophiqueaut}
%%Totaux pour les graphiques
\newcounter{sujetmetlitterairetotal}
\newcounter{sujetmetphilosophiquetotal}
%%% Autres (Ods et Odh) comptés ensemble
\newcounter{sujetmetallaut}
%%%%%%%%%%%%%%%%%%%%%%%%%%%%%%%%%%%%%%%%%%%%%%%%%%%%%%%%%%%%%%%%%%%%%%%
%%%%%%%%%%%%%%%%%%%%%%%%%%%%%%Compteurs%%%%%%%%%%%%%%%%%%%%%%%%%%%%%%%%
%%%%%%%%%%%%%%%%%%%%%%%%%%%%%%%%%%%%%%%%%%%%%%%%%%%%%%%%%%%%%%%%%%%%%%%
%%%Compter les sujets par thème%%%
%\newcounter{nosujetnca}
\newcounter{sujetncaall}%Total de questions pour le centre
\newcounter{sujetncatotal}%Total des questions en comptant en double les questions pour les graphiques
\newcounter{sujetncaHQ}
\newcounter{sujetncaRS}
%%%sujetncas simples%%%
\newcounter{sujetncaEXS}
\newcounter{sujetncaETA}
\newcounter{sujetncaMM}
\newcounter{sujetncaHV}
\newcounter{sujetncaCC}
\newcounter{sujetncaHL}
%%%sujetncas mixtes%%%
\newcounter{sujetncaEXSTA}
\newcounter{sujetncaEXSMM}
\newcounter{sujetncaETAMM}
\newcounter{sujetncaOdS}
\newcounter{sujetncaHVL}
\newcounter{sujetncaHVCC}
\newcounter{sujetncaHLCC}
\newcounter{sujetncaOdH}
%%%sujetncas mixtes semestres%%%
\newcounter{sujetncaETACC}
\newcounter{sujetncaETAHV}
\newcounter{sujetncaETAHL}
\newcounter{sujetncaEXSCC}
\newcounter{sujetncaEXSHV}
\newcounter{sujetncaEXSHL}
\newcounter{sujetncaMMCC}
\newcounter{sujetncaMMHV}
\newcounter{sujetncaMMHL}

%%%sujetncas mixtes années%%
\newcounter{sujetncapremterm}%%%Compteur général de sujetnca
\newcounter{sujetncaETAPdP}
\newcounter{sujetncaETADFI}
\newcounter{sujetncaETADMRC}
\newcounter{sujetncaETAHA}
\newcounter{sujetncapremtermETA}
\newcounter{sujetncaEXSPdP}
\newcounter{sujetncaEXSDFI}
\newcounter{sujetncaEXSDMRC}
\newcounter{sujetncaEXSHA}
\newcounter{sujetncapremtermEXS}
\newcounter{sujetncaMMPdP}
\newcounter{sujetncaMMDFI}
\newcounter{sujetncaMMDMRC}
\newcounter{sujetncaMMHA}
\newcounter{sujetncapremtermMM}
\newcounter{sujetncaCCPdP}
\newcounter{sujetncaCCDFI}
\newcounter{sujetncaCCDMRC}
\newcounter{sujetncaCCHA}
\newcounter{sujetncapremtermCC}
\newcounter{sujetncaHVPdP}
\newcounter{sujetncaHVDFI}
\newcounter{sujetncaHVDMRC}
\newcounter{sujetncaHVHA}
\newcounter{sujetncapremtermHV}
\newcounter{sujetncaHLPdP}
\newcounter{sujetncaHLDFI}
\newcounter{sujetncaHLDMRC}
\newcounter{sujetncaHLHA}
\newcounter{sujetncapremtermHL}
\newcounter{sujetncaRSPdP}
\newcounter{sujetncaRSRdM}
\newcounter{sujetncaRSDFI}
\newcounter{sujetncaRSDMRC}
\newcounter{sujetncaRSHA}
\newcounter{sujetncaHQPdP}
\newcounter{sujetncaHQRdM}
\newcounter{sujetncaHQDFI}
\newcounter{sujetncaHQDMRC}
\newcounter{sujetncaHQHA}

%%%Somme des compteurs%%%
\newcounter{sujetncaallETA}
\newcounter{sujetncaallEXS}
\newcounter{sujetncaallMM}
\newcounter{sujetncaallRS}
\newcounter{sujetncaallCC}
\newcounter{sujetncaallHV}
\newcounter{sujetncaallHL}
\newcounter{sujetncaallHQ}

%%%Compter les sujetncas par thème et types de questions%%%
%%%%%%%%%%%%%%%%%%%%%%%%%%%%%%%%%%%%%%%%%%%%%
%%%%%%%%%Question d'interprétation%%%%%%%%%%%
%%%%%%%%%%%%%%%%%%%%%%%%%%%%%%%%%%%%%%%%%%%%%

%%%%%%%%%%%%%%%Général%%%%%%%%%%%%%%%%%%%%%%%
\newcounter{nosujetncaint}
\newcounter{sujetncaHQint}
\newcounter{sujetncaRSint}
%%%sujetncas simples%%%
\newcounter{sujetncaintEXS}
\newcounter{sujetncaintETA}
\newcounter{sujetncaintMM}
\newcounter{sujetncaintHV}
\newcounter{sujetncaintCC}
\newcounter{sujetncaintHL}
%%%sujetncas mixtes%%%
\newcounter{sujetncaintEXSTA}
\newcounter{sujetncaintEXSMM}
\newcounter{sujetncaintETAMM}
\newcounter{sujetncaintOdS}
\newcounter{sujetncaintHVL}
\newcounter{sujetncaintHVCC}
\newcounter{sujetncaintHLCC}
\newcounter{sujetncaintOdH}
%%%sujetncas mixtes semestres%%%
\newcounter{sujetncaintETACC}
\newcounter{sujetncaintETAHV}
\newcounter{sujetncaintETAHL}
\newcounter{sujetncaintEXSCC}
\newcounter{sujetncaintEXSHV}
\newcounter{sujetncaintEXSHL}
\newcounter{sujetncaintMMCC}
\newcounter{sujetncaintMMHV}
\newcounter{sujetncaintMMHL}
%%%sujetncas mixtes années%%
\newcounter{sujetncaintpremterm}
\newcounter{sujetncaintpremtermRS}
\newcounter{sujetncaintpremtermHQ}
\newcounter{sujetncaintETAPdP}
\newcounter{sujetncaintETADFI}
\newcounter{sujetncaintETADMRC}
\newcounter{sujetncaintETAHA}
\newcounter{sujetncaintpremtermETA}
\newcounter{sujetncaintEXSPdP}
\newcounter{sujetncaintEXSDFI}
\newcounter{sujetncaintEXSDMRC}
\newcounter{sujetncaintEXSHA}
\newcounter{sujetncaintpremtermEXS}
\newcounter{sujetncaintMMPdP}
\newcounter{sujetncaintMMDFI}
\newcounter{sujetncaintMMDMRC}
\newcounter{sujetncaintMMHA}
\newcounter{sujetncaintpremtermMM}
\newcounter{sujetncaintCCPdP}
\newcounter{sujetncaintCCDFI}
\newcounter{sujetncaintCCDMRC}
\newcounter{sujetncaintCCHA}
\newcounter{sujetncaintpremtermCC}
\newcounter{sujetncaintHVPdP}
\newcounter{sujetncaintHVDFI}
\newcounter{sujetncaintHVDMRC}
\newcounter{sujetncaintHVHA}
\newcounter{sujetncaintpremtermHV}
\newcounter{sujetncaintHLPdP}
\newcounter{sujetncaintHLDFI}
\newcounter{sujetncaintHLDMRC}
\newcounter{sujetncaintHLHA}
\newcounter{sujetncaintpremtermHL}
\newcounter{sujetncaintRSPdP}
\newcounter{sujetncaintRSRdM}
\newcounter{sujetncaintRSDFI}
\newcounter{sujetncaintRSDMRC}
\newcounter{sujetncaintRSHA}
\newcounter{sujetncaintHQPdP}
\newcounter{sujetncaintHQRdM}
\newcounter{sujetncaintHQDFI}
\newcounter{sujetncaintHQDMRC}
\newcounter{sujetncaintHQHA}
%%%Somme des compteurs%%%
\newcounter{sujetncaallintETA}
\newcounter{sujetncaallintEXS}
\newcounter{sujetncaallintMM}
\newcounter{sujetncaallintRS}
\newcounter{sujetncaallintCC}
\newcounter{sujetncaallintHV}
\newcounter{sujetncaallintHL}
\newcounter{sujetncaallintHQ}

%%%%%%%%Interprétation philosophique%%%%%%%%%
\newcounter{nosujetncaintphilosophique}
\newcounter{sujetncaHQintphilosophique}
\newcounter{sujetncaRSintphilosophique}
%%%sujetncas simples%%%
\newcounter{sujetncaintphilosophiqueEXS}
\newcounter{sujetncaintphilosophiqueETA}
\newcounter{sujetncaintphilosophiqueMM}
\newcounter{sujetncaintphilosophiqueHV}
\newcounter{sujetncaintphilosophiqueCC}
\newcounter{sujetncaintphilosophiqueHL}
%%%sujetncas mixtes%%%
\newcounter{sujetncaintphilosophiqueEXSTA}
\newcounter{sujetncaintphilosophiqueEXSMM}
\newcounter{sujetncaintphilosophiqueETAMM}
\newcounter{sujetncaintphilosophiqueOdS}
\newcounter{sujetncaintphilosophiqueHVL}
\newcounter{sujetncaintphilosophiqueHVCC}
\newcounter{sujetncaintphilosophiqueHLCC}
\newcounter{sujetncaintphilosophiqueOdH}
%%%sujetncas mixtes semestres%%%
\newcounter{sujetncaintphilosophiqueETACC}
\newcounter{sujetncaintphilosophiqueETAHV}
\newcounter{sujetncaintphilosophiqueETAHL}
\newcounter{sujetncaintphilosophiqueEXSCC}
\newcounter{sujetncaintphilosophiqueEXSHV}
\newcounter{sujetncaintphilosophiqueEXSHL}
\newcounter{sujetncaintphilosophiqueMMCC}
\newcounter{sujetncaintphilosophiqueMMHV}
\newcounter{sujetncaintphilosophiqueMMHL}
%%%sujetncas mixtes années%%
\newcounter{sujetncaintphilosophiquepremterm}
\newcounter{sujetncaintphilosophiquepremtermRS}
\newcounter{sujetncaintphilosophiquepremtermHQ}
\newcounter{sujetncaintphilosophiqueETAPdP}
\newcounter{sujetncaintphilosophiqueETADFI}
\newcounter{sujetncaintphilosophiqueETADMRC}
\newcounter{sujetncaintphilosophiqueETAHA}
\newcounter{sujetncaintphilosophiquepremtermETA}
\newcounter{sujetncaintphilosophiqueEXSPdP}
\newcounter{sujetncaintphilosophiqueEXSDFI}
\newcounter{sujetncaintphilosophiqueEXSDMRC}
\newcounter{sujetncaintphilosophiqueEXSHA}
\newcounter{sujetncaintphilosophiquepremtermEXS}
\newcounter{sujetncaintphilosophiqueMMPdP}
\newcounter{sujetncaintphilosophiqueMMDFI}
\newcounter{sujetncaintphilosophiqueMMDMRC}
\newcounter{sujetncaintphilosophiqueMMHA}
\newcounter{sujetncaintphilosophiquepremtermMM}
\newcounter{sujetncaintphilosophiqueCCPdP}
\newcounter{sujetncaintphilosophiqueCCDFI}
\newcounter{sujetncaintphilosophiqueCCDMRC}
\newcounter{sujetncaintphilosophiqueCCHA}
\newcounter{sujetncaintphilosophiquepremtermCC}
\newcounter{sujetncaintphilosophiqueHVPdP}
\newcounter{sujetncaintphilosophiqueHVDFI}
\newcounter{sujetncaintphilosophiqueHVDMRC}
\newcounter{sujetncaintphilosophiqueHVHA}
\newcounter{sujetncaintphilosophiquepremtermHV}
\newcounter{sujetncaintphilosophiqueHLPdP}
\newcounter{sujetncaintphilosophiqueHLDFI}
\newcounter{sujetncaintphilosophiqueHLDMRC}
\newcounter{sujetncaintphilosophiqueHLHA}
\newcounter{sujetncaintphilosophiquepremtermHL}
\newcounter{sujetncaintphilosophiqueRSPdP}
\newcounter{sujetncaintphilosophiqueRSRdM}
\newcounter{sujetncaintphilosophiqueRSDFI}
\newcounter{sujetncaintphilosophiqueRSDMRC}
\newcounter{sujetncaintphilosophiqueRSHA}
\newcounter{sujetncaintphilosophiqueHQPdP}
\newcounter{sujetncaintphilosophiqueHQRdM}
\newcounter{sujetncaintphilosophiqueHQDFI}
\newcounter{sujetncaintphilosophiqueHQDMRC}
\newcounter{sujetncaintphilosophiqueHQHA}
%%%Somme des compteurs%%%
\newcounter{sujetncaallintphilosophiqueETA}
\newcounter{sujetncaallintphilosophiqueEXS}
\newcounter{sujetncaallintphilosophiqueMM}
\newcounter{sujetncaallintphilosophiqueRS}
\newcounter{sujetncaallintphilosophiqueCC}
\newcounter{sujetncaallintphilosophiqueHV}
\newcounter{sujetncaallintphilosophiqueHL}
\newcounter{sujetncaallintphilosophiqueHQ}

%%%%%%%%%Interprétation littéraire%%%%%%%%%%%
\newcounter{nosujetncaintlitteraire}
\newcounter{sujetncaHQintlitteraire}
\newcounter{sujetncaRSintlitteraire}
%%%sujetncas simples%%%
\newcounter{sujetncaintlitteraireEXS}
\newcounter{sujetncaintlitteraireETA}
\newcounter{sujetncaintlitteraireMM}
\newcounter{sujetncaintlitteraireHV}
\newcounter{sujetncaintlitteraireCC}
\newcounter{sujetncaintlitteraireHL}
%%%sujetncas mixtes%%%
\newcounter{sujetncaintlitteraireEXSTA}
\newcounter{sujetncaintlitteraireEXSMM}
\newcounter{sujetncaintlitteraireETAMM}
\newcounter{sujetncaintlitteraireOdS}
\newcounter{sujetncaintlitteraireHVL}
\newcounter{sujetncaintlitteraireHVCC}
\newcounter{sujetncaintlitteraireHLCC}
\newcounter{sujetncaintlitteraireOdH}
%%%sujetncas mixtes semestres%%%
\newcounter{sujetncaintlitteraireETACC}
\newcounter{sujetncaintlitteraireETAHV}
\newcounter{sujetncaintlitteraireETAHL}
\newcounter{sujetncaintlitteraireEXSCC}
\newcounter{sujetncaintlitteraireEXSHV}
\newcounter{sujetncaintlitteraireEXSHL}
\newcounter{sujetncaintlitteraireMMCC}
\newcounter{sujetncaintlitteraireMMHV}
\newcounter{sujetncaintlitteraireMMHL}
%%%sujetncas mixtes années%%
\newcounter{sujetncaintlitterairepremterm}
\newcounter{sujetncaintlitterairepremtermRS}
\newcounter{sujetncaintlitterairepremtermHQ}
\newcounter{sujetncaintlitteraireETAPdP}
\newcounter{sujetncaintlitteraireETADFI}
\newcounter{sujetncaintlitteraireETADMRC}
\newcounter{sujetncaintlitteraireETAHA}
\newcounter{sujetncaintlitterairepremtermETA}
\newcounter{sujetncaintlitteraireEXSPdP}
\newcounter{sujetncaintlitteraireEXSDFI}
\newcounter{sujetncaintlitteraireEXSDMRC}
\newcounter{sujetncaintlitteraireEXSHA}
\newcounter{sujetncaintlitterairepremtermEXS}
\newcounter{sujetncaintlitteraireMMPdP}
\newcounter{sujetncaintlitteraireMMDFI}
\newcounter{sujetncaintlitteraireMMDMRC}
\newcounter{sujetncaintlitteraireMMHA}
\newcounter{sujetncaintlitterairepremtermMM}
\newcounter{sujetncaintlitteraireCCPdP}
\newcounter{sujetncaintlitteraireCCDFI}
\newcounter{sujetncaintlitteraireCCDMRC}
\newcounter{sujetncaintlitteraireCCHA}
\newcounter{sujetncaintlitterairepremtermCC}
\newcounter{sujetncaintlitteraireHVPdP}
\newcounter{sujetncaintlitteraireHVDFI}
\newcounter{sujetncaintlitteraireHVDMRC}
\newcounter{sujetncaintlitteraireHVHA}
\newcounter{sujetncaintlitterairepremtermHV}
\newcounter{sujetncaintlitteraireHLPdP}
\newcounter{sujetncaintlitteraireHLDFI}
\newcounter{sujetncaintlitteraireHLDMRC}
\newcounter{sujetncaintlitteraireHLHA}
\newcounter{sujetncaintlitterairepremtermHL}
\newcounter{sujetncaintlitteraireRSPdP}
\newcounter{sujetncaintlitteraireRSRdM}
\newcounter{sujetncaintlitteraireRSDFI}
\newcounter{sujetncaintlitteraireRSDMRC}
\newcounter{sujetncaintlitteraireRSHA}
\newcounter{sujetncaintlitteraireHQPdP}
\newcounter{sujetncaintlitteraireHQRdM}
\newcounter{sujetncaintlitteraireHQDFI}
\newcounter{sujetncaintlitteraireHQDMRC}
\newcounter{sujetncaintlitteraireHQHA}
%%%Somme des compteurs%%%
\newcounter{sujetncaallintlitteraireETA}
\newcounter{sujetncaallintlitteraireEXS}
\newcounter{sujetncaallintlitteraireMM}
\newcounter{sujetncaallintlitteraireRS}
\newcounter{sujetncaallintlitteraireCC}
\newcounter{sujetncaallintlitteraireHV}
\newcounter{sujetncaallintlitteraireHL}
\newcounter{sujetncaallintlitteraireHQ}

%%%%%%%%%%%%%%%%%%%%%%%%%%%%%%%%%%%%%%%%%%%%%
%%%%%%%%%%%%Question de réflexion%%%%%%%%%%%%
%%%%%%%%%%%%%%%%%%%%%%%%%%%%%%%%%%%%%%%%%%%%%
\newcounter{nosujetncaref}
%%%sujetncas simples%%%
\newcounter{sujetncarefEXS}
\newcounter{sujetncarefETA}
\newcounter{sujetncarefMM}
\newcounter{sujetncarefHV}
\newcounter{sujetncarefCC}
\newcounter{sujetncarefHL}
%%%sujetncas mixtes%%%
\newcounter{sujetncarefEXSTA}
\newcounter{sujetncarefEXSMM}
\newcounter{sujetncarefETAMM}
\newcounter{sujetncarefOdS}
\newcounter{sujetncarefHVL}
\newcounter{sujetncarefHVCC}
\newcounter{sujetncarefHLCC}
\newcounter{sujetncarefOdH}
%%%sujetncas mixtes semestres%%%
\newcounter{sujetncarefETACC}
\newcounter{sujetncarefETAHV}
\newcounter{sujetncarefETAHL}
\newcounter{sujetncarefEXSCC}
\newcounter{sujetncarefEXSHV}
\newcounter{sujetncarefEXSHL}
\newcounter{sujetncarefMMCC}
\newcounter{sujetncarefMMHV}
\newcounter{sujetncarefMMHL}
%%%sujetncas mixtes années%%
\newcounter{sujetncarefpremterm}%%%Compteur général
\newcounter{sujetncarefETAPdP}
\newcounter{sujetncarefETADFI}
\newcounter{sujetncarefETADMRC}
\newcounter{sujetncarefETAHA}
\newcounter{sujetncarefpremtermETA}
\newcounter{sujetncarefEXSPdP}
\newcounter{sujetncarefEXSDFI}
\newcounter{sujetncarefEXSDMRC}
\newcounter{sujetncarefEXSHA}
\newcounter{sujetncarefpremtermEXS}
\newcounter{sujetncarefMMPdP}
\newcounter{sujetncarefMMDFI}
\newcounter{sujetncarefMMDMRC}
\newcounter{sujetncarefMMHA}
\newcounter{sujetncarefpremtermMM}
\newcounter{sujetncarefCCPdP}
\newcounter{sujetncarefCCDFI}
\newcounter{sujetncarefCCDMRC}
\newcounter{sujetncarefCCHA}
\newcounter{sujetncarefpremtermCC}
\newcounter{sujetncarefHVPdP}
\newcounter{sujetncarefHVDFI}
\newcounter{sujetncarefHVDMRC}
\newcounter{sujetncarefHVHA}
\newcounter{sujetncarefpremtermHV}
\newcounter{sujetncarefHLPdP}
\newcounter{sujetncarefHLDFI}
\newcounter{sujetncarefHLDMRC}
\newcounter{sujetncarefHLHA}
\newcounter{sujetncarefpremtermHL}
\newcounter{sujetncarefRSPdP}
\newcounter{sujetncarefRSRdM}
\newcounter{sujetncarefRSDFI}
\newcounter{sujetncarefRSDMRC}
\newcounter{sujetncarefRSHA}
\newcounter{sujetncarefHQPdP}
\newcounter{sujetncarefHQRdM}
\newcounter{sujetncarefHQDFI}
\newcounter{sujetncarefHQDMRC}
\newcounter{sujetncarefHQHA}
%%%Somme des compteurs%%%
\newcounter{sujetncaallrefETA}
\newcounter{sujetncaallrefEXS}
\newcounter{sujetncaallrefMM}
\newcounter{sujetncaallrefRS}
\newcounter{sujetncaallrefCC}
\newcounter{sujetncaallrefHV}
\newcounter{sujetncaallrefHL}
\newcounter{sujetncaallrefHQ}

%%%%%%%%Réflexion philosophique%%%%%%%%%
\newcounter{nosujetncarefphilosophique}
\newcounter{sujetncaHQrefphilosophique}
\newcounter{sujetncaRSrefphilosophique}
%%%sujetncas simples%%%
\newcounter{sujetncarefphilosophiqueEXS}
\newcounter{sujetncarefphilosophiqueETA}
\newcounter{sujetncarefphilosophiqueMM}
\newcounter{sujetncarefphilosophiqueHV}
\newcounter{sujetncarefphilosophiqueCC}
\newcounter{sujetncarefphilosophiqueHL}
%%%sujetncas mixtes%%%
\newcounter{sujetncarefphilosophiqueEXSTA}
\newcounter{sujetncarefphilosophiqueEXSMM}
\newcounter{sujetncarefphilosophiqueETAMM}
\newcounter{sujetncarefphilosophiqueOdS}
\newcounter{sujetncarefphilosophiqueHVL}
\newcounter{sujetncarefphilosophiqueHVCC}
\newcounter{sujetncarefphilosophiqueHLCC}
\newcounter{sujetncarefphilosophiqueOdH}
%%%sujetncas mixtes semestres%%%
\newcounter{sujetncarefphilosophiqueETACC}
\newcounter{sujetncarefphilosophiqueETAHV}
\newcounter{sujetncarefphilosophiqueETAHL}
\newcounter{sujetncarefphilosophiqueEXSCC}
\newcounter{sujetncarefphilosophiqueEXSHV}
\newcounter{sujetncarefphilosophiqueEXSHL}
\newcounter{sujetncarefphilosophiqueMMCC}
\newcounter{sujetncarefphilosophiqueMMHV}
\newcounter{sujetncarefphilosophiqueMMHL}
%%%sujetncas mixtes années%%
\newcounter{sujetncarefphilosophiquepremterm}
\newcounter{sujetncarefphilosophiquepremtermRS}
\newcounter{sujetncarefphilosophiquepremtermHQ}
\newcounter{sujetncarefphilosophiqueETAPdP}
\newcounter{sujetncarefphilosophiqueETADFI}
\newcounter{sujetncarefphilosophiqueETADMRC}
\newcounter{sujetncarefphilosophiqueETAHA}\newcounter{sujetncarefphilosophiquepremtermETA}
\newcounter{sujetncarefphilosophiqueEXSPdP}
\newcounter{sujetncarefphilosophiqueEXSDFI}
\newcounter{sujetncarefphilosophiqueEXSDMRC}
\newcounter{sujetncarefphilosophiqueEXSHA}
\newcounter{sujetncarefphilosophiquepremtermEXS}
\newcounter{sujetncarefphilosophiqueMMPdP}
\newcounter{sujetncarefphilosophiqueMMDFI}
\newcounter{sujetncarefphilosophiqueMMDMRC}
\newcounter{sujetncarefphilosophiqueMMHA}
\newcounter{sujetncarefphilosophiquepremtermMM}
\newcounter{sujetncarefphilosophiqueCCPdP}
\newcounter{sujetncarefphilosophiqueCCDFI}
\newcounter{sujetncarefphilosophiqueCCDMRC}
\newcounter{sujetncarefphilosophiqueCCHA}
\newcounter{sujetncarefphilosophiquepremtermCC}
\newcounter{sujetncarefphilosophiqueHVPdP}
\newcounter{sujetncarefphilosophiqueHVDFI}
\newcounter{sujetncarefphilosophiqueHVDMRC}
\newcounter{sujetncarefphilosophiqueHVHA}
\newcounter{sujetncarefphilosophiquepremtermHV}
\newcounter{sujetncarefphilosophiqueHLPdP}
\newcounter{sujetncarefphilosophiqueHLDFI}
\newcounter{sujetncarefphilosophiqueHLDMRC}
\newcounter{sujetncarefphilosophiqueHLHA}
\newcounter{sujetncarefphilosophiquepremtermHL}
\newcounter{sujetncarefphilosophiqueRSPdP}
\newcounter{sujetncarefphilosophiqueRSRdM}
\newcounter{sujetncarefphilosophiqueRSDFI}
\newcounter{sujetncarefphilosophiqueRSDMRC}
\newcounter{sujetncarefphilosophiqueRSHA}
\newcounter{sujetncarefphilosophiqueHQPdP}
\newcounter{sujetncarefphilosophiqueHQRdM}
\newcounter{sujetncarefphilosophiqueHQDFI}
\newcounter{sujetncarefphilosophiqueHQDMRC}
\newcounter{sujetncarefphilosophiqueHQHA}
%%%Somme des compteurs%%%
\newcounter{sujetncaallrefphilosophiqueETA}
\newcounter{sujetncaallrefphilosophiqueEXS}
\newcounter{sujetncaallrefphilosophiqueMM}
\newcounter{sujetncaallrefphilosophiqueRS}
\newcounter{sujetncaallrefphilosophiqueCC}
\newcounter{sujetncaallrefphilosophiqueHV}
\newcounter{sujetncaallrefphilosophiqueHL}
\newcounter{sujetncaallrefphilosophiqueHQ}

%%%%%%%%%Réflexion littéraire%%%%%%%%%%%
\newcounter{nosujetncareflitteraire}
\newcounter{sujetncaHQreflitteraire}
\newcounter{sujetncaRSreflitteraire}
%%%sujetncas simples%%%
\newcounter{sujetncareflitteraireEXS}
\newcounter{sujetncareflitteraireETA}
\newcounter{sujetncareflitteraireMM}
\newcounter{sujetncareflitteraireHV}
\newcounter{sujetncareflitteraireCC}
\newcounter{sujetncareflitteraireHL}
%%%sujetncas mixtes%%%
\newcounter{sujetncareflitteraireEXSTA}
\newcounter{sujetncareflitteraireEXSMM}
\newcounter{sujetncareflitteraireETAMM}
\newcounter{sujetncareflitteraireOdS}
\newcounter{sujetncareflitteraireHVL}
\newcounter{sujetncareflitteraireHVCC}
\newcounter{sujetncareflitteraireHLCC}
\newcounter{sujetncareflitteraireOdH}
%%%sujetncas mixtes semestres%%%
\newcounter{sujetncareflitteraireETACC}
\newcounter{sujetncareflitteraireETAHV}
\newcounter{sujetncareflitteraireETAHL}
\newcounter{sujetncareflitteraireEXSCC}
\newcounter{sujetncareflitteraireEXSHV}
\newcounter{sujetncareflitteraireEXSHL}
\newcounter{sujetncareflitteraireMMCC}
\newcounter{sujetncareflitteraireMMHV}
\newcounter{sujetncareflitteraireMMHL}
%%%sujetncas mixtes années%%
\newcounter{sujetncareflitterairepremterm}
\newcounter{sujetncareflitterairepremtermRS}
\newcounter{sujetncareflitterairepremtermHQ}
\newcounter{sujetncareflitteraireETAPdP}
\newcounter{sujetncareflitteraireETADFI}
\newcounter{sujetncareflitteraireETADMRC}
\newcounter{sujetncareflitteraireETAHA}
\newcounter{sujetncareflitterairepremtermETA}
\newcounter{sujetncareflitteraireEXSPdP}
\newcounter{sujetncareflitteraireEXSDFI}
\newcounter{sujetncareflitteraireEXSDMRC}
\newcounter{sujetncareflitteraireEXSHA}
\newcounter{sujetncareflitterairepremtermEXS}
\newcounter{sujetncareflitteraireMMPdP}
\newcounter{sujetncareflitteraireMMDFI}
\newcounter{sujetncareflitteraireMMDMRC}
\newcounter{sujetncareflitteraireMMHA}
\newcounter{sujetncareflitterairepremtermMM}
\newcounter{sujetncareflitteraireCCPdP}
\newcounter{sujetncareflitteraireCCDFI}
\newcounter{sujetncareflitteraireCCDMRC}
\newcounter{sujetncareflitteraireCCHA}
\newcounter{sujetncareflitterairepremtermCC}
\newcounter{sujetncareflitteraireHVPdP}
\newcounter{sujetncareflitteraireHVDFI}
\newcounter{sujetncareflitteraireHVDMRC}
\newcounter{sujetncareflitteraireHVHA}
\newcounter{sujetncareflitterairepremtermHV}
\newcounter{sujetncareflitteraireHLPdP}
\newcounter{sujetncareflitteraireHLDFI}
\newcounter{sujetncareflitteraireHLDMRC}
\newcounter{sujetncareflitteraireHLHA}
\newcounter{sujetncareflitterairepremtermHL}
\newcounter{sujetncareflitteraireRSPdP}
\newcounter{sujetncareflitteraireRSRdM}
\newcounter{sujetncareflitteraireRSDFI}
\newcounter{sujetncareflitteraireRSDMRC}
\newcounter{sujetncareflitteraireRSHA}
\newcounter{sujetncareflitteraireHQPdP}
\newcounter{sujetncareflitteraireHQRdM}
\newcounter{sujetncareflitteraireHQDFI}
\newcounter{sujetncareflitteraireHQDMRC}
\newcounter{sujetncareflitteraireHQHA}
%%%Somme des compteurs%%%
\newcounter{sujetncaallreflitteraireETA}
\newcounter{sujetncaallreflitteraireEXS}
\newcounter{sujetncaallreflitteraireMM}
\newcounter{sujetncaallreflitteraireCC}
\newcounter{sujetncaallreflitteraireRS}
\newcounter{sujetncaallreflitteraireHV}
\newcounter{sujetncaallreflitteraireHL}
\newcounter{sujetncaallreflitteraireHQ}
%%%Compteurs de statistiques globales sur les genres de question%%%
\newcounter{sujetncaexclRS}
\newcounter{sujetncamixtRS}
\newcounter{sujetncaexclHQ}
\newcounter{sujetncamixtHQ}
\newcounter{sujetncamixtHQRS}
%%%Compteurs de statistiques sur les transversaux (divers semestres, première et terminale)
\newcounter{sujettransncaintlitteraire}
\newcounter{sujettransncareflitteraire}
\newcounter{sujettransncaintphilosophique}
\newcounter{sujettransncarefphilosophique}
%%% Compteurs du nombre de sujets par type
\newcounter{sujetncaintlitteraire}
\newcounter{sujetncareflitteraire}
\newcounter{sujetncaintphilosophique}
\newcounter{sujetncarefphilosophique}
%%% Compteurs par discipline
%%Litt
\newcounter{sujetncalitteraireETA}
\newcounter{sujetncalitteraireEXS}
\newcounter{sujetncalitteraireMM}
\newcounter{sujetncalitteraireCC}
\newcounter{sujetncalitteraireHV}
\newcounter{sujetncalitteraireHL}
\newcounter{sujetncalitteraireEXSTA}
\newcounter{sujetncalitteraireETAMM}
\newcounter{sujetncalitteraireEXSMM}
\newcounter{sujetncalitteraireHVCC}
\newcounter{sujetncalitteraireHVL}
\newcounter{sujetncalitteraireaut}
%%Phil
\newcounter{sujetncaphilosophiqueETA}
\newcounter{sujetncaphilosophiqueEXS}
\newcounter{sujetncaphilosophiqueMM}
\newcounter{sujetncaphilosophiqueCC}
\newcounter{sujetncaphilosophiqueHV}
\newcounter{sujetncaphilosophiqueHL}
\newcounter{sujetncaphilosophiqueEXSTA}
\newcounter{sujetncaphilosophiqueETAMM}
\newcounter{sujetncaphilosophiqueEXSMM}
\newcounter{sujetncaphilosophiqueHVCC}
\newcounter{sujetncaphilosophiqueHVL}
\newcounter{sujetncaphilosophiqueaut}
%%Totaux pour les graphiques
\newcounter{sujetncalitterairetotal}
\newcounter{sujetncaphilosophiquetotal}
%%% Autres (Ods et Odh) comptés ensemble
\newcounter{sujetncaallaut}
%%%%%%%%%%%%%%%%%%%%%%%%%%%%%%%%%%%%%%%%%%%%%%%%%%%%%%%%%%%%%%%%%%%%%%%
%%%%%%%%%%%%%%%%%%%%%%%%%%%%%%Compteurs%%%%%%%%%%%%%%%%%%%%%%%%%%%%%%%%
%%%%%%%%%%%%%%%%%%%%%%%%%%%%%%%%%%%%%%%%%%%%%%%%%%%%%%%%%%%%%%%%%%%%%%%
%%%Compter les sujets par thème%%%
%\newcounter{nosujetpol}
\newcounter{sujetpolall}%Total de questions pour le centre
\newcounter{sujetpoltotal}%Total des questions en comptant en double les questions pour les graphiques
\newcounter{sujetpolHQ}
\newcounter{sujetpolRS}
%%%sujetpols simples%%%
\newcounter{sujetpolEXS}
\newcounter{sujetpolETA}
\newcounter{sujetpolMM}
\newcounter{sujetpolHV}
\newcounter{sujetpolCC}
\newcounter{sujetpolHL}
%%%sujetpols mixtes%%%
\newcounter{sujetpolEXSTA}
\newcounter{sujetpolEXSMM}
\newcounter{sujetpolETAMM}
\newcounter{sujetpolOdS}
\newcounter{sujetpolHVL}
\newcounter{sujetpolHVCC}
\newcounter{sujetpolHLCC}
\newcounter{sujetpolOdH}
%%%sujetpols mixtes semestres%%%
\newcounter{sujetpolETACC}
\newcounter{sujetpolETAHV}
\newcounter{sujetpolETAHL}
\newcounter{sujetpolEXSCC}
\newcounter{sujetpolEXSHV}
\newcounter{sujetpolEXSHL}
\newcounter{sujetpolMMCC}
\newcounter{sujetpolMMHV}
\newcounter{sujetpolMMHL}

%%%sujetpols mixtes années%%
\newcounter{sujetpolpremterm}%%%Compteur général de sujetpol
\newcounter{sujetpolETAPdP}
\newcounter{sujetpolETADFI}
\newcounter{sujetpolETADMRC}
\newcounter{sujetpolETAHA}
\newcounter{sujetpolpremtermETA}
\newcounter{sujetpolEXSPdP}
\newcounter{sujetpolEXSDFI}
\newcounter{sujetpolEXSDMRC}
\newcounter{sujetpolEXSHA}
\newcounter{sujetpolpremtermEXS}
\newcounter{sujetpolMMPdP}
\newcounter{sujetpolMMDFI}
\newcounter{sujetpolMMDMRC}
\newcounter{sujetpolMMHA}
\newcounter{sujetpolpremtermMM}
\newcounter{sujetpolCCPdP}
\newcounter{sujetpolCCDFI}
\newcounter{sujetpolCCDMRC}
\newcounter{sujetpolCCHA}
\newcounter{sujetpolpremtermCC}
\newcounter{sujetpolHVPdP}
\newcounter{sujetpolHVDFI}
\newcounter{sujetpolHVDMRC}
\newcounter{sujetpolHVHA}
\newcounter{sujetpolpremtermHV}
\newcounter{sujetpolHLPdP}
\newcounter{sujetpolHLDFI}
\newcounter{sujetpolHLDMRC}
\newcounter{sujetpolHLHA}
\newcounter{sujetpolpremtermHL}
\newcounter{sujetpolRSPdP}
\newcounter{sujetpolRSRdM}
\newcounter{sujetpolRSDFI}
\newcounter{sujetpolRSDMRC}
\newcounter{sujetpolRSHA}
\newcounter{sujetpolHQPdP}
\newcounter{sujetpolHQRdM}
\newcounter{sujetpolHQDFI}
\newcounter{sujetpolHQDMRC}
\newcounter{sujetpolHQHA}

%%%Somme des compteurs%%%
\newcounter{sujetpolallETA}
\newcounter{sujetpolallEXS}
\newcounter{sujetpolallMM}
\newcounter{sujetpolallRS}
\newcounter{sujetpolallCC}
\newcounter{sujetpolallHV}
\newcounter{sujetpolallHL}
\newcounter{sujetpolallHQ}

%%%Compter les sujetpols par thème et types de questions%%%
%%%%%%%%%%%%%%%%%%%%%%%%%%%%%%%%%%%%%%%%%%%%%
%%%%%%%%%Question d'interprétation%%%%%%%%%%%
%%%%%%%%%%%%%%%%%%%%%%%%%%%%%%%%%%%%%%%%%%%%%

%%%%%%%%%%%%%%%Général%%%%%%%%%%%%%%%%%%%%%%%
\newcounter{nosujetpolint}
\newcounter{sujetpolHQint}
\newcounter{sujetpolRSint}
%%%sujetpols simples%%%
\newcounter{sujetpolintEXS}
\newcounter{sujetpolintETA}
\newcounter{sujetpolintMM}
\newcounter{sujetpolintHV}
\newcounter{sujetpolintCC}
\newcounter{sujetpolintHL}
%%%sujetpols mixtes%%%
\newcounter{sujetpolintEXSTA}
\newcounter{sujetpolintEXSMM}
\newcounter{sujetpolintETAMM}
\newcounter{sujetpolintOdS}
\newcounter{sujetpolintHVL}
\newcounter{sujetpolintHVCC}
\newcounter{sujetpolintHLCC}
\newcounter{sujetpolintOdH}
%%%sujetpols mixtes semestres%%%
\newcounter{sujetpolintETACC}
\newcounter{sujetpolintETAHV}
\newcounter{sujetpolintETAHL}
\newcounter{sujetpolintEXSCC}
\newcounter{sujetpolintEXSHV}
\newcounter{sujetpolintEXSHL}
\newcounter{sujetpolintMMCC}
\newcounter{sujetpolintMMHV}
\newcounter{sujetpolintMMHL}
%%%sujetpols mixtes années%%
\newcounter{sujetpolintpremterm}
\newcounter{sujetpolintpremtermRS}
\newcounter{sujetpolintpremtermHQ}
\newcounter{sujetpolintETAPdP}
\newcounter{sujetpolintETADFI}
\newcounter{sujetpolintETADMRC}
\newcounter{sujetpolintETAHA}
\newcounter{sujetpolintpremtermETA}
\newcounter{sujetpolintEXSPdP}
\newcounter{sujetpolintEXSDFI}
\newcounter{sujetpolintEXSDMRC}
\newcounter{sujetpolintEXSHA}
\newcounter{sujetpolintpremtermEXS}
\newcounter{sujetpolintMMPdP}
\newcounter{sujetpolintMMDFI}
\newcounter{sujetpolintMMDMRC}
\newcounter{sujetpolintMMHA}
\newcounter{sujetpolintpremtermMM}
\newcounter{sujetpolintCCPdP}
\newcounter{sujetpolintCCDFI}
\newcounter{sujetpolintCCDMRC}
\newcounter{sujetpolintCCHA}
\newcounter{sujetpolintpremtermCC}
\newcounter{sujetpolintHVPdP}
\newcounter{sujetpolintHVDFI}
\newcounter{sujetpolintHVDMRC}
\newcounter{sujetpolintHVHA}
\newcounter{sujetpolintpremtermHV}
\newcounter{sujetpolintHLPdP}
\newcounter{sujetpolintHLDFI}
\newcounter{sujetpolintHLDMRC}
\newcounter{sujetpolintHLHA}
\newcounter{sujetpolintpremtermHL}
\newcounter{sujetpolintRSPdP}
\newcounter{sujetpolintRSRdM}
\newcounter{sujetpolintRSDFI}
\newcounter{sujetpolintRSDMRC}
\newcounter{sujetpolintRSHA}
\newcounter{sujetpolintHQPdP}
\newcounter{sujetpolintHQRdM}
\newcounter{sujetpolintHQDFI}
\newcounter{sujetpolintHQDMRC}
\newcounter{sujetpolintHQHA}
%%%Somme des compteurs%%%
\newcounter{sujetpolallintETA}
\newcounter{sujetpolallintEXS}
\newcounter{sujetpolallintMM}
\newcounter{sujetpolallintRS}
\newcounter{sujetpolallintCC}
\newcounter{sujetpolallintHV}
\newcounter{sujetpolallintHL}
\newcounter{sujetpolallintHQ}

%%%%%%%%Interprétation philosophique%%%%%%%%%
\newcounter{nosujetpolintphilosophique}
\newcounter{sujetpolHQintphilosophique}
\newcounter{sujetpolRSintphilosophique}
%%%sujetpols simples%%%
\newcounter{sujetpolintphilosophiqueEXS}
\newcounter{sujetpolintphilosophiqueETA}
\newcounter{sujetpolintphilosophiqueMM}
\newcounter{sujetpolintphilosophiqueHV}
\newcounter{sujetpolintphilosophiqueCC}
\newcounter{sujetpolintphilosophiqueHL}
%%%sujetpols mixtes%%%
\newcounter{sujetpolintphilosophiqueEXSTA}
\newcounter{sujetpolintphilosophiqueEXSMM}
\newcounter{sujetpolintphilosophiqueETAMM}
\newcounter{sujetpolintphilosophiqueOdS}
\newcounter{sujetpolintphilosophiqueHVL}
\newcounter{sujetpolintphilosophiqueHVCC}
\newcounter{sujetpolintphilosophiqueHLCC}
\newcounter{sujetpolintphilosophiqueOdH}
%%%sujetpols mixtes semestres%%%
\newcounter{sujetpolintphilosophiqueETACC}
\newcounter{sujetpolintphilosophiqueETAHV}
\newcounter{sujetpolintphilosophiqueETAHL}
\newcounter{sujetpolintphilosophiqueEXSCC}
\newcounter{sujetpolintphilosophiqueEXSHV}
\newcounter{sujetpolintphilosophiqueEXSHL}
\newcounter{sujetpolintphilosophiqueMMCC}
\newcounter{sujetpolintphilosophiqueMMHV}
\newcounter{sujetpolintphilosophiqueMMHL}
%%%sujetpols mixtes années%%
\newcounter{sujetpolintphilosophiquepremterm}
\newcounter{sujetpolintphilosophiquepremtermRS}
\newcounter{sujetpolintphilosophiquepremtermHQ}
\newcounter{sujetpolintphilosophiqueETAPdP}
\newcounter{sujetpolintphilosophiqueETADFI}
\newcounter{sujetpolintphilosophiqueETADMRC}
\newcounter{sujetpolintphilosophiqueETAHA}
\newcounter{sujetpolintphilosophiquepremtermETA}
\newcounter{sujetpolintphilosophiqueEXSPdP}
\newcounter{sujetpolintphilosophiqueEXSDFI}
\newcounter{sujetpolintphilosophiqueEXSDMRC}
\newcounter{sujetpolintphilosophiqueEXSHA}
\newcounter{sujetpolintphilosophiquepremtermEXS}
\newcounter{sujetpolintphilosophiqueMMPdP}
\newcounter{sujetpolintphilosophiqueMMDFI}
\newcounter{sujetpolintphilosophiqueMMDMRC}
\newcounter{sujetpolintphilosophiqueMMHA}
\newcounter{sujetpolintphilosophiquepremtermMM}
\newcounter{sujetpolintphilosophiqueCCPdP}
\newcounter{sujetpolintphilosophiqueCCDFI}
\newcounter{sujetpolintphilosophiqueCCDMRC}
\newcounter{sujetpolintphilosophiqueCCHA}
\newcounter{sujetpolintphilosophiquepremtermCC}
\newcounter{sujetpolintphilosophiqueHVPdP}
\newcounter{sujetpolintphilosophiqueHVDFI}
\newcounter{sujetpolintphilosophiqueHVDMRC}
\newcounter{sujetpolintphilosophiqueHVHA}
\newcounter{sujetpolintphilosophiquepremtermHV}
\newcounter{sujetpolintphilosophiqueHLPdP}
\newcounter{sujetpolintphilosophiqueHLDFI}
\newcounter{sujetpolintphilosophiqueHLDMRC}
\newcounter{sujetpolintphilosophiqueHLHA}
\newcounter{sujetpolintphilosophiquepremtermHL}
\newcounter{sujetpolintphilosophiqueRSPdP}
\newcounter{sujetpolintphilosophiqueRSRdM}
\newcounter{sujetpolintphilosophiqueRSDFI}
\newcounter{sujetpolintphilosophiqueRSDMRC}
\newcounter{sujetpolintphilosophiqueRSHA}
\newcounter{sujetpolintphilosophiqueHQPdP}
\newcounter{sujetpolintphilosophiqueHQRdM}
\newcounter{sujetpolintphilosophiqueHQDFI}
\newcounter{sujetpolintphilosophiqueHQDMRC}
\newcounter{sujetpolintphilosophiqueHQHA}
%%%Somme des compteurs%%%
\newcounter{sujetpolallintphilosophiqueETA}
\newcounter{sujetpolallintphilosophiqueEXS}
\newcounter{sujetpolallintphilosophiqueMM}
\newcounter{sujetpolallintphilosophiqueRS}
\newcounter{sujetpolallintphilosophiqueCC}
\newcounter{sujetpolallintphilosophiqueHV}
\newcounter{sujetpolallintphilosophiqueHL}
\newcounter{sujetpolallintphilosophiqueHQ}

%%%%%%%%%Interprétation littéraire%%%%%%%%%%%
\newcounter{nosujetpolintlitteraire}
\newcounter{sujetpolHQintlitteraire}
\newcounter{sujetpolRSintlitteraire}
%%%sujetpols simples%%%
\newcounter{sujetpolintlitteraireEXS}
\newcounter{sujetpolintlitteraireETA}
\newcounter{sujetpolintlitteraireMM}
\newcounter{sujetpolintlitteraireHV}
\newcounter{sujetpolintlitteraireCC}
\newcounter{sujetpolintlitteraireHL}
%%%sujetpols mixtes%%%
\newcounter{sujetpolintlitteraireEXSTA}
\newcounter{sujetpolintlitteraireEXSMM}
\newcounter{sujetpolintlitteraireETAMM}
\newcounter{sujetpolintlitteraireOdS}
\newcounter{sujetpolintlitteraireHVL}
\newcounter{sujetpolintlitteraireHVCC}
\newcounter{sujetpolintlitteraireHLCC}
\newcounter{sujetpolintlitteraireOdH}
%%%sujetpols mixtes semestres%%%
\newcounter{sujetpolintlitteraireETACC}
\newcounter{sujetpolintlitteraireETAHV}
\newcounter{sujetpolintlitteraireETAHL}
\newcounter{sujetpolintlitteraireEXSCC}
\newcounter{sujetpolintlitteraireEXSHV}
\newcounter{sujetpolintlitteraireEXSHL}
\newcounter{sujetpolintlitteraireMMCC}
\newcounter{sujetpolintlitteraireMMHV}
\newcounter{sujetpolintlitteraireMMHL}
%%%sujetpols mixtes années%%
\newcounter{sujetpolintlitterairepremterm}
\newcounter{sujetpolintlitterairepremtermRS}
\newcounter{sujetpolintlitterairepremtermHQ}
\newcounter{sujetpolintlitteraireETAPdP}
\newcounter{sujetpolintlitteraireETADFI}
\newcounter{sujetpolintlitteraireETADMRC}
\newcounter{sujetpolintlitteraireETAHA}
\newcounter{sujetpolintlitterairepremtermETA}
\newcounter{sujetpolintlitteraireEXSPdP}
\newcounter{sujetpolintlitteraireEXSDFI}
\newcounter{sujetpolintlitteraireEXSDMRC}
\newcounter{sujetpolintlitteraireEXSHA}
\newcounter{sujetpolintlitterairepremtermEXS}
\newcounter{sujetpolintlitteraireMMPdP}
\newcounter{sujetpolintlitteraireMMDFI}
\newcounter{sujetpolintlitteraireMMDMRC}
\newcounter{sujetpolintlitteraireMMHA}
\newcounter{sujetpolintlitterairepremtermMM}
\newcounter{sujetpolintlitteraireCCPdP}
\newcounter{sujetpolintlitteraireCCDFI}
\newcounter{sujetpolintlitteraireCCDMRC}
\newcounter{sujetpolintlitteraireCCHA}
\newcounter{sujetpolintlitterairepremtermCC}
\newcounter{sujetpolintlitteraireHVPdP}
\newcounter{sujetpolintlitteraireHVDFI}
\newcounter{sujetpolintlitteraireHVDMRC}
\newcounter{sujetpolintlitteraireHVHA}
\newcounter{sujetpolintlitterairepremtermHV}
\newcounter{sujetpolintlitteraireHLPdP}
\newcounter{sujetpolintlitteraireHLDFI}
\newcounter{sujetpolintlitteraireHLDMRC}
\newcounter{sujetpolintlitteraireHLHA}
\newcounter{sujetpolintlitterairepremtermHL}
\newcounter{sujetpolintlitteraireRSPdP}
\newcounter{sujetpolintlitteraireRSRdM}
\newcounter{sujetpolintlitteraireRSDFI}
\newcounter{sujetpolintlitteraireRSDMRC}
\newcounter{sujetpolintlitteraireRSHA}
\newcounter{sujetpolintlitteraireHQPdP}
\newcounter{sujetpolintlitteraireHQRdM}
\newcounter{sujetpolintlitteraireHQDFI}
\newcounter{sujetpolintlitteraireHQDMRC}
\newcounter{sujetpolintlitteraireHQHA}
%%%Somme des compteurs%%%
\newcounter{sujetpolallintlitteraireETA}
\newcounter{sujetpolallintlitteraireEXS}
\newcounter{sujetpolallintlitteraireMM}
\newcounter{sujetpolallintlitteraireRS}
\newcounter{sujetpolallintlitteraireCC}
\newcounter{sujetpolallintlitteraireHV}
\newcounter{sujetpolallintlitteraireHL}
\newcounter{sujetpolallintlitteraireHQ}

%%%%%%%%%%%%%%%%%%%%%%%%%%%%%%%%%%%%%%%%%%%%%
%%%%%%%%%%%%Question de réflexion%%%%%%%%%%%%
%%%%%%%%%%%%%%%%%%%%%%%%%%%%%%%%%%%%%%%%%%%%%
\newcounter{nosujetpolref}
%%%sujetpols simples%%%
\newcounter{sujetpolrefEXS}
\newcounter{sujetpolrefETA}
\newcounter{sujetpolrefMM}
\newcounter{sujetpolrefHV}
\newcounter{sujetpolrefCC}
\newcounter{sujetpolrefHL}
%%%sujetpols mixtes%%%
\newcounter{sujetpolrefEXSTA}
\newcounter{sujetpolrefEXSMM}
\newcounter{sujetpolrefETAMM}
\newcounter{sujetpolrefOdS}
\newcounter{sujetpolrefHVL}
\newcounter{sujetpolrefHVCC}
\newcounter{sujetpolrefHLCC}
\newcounter{sujetpolrefOdH}
%%%sujetpols mixtes semestres%%%
\newcounter{sujetpolrefETACC}
\newcounter{sujetpolrefETAHV}
\newcounter{sujetpolrefETAHL}
\newcounter{sujetpolrefEXSCC}
\newcounter{sujetpolrefEXSHV}
\newcounter{sujetpolrefEXSHL}
\newcounter{sujetpolrefMMCC}
\newcounter{sujetpolrefMMHV}
\newcounter{sujetpolrefMMHL}
%%%sujetpols mixtes années%%
\newcounter{sujetpolrefpremterm}%%%Compteur général
\newcounter{sujetpolrefETAPdP}
\newcounter{sujetpolrefETADFI}
\newcounter{sujetpolrefETADMRC}
\newcounter{sujetpolrefETAHA}
\newcounter{sujetpolrefpremtermETA}
\newcounter{sujetpolrefEXSPdP}
\newcounter{sujetpolrefEXSDFI}
\newcounter{sujetpolrefEXSDMRC}
\newcounter{sujetpolrefEXSHA}
\newcounter{sujetpolrefpremtermEXS}
\newcounter{sujetpolrefMMPdP}
\newcounter{sujetpolrefMMDFI}
\newcounter{sujetpolrefMMDMRC}
\newcounter{sujetpolrefMMHA}
\newcounter{sujetpolrefpremtermMM}
\newcounter{sujetpolrefCCPdP}
\newcounter{sujetpolrefCCDFI}
\newcounter{sujetpolrefCCDMRC}
\newcounter{sujetpolrefCCHA}
\newcounter{sujetpolrefpremtermCC}
\newcounter{sujetpolrefHVPdP}
\newcounter{sujetpolrefHVDFI}
\newcounter{sujetpolrefHVDMRC}
\newcounter{sujetpolrefHVHA}
\newcounter{sujetpolrefpremtermHV}
\newcounter{sujetpolrefHLPdP}
\newcounter{sujetpolrefHLDFI}
\newcounter{sujetpolrefHLDMRC}
\newcounter{sujetpolrefHLHA}
\newcounter{sujetpolrefpremtermHL}
\newcounter{sujetpolrefRSPdP}
\newcounter{sujetpolrefRSRdM}
\newcounter{sujetpolrefRSDFI}
\newcounter{sujetpolrefRSDMRC}
\newcounter{sujetpolrefRSHA}
\newcounter{sujetpolrefHQPdP}
\newcounter{sujetpolrefHQRdM}
\newcounter{sujetpolrefHQDFI}
\newcounter{sujetpolrefHQDMRC}
\newcounter{sujetpolrefHQHA}
%%%Somme des compteurs%%%
\newcounter{sujetpolallrefETA}
\newcounter{sujetpolallrefEXS}
\newcounter{sujetpolallrefMM}
\newcounter{sujetpolallrefRS}
\newcounter{sujetpolallrefCC}
\newcounter{sujetpolallrefHV}
\newcounter{sujetpolallrefHL}
\newcounter{sujetpolallrefHQ}

%%%%%%%%Réflexion philosophique%%%%%%%%%
\newcounter{nosujetpolrefphilosophique}
\newcounter{sujetpolHQrefphilosophique}
\newcounter{sujetpolRSrefphilosophique}
%%%sujetpols simples%%%
\newcounter{sujetpolrefphilosophiqueEXS}
\newcounter{sujetpolrefphilosophiqueETA}
\newcounter{sujetpolrefphilosophiqueMM}
\newcounter{sujetpolrefphilosophiqueHV}
\newcounter{sujetpolrefphilosophiqueCC}
\newcounter{sujetpolrefphilosophiqueHL}
%%%sujetpols mixtes%%%
\newcounter{sujetpolrefphilosophiqueEXSTA}
\newcounter{sujetpolrefphilosophiqueEXSMM}
\newcounter{sujetpolrefphilosophiqueETAMM}
\newcounter{sujetpolrefphilosophiqueOdS}
\newcounter{sujetpolrefphilosophiqueHVL}
\newcounter{sujetpolrefphilosophiqueHVCC}
\newcounter{sujetpolrefphilosophiqueHLCC}
\newcounter{sujetpolrefphilosophiqueOdH}
%%%sujetpols mixtes semestres%%%
\newcounter{sujetpolrefphilosophiqueETACC}
\newcounter{sujetpolrefphilosophiqueETAHV}
\newcounter{sujetpolrefphilosophiqueETAHL}
\newcounter{sujetpolrefphilosophiqueEXSCC}
\newcounter{sujetpolrefphilosophiqueEXSHV}
\newcounter{sujetpolrefphilosophiqueEXSHL}
\newcounter{sujetpolrefphilosophiqueMMCC}
\newcounter{sujetpolrefphilosophiqueMMHV}
\newcounter{sujetpolrefphilosophiqueMMHL}
%%%sujetpols mixtes années%%
\newcounter{sujetpolrefphilosophiquepremterm}
\newcounter{sujetpolrefphilosophiquepremtermRS}
\newcounter{sujetpolrefphilosophiquepremtermHQ}
\newcounter{sujetpolrefphilosophiqueETAPdP}
\newcounter{sujetpolrefphilosophiqueETADFI}
\newcounter{sujetpolrefphilosophiqueETADMRC}
\newcounter{sujetpolrefphilosophiqueETAHA}\newcounter{sujetpolrefphilosophiquepremtermETA}
\newcounter{sujetpolrefphilosophiqueEXSPdP}
\newcounter{sujetpolrefphilosophiqueEXSDFI}
\newcounter{sujetpolrefphilosophiqueEXSDMRC}
\newcounter{sujetpolrefphilosophiqueEXSHA}
\newcounter{sujetpolrefphilosophiquepremtermEXS}
\newcounter{sujetpolrefphilosophiqueMMPdP}
\newcounter{sujetpolrefphilosophiqueMMDFI}
\newcounter{sujetpolrefphilosophiqueMMDMRC}
\newcounter{sujetpolrefphilosophiqueMMHA}
\newcounter{sujetpolrefphilosophiquepremtermMM}
\newcounter{sujetpolrefphilosophiqueCCPdP}
\newcounter{sujetpolrefphilosophiqueCCDFI}
\newcounter{sujetpolrefphilosophiqueCCDMRC}
\newcounter{sujetpolrefphilosophiqueCCHA}
\newcounter{sujetpolrefphilosophiquepremtermCC}
\newcounter{sujetpolrefphilosophiqueHVPdP}
\newcounter{sujetpolrefphilosophiqueHVDFI}
\newcounter{sujetpolrefphilosophiqueHVDMRC}
\newcounter{sujetpolrefphilosophiqueHVHA}
\newcounter{sujetpolrefphilosophiquepremtermHV}
\newcounter{sujetpolrefphilosophiqueHLPdP}
\newcounter{sujetpolrefphilosophiqueHLDFI}
\newcounter{sujetpolrefphilosophiqueHLDMRC}
\newcounter{sujetpolrefphilosophiqueHLHA}
\newcounter{sujetpolrefphilosophiquepremtermHL}
\newcounter{sujetpolrefphilosophiqueRSPdP}
\newcounter{sujetpolrefphilosophiqueRSRdM}
\newcounter{sujetpolrefphilosophiqueRSDFI}
\newcounter{sujetpolrefphilosophiqueRSDMRC}
\newcounter{sujetpolrefphilosophiqueRSHA}
\newcounter{sujetpolrefphilosophiqueHQPdP}
\newcounter{sujetpolrefphilosophiqueHQRdM}
\newcounter{sujetpolrefphilosophiqueHQDFI}
\newcounter{sujetpolrefphilosophiqueHQDMRC}
\newcounter{sujetpolrefphilosophiqueHQHA}
%%%Somme des compteurs%%%
\newcounter{sujetpolallrefphilosophiqueETA}
\newcounter{sujetpolallrefphilosophiqueEXS}
\newcounter{sujetpolallrefphilosophiqueMM}
\newcounter{sujetpolallrefphilosophiqueRS}
\newcounter{sujetpolallrefphilosophiqueCC}
\newcounter{sujetpolallrefphilosophiqueHV}
\newcounter{sujetpolallrefphilosophiqueHL}
\newcounter{sujetpolallrefphilosophiqueHQ}

%%%%%%%%%Réflexion littéraire%%%%%%%%%%%
\newcounter{nosujetpolreflitteraire}
\newcounter{sujetpolHQreflitteraire}
\newcounter{sujetpolRSreflitteraire}
%%%sujetpols simples%%%
\newcounter{sujetpolreflitteraireEXS}
\newcounter{sujetpolreflitteraireETA}
\newcounter{sujetpolreflitteraireMM}
\newcounter{sujetpolreflitteraireHV}
\newcounter{sujetpolreflitteraireCC}
\newcounter{sujetpolreflitteraireHL}
%%%sujetpols mixtes%%%
\newcounter{sujetpolreflitteraireEXSTA}
\newcounter{sujetpolreflitteraireEXSMM}
\newcounter{sujetpolreflitteraireETAMM}
\newcounter{sujetpolreflitteraireOdS}
\newcounter{sujetpolreflitteraireHVL}
\newcounter{sujetpolreflitteraireHVCC}
\newcounter{sujetpolreflitteraireHLCC}
\newcounter{sujetpolreflitteraireOdH}
%%%sujetpols mixtes semestres%%%
\newcounter{sujetpolreflitteraireETACC}
\newcounter{sujetpolreflitteraireETAHV}
\newcounter{sujetpolreflitteraireETAHL}
\newcounter{sujetpolreflitteraireEXSCC}
\newcounter{sujetpolreflitteraireEXSHV}
\newcounter{sujetpolreflitteraireEXSHL}
\newcounter{sujetpolreflitteraireMMCC}
\newcounter{sujetpolreflitteraireMMHV}
\newcounter{sujetpolreflitteraireMMHL}
%%%sujetpols mixtes années%%
\newcounter{sujetpolreflitterairepremterm}
\newcounter{sujetpolreflitterairepremtermRS}
\newcounter{sujetpolreflitterairepremtermHQ}
\newcounter{sujetpolreflitteraireETAPdP}
\newcounter{sujetpolreflitteraireETADFI}
\newcounter{sujetpolreflitteraireETADMRC}
\newcounter{sujetpolreflitteraireETAHA}
\newcounter{sujetpolreflitterairepremtermETA}
\newcounter{sujetpolreflitteraireEXSPdP}
\newcounter{sujetpolreflitteraireEXSDFI}
\newcounter{sujetpolreflitteraireEXSDMRC}
\newcounter{sujetpolreflitteraireEXSHA}
\newcounter{sujetpolreflitterairepremtermEXS}
\newcounter{sujetpolreflitteraireMMPdP}
\newcounter{sujetpolreflitteraireMMDFI}
\newcounter{sujetpolreflitteraireMMDMRC}
\newcounter{sujetpolreflitteraireMMHA}
\newcounter{sujetpolreflitterairepremtermMM}
\newcounter{sujetpolreflitteraireCCPdP}
\newcounter{sujetpolreflitteraireCCDFI}
\newcounter{sujetpolreflitteraireCCDMRC}
\newcounter{sujetpolreflitteraireCCHA}
\newcounter{sujetpolreflitterairepremtermCC}
\newcounter{sujetpolreflitteraireHVPdP}
\newcounter{sujetpolreflitteraireHVDFI}
\newcounter{sujetpolreflitteraireHVDMRC}
\newcounter{sujetpolreflitteraireHVHA}
\newcounter{sujetpolreflitterairepremtermHV}
\newcounter{sujetpolreflitteraireHLPdP}
\newcounter{sujetpolreflitteraireHLDFI}
\newcounter{sujetpolreflitteraireHLDMRC}
\newcounter{sujetpolreflitteraireHLHA}
\newcounter{sujetpolreflitterairepremtermHL}
\newcounter{sujetpolreflitteraireRSPdP}
\newcounter{sujetpolreflitteraireRSRdM}
\newcounter{sujetpolreflitteraireRSDFI}
\newcounter{sujetpolreflitteraireRSDMRC}
\newcounter{sujetpolreflitteraireRSHA}
\newcounter{sujetpolreflitteraireHQPdP}
\newcounter{sujetpolreflitteraireHQRdM}
\newcounter{sujetpolreflitteraireHQDFI}
\newcounter{sujetpolreflitteraireHQDMRC}
\newcounter{sujetpolreflitteraireHQHA}
%%%Somme des compteurs%%%
\newcounter{sujetpolallreflitteraireETA}
\newcounter{sujetpolallreflitteraireEXS}
\newcounter{sujetpolallreflitteraireMM}
\newcounter{sujetpolallreflitteraireCC}
\newcounter{sujetpolallreflitteraireRS}
\newcounter{sujetpolallreflitteraireHV}
\newcounter{sujetpolallreflitteraireHL}
\newcounter{sujetpolallreflitteraireHQ}
%%%Compteurs de statistiques globales sur les genres de question%%%
\newcounter{sujetpolexclRS}
\newcounter{sujetpolmixtRS}
\newcounter{sujetpolexclHQ}
\newcounter{sujetpolmixtHQ}
\newcounter{sujetpolmixtHQRS}
%%%Compteurs de statistiques sur les transversaux (divers semestres, première et terminale)
\newcounter{sujettranspolintlitteraire}
\newcounter{sujettranspolreflitteraire}
\newcounter{sujettranspolintphilosophique}
\newcounter{sujettranspolrefphilosophique}
%%% Compteurs du nombre de sujets par type
\newcounter{sujetpolintlitteraire}
\newcounter{sujetpolreflitteraire}
\newcounter{sujetpolintphilosophique}
\newcounter{sujetpolrefphilosophique}
%%% Compteurs par discipline
%%Litt
\newcounter{sujetpollitteraireETA}
\newcounter{sujetpollitteraireEXS}
\newcounter{sujetpollitteraireMM}
\newcounter{sujetpollitteraireCC}
\newcounter{sujetpollitteraireHV}
\newcounter{sujetpollitteraireHL}
\newcounter{sujetpollitteraireEXSTA}
\newcounter{sujetpollitteraireETAMM}
\newcounter{sujetpollitteraireEXSMM}
\newcounter{sujetpollitteraireHVCC}
\newcounter{sujetpollitteraireHVL}
\newcounter{sujetpollitteraireaut}
%%Phil
\newcounter{sujetpolphilosophiqueETA}
\newcounter{sujetpolphilosophiqueEXS}
\newcounter{sujetpolphilosophiqueMM}
\newcounter{sujetpolphilosophiqueCC}
\newcounter{sujetpolphilosophiqueHV}
\newcounter{sujetpolphilosophiqueHL}
\newcounter{sujetpolphilosophiqueEXSTA}
\newcounter{sujetpolphilosophiqueETAMM}
\newcounter{sujetpolphilosophiqueEXSMM}
\newcounter{sujetpolphilosophiqueHVCC}
\newcounter{sujetpolphilosophiqueHVL}
\newcounter{sujetpolphilosophiqueaut}
%%Totaux pour les graphiques
\newcounter{sujetpollitterairetotal}
\newcounter{sujetpolphilosophiquetotal}
%%% Autres (Ods et Odh) comptés ensemble
\newcounter{sujetpolallaut}
%%%%%%%%%%%%%%%%%%%%%%%%%%%%%%%%%%%%%%%%%%%%%%%%%%%%%%%%%%%%%%%%%%%%%%%
%%%%%%%%%%%%%%%%%%%%%%%%%%%%%%Compteurs%%%%%%%%%%%%%%%%%%%%%%%%%%%%%%%%
%%%%%%%%%%%%%%%%%%%%%%%%%%%%%%%%%%%%%%%%%%%%%%%%%%%%%%%%%%%%%%%%%%%%%%%
%%%Compter les sujets par thème%%%
%\newcounter{nosujetreu}
\newcounter{sujetreuall}%Total de questions pour le centre
\newcounter{sujetreutotal}%Total des questions en comptant en double les questions pour les graphiques
\newcounter{sujetreuHQ}
\newcounter{sujetreuRS}
%%%sujetreus simples%%%
\newcounter{sujetreuEXS}
\newcounter{sujetreuETA}
\newcounter{sujetreuMM}
\newcounter{sujetreuHV}
\newcounter{sujetreuCC}
\newcounter{sujetreuHL}
%%%sujetreus mixtes%%%
\newcounter{sujetreuEXSTA}
\newcounter{sujetreuEXSMM}
\newcounter{sujetreuETAMM}
\newcounter{sujetreuOdS}
\newcounter{sujetreuHVL}
\newcounter{sujetreuHVCC}
\newcounter{sujetreuHLCC}
\newcounter{sujetreuOdH}
%%%sujetreus mixtes semestres%%%
\newcounter{sujetreuETACC}
\newcounter{sujetreuETAHV}
\newcounter{sujetreuETAHL}
\newcounter{sujetreuEXSCC}
\newcounter{sujetreuEXSHV}
\newcounter{sujetreuEXSHL}
\newcounter{sujetreuMMCC}
\newcounter{sujetreuMMHV}
\newcounter{sujetreuMMHL}

%%%sujetreus mixtes années%%
\newcounter{sujetreupremterm}%%%Compteur général de sujetreu
\newcounter{sujetreuETAPdP}
\newcounter{sujetreuETADFI}
\newcounter{sujetreuETADMRC}
\newcounter{sujetreuETAHA}
\newcounter{sujetreupremtermETA}
\newcounter{sujetreuEXSPdP}
\newcounter{sujetreuEXSDFI}
\newcounter{sujetreuEXSDMRC}
\newcounter{sujetreuEXSHA}
\newcounter{sujetreupremtermEXS}
\newcounter{sujetreuMMPdP}
\newcounter{sujetreuMMDFI}
\newcounter{sujetreuMMDMRC}
\newcounter{sujetreuMMHA}
\newcounter{sujetreupremtermMM}
\newcounter{sujetreuCCPdP}
\newcounter{sujetreuCCDFI}
\newcounter{sujetreuCCDMRC}
\newcounter{sujetreuCCHA}
\newcounter{sujetreupremtermCC}
\newcounter{sujetreuHVPdP}
\newcounter{sujetreuHVDFI}
\newcounter{sujetreuHVDMRC}
\newcounter{sujetreuHVHA}
\newcounter{sujetreupremtermHV}
\newcounter{sujetreuHLPdP}
\newcounter{sujetreuHLDFI}
\newcounter{sujetreuHLDMRC}
\newcounter{sujetreuHLHA}
\newcounter{sujetreupremtermHL}
\newcounter{sujetreuRSPdP}
\newcounter{sujetreuRSRdM}
\newcounter{sujetreuRSDFI}
\newcounter{sujetreuRSDMRC}
\newcounter{sujetreuRSHA}
\newcounter{sujetreuHQPdP}
\newcounter{sujetreuHQRdM}
\newcounter{sujetreuHQDFI}
\newcounter{sujetreuHQDMRC}
\newcounter{sujetreuHQHA}

%%%Somme des compteurs%%%
\newcounter{sujetreuallETA}
\newcounter{sujetreuallEXS}
\newcounter{sujetreuallMM}
\newcounter{sujetreuallRS}
\newcounter{sujetreuallCC}
\newcounter{sujetreuallHV}
\newcounter{sujetreuallHL}
\newcounter{sujetreuallHQ}

%%%Compter les sujetreus par thème et types de questions%%%
%%%%%%%%%%%%%%%%%%%%%%%%%%%%%%%%%%%%%%%%%%%%%
%%%%%%%%%Question d'interprétation%%%%%%%%%%%
%%%%%%%%%%%%%%%%%%%%%%%%%%%%%%%%%%%%%%%%%%%%%

%%%%%%%%%%%%%%%Général%%%%%%%%%%%%%%%%%%%%%%%
\newcounter{nosujetreuint}
\newcounter{sujetreuHQint}
\newcounter{sujetreuRSint}
%%%sujetreus simples%%%
\newcounter{sujetreuintEXS}
\newcounter{sujetreuintETA}
\newcounter{sujetreuintMM}
\newcounter{sujetreuintHV}
\newcounter{sujetreuintCC}
\newcounter{sujetreuintHL}
%%%sujetreus mixtes%%%
\newcounter{sujetreuintEXSTA}
\newcounter{sujetreuintEXSMM}
\newcounter{sujetreuintETAMM}
\newcounter{sujetreuintOdS}
\newcounter{sujetreuintHVL}
\newcounter{sujetreuintHVCC}
\newcounter{sujetreuintHLCC}
\newcounter{sujetreuintOdH}
%%%sujetreus mixtes semestres%%%
\newcounter{sujetreuintETACC}
\newcounter{sujetreuintETAHV}
\newcounter{sujetreuintETAHL}
\newcounter{sujetreuintEXSCC}
\newcounter{sujetreuintEXSHV}
\newcounter{sujetreuintEXSHL}
\newcounter{sujetreuintMMCC}
\newcounter{sujetreuintMMHV}
\newcounter{sujetreuintMMHL}
%%%sujetreus mixtes années%%
\newcounter{sujetreuintpremterm}
\newcounter{sujetreuintpremtermRS}
\newcounter{sujetreuintpremtermHQ}
\newcounter{sujetreuintETAPdP}
\newcounter{sujetreuintETADFI}
\newcounter{sujetreuintETADMRC}
\newcounter{sujetreuintETAHA}
\newcounter{sujetreuintpremtermETA}
\newcounter{sujetreuintEXSPdP}
\newcounter{sujetreuintEXSDFI}
\newcounter{sujetreuintEXSDMRC}
\newcounter{sujetreuintEXSHA}
\newcounter{sujetreuintpremtermEXS}
\newcounter{sujetreuintMMPdP}
\newcounter{sujetreuintMMDFI}
\newcounter{sujetreuintMMDMRC}
\newcounter{sujetreuintMMHA}
\newcounter{sujetreuintpremtermMM}
\newcounter{sujetreuintCCPdP}
\newcounter{sujetreuintCCDFI}
\newcounter{sujetreuintCCDMRC}
\newcounter{sujetreuintCCHA}
\newcounter{sujetreuintpremtermCC}
\newcounter{sujetreuintHVPdP}
\newcounter{sujetreuintHVDFI}
\newcounter{sujetreuintHVDMRC}
\newcounter{sujetreuintHVHA}
\newcounter{sujetreuintpremtermHV}
\newcounter{sujetreuintHLPdP}
\newcounter{sujetreuintHLDFI}
\newcounter{sujetreuintHLDMRC}
\newcounter{sujetreuintHLHA}
\newcounter{sujetreuintpremtermHL}
\newcounter{sujetreuintRSPdP}
\newcounter{sujetreuintRSRdM}
\newcounter{sujetreuintRSDFI}
\newcounter{sujetreuintRSDMRC}
\newcounter{sujetreuintRSHA}
\newcounter{sujetreuintHQPdP}
\newcounter{sujetreuintHQRdM}
\newcounter{sujetreuintHQDFI}
\newcounter{sujetreuintHQDMRC}
\newcounter{sujetreuintHQHA}
%%%Somme des compteurs%%%
\newcounter{sujetreuallintETA}
\newcounter{sujetreuallintEXS}
\newcounter{sujetreuallintMM}
\newcounter{sujetreuallintRS}
\newcounter{sujetreuallintCC}
\newcounter{sujetreuallintHV}
\newcounter{sujetreuallintHL}
\newcounter{sujetreuallintHQ}

%%%%%%%%Interprétation philosophique%%%%%%%%%
\newcounter{nosujetreuintphilosophique}
\newcounter{sujetreuHQintphilosophique}
\newcounter{sujetreuRSintphilosophique}
%%%sujetreus simples%%%
\newcounter{sujetreuintphilosophiqueEXS}
\newcounter{sujetreuintphilosophiqueETA}
\newcounter{sujetreuintphilosophiqueMM}
\newcounter{sujetreuintphilosophiqueHV}
\newcounter{sujetreuintphilosophiqueCC}
\newcounter{sujetreuintphilosophiqueHL}
%%%sujetreus mixtes%%%
\newcounter{sujetreuintphilosophiqueEXSTA}
\newcounter{sujetreuintphilosophiqueEXSMM}
\newcounter{sujetreuintphilosophiqueETAMM}
\newcounter{sujetreuintphilosophiqueOdS}
\newcounter{sujetreuintphilosophiqueHVL}
\newcounter{sujetreuintphilosophiqueHVCC}
\newcounter{sujetreuintphilosophiqueHLCC}
\newcounter{sujetreuintphilosophiqueOdH}
%%%sujetreus mixtes semestres%%%
\newcounter{sujetreuintphilosophiqueETACC}
\newcounter{sujetreuintphilosophiqueETAHV}
\newcounter{sujetreuintphilosophiqueETAHL}
\newcounter{sujetreuintphilosophiqueEXSCC}
\newcounter{sujetreuintphilosophiqueEXSHV}
\newcounter{sujetreuintphilosophiqueEXSHL}
\newcounter{sujetreuintphilosophiqueMMCC}
\newcounter{sujetreuintphilosophiqueMMHV}
\newcounter{sujetreuintphilosophiqueMMHL}
%%%sujetreus mixtes années%%
\newcounter{sujetreuintphilosophiquepremterm}
\newcounter{sujetreuintphilosophiquepremtermRS}
\newcounter{sujetreuintphilosophiquepremtermHQ}
\newcounter{sujetreuintphilosophiqueETAPdP}
\newcounter{sujetreuintphilosophiqueETADFI}
\newcounter{sujetreuintphilosophiqueETADMRC}
\newcounter{sujetreuintphilosophiqueETAHA}
\newcounter{sujetreuintphilosophiquepremtermETA}
\newcounter{sujetreuintphilosophiqueEXSPdP}
\newcounter{sujetreuintphilosophiqueEXSDFI}
\newcounter{sujetreuintphilosophiqueEXSDMRC}
\newcounter{sujetreuintphilosophiqueEXSHA}
\newcounter{sujetreuintphilosophiquepremtermEXS}
\newcounter{sujetreuintphilosophiqueMMPdP}
\newcounter{sujetreuintphilosophiqueMMDFI}
\newcounter{sujetreuintphilosophiqueMMDMRC}
\newcounter{sujetreuintphilosophiqueMMHA}
\newcounter{sujetreuintphilosophiquepremtermMM}
\newcounter{sujetreuintphilosophiqueCCPdP}
\newcounter{sujetreuintphilosophiqueCCDFI}
\newcounter{sujetreuintphilosophiqueCCDMRC}
\newcounter{sujetreuintphilosophiqueCCHA}
\newcounter{sujetreuintphilosophiquepremtermCC}
\newcounter{sujetreuintphilosophiqueHVPdP}
\newcounter{sujetreuintphilosophiqueHVDFI}
\newcounter{sujetreuintphilosophiqueHVDMRC}
\newcounter{sujetreuintphilosophiqueHVHA}
\newcounter{sujetreuintphilosophiquepremtermHV}
\newcounter{sujetreuintphilosophiqueHLPdP}
\newcounter{sujetreuintphilosophiqueHLDFI}
\newcounter{sujetreuintphilosophiqueHLDMRC}
\newcounter{sujetreuintphilosophiqueHLHA}
\newcounter{sujetreuintphilosophiquepremtermHL}
\newcounter{sujetreuintphilosophiqueRSPdP}
\newcounter{sujetreuintphilosophiqueRSRdM}
\newcounter{sujetreuintphilosophiqueRSDFI}
\newcounter{sujetreuintphilosophiqueRSDMRC}
\newcounter{sujetreuintphilosophiqueRSHA}
\newcounter{sujetreuintphilosophiqueHQPdP}
\newcounter{sujetreuintphilosophiqueHQRdM}
\newcounter{sujetreuintphilosophiqueHQDFI}
\newcounter{sujetreuintphilosophiqueHQDMRC}
\newcounter{sujetreuintphilosophiqueHQHA}
%%%Somme des compteurs%%%
\newcounter{sujetreuallintphilosophiqueETA}
\newcounter{sujetreuallintphilosophiqueEXS}
\newcounter{sujetreuallintphilosophiqueMM}
\newcounter{sujetreuallintphilosophiqueRS}
\newcounter{sujetreuallintphilosophiqueCC}
\newcounter{sujetreuallintphilosophiqueHV}
\newcounter{sujetreuallintphilosophiqueHL}
\newcounter{sujetreuallintphilosophiqueHQ}

%%%%%%%%%Interprétation littéraire%%%%%%%%%%%
\newcounter{nosujetreuintlitteraire}
\newcounter{sujetreuHQintlitteraire}
\newcounter{sujetreuRSintlitteraire}
%%%sujetreus simples%%%
\newcounter{sujetreuintlitteraireEXS}
\newcounter{sujetreuintlitteraireETA}
\newcounter{sujetreuintlitteraireMM}
\newcounter{sujetreuintlitteraireHV}
\newcounter{sujetreuintlitteraireCC}
\newcounter{sujetreuintlitteraireHL}
%%%sujetreus mixtes%%%
\newcounter{sujetreuintlitteraireEXSTA}
\newcounter{sujetreuintlitteraireEXSMM}
\newcounter{sujetreuintlitteraireETAMM}
\newcounter{sujetreuintlitteraireOdS}
\newcounter{sujetreuintlitteraireHVL}
\newcounter{sujetreuintlitteraireHVCC}
\newcounter{sujetreuintlitteraireHLCC}
\newcounter{sujetreuintlitteraireOdH}
%%%sujetreus mixtes semestres%%%
\newcounter{sujetreuintlitteraireETACC}
\newcounter{sujetreuintlitteraireETAHV}
\newcounter{sujetreuintlitteraireETAHL}
\newcounter{sujetreuintlitteraireEXSCC}
\newcounter{sujetreuintlitteraireEXSHV}
\newcounter{sujetreuintlitteraireEXSHL}
\newcounter{sujetreuintlitteraireMMCC}
\newcounter{sujetreuintlitteraireMMHV}
\newcounter{sujetreuintlitteraireMMHL}
%%%sujetreus mixtes années%%
\newcounter{sujetreuintlitterairepremterm}
\newcounter{sujetreuintlitterairepremtermRS}
\newcounter{sujetreuintlitterairepremtermHQ}
\newcounter{sujetreuintlitteraireETAPdP}
\newcounter{sujetreuintlitteraireETADFI}
\newcounter{sujetreuintlitteraireETADMRC}
\newcounter{sujetreuintlitteraireETAHA}
\newcounter{sujetreuintlitterairepremtermETA}
\newcounter{sujetreuintlitteraireEXSPdP}
\newcounter{sujetreuintlitteraireEXSDFI}
\newcounter{sujetreuintlitteraireEXSDMRC}
\newcounter{sujetreuintlitteraireEXSHA}
\newcounter{sujetreuintlitterairepremtermEXS}
\newcounter{sujetreuintlitteraireMMPdP}
\newcounter{sujetreuintlitteraireMMDFI}
\newcounter{sujetreuintlitteraireMMDMRC}
\newcounter{sujetreuintlitteraireMMHA}
\newcounter{sujetreuintlitterairepremtermMM}
\newcounter{sujetreuintlitteraireCCPdP}
\newcounter{sujetreuintlitteraireCCDFI}
\newcounter{sujetreuintlitteraireCCDMRC}
\newcounter{sujetreuintlitteraireCCHA}
\newcounter{sujetreuintlitterairepremtermCC}
\newcounter{sujetreuintlitteraireHVPdP}
\newcounter{sujetreuintlitteraireHVDFI}
\newcounter{sujetreuintlitteraireHVDMRC}
\newcounter{sujetreuintlitteraireHVHA}
\newcounter{sujetreuintlitterairepremtermHV}
\newcounter{sujetreuintlitteraireHLPdP}
\newcounter{sujetreuintlitteraireHLDFI}
\newcounter{sujetreuintlitteraireHLDMRC}
\newcounter{sujetreuintlitteraireHLHA}
\newcounter{sujetreuintlitterairepremtermHL}
\newcounter{sujetreuintlitteraireRSPdP}
\newcounter{sujetreuintlitteraireRSRdM}
\newcounter{sujetreuintlitteraireRSDFI}
\newcounter{sujetreuintlitteraireRSDMRC}
\newcounter{sujetreuintlitteraireRSHA}
\newcounter{sujetreuintlitteraireHQPdP}
\newcounter{sujetreuintlitteraireHQRdM}
\newcounter{sujetreuintlitteraireHQDFI}
\newcounter{sujetreuintlitteraireHQDMRC}
\newcounter{sujetreuintlitteraireHQHA}
%%%Somme des compteurs%%%
\newcounter{sujetreuallintlitteraireETA}
\newcounter{sujetreuallintlitteraireEXS}
\newcounter{sujetreuallintlitteraireMM}
\newcounter{sujetreuallintlitteraireRS}
\newcounter{sujetreuallintlitteraireCC}
\newcounter{sujetreuallintlitteraireHV}
\newcounter{sujetreuallintlitteraireHL}
\newcounter{sujetreuallintlitteraireHQ}

%%%%%%%%%%%%%%%%%%%%%%%%%%%%%%%%%%%%%%%%%%%%%
%%%%%%%%%%%%Question de réflexion%%%%%%%%%%%%
%%%%%%%%%%%%%%%%%%%%%%%%%%%%%%%%%%%%%%%%%%%%%
\newcounter{nosujetreuref}
%%%sujetreus simples%%%
\newcounter{sujetreurefEXS}
\newcounter{sujetreurefETA}
\newcounter{sujetreurefMM}
\newcounter{sujetreurefHV}
\newcounter{sujetreurefCC}
\newcounter{sujetreurefHL}
%%%sujetreus mixtes%%%
\newcounter{sujetreurefEXSTA}
\newcounter{sujetreurefEXSMM}
\newcounter{sujetreurefETAMM}
\newcounter{sujetreurefOdS}
\newcounter{sujetreurefHVL}
\newcounter{sujetreurefHVCC}
\newcounter{sujetreurefHLCC}
\newcounter{sujetreurefOdH}
%%%sujetreus mixtes semestres%%%
\newcounter{sujetreurefETACC}
\newcounter{sujetreurefETAHV}
\newcounter{sujetreurefETAHL}
\newcounter{sujetreurefEXSCC}
\newcounter{sujetreurefEXSHV}
\newcounter{sujetreurefEXSHL}
\newcounter{sujetreurefMMCC}
\newcounter{sujetreurefMMHV}
\newcounter{sujetreurefMMHL}
%%%sujetreus mixtes années%%
\newcounter{sujetreurefpremterm}%%%Compteur général
\newcounter{sujetreurefETAPdP}
\newcounter{sujetreurefETADFI}
\newcounter{sujetreurefETADMRC}
\newcounter{sujetreurefETAHA}
\newcounter{sujetreurefpremtermETA}
\newcounter{sujetreurefEXSPdP}
\newcounter{sujetreurefEXSDFI}
\newcounter{sujetreurefEXSDMRC}
\newcounter{sujetreurefEXSHA}
\newcounter{sujetreurefpremtermEXS}
\newcounter{sujetreurefMMPdP}
\newcounter{sujetreurefMMDFI}
\newcounter{sujetreurefMMDMRC}
\newcounter{sujetreurefMMHA}
\newcounter{sujetreurefpremtermMM}
\newcounter{sujetreurefCCPdP}
\newcounter{sujetreurefCCDFI}
\newcounter{sujetreurefCCDMRC}
\newcounter{sujetreurefCCHA}
\newcounter{sujetreurefpremtermCC}
\newcounter{sujetreurefHVPdP}
\newcounter{sujetreurefHVDFI}
\newcounter{sujetreurefHVDMRC}
\newcounter{sujetreurefHVHA}
\newcounter{sujetreurefpremtermHV}
\newcounter{sujetreurefHLPdP}
\newcounter{sujetreurefHLDFI}
\newcounter{sujetreurefHLDMRC}
\newcounter{sujetreurefHLHA}
\newcounter{sujetreurefpremtermHL}
\newcounter{sujetreurefRSPdP}
\newcounter{sujetreurefRSRdM}
\newcounter{sujetreurefRSDFI}
\newcounter{sujetreurefRSDMRC}
\newcounter{sujetreurefRSHA}
\newcounter{sujetreurefHQPdP}
\newcounter{sujetreurefHQRdM}
\newcounter{sujetreurefHQDFI}
\newcounter{sujetreurefHQDMRC}
\newcounter{sujetreurefHQHA}
%%%Somme des compteurs%%%
\newcounter{sujetreuallrefETA}
\newcounter{sujetreuallrefEXS}
\newcounter{sujetreuallrefMM}
\newcounter{sujetreuallrefRS}
\newcounter{sujetreuallrefCC}
\newcounter{sujetreuallrefHV}
\newcounter{sujetreuallrefHL}
\newcounter{sujetreuallrefHQ}

%%%%%%%%Réflexion philosophique%%%%%%%%%
\newcounter{nosujetreurefphilosophique}
\newcounter{sujetreuHQrefphilosophique}
\newcounter{sujetreuRSrefphilosophique}
%%%sujetreus simples%%%
\newcounter{sujetreurefphilosophiqueEXS}
\newcounter{sujetreurefphilosophiqueETA}
\newcounter{sujetreurefphilosophiqueMM}
\newcounter{sujetreurefphilosophiqueHV}
\newcounter{sujetreurefphilosophiqueCC}
\newcounter{sujetreurefphilosophiqueHL}
%%%sujetreus mixtes%%%
\newcounter{sujetreurefphilosophiqueEXSTA}
\newcounter{sujetreurefphilosophiqueEXSMM}
\newcounter{sujetreurefphilosophiqueETAMM}
\newcounter{sujetreurefphilosophiqueOdS}
\newcounter{sujetreurefphilosophiqueHVL}
\newcounter{sujetreurefphilosophiqueHVCC}
\newcounter{sujetreurefphilosophiqueHLCC}
\newcounter{sujetreurefphilosophiqueOdH}
%%%sujetreus mixtes semestres%%%
\newcounter{sujetreurefphilosophiqueETACC}
\newcounter{sujetreurefphilosophiqueETAHV}
\newcounter{sujetreurefphilosophiqueETAHL}
\newcounter{sujetreurefphilosophiqueEXSCC}
\newcounter{sujetreurefphilosophiqueEXSHV}
\newcounter{sujetreurefphilosophiqueEXSHL}
\newcounter{sujetreurefphilosophiqueMMCC}
\newcounter{sujetreurefphilosophiqueMMHV}
\newcounter{sujetreurefphilosophiqueMMHL}
%%%sujetreus mixtes années%%
\newcounter{sujetreurefphilosophiquepremterm}
\newcounter{sujetreurefphilosophiquepremtermRS}
\newcounter{sujetreurefphilosophiquepremtermHQ}
\newcounter{sujetreurefphilosophiqueETAPdP}
\newcounter{sujetreurefphilosophiqueETADFI}
\newcounter{sujetreurefphilosophiqueETADMRC}
\newcounter{sujetreurefphilosophiqueETAHA}\newcounter{sujetreurefphilosophiquepremtermETA}
\newcounter{sujetreurefphilosophiqueEXSPdP}
\newcounter{sujetreurefphilosophiqueEXSDFI}
\newcounter{sujetreurefphilosophiqueEXSDMRC}
\newcounter{sujetreurefphilosophiqueEXSHA}
\newcounter{sujetreurefphilosophiquepremtermEXS}
\newcounter{sujetreurefphilosophiqueMMPdP}
\newcounter{sujetreurefphilosophiqueMMDFI}
\newcounter{sujetreurefphilosophiqueMMDMRC}
\newcounter{sujetreurefphilosophiqueMMHA}
\newcounter{sujetreurefphilosophiquepremtermMM}
\newcounter{sujetreurefphilosophiqueCCPdP}
\newcounter{sujetreurefphilosophiqueCCDFI}
\newcounter{sujetreurefphilosophiqueCCDMRC}
\newcounter{sujetreurefphilosophiqueCCHA}
\newcounter{sujetreurefphilosophiquepremtermCC}
\newcounter{sujetreurefphilosophiqueHVPdP}
\newcounter{sujetreurefphilosophiqueHVDFI}
\newcounter{sujetreurefphilosophiqueHVDMRC}
\newcounter{sujetreurefphilosophiqueHVHA}
\newcounter{sujetreurefphilosophiquepremtermHV}
\newcounter{sujetreurefphilosophiqueHLPdP}
\newcounter{sujetreurefphilosophiqueHLDFI}
\newcounter{sujetreurefphilosophiqueHLDMRC}
\newcounter{sujetreurefphilosophiqueHLHA}
\newcounter{sujetreurefphilosophiquepremtermHL}
\newcounter{sujetreurefphilosophiqueRSPdP}
\newcounter{sujetreurefphilosophiqueRSRdM}
\newcounter{sujetreurefphilosophiqueRSDFI}
\newcounter{sujetreurefphilosophiqueRSDMRC}
\newcounter{sujetreurefphilosophiqueRSHA}
\newcounter{sujetreurefphilosophiqueHQPdP}
\newcounter{sujetreurefphilosophiqueHQRdM}
\newcounter{sujetreurefphilosophiqueHQDFI}
\newcounter{sujetreurefphilosophiqueHQDMRC}
\newcounter{sujetreurefphilosophiqueHQHA}
%%%Somme des compteurs%%%
\newcounter{sujetreuallrefphilosophiqueETA}
\newcounter{sujetreuallrefphilosophiqueEXS}
\newcounter{sujetreuallrefphilosophiqueMM}
\newcounter{sujetreuallrefphilosophiqueRS}
\newcounter{sujetreuallrefphilosophiqueCC}
\newcounter{sujetreuallrefphilosophiqueHV}
\newcounter{sujetreuallrefphilosophiqueHL}
\newcounter{sujetreuallrefphilosophiqueHQ}

%%%%%%%%%Réflexion littéraire%%%%%%%%%%%
\newcounter{nosujetreureflitteraire}
\newcounter{sujetreuHQreflitteraire}
\newcounter{sujetreuRSreflitteraire}
%%%sujetreus simples%%%
\newcounter{sujetreureflitteraireEXS}
\newcounter{sujetreureflitteraireETA}
\newcounter{sujetreureflitteraireMM}
\newcounter{sujetreureflitteraireHV}
\newcounter{sujetreureflitteraireCC}
\newcounter{sujetreureflitteraireHL}
%%%sujetreus mixtes%%%
\newcounter{sujetreureflitteraireEXSTA}
\newcounter{sujetreureflitteraireEXSMM}
\newcounter{sujetreureflitteraireETAMM}
\newcounter{sujetreureflitteraireOdS}
\newcounter{sujetreureflitteraireHVL}
\newcounter{sujetreureflitteraireHVCC}
\newcounter{sujetreureflitteraireHLCC}
\newcounter{sujetreureflitteraireOdH}
%%%sujetreus mixtes semestres%%%
\newcounter{sujetreureflitteraireETACC}
\newcounter{sujetreureflitteraireETAHV}
\newcounter{sujetreureflitteraireETAHL}
\newcounter{sujetreureflitteraireEXSCC}
\newcounter{sujetreureflitteraireEXSHV}
\newcounter{sujetreureflitteraireEXSHL}
\newcounter{sujetreureflitteraireMMCC}
\newcounter{sujetreureflitteraireMMHV}
\newcounter{sujetreureflitteraireMMHL}
%%%sujetreus mixtes années%%
\newcounter{sujetreureflitterairepremterm}
\newcounter{sujetreureflitterairepremtermRS}
\newcounter{sujetreureflitterairepremtermHQ}
\newcounter{sujetreureflitteraireETAPdP}
\newcounter{sujetreureflitteraireETADFI}
\newcounter{sujetreureflitteraireETADMRC}
\newcounter{sujetreureflitteraireETAHA}
\newcounter{sujetreureflitterairepremtermETA}
\newcounter{sujetreureflitteraireEXSPdP}
\newcounter{sujetreureflitteraireEXSDFI}
\newcounter{sujetreureflitteraireEXSDMRC}
\newcounter{sujetreureflitteraireEXSHA}
\newcounter{sujetreureflitterairepremtermEXS}
\newcounter{sujetreureflitteraireMMPdP}
\newcounter{sujetreureflitteraireMMDFI}
\newcounter{sujetreureflitteraireMMDMRC}
\newcounter{sujetreureflitteraireMMHA}
\newcounter{sujetreureflitterairepremtermMM}
\newcounter{sujetreureflitteraireCCPdP}
\newcounter{sujetreureflitteraireCCDFI}
\newcounter{sujetreureflitteraireCCDMRC}
\newcounter{sujetreureflitteraireCCHA}
\newcounter{sujetreureflitterairepremtermCC}
\newcounter{sujetreureflitteraireHVPdP}
\newcounter{sujetreureflitteraireHVDFI}
\newcounter{sujetreureflitteraireHVDMRC}
\newcounter{sujetreureflitteraireHVHA}
\newcounter{sujetreureflitterairepremtermHV}
\newcounter{sujetreureflitteraireHLPdP}
\newcounter{sujetreureflitteraireHLDFI}
\newcounter{sujetreureflitteraireHLDMRC}
\newcounter{sujetreureflitteraireHLHA}
\newcounter{sujetreureflitterairepremtermHL}
\newcounter{sujetreureflitteraireRSPdP}
\newcounter{sujetreureflitteraireRSRdM}
\newcounter{sujetreureflitteraireRSDFI}
\newcounter{sujetreureflitteraireRSDMRC}
\newcounter{sujetreureflitteraireRSHA}
\newcounter{sujetreureflitteraireHQPdP}
\newcounter{sujetreureflitteraireHQRdM}
\newcounter{sujetreureflitteraireHQDFI}
\newcounter{sujetreureflitteraireHQDMRC}
\newcounter{sujetreureflitteraireHQHA}
%%%Somme des compteurs%%%
\newcounter{sujetreuallreflitteraireETA}
\newcounter{sujetreuallreflitteraireEXS}
\newcounter{sujetreuallreflitteraireMM}
\newcounter{sujetreuallreflitteraireCC}
\newcounter{sujetreuallreflitteraireRS}
\newcounter{sujetreuallreflitteraireHV}
\newcounter{sujetreuallreflitteraireHL}
\newcounter{sujetreuallreflitteraireHQ}
%%%Compteurs de statistiques globales sur les genres de question%%%
\newcounter{sujetreuexclRS}
\newcounter{sujetreumixtRS}
\newcounter{sujetreuexclHQ}
\newcounter{sujetreumixtHQ}
\newcounter{sujetreumixtHQRS}
%%%Compteurs de statistiques sur les transversaux (divers semestres, première et terminale)
\newcounter{sujettransreuintlitteraire}
\newcounter{sujettransreureflitteraire}
\newcounter{sujettransreuintphilosophique}
\newcounter{sujettransreurefphilosophique}
%%% Compteurs du nombre de sujets par type
\newcounter{sujetreuintlitteraire}
\newcounter{sujetreureflitteraire}
\newcounter{sujetreuintphilosophique}
\newcounter{sujetreurefphilosophique}
%%% Compteurs par discipline
%%Litt
\newcounter{sujetreulitteraireETA}
\newcounter{sujetreulitteraireEXS}
\newcounter{sujetreulitteraireMM}
\newcounter{sujetreulitteraireCC}
\newcounter{sujetreulitteraireHV}
\newcounter{sujetreulitteraireHL}
\newcounter{sujetreulitteraireEXSTA}
\newcounter{sujetreulitteraireETAMM}
\newcounter{sujetreulitteraireEXSMM}
\newcounter{sujetreulitteraireHVCC}
\newcounter{sujetreulitteraireHVL}
\newcounter{sujetreulitteraireaut}
%%Phil
\newcounter{sujetreuphilosophiqueETA}
\newcounter{sujetreuphilosophiqueEXS}
\newcounter{sujetreuphilosophiqueMM}
\newcounter{sujetreuphilosophiqueCC}
\newcounter{sujetreuphilosophiqueHV}
\newcounter{sujetreuphilosophiqueHL}
\newcounter{sujetreuphilosophiqueEXSTA}
\newcounter{sujetreuphilosophiqueETAMM}
\newcounter{sujetreuphilosophiqueEXSMM}
\newcounter{sujetreuphilosophiqueHVCC}
\newcounter{sujetreuphilosophiqueHVL}
\newcounter{sujetreuphilosophiqueaut}
%%Totaux pour les graphiques
\newcounter{sujetreulitterairetotal}
\newcounter{sujetreuphilosophiquetotal}
%%% Autres (Ods et Odh) comptés ensemble
\newcounter{sujetreuallaut}
%%%%%%%%%%%%%%%%%%%%%%%%%%%%%%%%%%%%%%%%%%%%%%%%%%%%%%%%%%%%%%%%%%%%%%%
%%%%%%%%%%%%%%%%%%%%%%%%%%%%%%Compteurs%%%%%%%%%%%%%%%%%%%%%%%%%%%%%%%%
%%%%%%%%%%%%%%%%%%%%%%%%%%%%%%%%%%%%%%%%%%%%%%%%%%%%%%%%%%%%%%%%%%%%%%%
%%%Compter les sujets par thème%%%
%\newcounter{nosujetsujetzero}
\newcounter{sujetsujall}%Total de questions pour le centre
\newcounter{sujetsujtotal}%Total des questions en comptant en double les questions pour les graphiques
\newcounter{sujetsujHQ}
\newcounter{sujetsujRS}
%%%sujetsujs simples%%%
\newcounter{sujetsujEXS}
\newcounter{sujetsujETA}
\newcounter{sujetsujMM}
\newcounter{sujetsujHV}
\newcounter{sujetsujCC}
\newcounter{sujetsujHL}
%%%sujetsujs mixtes%%%
\newcounter{sujetsujEXSTA}
\newcounter{sujetsujEXSMM}
\newcounter{sujetsujETAMM}
\newcounter{sujetsujOdS}
\newcounter{sujetsujHVL}
\newcounter{sujetsujHVCC}
\newcounter{sujetsujHLCC}
\newcounter{sujetsujOdH}
%%%sujetsujs mixtes semestres%%%
\newcounter{sujetsujETACC}
\newcounter{sujetsujETAHV}
\newcounter{sujetsujETAHL}
\newcounter{sujetsujEXSCC}
\newcounter{sujetsujEXSHV}
\newcounter{sujetsujEXSHL}
\newcounter{sujetsujMMCC}
\newcounter{sujetsujMMHV}
\newcounter{sujetsujMMHL}

%%%sujetsujs mixtes années%%
\newcounter{sujetsujpremterm}%%%Compteur général de sujetsuj
\newcounter{sujetsujETAPdP}
\newcounter{sujetsujETADFI}
\newcounter{sujetsujETADMRC}
\newcounter{sujetsujETAHA}
\newcounter{sujetsujpremtermETA}
\newcounter{sujetsujEXSPdP}
\newcounter{sujetsujEXSDFI}
\newcounter{sujetsujEXSDMRC}
\newcounter{sujetsujEXSHA}
\newcounter{sujetsujpremtermEXS}
\newcounter{sujetsujMMPdP}
\newcounter{sujetsujMMDFI}
\newcounter{sujetsujMMDMRC}
\newcounter{sujetsujMMHA}
\newcounter{sujetsujpremtermMM}
\newcounter{sujetsujCCPdP}
\newcounter{sujetsujCCDFI}
\newcounter{sujetsujCCDMRC}
\newcounter{sujetsujCCHA}
\newcounter{sujetsujpremtermCC}
\newcounter{sujetsujHVPdP}
\newcounter{sujetsujHVDFI}
\newcounter{sujetsujHVDMRC}
\newcounter{sujetsujHVHA}
\newcounter{sujetsujpremtermHV}
\newcounter{sujetsujHLPdP}
\newcounter{sujetsujHLDFI}
\newcounter{sujetsujHLDMRC}
\newcounter{sujetsujHLHA}
\newcounter{sujetsujpremtermHL}
\newcounter{sujetsujRSPdP}
\newcounter{sujetsujRSRdM}
\newcounter{sujetsujRSDFI}
\newcounter{sujetsujRSDMRC}
\newcounter{sujetsujRSHA}
\newcounter{sujetsujHQPdP}
\newcounter{sujetsujHQRdM}
\newcounter{sujetsujHQDFI}
\newcounter{sujetsujHQDMRC}
\newcounter{sujetsujHQHA}

%%%Somme des compteurs%%%
\newcounter{sujetsujallETA}
\newcounter{sujetsujallEXS}
\newcounter{sujetsujallMM}
\newcounter{sujetsujallRS}
\newcounter{sujetsujallCC}
\newcounter{sujetsujallHV}
\newcounter{sujetsujallHL}
\newcounter{sujetsujallHQ}

%%%Compter les sujetsujs par thème et types de questions%%%
%%%%%%%%%%%%%%%%%%%%%%%%%%%%%%%%%%%%%%%%%%%%%
%%%%%%%%%Question d'interprétation%%%%%%%%%%%
%%%%%%%%%%%%%%%%%%%%%%%%%%%%%%%%%%%%%%%%%%%%%

%%%%%%%%%%%%%%%Général%%%%%%%%%%%%%%%%%%%%%%%
\newcounter{nosujetsujint}
\newcounter{sujetsujHQint}
\newcounter{sujetsujRSint}
%%%sujetsujs simples%%%
\newcounter{sujetsujintEXS}
\newcounter{sujetsujintETA}
\newcounter{sujetsujintMM}
\newcounter{sujetsujintHV}
\newcounter{sujetsujintCC}
\newcounter{sujetsujintHL}
%%%sujetsujs mixtes%%%
\newcounter{sujetsujintEXSTA}
\newcounter{sujetsujintEXSMM}
\newcounter{sujetsujintETAMM}
\newcounter{sujetsujintOdS}
\newcounter{sujetsujintHVL}
\newcounter{sujetsujintHVCC}
\newcounter{sujetsujintHLCC}
\newcounter{sujetsujintOdH}
%%%sujetsujs mixtes semestres%%%
\newcounter{sujetsujintETACC}
\newcounter{sujetsujintETAHV}
\newcounter{sujetsujintETAHL}
\newcounter{sujetsujintEXSCC}
\newcounter{sujetsujintEXSHV}
\newcounter{sujetsujintEXSHL}
\newcounter{sujetsujintMMCC}
\newcounter{sujetsujintMMHV}
\newcounter{sujetsujintMMHL}
%%%sujetsujs mixtes années%%
\newcounter{sujetsujintpremterm}
\newcounter{sujetsujintpremtermRS}
\newcounter{sujetsujintpremtermHQ}
\newcounter{sujetsujintETAPdP}
\newcounter{sujetsujintETADFI}
\newcounter{sujetsujintETADMRC}
\newcounter{sujetsujintETAHA}
\newcounter{sujetsujintpremtermETA}
\newcounter{sujetsujintEXSPdP}
\newcounter{sujetsujintEXSDFI}
\newcounter{sujetsujintEXSDMRC}
\newcounter{sujetsujintEXSHA}
\newcounter{sujetsujintpremtermEXS}
\newcounter{sujetsujintMMPdP}
\newcounter{sujetsujintMMDFI}
\newcounter{sujetsujintMMDMRC}
\newcounter{sujetsujintMMHA}
\newcounter{sujetsujintpremtermMM}
\newcounter{sujetsujintCCPdP}
\newcounter{sujetsujintCCDFI}
\newcounter{sujetsujintCCDMRC}
\newcounter{sujetsujintCCHA}
\newcounter{sujetsujintpremtermCC}
\newcounter{sujetsujintHVPdP}
\newcounter{sujetsujintHVDFI}
\newcounter{sujetsujintHVDMRC}
\newcounter{sujetsujintHVHA}
\newcounter{sujetsujintpremtermHV}
\newcounter{sujetsujintHLPdP}
\newcounter{sujetsujintHLDFI}
\newcounter{sujetsujintHLDMRC}
\newcounter{sujetsujintHLHA}
\newcounter{sujetsujintpremtermHL}
\newcounter{sujetsujintRSPdP}
\newcounter{sujetsujintRSRdM}
\newcounter{sujetsujintRSDFI}
\newcounter{sujetsujintRSDMRC}
\newcounter{sujetsujintRSHA}
\newcounter{sujetsujintHQPdP}
\newcounter{sujetsujintHQRdM}
\newcounter{sujetsujintHQDFI}
\newcounter{sujetsujintHQDMRC}
\newcounter{sujetsujintHQHA}
%%%Somme des compteurs%%%
\newcounter{sujetsujallintETA}
\newcounter{sujetsujallintEXS}
\newcounter{sujetsujallintMM}
\newcounter{sujetsujallintRS}
\newcounter{sujetsujallintCC}
\newcounter{sujetsujallintHV}
\newcounter{sujetsujallintHL}
\newcounter{sujetsujallintHQ}

%%%%%%%%Interprétation philosophique%%%%%%%%%
\newcounter{nosujetsujintphilosophique}
\newcounter{sujetsujHQintphilosophique}
\newcounter{sujetsujRSintphilosophique}
%%%sujetsujs simples%%%
\newcounter{sujetsujintphilosophiqueEXS}
\newcounter{sujetsujintphilosophiqueETA}
\newcounter{sujetsujintphilosophiqueMM}
\newcounter{sujetsujintphilosophiqueHV}
\newcounter{sujetsujintphilosophiqueCC}
\newcounter{sujetsujintphilosophiqueHL}
%%%sujetsujs mixtes%%%
\newcounter{sujetsujintphilosophiqueEXSTA}
\newcounter{sujetsujintphilosophiqueEXSMM}
\newcounter{sujetsujintphilosophiqueETAMM}
\newcounter{sujetsujintphilosophiqueOdS}
\newcounter{sujetsujintphilosophiqueHVL}
\newcounter{sujetsujintphilosophiqueHVCC}
\newcounter{sujetsujintphilosophiqueHLCC}
\newcounter{sujetsujintphilosophiqueOdH}
%%%sujetsujs mixtes semestres%%%
\newcounter{sujetsujintphilosophiqueETACC}
\newcounter{sujetsujintphilosophiqueETAHV}
\newcounter{sujetsujintphilosophiqueETAHL}
\newcounter{sujetsujintphilosophiqueEXSCC}
\newcounter{sujetsujintphilosophiqueEXSHV}
\newcounter{sujetsujintphilosophiqueEXSHL}
\newcounter{sujetsujintphilosophiqueMMCC}
\newcounter{sujetsujintphilosophiqueMMHV}
\newcounter{sujetsujintphilosophiqueMMHL}
%%%sujetsujs mixtes années%%
\newcounter{sujetsujintphilosophiquepremterm}
\newcounter{sujetsujintphilosophiquepremtermRS}
\newcounter{sujetsujintphilosophiquepremtermHQ}
\newcounter{sujetsujintphilosophiqueETAPdP}
\newcounter{sujetsujintphilosophiqueETADFI}
\newcounter{sujetsujintphilosophiqueETADMRC}
\newcounter{sujetsujintphilosophiqueETAHA}
\newcounter{sujetsujintphilosophiquepremtermETA}
\newcounter{sujetsujintphilosophiqueEXSPdP}
\newcounter{sujetsujintphilosophiqueEXSDFI}
\newcounter{sujetsujintphilosophiqueEXSDMRC}
\newcounter{sujetsujintphilosophiqueEXSHA}
\newcounter{sujetsujintphilosophiquepremtermEXS}
\newcounter{sujetsujintphilosophiqueMMPdP}
\newcounter{sujetsujintphilosophiqueMMDFI}
\newcounter{sujetsujintphilosophiqueMMDMRC}
\newcounter{sujetsujintphilosophiqueMMHA}
\newcounter{sujetsujintphilosophiquepremtermMM}
\newcounter{sujetsujintphilosophiqueCCPdP}
\newcounter{sujetsujintphilosophiqueCCDFI}
\newcounter{sujetsujintphilosophiqueCCDMRC}
\newcounter{sujetsujintphilosophiqueCCHA}
\newcounter{sujetsujintphilosophiquepremtermCC}
\newcounter{sujetsujintphilosophiqueHVPdP}
\newcounter{sujetsujintphilosophiqueHVDFI}
\newcounter{sujetsujintphilosophiqueHVDMRC}
\newcounter{sujetsujintphilosophiqueHVHA}
\newcounter{sujetsujintphilosophiquepremtermHV}
\newcounter{sujetsujintphilosophiqueHLPdP}
\newcounter{sujetsujintphilosophiqueHLDFI}
\newcounter{sujetsujintphilosophiqueHLDMRC}
\newcounter{sujetsujintphilosophiqueHLHA}
\newcounter{sujetsujintphilosophiquepremtermHL}
\newcounter{sujetsujintphilosophiqueRSPdP}
\newcounter{sujetsujintphilosophiqueRSRdM}
\newcounter{sujetsujintphilosophiqueRSDFI}
\newcounter{sujetsujintphilosophiqueRSDMRC}
\newcounter{sujetsujintphilosophiqueRSHA}
\newcounter{sujetsujintphilosophiqueHQPdP}
\newcounter{sujetsujintphilosophiqueHQRdM}
\newcounter{sujetsujintphilosophiqueHQDFI}
\newcounter{sujetsujintphilosophiqueHQDMRC}
\newcounter{sujetsujintphilosophiqueHQHA}
%%%Somme des compteurs%%%
\newcounter{sujetsujallintphilosophiqueETA}
\newcounter{sujetsujallintphilosophiqueEXS}
\newcounter{sujetsujallintphilosophiqueMM}
\newcounter{sujetsujallintphilosophiqueRS}
\newcounter{sujetsujallintphilosophiqueCC}
\newcounter{sujetsujallintphilosophiqueHV}
\newcounter{sujetsujallintphilosophiqueHL}
\newcounter{sujetsujallintphilosophiqueHQ}

%%%%%%%%%Interprétation littéraire%%%%%%%%%%%
\newcounter{nosujetsujintlitteraire}
\newcounter{sujetsujHQintlitteraire}
\newcounter{sujetsujRSintlitteraire}
%%%sujetsujs simples%%%
\newcounter{sujetsujintlitteraireEXS}
\newcounter{sujetsujintlitteraireETA}
\newcounter{sujetsujintlitteraireMM}
\newcounter{sujetsujintlitteraireHV}
\newcounter{sujetsujintlitteraireCC}
\newcounter{sujetsujintlitteraireHL}
%%%sujetsujs mixtes%%%
\newcounter{sujetsujintlitteraireEXSTA}
\newcounter{sujetsujintlitteraireEXSMM}
\newcounter{sujetsujintlitteraireETAMM}
\newcounter{sujetsujintlitteraireOdS}
\newcounter{sujetsujintlitteraireHVL}
\newcounter{sujetsujintlitteraireHVCC}
\newcounter{sujetsujintlitteraireHLCC}
\newcounter{sujetsujintlitteraireOdH}
%%%sujetsujs mixtes semestres%%%
\newcounter{sujetsujintlitteraireETACC}
\newcounter{sujetsujintlitteraireETAHV}
\newcounter{sujetsujintlitteraireETAHL}
\newcounter{sujetsujintlitteraireEXSCC}
\newcounter{sujetsujintlitteraireEXSHV}
\newcounter{sujetsujintlitteraireEXSHL}
\newcounter{sujetsujintlitteraireMMCC}
\newcounter{sujetsujintlitteraireMMHV}
\newcounter{sujetsujintlitteraireMMHL}
%%%sujetsujs mixtes années%%
\newcounter{sujetsujintlitterairepremterm}
\newcounter{sujetsujintlitterairepremtermRS}
\newcounter{sujetsujintlitterairepremtermHQ}
\newcounter{sujetsujintlitteraireETAPdP}
\newcounter{sujetsujintlitteraireETADFI}
\newcounter{sujetsujintlitteraireETADMRC}
\newcounter{sujetsujintlitteraireETAHA}
\newcounter{sujetsujintlitterairepremtermETA}
\newcounter{sujetsujintlitteraireEXSPdP}
\newcounter{sujetsujintlitteraireEXSDFI}
\newcounter{sujetsujintlitteraireEXSDMRC}
\newcounter{sujetsujintlitteraireEXSHA}
\newcounter{sujetsujintlitterairepremtermEXS}
\newcounter{sujetsujintlitteraireMMPdP}
\newcounter{sujetsujintlitteraireMMDFI}
\newcounter{sujetsujintlitteraireMMDMRC}
\newcounter{sujetsujintlitteraireMMHA}
\newcounter{sujetsujintlitterairepremtermMM}
\newcounter{sujetsujintlitteraireCCPdP}
\newcounter{sujetsujintlitteraireCCDFI}
\newcounter{sujetsujintlitteraireCCDMRC}
\newcounter{sujetsujintlitteraireCCHA}
\newcounter{sujetsujintlitterairepremtermCC}
\newcounter{sujetsujintlitteraireHVPdP}
\newcounter{sujetsujintlitteraireHVDFI}
\newcounter{sujetsujintlitteraireHVDMRC}
\newcounter{sujetsujintlitteraireHVHA}
\newcounter{sujetsujintlitterairepremtermHV}
\newcounter{sujetsujintlitteraireHLPdP}
\newcounter{sujetsujintlitteraireHLDFI}
\newcounter{sujetsujintlitteraireHLDMRC}
\newcounter{sujetsujintlitteraireHLHA}
\newcounter{sujetsujintlitterairepremtermHL}
\newcounter{sujetsujintlitteraireRSPdP}
\newcounter{sujetsujintlitteraireRSRdM}
\newcounter{sujetsujintlitteraireRSDFI}
\newcounter{sujetsujintlitteraireRSDMRC}
\newcounter{sujetsujintlitteraireRSHA}
\newcounter{sujetsujintlitteraireHQPdP}
\newcounter{sujetsujintlitteraireHQRdM}
\newcounter{sujetsujintlitteraireHQDFI}
\newcounter{sujetsujintlitteraireHQDMRC}
\newcounter{sujetsujintlitteraireHQHA}
%%%Somme des compteurs%%%
\newcounter{sujetsujallintlitteraireETA}
\newcounter{sujetsujallintlitteraireEXS}
\newcounter{sujetsujallintlitteraireMM}
\newcounter{sujetsujallintlitteraireRS}
\newcounter{sujetsujallintlitteraireCC}
\newcounter{sujetsujallintlitteraireHV}
\newcounter{sujetsujallintlitteraireHL}
\newcounter{sujetsujallintlitteraireHQ}

%%%%%%%%%%%%%%%%%%%%%%%%%%%%%%%%%%%%%%%%%%%%%
%%%%%%%%%%%%Question de réflexion%%%%%%%%%%%%
%%%%%%%%%%%%%%%%%%%%%%%%%%%%%%%%%%%%%%%%%%%%%
\newcounter{nosujetsujref}
%%%sujetsujs simples%%%
\newcounter{sujetsujrefEXS}
\newcounter{sujetsujrefETA}
\newcounter{sujetsujrefMM}
\newcounter{sujetsujrefHV}
\newcounter{sujetsujrefCC}
\newcounter{sujetsujrefHL}
%%%sujetsujs mixtes%%%
\newcounter{sujetsujrefEXSTA}
\newcounter{sujetsujrefEXSMM}
\newcounter{sujetsujrefETAMM}
\newcounter{sujetsujrefOdS}
\newcounter{sujetsujrefHVL}
\newcounter{sujetsujrefHVCC}
\newcounter{sujetsujrefHLCC}
\newcounter{sujetsujrefOdH}
%%%sujetsujs mixtes semestres%%%
\newcounter{sujetsujrefETACC}
\newcounter{sujetsujrefETAHV}
\newcounter{sujetsujrefETAHL}
\newcounter{sujetsujrefEXSCC}
\newcounter{sujetsujrefEXSHV}
\newcounter{sujetsujrefEXSHL}
\newcounter{sujetsujrefMMCC}
\newcounter{sujetsujrefMMHV}
\newcounter{sujetsujrefMMHL}
%%%sujetsujs mixtes années%%
\newcounter{sujetsujrefpremterm}%%%Compteur général
\newcounter{sujetsujrefETAPdP}
\newcounter{sujetsujrefETADFI}
\newcounter{sujetsujrefETADMRC}
\newcounter{sujetsujrefETAHA}
\newcounter{sujetsujrefpremtermETA}
\newcounter{sujetsujrefEXSPdP}
\newcounter{sujetsujrefEXSDFI}
\newcounter{sujetsujrefEXSDMRC}
\newcounter{sujetsujrefEXSHA}
\newcounter{sujetsujrefpremtermEXS}
\newcounter{sujetsujrefMMPdP}
\newcounter{sujetsujrefMMDFI}
\newcounter{sujetsujrefMMDMRC}
\newcounter{sujetsujrefMMHA}
\newcounter{sujetsujrefpremtermMM}
\newcounter{sujetsujrefCCPdP}
\newcounter{sujetsujrefCCDFI}
\newcounter{sujetsujrefCCDMRC}
\newcounter{sujetsujrefCCHA}
\newcounter{sujetsujrefpremtermCC}
\newcounter{sujetsujrefHVPdP}
\newcounter{sujetsujrefHVDFI}
\newcounter{sujetsujrefHVDMRC}
\newcounter{sujetsujrefHVHA}
\newcounter{sujetsujrefpremtermHV}
\newcounter{sujetsujrefHLPdP}
\newcounter{sujetsujrefHLDFI}
\newcounter{sujetsujrefHLDMRC}
\newcounter{sujetsujrefHLHA}
\newcounter{sujetsujrefpremtermHL}
\newcounter{sujetsujrefRSPdP}
\newcounter{sujetsujrefRSRdM}
\newcounter{sujetsujrefRSDFI}
\newcounter{sujetsujrefRSDMRC}
\newcounter{sujetsujrefRSHA}
\newcounter{sujetsujrefHQPdP}
\newcounter{sujetsujrefHQRdM}
\newcounter{sujetsujrefHQDFI}
\newcounter{sujetsujrefHQDMRC}
\newcounter{sujetsujrefHQHA}
%%%Somme des compteurs%%%
\newcounter{sujetsujallrefETA}
\newcounter{sujetsujallrefEXS}
\newcounter{sujetsujallrefMM}
\newcounter{sujetsujallrefRS}
\newcounter{sujetsujallrefCC}
\newcounter{sujetsujallrefHV}
\newcounter{sujetsujallrefHL}
\newcounter{sujetsujallrefHQ}

%%%%%%%%Réflexion philosophique%%%%%%%%%
\newcounter{nosujetsujrefphilosophique}
\newcounter{sujetsujHQrefphilosophique}
\newcounter{sujetsujRSrefphilosophique}
%%%sujetsujs simples%%%
\newcounter{sujetsujrefphilosophiqueEXS}
\newcounter{sujetsujrefphilosophiqueETA}
\newcounter{sujetsujrefphilosophiqueMM}
\newcounter{sujetsujrefphilosophiqueHV}
\newcounter{sujetsujrefphilosophiqueCC}
\newcounter{sujetsujrefphilosophiqueHL}
%%%sujetsujs mixtes%%%
\newcounter{sujetsujrefphilosophiqueEXSTA}
\newcounter{sujetsujrefphilosophiqueEXSMM}
\newcounter{sujetsujrefphilosophiqueETAMM}
\newcounter{sujetsujrefphilosophiqueOdS}
\newcounter{sujetsujrefphilosophiqueHVL}
\newcounter{sujetsujrefphilosophiqueHVCC}
\newcounter{sujetsujrefphilosophiqueHLCC}
\newcounter{sujetsujrefphilosophiqueOdH}
%%%sujetsujs mixtes semestres%%%
\newcounter{sujetsujrefphilosophiqueETACC}
\newcounter{sujetsujrefphilosophiqueETAHV}
\newcounter{sujetsujrefphilosophiqueETAHL}
\newcounter{sujetsujrefphilosophiqueEXSCC}
\newcounter{sujetsujrefphilosophiqueEXSHV}
\newcounter{sujetsujrefphilosophiqueEXSHL}
\newcounter{sujetsujrefphilosophiqueMMCC}
\newcounter{sujetsujrefphilosophiqueMMHV}
\newcounter{sujetsujrefphilosophiqueMMHL}
%%%sujetsujs mixtes années%%
\newcounter{sujetsujrefphilosophiquepremterm}
\newcounter{sujetsujrefphilosophiquepremtermRS}
\newcounter{sujetsujrefphilosophiquepremtermHQ}
\newcounter{sujetsujrefphilosophiqueETAPdP}
\newcounter{sujetsujrefphilosophiqueETADFI}
\newcounter{sujetsujrefphilosophiqueETADMRC}
\newcounter{sujetsujrefphilosophiqueETAHA}\newcounter{sujetsujrefphilosophiquepremtermETA}
\newcounter{sujetsujrefphilosophiqueEXSPdP}
\newcounter{sujetsujrefphilosophiqueEXSDFI}
\newcounter{sujetsujrefphilosophiqueEXSDMRC}
\newcounter{sujetsujrefphilosophiqueEXSHA}
\newcounter{sujetsujrefphilosophiquepremtermEXS}
\newcounter{sujetsujrefphilosophiqueMMPdP}
\newcounter{sujetsujrefphilosophiqueMMDFI}
\newcounter{sujetsujrefphilosophiqueMMDMRC}
\newcounter{sujetsujrefphilosophiqueMMHA}
\newcounter{sujetsujrefphilosophiquepremtermMM}
\newcounter{sujetsujrefphilosophiqueCCPdP}
\newcounter{sujetsujrefphilosophiqueCCDFI}
\newcounter{sujetsujrefphilosophiqueCCDMRC}
\newcounter{sujetsujrefphilosophiqueCCHA}
\newcounter{sujetsujrefphilosophiquepremtermCC}
\newcounter{sujetsujrefphilosophiqueHVPdP}
\newcounter{sujetsujrefphilosophiqueHVDFI}
\newcounter{sujetsujrefphilosophiqueHVDMRC}
\newcounter{sujetsujrefphilosophiqueHVHA}
\newcounter{sujetsujrefphilosophiquepremtermHV}
\newcounter{sujetsujrefphilosophiqueHLPdP}
\newcounter{sujetsujrefphilosophiqueHLDFI}
\newcounter{sujetsujrefphilosophiqueHLDMRC}
\newcounter{sujetsujrefphilosophiqueHLHA}
\newcounter{sujetsujrefphilosophiquepremtermHL}
\newcounter{sujetsujrefphilosophiqueRSPdP}
\newcounter{sujetsujrefphilosophiqueRSRdM}
\newcounter{sujetsujrefphilosophiqueRSDFI}
\newcounter{sujetsujrefphilosophiqueRSDMRC}
\newcounter{sujetsujrefphilosophiqueRSHA}
\newcounter{sujetsujrefphilosophiqueHQPdP}
\newcounter{sujetsujrefphilosophiqueHQRdM}
\newcounter{sujetsujrefphilosophiqueHQDFI}
\newcounter{sujetsujrefphilosophiqueHQDMRC}
\newcounter{sujetsujrefphilosophiqueHQHA}
%%%Somme des compteurs%%%
\newcounter{sujetsujallrefphilosophiqueETA}
\newcounter{sujetsujallrefphilosophiqueEXS}
\newcounter{sujetsujallrefphilosophiqueMM}
\newcounter{sujetsujallrefphilosophiqueRS}
\newcounter{sujetsujallrefphilosophiqueCC}
\newcounter{sujetsujallrefphilosophiqueHV}
\newcounter{sujetsujallrefphilosophiqueHL}
\newcounter{sujetsujallrefphilosophiqueHQ}

%%%%%%%%%Réflexion littéraire%%%%%%%%%%%
\newcounter{nosujetsujreflitteraire}
\newcounter{sujetsujHQreflitteraire}
\newcounter{sujetsujRSreflitteraire}
%%%sujetsujs simples%%%
\newcounter{sujetsujreflitteraireEXS}
\newcounter{sujetsujreflitteraireETA}
\newcounter{sujetsujreflitteraireMM}
\newcounter{sujetsujreflitteraireHV}
\newcounter{sujetsujreflitteraireCC}
\newcounter{sujetsujreflitteraireHL}
%%%sujetsujs mixtes%%%
\newcounter{sujetsujreflitteraireEXSTA}
\newcounter{sujetsujreflitteraireEXSMM}
\newcounter{sujetsujreflitteraireETAMM}
\newcounter{sujetsujreflitteraireOdS}
\newcounter{sujetsujreflitteraireHVL}
\newcounter{sujetsujreflitteraireHVCC}
\newcounter{sujetsujreflitteraireHLCC}
\newcounter{sujetsujreflitteraireOdH}
%%%sujetsujs mixtes semestres%%%
\newcounter{sujetsujreflitteraireETACC}
\newcounter{sujetsujreflitteraireETAHV}
\newcounter{sujetsujreflitteraireETAHL}
\newcounter{sujetsujreflitteraireEXSCC}
\newcounter{sujetsujreflitteraireEXSHV}
\newcounter{sujetsujreflitteraireEXSHL}
\newcounter{sujetsujreflitteraireMMCC}
\newcounter{sujetsujreflitteraireMMHV}
\newcounter{sujetsujreflitteraireMMHL}
%%%sujetsujs mixtes années%%
\newcounter{sujetsujreflitterairepremterm}
\newcounter{sujetsujreflitterairepremtermRS}
\newcounter{sujetsujreflitterairepremtermHQ}
\newcounter{sujetsujreflitteraireETAPdP}
\newcounter{sujetsujreflitteraireETADFI}
\newcounter{sujetsujreflitteraireETADMRC}
\newcounter{sujetsujreflitteraireETAHA}
\newcounter{sujetsujreflitterairepremtermETA}
\newcounter{sujetsujreflitteraireEXSPdP}
\newcounter{sujetsujreflitteraireEXSDFI}
\newcounter{sujetsujreflitteraireEXSDMRC}
\newcounter{sujetsujreflitteraireEXSHA}
\newcounter{sujetsujreflitterairepremtermEXS}
\newcounter{sujetsujreflitteraireMMPdP}
\newcounter{sujetsujreflitteraireMMDFI}
\newcounter{sujetsujreflitteraireMMDMRC}
\newcounter{sujetsujreflitteraireMMHA}
\newcounter{sujetsujreflitterairepremtermMM}
\newcounter{sujetsujreflitteraireCCPdP}
\newcounter{sujetsujreflitteraireCCDFI}
\newcounter{sujetsujreflitteraireCCDMRC}
\newcounter{sujetsujreflitteraireCCHA}
\newcounter{sujetsujreflitterairepremtermCC}
\newcounter{sujetsujreflitteraireHVPdP}
\newcounter{sujetsujreflitteraireHVDFI}
\newcounter{sujetsujreflitteraireHVDMRC}
\newcounter{sujetsujreflitteraireHVHA}
\newcounter{sujetsujreflitterairepremtermHV}
\newcounter{sujetsujreflitteraireHLPdP}
\newcounter{sujetsujreflitteraireHLDFI}
\newcounter{sujetsujreflitteraireHLDMRC}
\newcounter{sujetsujreflitteraireHLHA}
\newcounter{sujetsujreflitterairepremtermHL}
\newcounter{sujetsujreflitteraireRSPdP}
\newcounter{sujetsujreflitteraireRSRdM}
\newcounter{sujetsujreflitteraireRSDFI}
\newcounter{sujetsujreflitteraireRSDMRC}
\newcounter{sujetsujreflitteraireRSHA}
\newcounter{sujetsujreflitteraireHQPdP}
\newcounter{sujetsujreflitteraireHQRdM}
\newcounter{sujetsujreflitteraireHQDFI}
\newcounter{sujetsujreflitteraireHQDMRC}
\newcounter{sujetsujreflitteraireHQHA}
%%%Somme des compteurs%%%
\newcounter{sujetsujallreflitteraireETA}
\newcounter{sujetsujallreflitteraireEXS}
\newcounter{sujetsujallreflitteraireMM}
\newcounter{sujetsujallreflitteraireCC}
\newcounter{sujetsujallreflitteraireRS}
\newcounter{sujetsujallreflitteraireHV}
\newcounter{sujetsujallreflitteraireHL}
\newcounter{sujetsujallreflitteraireHQ}
%%%Compteurs de statistiques globales sur les genres de question%%%
\newcounter{sujetsujexclRS}
\newcounter{sujetsujmixtRS}
\newcounter{sujetsujexclHQ}
\newcounter{sujetsujmixtHQ}
\newcounter{sujetsujmixtHQRS}
%%%Compteurs de statistiques sur les transversaux (divers semestres, première et terminale)
\newcounter{sujettranssujintlitteraire}
\newcounter{sujettranssujreflitteraire}
\newcounter{sujettranssujintphilosophique}
\newcounter{sujettranssujrefphilosophique}
%%% Compteurs du nombre de sujets par type
\newcounter{sujetsujintlitteraire}
\newcounter{sujetsujreflitteraire}
\newcounter{sujetsujintphilosophique}
\newcounter{sujetsujrefphilosophique}
%%% Compteurs par discipline
%%Litt
\newcounter{sujetsujlitteraireETA}
\newcounter{sujetsujlitteraireEXS}
\newcounter{sujetsujlitteraireMM}
\newcounter{sujetsujlitteraireCC}
\newcounter{sujetsujlitteraireHV}
\newcounter{sujetsujlitteraireHL}
\newcounter{sujetsujlitteraireEXSTA}
\newcounter{sujetsujlitteraireETAMM}
\newcounter{sujetsujlitteraireEXSMM}
\newcounter{sujetsujlitteraireHVCC}
\newcounter{sujetsujlitteraireHVL}
\newcounter{sujetsujlitteraireaut}
%%Phil
\newcounter{sujetsujphilosophiqueETA}
\newcounter{sujetsujphilosophiqueEXS}
\newcounter{sujetsujphilosophiqueMM}
\newcounter{sujetsujphilosophiqueCC}
\newcounter{sujetsujphilosophiqueHV}
\newcounter{sujetsujphilosophiqueHL}
\newcounter{sujetsujphilosophiqueEXSTA}
\newcounter{sujetsujphilosophiqueETAMM}
\newcounter{sujetsujphilosophiqueEXSMM}
\newcounter{sujetsujphilosophiqueHVCC}
\newcounter{sujetsujphilosophiqueHVL}
\newcounter{sujetsujphilosophiqueaut}
%%Totaux pour les graphiques
\newcounter{sujetsujlitterairetotal}
\newcounter{sujetsujphilosophiquetotal}
%%% Autres (Ods et Odh) comptés ensemble
\newcounter{sujetsujallaut}
\newcommand{\compteursfinauxexamens}{
\sommecompteursexamens[amnord]
\sommecompteursexamens[amsud]
\sommecompteursexamens[asi]
\sommecompteursexamens[cea]
\sommecompteursexamens[lib]
\sommecompteursexamens[unk]
\sommecompteursexamens[met]
\sommecompteursexamens[nca]
\sommecompteursexamens[pol]
\sommecompteursexamens[reu]
\sommecompteursexamens[suj]
\sommecompteursexamens
}

\NewDocumentCommand{\sommecompteursexamens}{O{}}{

\setcounter{sujet#1allETA}{\numexpr\value{sujet#1ETA}+\value{sujet#1EXSTA}+\value{sujet#1ETAMM}+\value{sujet#1ETACC}+\value{sujet#1ETAHV}+\value{sujet#1ETAHL}+\value{sujet#1ETAPdP}+\value{sujet#1ETADFI}+\value{sujet#1ETADMRC}+\value{sujet#1ETAHA}\relax}

\setcounter{sujet#1premtermETA}{\numexpr%
\value{sujet#1ETAPdP}+\value{sujet#1ETADFI}+\value{sujet#1ETADMRC}+\value{sujet#1ETAHA}%
\relax}

\setcounter{sujet#1allEXS}{\numexpr%
\value{sujet#1EXS}+\value{sujet#1EXSTA}+\value{sujet#1EXSMM}+\value{sujet#1EXSHV}+\value{sujet#1EXSHL}+\value{sujet#1EXSCC}+\value{sujet#1EXSPdP}+\value{sujet#1EXSDFI}+\value{sujet#1EXSDMRC}+\value{sujet#1EXSHA}%
\relax}

\setcounter{sujet#1premtermEXS}{\numexpr%
\value{sujet#1EXSPdP}+\value{sujet#1EXSDFI}+\value{sujet#1EXSDMRC}+\value{sujet#1EXSHA}%
\relax}

\setcounter{sujet#1allMM}{\numexpr%
\value{sujet#1MM}+\value{sujet#1EXSMM}+\value{sujet#1ETAMM}+\value{sujet#1MMCC}+\value{sujet#1MMHV}+\value{sujet#1MMHL}+\value{sujet#1MMPdP}+\value{sujet#1MMDFI}+\value{sujet#1MMDMRC}+\value{sujet#1MMHA}%
\relax}

\setcounter{sujet#1premtermMM}{\numexpr%
\value{sujet#1MMPdP}+\value{sujet#1MMDFI}+\value{sujet#1MMDMRC}+\value{sujet#1MMHA}%
\relax}

\setcounter{sujet#1allCC}{\numexpr%
\value{sujet#1CC}+\value{sujet#1HVCC}+\value{sujet#1HLCC}+\value{sujet#1ETACC}+\value{sujet#1EXSCC}+\value{sujet#1MMCC}+\value{sujet#1CCPdP}+\value{sujet#1CCDFI}+\value{sujet#1CCDMRC}+\value{sujet#1CCHA}%
\relax}

\setcounter{sujet#1premtermCC}{\numexpr%
\value{sujet#1CCPdP}+\value{sujet#1CCDFI}+\value{sujet#1CCDMRC}+\value{sujet#1CCHA}%
\relax}

\setcounter{sujet#1allHV}{\numexpr%
\value{sujet#1HV}+\value{sujet#1HVL}+\value{sujet#1HVCC}+\value{sujet#1ETAHV}+\value{sujet#1EXSHV}+\value{sujet#1MMHV}+\value{sujet#1HVPdP}+\value{sujet#1HVDFI}+\value{sujet#1HVDMRC}+\value{sujet#1HVHA}%
\relax}

\setcounter{sujet#1premtermHV}{\numexpr%
\value{sujet#1HVPdP}+\value{sujet#1HVDFI}+\value{sujet#1HVDMRC}+\value{sujet#1HVHA}%
\relax}

\setcounter{sujet#1allHL}{\numexpr%
\value{sujet#1HL}+\value{sujet#1HVL}+\value{sujet#1HLCC}+\value{sujet#1ETAHL}+\value{sujet#1EXSHL}+\value{sujet#1MMHL}+\value{sujet#1HLPdP}+\value{sujet#1HLDFI}+\value{sujet#1HLDMRC}+\value{sujet#1HLHA}%
\relax}

\setcounter{sujet#1premtermHL}{\numexpr%
\value{sujet#1HLPdP}+\value{sujet#1HLDFI}+\value{sujet#1HLDMRC}+\value{sujet#1HLHA}%
\relax}

\setcounter{sujet#1allintETA}{\numexpr\value{sujet#1intETA}+\value{sujet#1intEXSTA}+\value{sujet#1intETAMM}+\value{sujet#1intETACC}+\value{sujet#1intETAHV}+\value{sujet#1intETAHL}+\value{sujet#1intETAPdP}+\value{sujet#1intETADFI}+\value{sujet#1intETADMRC}+\value{sujet#1intETAHA}\relax}

\setcounter{sujet#1intpremtermETA}{\numexpr%
\value{sujet#1intETAPdP}+\value{sujet#1intETADFI}+\value{sujet#1intETADMRC}+\value{sujet#1intETAHA}%
\relax}

\setcounter{sujet#1allintEXS}{\numexpr%
\value{sujet#1intEXS}+\value{sujet#1intEXSTA}+\value{sujet#1intEXSMM}+\value{sujet#1intEXSHV}+\value{sujet#1intEXSHL}+\value{sujet#1intEXSCC}+\value{sujet#1intEXSPdP}+\value{sujet#1intEXSDFI}+\value{sujet#1intEXSDMRC}+\value{sujet#1intEXSHA}%
\relax}

\setcounter{sujet#1intpremtermEXS}{\numexpr%
\value{sujet#1intEXSPdP}+\value{sujet#1intEXSDFI}+\value{sujet#1intEXSDMRC}+\value{sujet#1intEXSHA}%
\relax}

\setcounter{sujet#1allintMM}{\numexpr%
\value{sujet#1intMM}+\value{sujet#1intEXSMM}+\value{sujet#1intETAMM}+\value{sujet#1intMMCC}+\value{sujet#1intMMHV}+\value{sujet#1intMMHL}+\value{sujet#1intMMPdP}+\value{sujet#1intMMDFI}+\value{sujet#1intMMDMRC}+\value{sujet#1intMMHA}%
\relax}

\setcounter{sujet#1intpremtermMM}{\numexpr%
\value{sujet#1intMMPdP}+\value{sujet#1intMMDFI}+\value{sujet#1intMMDMRC}+\value{sujet#1intMMHA}%
\relax}

\setcounter{sujet#1allintCC}{\numexpr%
\value{sujet#1intCC}+\value{sujet#1intHVCC}+\value{sujet#1intHLCC}+\value{sujet#1intETACC}+\value{sujet#1intEXSCC}+\value{sujet#1intMMCC}+\value{sujet#1intCCPdP}+\value{sujet#1intCCDFI}+\value{sujet#1intCCDMRC}+\value{sujet#1intCCHA}%
\relax}

\setcounter{sujet#1intpremtermCC}{\numexpr%
\value{sujet#1intCCPdP}+\value{sujet#1intCCDFI}+\value{sujet#1intCCDMRC}+\value{sujet#1intCCHA}%
\relax}

\setcounter{sujet#1allintHV}{\numexpr%
\value{sujet#1intHV}+\value{sujet#1intHVL}+\value{sujet#1intHVCC}+\value{sujet#1intETAHV}+\value{sujet#1intEXSHV}+\value{sujet#1intMMHV}+\value{sujet#1intHVPdP}+\value{sujet#1intHVDFI}+\value{sujet#1intHVDMRC}+\value{sujet#1intHVHA}%
\relax}

\setcounter{sujet#1intpremtermHV}{\numexpr%
\value{sujet#1intHVPdP}+\value{sujet#1intHVDFI}+\value{sujet#1intHVDMRC}+\value{sujet#1intHVHA}%
\relax}

\setcounter{sujet#1allintHL}{\numexpr%
\value{sujet#1intHL}+\value{sujet#1intHVL}+\value{sujet#1intHLCC}+\value{sujet#1intETAHL}+\value{sujet#1intEXSHL}+\value{sujet#1intMMHL}+\value{sujet#1intHLPdP}+\value{sujet#1intHLDFI}+\value{sujet#1intHLDMRC}+\value{sujet#1intHLHA}%
\relax}

\setcounter{sujet#1intpremtermHL}{\numexpr%
\value{sujet#1intHLPdP}+\value{sujet#1intHLDFI}+\value{sujet#1intHLDMRC}+\value{sujet#1intHLHA}%
\relax}

%%%%%%%%%%%%%Int Litt.%%%%%%%%%%%%%%%%%%%%%%%
\setcounter{sujet#1allintlitteraireETA}{\numexpr\value{sujet#1intlitteraireETA}+\value{sujet#1intlitteraireEXSTA}+\value{sujet#1intlitteraireETAMM}+\value{sujet#1intlitteraireETACC}+\value{sujet#1intlitteraireETAHV}+\value{sujet#1intlitteraireETAHL}+\value{sujet#1intlitteraireETAPdP}+\value{sujet#1intlitteraireETADFI}+\value{sujet#1intlitteraireETADMRC}+\value{sujet#1intlitteraireETAHA}\relax}

\setcounter{sujet#1intlitterairepremtermETA}{\numexpr%
\value{sujet#1intlitteraireETAPdP}+\value{sujet#1intlitteraireETADFI}+\value{sujet#1intlitteraireETADMRC}+\value{sujet#1intlitteraireETAHA}%
\relax}

\setcounter{sujet#1allintlitteraireEXS}{\numexpr%
\value{sujet#1intlitteraireEXS}+\value{sujet#1intlitteraireEXSTA}+\value{sujet#1intlitteraireEXSMM}+\value{sujet#1intlitteraireEXSHV}+\value{sujet#1intlitteraireEXSHL}+\value{sujet#1intlitteraireEXSCC}+\value{sujet#1intlitteraireEXSPdP}+\value{sujet#1intlitteraireEXSDFI}+\value{sujet#1intlitteraireEXSDMRC}+\value{sujet#1intlitteraireEXSHA}%
\relax}

\setcounter{sujet#1intlitterairepremtermEXS}{\numexpr%
\value{sujet#1intlitteraireEXSPdP}+\value{sujet#1intlitteraireEXSDFI}+\value{sujet#1intlitteraireEXSDMRC}+\value{sujet#1intlitteraireEXSHA}%
\relax}

\setcounter{sujet#1allintlitteraireMM}{\numexpr%
\value{sujet#1intlitteraireMM}+\value{sujet#1intlitteraireEXSMM}+\value{sujet#1intlitteraireETAMM}+\value{sujet#1intlitteraireMMCC}+\value{sujet#1intlitteraireMMHV}+\value{sujet#1intlitteraireMMHL}+\value{sujet#1intlitteraireMMPdP}+\value{sujet#1intlitteraireMMDFI}+\value{sujet#1intlitteraireMMDMRC}+\value{sujet#1intlitteraireMMHA}%
\relax}

\setcounter{sujet#1intlitterairepremtermMM}{\numexpr%
\value{sujet#1intlitteraireMMPdP}+\value{sujet#1intlitteraireMMDFI}+\value{sujet#1intlitteraireMMDMRC}+\value{sujet#1intlitteraireMMHA}%
\relax}

\setcounter{sujet#1allintlitteraireCC}{\numexpr%
\value{sujet#1intlitteraireCC}+\value{sujet#1intlitteraireHVCC}+\value{sujet#1intlitteraireHLCC}+\value{sujet#1intlitteraireETACC}+\value{sujet#1intlitteraireEXSCC}+\value{sujet#1intlitteraireMMCC}+\value{sujet#1intlitteraireCCPdP}+\value{sujet#1intlitteraireCCDFI}+\value{sujet#1intlitteraireCCDMRC}+\value{sujet#1intlitteraireCCHA}%
\relax}

\setcounter{sujet#1intlitterairepremtermCC}{\numexpr%
\value{sujet#1intlitteraireCCPdP}+\value{sujet#1intlitteraireCCDFI}+\value{sujet#1intlitteraireCCDMRC}+\value{sujet#1intlitteraireCCHA}%
\relax}

\setcounter{sujet#1allintlitteraireHV}{\numexpr%
\value{sujet#1intlitteraireHV}+\value{sujet#1intlitteraireHVL}+\value{sujet#1intlitteraireHVCC}+\value{sujet#1intlitteraireETAHV}+\value{sujet#1intlitteraireEXSHV}+\value{sujet#1intlitteraireMMHV}+\value{sujet#1intlitteraireHVPdP}+\value{sujet#1intlitteraireHVDFI}+\value{sujet#1intlitteraireHVDMRC}+\value{sujet#1intlitteraireHVHA}%
\relax}

\setcounter{sujet#1intlitterairepremtermHV}{\numexpr%
\value{sujet#1intlitteraireHVPdP}+\value{sujet#1intlitteraireHVDFI}+\value{sujet#1intlitteraireHVDMRC}+\value{sujet#1intlitteraireHVHA}%
\relax}

\setcounter{sujet#1allintlitteraireHL}{\numexpr%
\value{sujet#1intlitteraireHL}+\value{sujet#1intlitteraireHVL}+\value{sujet#1intlitteraireHLCC}+\value{sujet#1intlitteraireETAHL}+\value{sujet#1intlitteraireEXSHL}+\value{sujet#1intlitteraireMMHL}+\value{sujet#1intlitteraireHLPdP}+\value{sujet#1intlitteraireHLDFI}+\value{sujet#1intlitteraireHLDMRC}+\value{sujet#1intlitteraireHLHA}%
\relax}

\setcounter{sujet#1intlitterairepremtermHL}{\numexpr%
\value{sujet#1intlitteraireHLPdP}+\value{sujet#1intlitteraireHLDFI}+\value{sujet#1intlitteraireHLDMRC}+\value{sujet#1intlitteraireHLHA}%
\relax}

%%%%%%%%%%%%%Int Phil.%%%%%%%%%%%%%%%%%%%%%%%
\setcounter{sujet#1allintphilosophiqueETA}{\numexpr\value{sujet#1intphilosophiqueETA}+\value{sujet#1intphilosophiqueEXSTA}+\value{sujet#1intphilosophiqueETAMM}+\value{sujet#1intphilosophiqueETACC}+\value{sujet#1intphilosophiqueETAHV}+\value{sujet#1intphilosophiqueETAHL}+\value{sujet#1intphilosophiqueETAPdP}+\value{sujet#1intphilosophiqueETADFI}+\value{sujet#1intphilosophiqueETADMRC}+\value{sujet#1intphilosophiqueETAHA}\relax}

\setcounter{sujet#1intphilosophiquepremtermETA}{\numexpr%
\value{sujet#1intphilosophiqueETAPdP}+\value{sujet#1intphilosophiqueETADFI}+\value{sujet#1intphilosophiqueETADMRC}+\value{sujet#1intphilosophiqueETAHA}%
\relax}

\setcounter{sujet#1allintphilosophiqueEXS}{\numexpr%
\value{sujet#1intphilosophiqueEXS}+\value{sujet#1intphilosophiqueEXSTA}+\value{sujet#1intphilosophiqueEXSMM}+\value{sujet#1intphilosophiqueEXSHV}+\value{sujet#1intphilosophiqueEXSHL}+\value{sujet#1intphilosophiqueEXSCC}+\value{sujet#1intphilosophiqueEXSPdP}+\value{sujet#1intphilosophiqueEXSDFI}+\value{sujet#1intphilosophiqueEXSDMRC}+\value{sujet#1intphilosophiqueEXSHA}%
\relax}

\setcounter{sujet#1intphilosophiquepremtermEXS}{\numexpr%
\value{sujet#1intphilosophiqueEXSPdP}+\value{sujet#1intphilosophiqueEXSDFI}+\value{sujet#1intphilosophiqueEXSDMRC}+\value{sujet#1intphilosophiqueEXSHA}%
\relax}

\setcounter{sujet#1allintphilosophiqueMM}{\numexpr%
\value{sujet#1intphilosophiqueMM}+\value{sujet#1intphilosophiqueEXSMM}+\value{sujet#1intphilosophiqueETAMM}+\value{sujet#1intphilosophiqueMMCC}+\value{sujet#1intphilosophiqueMMHV}+\value{sujet#1intphilosophiqueMMHL}+\value{sujet#1intphilosophiqueMMPdP}+\value{sujet#1intphilosophiqueMMDFI}+\value{sujet#1intphilosophiqueMMDMRC}+\value{sujet#1intphilosophiqueMMHA}%
\relax}

\setcounter{sujet#1intphilosophiquepremtermMM}{\numexpr%
\value{sujet#1intphilosophiqueMMPdP}+\value{sujet#1intphilosophiqueMMDFI}+\value{sujet#1intphilosophiqueMMDMRC}+\value{sujet#1intphilosophiqueMMHA}%
\relax}

\setcounter{sujet#1allintphilosophiqueCC}{\numexpr%
\value{sujet#1intphilosophiqueCC}+\value{sujet#1intphilosophiqueHVCC}+\value{sujet#1intphilosophiqueHLCC}+\value{sujet#1intphilosophiqueETACC}+\value{sujet#1intphilosophiqueEXSCC}+\value{sujet#1intphilosophiqueMMCC}+\value{sujet#1intphilosophiqueCCPdP}+\value{sujet#1intphilosophiqueCCDFI}+\value{sujet#1intphilosophiqueCCDMRC}+\value{sujet#1intphilosophiqueCCHA}%
\relax}

\setcounter{sujet#1intphilosophiquepremtermCC}{\numexpr%
\value{sujet#1intphilosophiqueCCPdP}+\value{sujet#1intphilosophiqueCCDFI}+\value{sujet#1intphilosophiqueCCDMRC}+\value{sujet#1intphilosophiqueCCHA}%
\relax}

\setcounter{sujet#1allintphilosophiqueHV}{\numexpr%
\value{sujet#1intphilosophiqueHV}+\value{sujet#1intphilosophiqueHVL}+\value{sujet#1intphilosophiqueHVCC}+\value{sujet#1intphilosophiqueETAHV}+\value{sujet#1intphilosophiqueEXSHV}+\value{sujet#1intphilosophiqueMMHV}+\value{sujet#1intphilosophiqueHVPdP}+\value{sujet#1intphilosophiqueHVDFI}+\value{sujet#1intphilosophiqueHVDMRC}+\value{sujet#1intphilosophiqueHVHA}%
\relax}

\setcounter{sujet#1intphilosophiquepremtermHV}{\numexpr%
\value{sujet#1intphilosophiqueHVPdP}+\value{sujet#1intphilosophiqueHVDFI}+\value{sujet#1intphilosophiqueHVDMRC}+\value{sujet#1intphilosophiqueHVHA}%
\relax}

\setcounter{sujet#1allintphilosophiqueHL}{\numexpr%
\value{sujet#1intphilosophiqueHL}+\value{sujet#1intphilosophiqueHVL}+\value{sujet#1intphilosophiqueHLCC}+\value{sujet#1intphilosophiqueETAHL}+\value{sujet#1intphilosophiqueEXSHL}+\value{sujet#1intphilosophiqueMMHL}+\value{sujet#1intphilosophiqueHLPdP}+\value{sujet#1intphilosophiqueHLDFI}+\value{sujet#1intphilosophiqueHLDMRC}+\value{sujet#1intphilosophiqueHLHA}%
\relax}

\setcounter{sujet#1intphilosophiquepremtermHL}{\numexpr%
\value{sujet#1intphilosophiqueHLPdP}+\value{sujet#1intphilosophiqueHLDFI}+\value{sujet#1intphilosophiqueHLDMRC}+\value{sujet#1intphilosophiqueHLHA}%
\relax}

%%%%%%%%%%%%%%%Questions de réflexion%%%%%%%%%%
\setcounter{sujet#1allrefETA}{\numexpr\value{sujet#1refETA}+\value{sujet#1refEXSTA}+\value{sujet#1refETAMM}+\value{sujet#1refETACC}+\value{sujet#1refETAHV}+\value{sujet#1refETAHL}+\value{sujet#1refETAPdP}+\value{sujet#1refETADFI}+\value{sujet#1refETADMRC}+\value{sujet#1refETAHA}\relax}

\setcounter{sujet#1refpremtermETA}{\numexpr%
\value{sujet#1refETAPdP}+\value{sujet#1refETADFI}+\value{sujet#1refETADMRC}+\value{sujet#1refETAHA}%
\relax}

\setcounter{sujet#1allrefEXS}{\numexpr%
\value{sujet#1refEXS}+\value{sujet#1refEXSTA}+\value{sujet#1refEXSMM}+\value{sujet#1refEXSHV}+\value{sujet#1refEXSHL}+\value{sujet#1refEXSCC}+\value{sujet#1refEXSPdP}+\value{sujet#1refEXSDFI}+\value{sujet#1refEXSDMRC}+\value{sujet#1refEXSHA}%
\relax}

\setcounter{sujet#1refpremtermEXS}{\numexpr%
\value{sujet#1refEXSPdP}+\value{sujet#1refEXSDFI}+\value{sujet#1refEXSDMRC}+\value{sujet#1refEXSHA}%
\relax}

\setcounter{sujet#1allrefMM}{\numexpr%
\value{sujet#1refMM}+\value{sujet#1refEXSMM}+\value{sujet#1refETAMM}+\value{sujet#1refMMCC}+\value{sujet#1refMMHV}+\value{sujet#1refMMHL}+\value{sujet#1refMMPdP}+\value{sujet#1refMMDFI}+\value{sujet#1refMMDMRC}+\value{sujet#1refMMHA}%
\relax}

\setcounter{sujet#1refpremtermMM}{\numexpr%
\value{sujet#1refMMPdP}+\value{sujet#1refMMDFI}+\value{sujet#1refMMDMRC}+\value{sujet#1refMMHA}%
\relax}

\setcounter{sujet#1allrefCC}{\numexpr%
\value{sujet#1refCC}+\value{sujet#1refHVCC}+\value{sujet#1refHLCC}+\value{sujet#1refETACC}+\value{sujet#1refEXSCC}+\value{sujet#1refMMCC}+\value{sujet#1refCCPdP}+\value{sujet#1refCCDFI}+\value{sujet#1refCCDMRC}+\value{sujet#1refCCHA}%
\relax}

\setcounter{sujet#1refpremtermCC}{\numexpr%
\value{sujet#1refCCPdP}+\value{sujet#1refCCDFI}+\value{sujet#1refCCDMRC}+\value{sujet#1refCCHA}%
\relax}

\setcounter{sujet#1allrefHV}{\numexpr%
\value{sujet#1refHV}+\value{sujet#1refHVL}+\value{sujet#1refHVCC}+\value{sujet#1refETAHV}+\value{sujet#1refEXSHV}+\value{sujet#1refMMHV}+\value{sujet#1refHVPdP}+\value{sujet#1refHVDFI}+\value{sujet#1refHVDMRC}+\value{sujet#1refHVHA}%
\relax}

\setcounter{sujet#1refpremtermHV}{\numexpr%
\value{sujet#1refHVPdP}+\value{sujet#1refHVDFI}+\value{sujet#1refHVDMRC}+\value{sujet#1refHVHA}%
\relax}

\setcounter{sujet#1allrefHL}{\numexpr%
\value{sujet#1refHL}+\value{sujet#1refHVL}+\value{sujet#1refHLCC}+\value{sujet#1refETAHL}+\value{sujet#1refEXSHL}+\value{sujet#1refMMHL}+\value{sujet#1refHLPdP}+\value{sujet#1refHLDFI}+\value{sujet#1refHLDMRC}+\value{sujet#1refHLHA}%
\relax}

\setcounter{sujet#1refpremtermHL}{\numexpr%
\value{sujet#1refHLPdP}+\value{sujet#1refHLDFI}+\value{sujet#1refHLDMRC}+\value{sujet#1refHLHA}%
\relax}

\setcounter{sujet#1allreflitteraireETA}{\numexpr\value{sujet#1reflitteraireETA}+\value{sujet#1reflitteraireEXSTA}+\value{sujet#1reflitteraireETAMM}+\value{sujet#1reflitteraireETACC}+\value{sujet#1reflitteraireETAHV}+\value{sujet#1reflitteraireETAHL}+\value{sujet#1reflitteraireETAPdP}+\value{sujet#1reflitteraireETADFI}+\value{sujet#1reflitteraireETADMRC}+\value{sujet#1reflitteraireETAHA}\relax}

\setcounter{sujet#1reflitterairepremtermETA}{\numexpr%
\value{sujet#1reflitteraireETAPdP}+\value{sujet#1reflitteraireETADFI}+\value{sujet#1reflitteraireETADMRC}+\value{sujet#1reflitteraireETAHA}%
\relax}

\setcounter{sujet#1allreflitteraireEXS}{\numexpr%
\value{sujet#1reflitteraireEXS}+\value{sujet#1reflitteraireEXSTA}+\value{sujet#1reflitteraireEXSMM}+\value{sujet#1reflitteraireEXSHV}+\value{sujet#1reflitteraireEXSHL}+\value{sujet#1reflitteraireEXSCC}+\value{sujet#1reflitteraireEXSPdP}+\value{sujet#1reflitteraireEXSDFI}+\value{sujet#1reflitteraireEXSDMRC}+\value{sujet#1reflitteraireEXSHA}%
\relax}

\setcounter{sujet#1reflitterairepremtermEXS}{\numexpr%
\value{sujet#1reflitteraireEXSPdP}+\value{sujet#1reflitteraireEXSDFI}+\value{sujet#1reflitteraireEXSDMRC}+\value{sujet#1reflitteraireEXSHA}%
\relax}

\setcounter{sujet#1allreflitteraireMM}{\numexpr%
\value{sujet#1reflitteraireMM}+\value{sujet#1reflitteraireEXSMM}+\value{sujet#1reflitteraireETAMM}+\value{sujet#1reflitteraireMMCC}+\value{sujet#1reflitteraireMMHV}+\value{sujet#1reflitteraireMMHL}+\value{sujet#1reflitteraireMMPdP}+\value{sujet#1reflitteraireMMDFI}+\value{sujet#1reflitteraireMMDMRC}+\value{sujet#1reflitteraireMMHA}%
\relax}

\setcounter{sujet#1reflitterairepremtermMM}{\numexpr%
\value{sujet#1reflitteraireMMPdP}+\value{sujet#1reflitteraireMMDFI}+\value{sujet#1reflitteraireMMDMRC}+\value{sujet#1reflitteraireMMHA}%
\relax}

\setcounter{sujet#1allreflitteraireCC}{\numexpr%
\value{sujet#1reflitteraireCC}+\value{sujet#1reflitteraireHVCC}+\value{sujet#1reflitteraireHLCC}+\value{sujet#1reflitteraireETACC}+\value{sujet#1reflitteraireEXSCC}+\value{sujet#1reflitteraireMMCC}+\value{sujet#1reflitteraireCCPdP}+\value{sujet#1reflitteraireCCDFI}+\value{sujet#1reflitteraireCCDMRC}+\value{sujet#1reflitteraireCCHA}%
\relax}

\setcounter{sujet#1reflitterairepremtermCC}{\numexpr%
\value{sujet#1reflitteraireCCPdP}+\value{sujet#1reflitteraireCCDFI}+\value{sujet#1reflitteraireCCDMRC}+\value{sujet#1reflitteraireCCHA}%
\relax}

\setcounter{sujet#1allreflitteraireHV}{\numexpr%
\value{sujet#1reflitteraireHV}+\value{sujet#1reflitteraireHVL}+\value{sujet#1reflitteraireHVCC}+\value{sujet#1reflitteraireETAHV}+\value{sujet#1reflitteraireEXSHV}+\value{sujet#1reflitteraireMMHV}+\value{sujet#1reflitteraireHVPdP}+\value{sujet#1reflitteraireHVDFI}+\value{sujet#1reflitteraireHVDMRC}+\value{sujet#1reflitteraireHVHA}%
\relax}

\setcounter{sujet#1reflitterairepremtermHV}{\numexpr%
\value{sujet#1reflitteraireHVPdP}+\value{sujet#1reflitteraireHVDFI}+\value{sujet#1reflitteraireHVDMRC}+\value{sujet#1reflitteraireHVHA}%
\relax}

\setcounter{sujet#1allreflitteraireHL}{\numexpr%
\value{sujet#1reflitteraireHL}+\value{sujet#1reflitteraireHVL}+\value{sujet#1reflitteraireHLCC}+\value{sujet#1reflitteraireETAHL}+\value{sujet#1reflitteraireEXSHL}+\value{sujet#1reflitteraireMMHL}+\value{sujet#1reflitteraireHLPdP}+\value{sujet#1reflitteraireHLDFI}+\value{sujet#1reflitteraireHLDMRC}+\value{sujet#1reflitteraireHLHA}%
\relax}

\setcounter{sujet#1reflitterairepremtermHL}{\numexpr%
\value{sujet#1reflitteraireHLPdP}+\value{sujet#1reflitteraireHLDFI}+\value{sujet#1reflitteraireHLDMRC}+\value{sujet#1reflitteraireHLHA}%
\relax}

%%%%%%%%%Réfl. philosophique%%%%
\setcounter{sujet#1allrefphilosophiqueETA}{\numexpr\value{sujet#1refphilosophiqueETA}+\value{sujet#1refphilosophiqueEXSTA}+\value{sujet#1refphilosophiqueETAMM}+\value{sujet#1refphilosophiqueETACC}+\value{sujet#1refphilosophiqueETAHV}+\value{sujet#1refphilosophiqueETAHL}+\value{sujet#1refphilosophiqueETAPdP}+\value{sujet#1refphilosophiqueETADFI}+\value{sujet#1refphilosophiqueETADMRC}+\value{sujet#1refphilosophiqueETAHA}\relax}

\setcounter{sujet#1refphilosophiquepremtermETA}{\numexpr%
\value{sujet#1refphilosophiqueETAPdP}+\value{sujet#1refphilosophiqueETADFI}+\value{sujet#1refphilosophiqueETADMRC}+\value{sujet#1refphilosophiqueETAHA}%
\relax}

\setcounter{sujet#1allrefphilosophiqueEXS}{\numexpr%
\value{sujet#1refphilosophiqueEXS}+\value{sujet#1refphilosophiqueEXSTA}+\value{sujet#1refphilosophiqueEXSMM}+\value{sujet#1refphilosophiqueEXSHV}+\value{sujet#1refphilosophiqueEXSHL}+\value{sujet#1refphilosophiqueEXSCC}+\value{sujet#1refphilosophiqueEXSPdP}+\value{sujet#1refphilosophiqueEXSDFI}+\value{sujet#1refphilosophiqueEXSDMRC}+\value{sujet#1refphilosophiqueEXSHA}%
\relax}

\setcounter{sujet#1refphilosophiquepremtermEXS}{\numexpr%
\value{sujet#1refphilosophiqueEXSPdP}+\value{sujet#1refphilosophiqueEXSDFI}+\value{sujet#1refphilosophiqueEXSDMRC}+\value{sujet#1refphilosophiqueEXSHA}%
\relax}

\setcounter{sujet#1allrefphilosophiqueMM}{\numexpr%
\value{sujet#1refphilosophiqueMM}+\value{sujet#1refphilosophiqueEXSMM}+\value{sujet#1refphilosophiqueETAMM}+\value{sujet#1refphilosophiqueMMCC}+\value{sujet#1refphilosophiqueMMHV}+\value{sujet#1refphilosophiqueMMHL}+\value{sujet#1refphilosophiqueMMPdP}+\value{sujet#1refphilosophiqueMMDFI}+\value{sujet#1refphilosophiqueMMDMRC}+\value{sujet#1refphilosophiqueMMHA}%
\relax}

\setcounter{sujet#1refphilosophiquepremtermMM}{\numexpr%
\value{sujet#1refphilosophiqueMMPdP}+\value{sujet#1refphilosophiqueMMDFI}+\value{sujet#1refphilosophiqueMMDMRC}+\value{sujet#1refphilosophiqueMMHA}%
\relax}

\setcounter{sujet#1allrefphilosophiqueCC}{\numexpr%
\value{sujet#1refphilosophiqueCC}+\value{sujet#1refphilosophiqueHVCC}+\value{sujet#1refphilosophiqueHLCC}+\value{sujet#1refphilosophiqueETACC}+\value{sujet#1refphilosophiqueEXSCC}+\value{sujet#1refphilosophiqueMMCC}+\value{sujet#1refphilosophiqueCCPdP}+\value{sujet#1refphilosophiqueCCDFI}+\value{sujet#1refphilosophiqueCCDMRC}+\value{sujet#1refphilosophiqueCCHA}%
\relax}

\setcounter{sujet#1refphilosophiquepremtermCC}{\numexpr%
\value{sujet#1refphilosophiqueCCPdP}+\value{sujet#1refphilosophiqueCCDFI}+\value{sujet#1refphilosophiqueCCDMRC}+\value{sujet#1refphilosophiqueCCHA}%
\relax}

\setcounter{sujet#1allrefphilosophiqueHV}{\numexpr%
\value{sujet#1refphilosophiqueHV}+\value{sujet#1refphilosophiqueHVL}+\value{sujet#1refphilosophiqueHVCC}+\value{sujet#1refphilosophiqueETAHV}+\value{sujet#1refphilosophiqueEXSHV}+\value{sujet#1refphilosophiqueMMHV}+\value{sujet#1refphilosophiqueHVPdP}+\value{sujet#1refphilosophiqueHVDFI}+\value{sujet#1refphilosophiqueHVDMRC}+\value{sujet#1refphilosophiqueHVHA}%
\relax}

\setcounter{sujet#1refphilosophiquepremtermHV}{\numexpr%
\value{sujet#1refphilosophiqueHVPdP}+\value{sujet#1refphilosophiqueHVDFI}+\value{sujet#1refphilosophiqueHVDMRC}+\value{sujet#1refphilosophiqueHVHA}%
\relax}

\setcounter{sujet#1allrefphilosophiqueHL}{\numexpr%
\value{sujet#1refphilosophiqueHL}+\value{sujet#1refphilosophiqueHVL}+\value{sujet#1refphilosophiqueHLCC}+\value{sujet#1refphilosophiqueETAHL}+\value{sujet#1refphilosophiqueEXSHL}+\value{sujet#1refphilosophiqueMMHL}+\value{sujet#1refphilosophiqueHLPdP}+\value{sujet#1refphilosophiqueHLDFI}+\value{sujet#1refphilosophiqueHLDMRC}+\value{sujet#1refphilosophiqueHLHA}%
\relax}

\setcounter{sujet#1refphilosophiquepremtermHL}{\numexpr%
\value{sujet#1refphilosophiqueHLPdP}+\value{sujet#1refphilosophiqueHLDFI}+\value{sujet#1refphilosophiqueHLDMRC}+\value{sujet#1refphilosophiqueHLHA}%
\relax}

%%%%%%%%%%%%%%%%%%%%%%%%%%%%%%%%%%%%%%%%%%%%%%%%%%%%%%%%%%%%%%%%%%%%
%%%%%%%%%%%%%%%%%Compteurs des sujet#1s sur plusieurs années%%%%%%%%%%
%%%%%%%%%%%%%%%%%%%%%%%%%%%%%%%%%%%%%%%%%%%%%%%%%%%%%%%%%%%%%%%%%%%%

\setcounter{sujet#1RSPdP}{\numexpr\value{sujet#1ETAPdP}+\value{sujet#1EXSPdP}+\value{sujet#1MMPdP}\relax}

\setcounter{sujet#1HQPdP}{\numexpr\value{sujet#1CCPdP}+\value{sujet#1HVPdP}+\value{sujet#1HLPdP}\relax}

\setcounter{sujet#1RSDFI}{\numexpr\value{sujet#1ETADFI}+\value{sujet#1EXSDFI}+\value{sujet#1MMDFI}\relax}

\setcounter{sujet#1HQDFI}{\numexpr\value{sujet#1CCDFI}+\value{sujet#1HVDFI}+\value{sujet#1HLDFI}\relax}

\setcounter{sujet#1RSDMRC}{\numexpr\value{sujet#1ETADMRC}+\value{sujet#1EXSDMRC}+\value{sujet#1MMDMRC}\relax}

\setcounter{sujet#1HQDMRC}{\numexpr\value{sujet#1CCDMRC}+\value{sujet#1HVDMRC}+\value{sujet#1HLDMRC}\relax}

\setcounter{sujet#1RSHA}{\numexpr\value{sujet#1ETAHA}+\value{sujet#1EXSHA}+\value{sujet#1ETAHA}\relax}

\setcounter{sujet#1HQHA}{\numexpr\value{sujet#1CCHA}+\value{sujet#1HVHA}+\value{sujet#1HLHA}\relax}

%%%%%%%%%Interprétations
\setcounter{sujet#1intRSPdP}{\numexpr\value{sujet#1intETAPdP}+\value{sujet#1intEXSPdP}+\value{sujet#1intMMPdP}\relax}

\setcounter{sujet#1intHQPdP}{\numexpr\value{sujet#1intCCPdP}+\value{sujet#1intHVPdP}+\value{sujet#1intHLPdP}\relax}

\setcounter{sujet#1intRSDFI}{\numexpr\value{sujet#1intETADFI}+\value{sujet#1intEXSDFI}+\value{sujet#1intMMDFI}\relax}

\setcounter{sujet#1intHQDFI}{\numexpr\value{sujet#1intCCDFI}+\value{sujet#1intHVDFI}+\value{sujet#1intHLDFI}\relax}

\setcounter{sujet#1intRSDMRC}{\numexpr\value{sujet#1intETADMRC}+\value{sujet#1intEXSDMRC}+\value{sujet#1intMMDMRC}\relax}

\setcounter{sujet#1intHQDMRC}{\numexpr\value{sujet#1intCCDMRC}+\value{sujet#1intHVDMRC}+\value{sujet#1intHLDMRC}\relax}

\setcounter{sujet#1intRSHA}{\numexpr\value{sujet#1intETAHA}+\value{sujet#1intEXSHA}+\value{sujet#1intETAHA}\relax}

\setcounter{sujet#1intHQHA}{\numexpr\value{sujet#1intCCHA}+\value{sujet#1intHVHA}+\value{sujet#1intHLHA}\relax}

%%%%%%%%%%%%%%Littéraire%%%%%%%%%%%%%%%%%%%%%
\setcounter{sujet#1intlitteraireRSPdP}{\numexpr\value{sujet#1intlitteraireETAPdP}+\value{sujet#1intlitteraireEXSPdP}+\value{sujet#1intlitteraireMMPdP}\relax}

\setcounter{sujet#1intlitteraireHQPdP}{\numexpr\value{sujet#1intlitteraireCCPdP}+\value{sujet#1intlitteraireHVPdP}+\value{sujet#1intlitteraireHLPdP}\relax}

\setcounter{sujet#1intlitteraireRSDFI}{\numexpr\value{sujet#1intlitteraireETADFI}+\value{sujet#1intlitteraireEXSDFI}+\value{sujet#1intlitteraireMMDFI}\relax}

\setcounter{sujet#1intlitteraireHQDFI}{\numexpr\value{sujet#1intlitteraireCCDFI}+\value{sujet#1intlitteraireHVDFI}+\value{sujet#1intlitteraireHLDFI}\relax}

\setcounter{sujet#1intlitteraireRSDMRC}{\numexpr\value{sujet#1intlitteraireETADMRC}+\value{sujet#1intlitteraireEXSDMRC}+\value{sujet#1intlitteraireMMDMRC}\relax}

\setcounter{sujet#1intlitteraireHQDMRC}{\numexpr\value{sujet#1intlitteraireCCDMRC}+\value{sujet#1intlitteraireHVDMRC}+\value{sujet#1intlitteraireHLDMRC}\relax}

\setcounter{sujet#1intlitteraireRSHA}{\numexpr\value{sujet#1intlitteraireETAHA}+\value{sujet#1intlitteraireEXSHA}+\value{sujet#1intlitteraireETAHA}\relax}

\setcounter{sujet#1intlitteraireHQHA}{\numexpr\value{sujet#1intlitteraireCCHA}+\value{sujet#1intlitteraireHVHA}+\value{sujet#1intlitteraireHLHA}\relax}
%%%%%%%%%%%%%%Philosophique%%%%%%%%%%%%%%%%%%
\setcounter{sujet#1intphilosophiqueRSPdP}{\numexpr\value{sujet#1intphilosophiqueETAPdP}+\value{sujet#1intphilosophiqueEXSPdP}+\value{sujet#1intphilosophiqueMMPdP}\relax}

\setcounter{sujet#1intphilosophiqueHQPdP}{\numexpr\value{sujet#1intphilosophiqueCCPdP}+\value{sujet#1intphilosophiqueHVPdP}+\value{sujet#1intphilosophiqueHLPdP}\relax}

\setcounter{sujet#1intphilosophiqueRSDFI}{\numexpr\value{sujet#1intphilosophiqueETADFI}+\value{sujet#1intphilosophiqueEXSDFI}+\value{sujet#1intphilosophiqueMMDFI}\relax}

\setcounter{sujet#1intphilosophiqueHQDFI}{\numexpr\value{sujet#1intphilosophiqueCCDFI}+\value{sujet#1intphilosophiqueHVDFI}+\value{sujet#1intphilosophiqueHLDFI}\relax}

\setcounter{sujet#1intphilosophiqueRSDMRC}{\numexpr\value{sujet#1intphilosophiqueETADMRC}+\value{sujet#1intphilosophiqueEXSDMRC}+\value{sujet#1intphilosophiqueMMDMRC}\relax}

\setcounter{sujet#1intphilosophiqueHQDMRC}{\numexpr\value{sujet#1intphilosophiqueCCDMRC}+\value{sujet#1intphilosophiqueHVDMRC}+\value{sujet#1intphilosophiqueHLDMRC}\relax}

\setcounter{sujet#1intphilosophiqueRSHA}{\numexpr\value{sujet#1intphilosophiqueETAHA}+\value{sujet#1intphilosophiqueEXSHA}+\value{sujet#1intphilosophiqueETAHA}\relax}

\setcounter{sujet#1intphilosophiqueHQHA}{\numexpr\value{sujet#1intphilosophiqueCCHA}+\value{sujet#1intphilosophiqueHVHA}+\value{sujet#1intphilosophiqueHLHA}\relax}
%%%%%%%%%%%%%%Réflexions%%%%%%%%%%%%%%%%%%%%%
\setcounter{sujet#1refRSPdP}{\numexpr\value{sujet#1refETAPdP}+\value{sujet#1refEXSPdP}+\value{sujet#1refMMPdP}\relax}

\setcounter{sujet#1refHQPdP}{\numexpr\value{sujet#1refCCPdP}+\value{sujet#1refHVPdP}+\value{sujet#1refHLPdP}\relax}

\setcounter{sujet#1refRSDFI}{\numexpr\value{sujet#1refETADFI}+\value{sujet#1refEXSDFI}+\value{sujet#1refMMDFI}\relax}

\setcounter{sujet#1refHQDFI}{\numexpr\value{sujet#1refCCDFI}+\value{sujet#1refHVDFI}+\value{sujet#1refHLDFI}\relax}

\setcounter{sujet#1refRSDMRC}{\numexpr\value{sujet#1refETADMRC}+\value{sujet#1refEXSDMRC}+\value{sujet#1refMMDMRC}\relax}

\setcounter{sujet#1refHQDMRC}{\numexpr\value{sujet#1refCCDMRC}+\value{sujet#1refHVDMRC}+\value{sujet#1refHLDMRC}\relax}

\setcounter{sujet#1refRSHA}{\numexpr\value{sujet#1refETAHA}+\value{sujet#1refEXSHA}+\value{sujet#1refETAHA}\relax}

\setcounter{sujet#1refHQHA}{\numexpr\value{sujet#1refCCHA}+\value{sujet#1refHVHA}+\value{sujet#1refHLHA}\relax}

%%%%%%%%%%%%%%Littéraire%%%%%%%%%%%%%%%%%%%%%
\setcounter{sujet#1reflitteraireRSPdP}{\numexpr\value{sujet#1reflitteraireETAPdP}+\value{sujet#1reflitteraireEXSPdP}+\value{sujet#1reflitteraireMMPdP}\relax}

\setcounter{sujet#1reflitteraireHQPdP}{\numexpr\value{sujet#1reflitteraireCCPdP}+\value{sujet#1reflitteraireHVPdP}+\value{sujet#1reflitteraireHLPdP}\relax}

\setcounter{sujet#1reflitteraireRSDFI}{\numexpr\value{sujet#1reflitteraireETADFI}+\value{sujet#1reflitteraireEXSDFI}+\value{sujet#1reflitteraireMMDFI}\relax}

\setcounter{sujet#1reflitteraireHQDFI}{\numexpr\value{sujet#1reflitteraireCCDFI}+\value{sujet#1reflitteraireHVDFI}+\value{sujet#1reflitteraireHLDFI}\relax}

\setcounter{sujet#1reflitteraireRSDMRC}{\numexpr\value{sujet#1reflitteraireETADMRC}+\value{sujet#1reflitteraireEXSDMRC}+\value{sujet#1reflitteraireMMDMRC}\relax}

\setcounter{sujet#1reflitteraireHQDMRC}{\numexpr\value{sujet#1reflitteraireCCDMRC}+\value{sujet#1reflitteraireHVDMRC}+\value{sujet#1reflitteraireHLDMRC}\relax}

\setcounter{sujet#1reflitteraireRSHA}{\numexpr\value{sujet#1reflitteraireETAHA}+\value{sujet#1reflitteraireEXSHA}+\value{sujet#1reflitteraireETAHA}\relax}

\setcounter{sujet#1reflitteraireHQHA}{\numexpr\value{sujet#1reflitteraireCCHA}+\value{sujet#1reflitteraireHVHA}+\value{sujet#1reflitteraireHLHA}\relax}

%%%%%%%%%%%%%%Philosophique%%%%%%%%%%%%%%%%%%
\setcounter{sujet#1refphilosophiqueRSPdP}{\numexpr\value{sujet#1refphilosophiqueETAPdP}+\value{sujet#1refphilosophiqueEXSPdP}+\value{sujet#1refphilosophiqueMMPdP}\relax}

\setcounter{sujet#1refphilosophiqueHQPdP}{\numexpr\value{sujet#1refphilosophiqueCCPdP}+\value{sujet#1refphilosophiqueHVPdP}+\value{sujet#1refphilosophiqueHLPdP}\relax}

\setcounter{sujet#1refphilosophiqueRSDFI}{\numexpr\value{sujet#1refphilosophiqueETADFI}+\value{sujet#1refphilosophiqueEXSDFI}+\value{sujet#1refphilosophiqueMMDFI}\relax}

\setcounter{sujet#1refphilosophiqueHQDFI}{\numexpr\value{sujet#1refphilosophiqueCCDFI}+\value{sujet#1refphilosophiqueHVDFI}+\value{sujet#1refphilosophiqueHLDFI}\relax}

\setcounter{sujet#1refphilosophiqueRSDMRC}{\numexpr\value{sujet#1refphilosophiqueETADMRC}+\value{sujet#1refphilosophiqueEXSDMRC}+\value{sujet#1refphilosophiqueMMDMRC}\relax}

\setcounter{sujet#1refphilosophiqueHQDMRC}{\numexpr\value{sujet#1refphilosophiqueCCDMRC}+\value{sujet#1refphilosophiqueHVDMRC}+\value{sujet#1refphilosophiqueHLDMRC}\relax}

\setcounter{sujet#1refphilosophiqueRSHA}{\numexpr\value{sujet#1refphilosophiqueETAHA}+\value{sujet#1refphilosophiqueEXSHA}+\value{sujet#1refphilosophiqueETAHA}\relax}

\setcounter{sujet#1refphilosophiqueHQHA}{\numexpr\value{sujet#1refphilosophiqueCCHA}+\value{sujet#1refphilosophiqueHVHA}+\value{sujet#1refphilosophiqueHLHA}\relax}

%%%%Compteurs globaux%%%
\setcounter{sujet#1RSintlitteraire}{\numexpr\value{sujet#1intlitteraireETA}+\value{sujet#1intlitteraireEXS}+\value{sujet#1intlitteraireMM}+\value{sujet#1intlitteraireEXSTA}+\value{sujet#1intlitteraireETAMM}+\value{sujet#1intlitteraireETACC}+\value{sujet#1intlitteraireEXSMM}+\value{sujet#1intlitteraireETAHV}+\value{sujet#1intlitteraireETAHL}+\value{sujet#1intlitteraireETAPdP}+\value{sujet#1intlitteraireETADFI}+\value{sujet#1intlitteraireETADMRC}+\value{sujet#1intlitteraireETAHA}+\value{sujet#1intlitteraireEXSHV}+\value{sujet#1intlitteraireEXSHL}+\value{sujet#1intlitteraireEXSCC}+\value{sujet#1intlitteraireEXSPdP}+\value{sujet#1intlitteraireEXSDFI}+\value{sujet#1intlitteraireEXSDMRC}+\value{sujet#1intlitteraireEXSHA}+\value{sujet#1intlitteraireMMCC}+\value{sujet#1intlitteraireMMHV}+\value{sujet#1intlitteraireMMHL}+\value{sujet#1intlitteraireMMPdP}+\value{sujet#1intlitteraireMMDFI}+\value{sujet#1intlitteraireMMDMRC}+\value{sujet#1intlitteraireMMHA}+\value{sujet#1intlitteraireOdS}\relax}

\setcounter{sujet#1HQintlitteraire}{%
\numexpr\value{sujet#1intlitteraireCC}+\value{sujet#1intlitteraireHV}+\value{sujet#1intlitteraireHL}+\value{sujet#1intlitteraireHVCC}+\value{sujet#1intlitteraireHLCC}+\value{sujet#1intlitteraireHVL}+\value{sujet#1intlitteraireETACC}+\value{sujet#1intlitteraireEXSCC}+\value{sujet#1intlitteraireMMCC}+\value{sujet#1intlitteraireCCPdP}+\value{sujet#1intlitteraireCCDFI}+\value{sujet#1intlitteraireCCDMRC}+\value{sujet#1intlitteraireCCHA}+\value{sujet#1intlitteraireETAHV}+\value{sujet#1intlitteraireEXSHV}+\value{sujet#1intlitteraireMMHV}+\value{sujet#1intlitteraireHVPdP}+\value{sujet#1intlitteraireHVDFI}+\value{sujet#1intlitteraireHVDMRC}+\value{sujet#1intlitteraireHVHA}+\value{sujet#1intlitteraireETAHL}+\value{sujet#1intlitteraireEXSHL}+\value{sujet#1intlitteraireMMHL}+\value{sujet#1intlitteraireHLPdP}+\value{sujet#1intlitteraireHLDFI}+\value{sujet#1intlitteraireHLDMRC}+\value{sujet#1intlitteraireHLHA}+\value{sujet#1intlitteraireOdH}\relax}

\setcounter{sujet#1RSintphilosophique}{\numexpr\value{sujet#1intphilosophiqueETA}+\value{sujet#1intphilosophiqueEXS}+\value{sujet#1intphilosophiqueMM}+\value{sujet#1intphilosophiqueEXSTA}+\value{sujet#1intphilosophiqueETAMM}+\value{sujet#1intphilosophiqueETACC}+\value{sujet#1intphilosophiqueEXSMM}+\value{sujet#1intphilosophiqueETAHV}+\value{sujet#1intphilosophiqueETAHL}+\value{sujet#1intphilosophiqueETAPdP}+\value{sujet#1intphilosophiqueETADFI}+\value{sujet#1intphilosophiqueETADMRC}+\value{sujet#1intphilosophiqueETAHA}+\value{sujet#1intphilosophiqueEXSHV}+\value{sujet#1intphilosophiqueEXSHL}+\value{sujet#1intphilosophiqueEXSCC}+\value{sujet#1intphilosophiqueEXSPdP}+\value{sujet#1intphilosophiqueEXSDFI}+\value{sujet#1intphilosophiqueEXSDMRC}+\value{sujet#1intphilosophiqueEXSHA}+\value{sujet#1intphilosophiqueMMCC}+\value{sujet#1intphilosophiqueMMHV}+\value{sujet#1intphilosophiqueMMHL}+\value{sujet#1intphilosophiqueMMPdP}+\value{sujet#1intphilosophiqueMMDFI}+\value{sujet#1intphilosophiqueMMDMRC}+\value{sujet#1intphilosophiqueMMHA}+\value{sujet#1intphilosophiqueOdS}\relax}

\setcounter{sujet#1HQintphilosophique}{%
\numexpr\value{sujet#1intphilosophiqueCC}+\value{sujet#1intphilosophiqueHV}+\value{sujet#1intphilosophiqueHL}+\value{sujet#1intphilosophiqueHVCC}+\value{sujet#1intphilosophiqueHLCC}+\value{sujet#1intphilosophiqueHVL}+\value{sujet#1intphilosophiqueETACC}+\value{sujet#1intphilosophiqueEXSCC}+\value{sujet#1intphilosophiqueMMCC}+\value{sujet#1intphilosophiqueCCPdP}+\value{sujet#1intphilosophiqueCCDFI}+\value{sujet#1intphilosophiqueCCDMRC}+\value{sujet#1intphilosophiqueCCHA}+\value{sujet#1intphilosophiqueETAHV}+\value{sujet#1intphilosophiqueEXSHV}+\value{sujet#1intphilosophiqueMMHV}+\value{sujet#1intphilosophiqueHVPdP}+\value{sujet#1intphilosophiqueHVDFI}+\value{sujet#1intphilosophiqueHVDMRC}+\value{sujet#1intphilosophiqueHVHA}+\value{sujet#1intphilosophiqueETAHL}+\value{sujet#1intphilosophiqueEXSHL}+\value{sujet#1intphilosophiqueMMHL}+\value{sujet#1intphilosophiqueHLPdP}+\value{sujet#1intphilosophiqueHLDFI}+\value{sujet#1intphilosophiqueHLDMRC}+\value{sujet#1intphilosophiqueHLHA}+\value{sujet#1intphilosophiqueOdH}\relax}

\setcounter{sujet#1RSreflitteraire}{\numexpr\value{sujet#1reflitteraireETA}+\value{sujet#1reflitteraireEXS}+\value{sujet#1reflitteraireMM}+\value{sujet#1reflitteraireEXSTA}+\value{sujet#1reflitteraireETAMM}+\value{sujet#1reflitteraireETACC}+\value{sujet#1reflitteraireEXSMM}+\value{sujet#1reflitteraireETAHV}+\value{sujet#1reflitteraireETAHL}+\value{sujet#1reflitteraireETAPdP}+\value{sujet#1reflitteraireETADFI}+\value{sujet#1reflitteraireETADMRC}+\value{sujet#1reflitteraireETAHA}+\value{sujet#1reflitteraireEXSHV}+\value{sujet#1reflitteraireEXSHL}+\value{sujet#1reflitteraireEXSCC}+\value{sujet#1reflitteraireEXSPdP}+\value{sujet#1reflitteraireEXSDFI}+\value{sujet#1reflitteraireEXSDMRC}+\value{sujet#1reflitteraireEXSHA}+\value{sujet#1reflitteraireMMCC}+\value{sujet#1reflitteraireMMHV}+\value{sujet#1reflitteraireMMHL}+\value{sujet#1reflitteraireMMPdP}+\value{sujet#1reflitteraireMMDFI}+\value{sujet#1reflitteraireMMDMRC}+\value{sujet#1reflitteraireMMHA}+\value{sujet#1reflitteraireOdS}\relax}

\setcounter{sujet#1HQreflitteraire}{%
\numexpr\value{sujet#1reflitteraireCC}+\value{sujet#1reflitteraireHV}+\value{sujet#1reflitteraireHL}+\value{sujet#1reflitteraireHVCC}+\value{sujet#1reflitteraireHLCC}+\value{sujet#1reflitteraireHVL}+\value{sujet#1reflitteraireETACC}+\value{sujet#1reflitteraireEXSCC}+\value{sujet#1reflitteraireMMCC}+\value{sujet#1reflitteraireCCPdP}+\value{sujet#1reflitteraireCCDFI}+\value{sujet#1reflitteraireCCDMRC}+\value{sujet#1reflitteraireCCHA}+\value{sujet#1reflitteraireETAHV}+\value{sujet#1reflitteraireEXSHV}+\value{sujet#1reflitteraireMMHV}+\value{sujet#1reflitteraireHVPdP}+\value{sujet#1reflitteraireHVDFI}+\value{sujet#1reflitteraireHVDMRC}+\value{sujet#1reflitteraireHVHA}+\value{sujet#1reflitteraireETAHL}+\value{sujet#1reflitteraireEXSHL}+\value{sujet#1reflitteraireMMHL}+\value{sujet#1reflitteraireHLPdP}+\value{sujet#1reflitteraireHLDFI}+\value{sujet#1reflitteraireHLDMRC}+\value{sujet#1reflitteraireHLHA}+\value{sujet#1reflitteraireOdH}\relax}

\setcounter{sujet#1RSrefphilosophique}{\numexpr\value{sujet#1refphilosophiqueETA}+\value{sujet#1refphilosophiqueEXS}+\value{sujet#1refphilosophiqueMM}+\value{sujet#1refphilosophiqueEXSTA}+\value{sujet#1refphilosophiqueETAMM}+\value{sujet#1refphilosophiqueETACC}+\value{sujet#1refphilosophiqueEXSMM}+\value{sujet#1refphilosophiqueETAHV}+\value{sujet#1refphilosophiqueETAHL}+\value{sujet#1refphilosophiqueETAPdP}+\value{sujet#1refphilosophiqueETADFI}+\value{sujet#1refphilosophiqueETADMRC}+\value{sujet#1refphilosophiqueETAHA}+\value{sujet#1refphilosophiqueEXSHV}+\value{sujet#1refphilosophiqueEXSHL}+\value{sujet#1refphilosophiqueEXSCC}+\value{sujet#1refphilosophiqueEXSPdP}+\value{sujet#1refphilosophiqueEXSDFI}+\value{sujet#1refphilosophiqueEXSDMRC}+\value{sujet#1refphilosophiqueEXSHA}+\value{sujet#1refphilosophiqueMMCC}+\value{sujet#1refphilosophiqueMMHV}+\value{sujet#1refphilosophiqueMMHL}+\value{sujet#1refphilosophiqueMMPdP}+\value{sujet#1refphilosophiqueMMDFI}+\value{sujet#1refphilosophiqueMMDMRC}+\value{sujet#1refphilosophiqueMMHA}+\value{sujet#1refphilosophiqueOdS}\relax}

\setcounter{sujet#1HQrefphilosophique}{%
\numexpr\value{sujet#1refphilosophiqueCC}+\value{sujet#1refphilosophiqueHV}+\value{sujet#1refphilosophiqueHL}+\value{sujet#1refphilosophiqueHVCC}+\value{sujet#1refphilosophiqueHLCC}+\value{sujet#1refphilosophiqueHVL}+\value{sujet#1refphilosophiqueETACC}+\value{sujet#1refphilosophiqueEXSCC}+\value{sujet#1refphilosophiqueMMCC}+\value{sujet#1refphilosophiqueCCPdP}+\value{sujet#1refphilosophiqueCCDFI}+\value{sujet#1refphilosophiqueCCDMRC}+\value{sujet#1refphilosophiqueCCHA}+\value{sujet#1refphilosophiqueETAHV}+\value{sujet#1refphilosophiqueEXSHV}+\value{sujet#1refphilosophiqueMMHV}+\value{sujet#1refphilosophiqueHVPdP}+\value{sujet#1refphilosophiqueHVDFI}+\value{sujet#1refphilosophiqueHVDMRC}+\value{sujet#1refphilosophiqueHVHA}+\value{sujet#1refphilosophiqueETAHL}+\value{sujet#1refphilosophiqueEXSHL}+\value{sujet#1refphilosophiqueMMHL}+\value{sujet#1refphilosophiqueHLPdP}+\value{sujet#1refphilosophiqueHLDFI}+\value{sujet#1refphilosophiqueHLDMRC}+\value{sujet#1refphilosophiqueHLHA}+\value{sujet#1refphilosophiqueOdH}\relax}

\setcounter{sujet#1intlitterairepremtermRS}{\numexpr%
\value{sujet#1intlitterairepremtermETA}+\value{sujet#1intlitterairepremtermEXS}+\value{sujet#1intlitterairepremtermMM}\relax}

\setcounter{sujet#1intlitterairepremtermHQ}{\numexpr%
\value{sujet#1intlitterairepremtermCC}+\value{sujet#1intlitterairepremtermHV}+\value{sujet#1intlitterairepremtermHL}\relax}

\setcounter{sujet#1intlitterairepremterm}{\numexpr%
\value{sujet#1intlitterairepremtermRS}+\value{sujet#1intlitterairepremtermHQ}\relax}

\setcounter{sujet#1reflitterairepremtermRS}{\numexpr%
\value{sujet#1reflitterairepremtermETA}+\value{sujet#1reflitterairepremtermEXS}+\value{sujet#1reflitterairepremtermMM}\relax}

\setcounter{sujet#1reflitterairepremtermHQ}{\numexpr%
\value{sujet#1reflitterairepremtermCC}+\value{sujet#1reflitterairepremtermHV}+\value{sujet#1reflitterairepremtermHL}\relax}

\setcounter{sujet#1reflitterairepremterm}{\numexpr%
\value{sujet#1reflitterairepremtermRS}+\value{sujet#1reflitterairepremtermHQ}\relax}

\setcounter{sujet#1intphilosophiquepremtermRS}{\numexpr%
\value{sujet#1intphilosophiquepremtermETA}+\value{sujet#1intphilosophiquepremtermEXS}+\value{sujet#1intphilosophiquepremtermMM}\relax}

\setcounter{sujet#1intphilosophiquepremtermHQ}{\numexpr%
\value{sujet#1intphilosophiquepremtermCC}+\value{sujet#1intphilosophiquepremtermHV}+\value{sujet#1intphilosophiquepremtermHL}\relax}

\setcounter{sujet#1intphilosophiquepremterm}{\numexpr%
\value{sujet#1intphilosophiquepremtermRS}+\value{sujet#1intphilosophiquepremtermHQ}\relax}

\setcounter{sujet#1refphilosophiquepremtermRS}{\numexpr%
\value{sujet#1refphilosophiquepremtermETA}+\value{sujet#1refphilosophiquepremtermEXS}+\value{sujet#1refphilosophiquepremtermMM}\relax}

\setcounter{sujet#1refphilosophiquepremtermHQ}{\numexpr%
\value{sujet#1refphilosophiquepremtermCC}+\value{sujet#1refphilosophiquepremtermHV}+\value{sujet#1refphilosophiquepremtermHL}\relax}

\setcounter{sujet#1refphilosophiquepremterm}{\numexpr%
\value{sujet#1refphilosophiquepremtermRS}+\value{sujet#1refphilosophiquepremtermHQ}\relax}

\setcounter{sujet#1premterm}{\numexpr\value{sujet#1intlitterairepremterm}+\value{sujet#1reflitterairepremterm}+\value{sujet#1intphilosophiquepremterm}+\value{sujet#1refphilosophiquepremterm}\relax}

\setcounter{sujet#1all}{
\numexpr\value{sujet#1RSintlitteraire}+\value{sujet#1HQintlitteraire}+\value{sujet#1HQintphilosophique}+\value{sujet#1RSintphilosophique}-\value{sujet#1intlitteraireETACC}-\value{sujet#1intlitteraireEXSCC}-\value{sujet#1intlitteraireMMCC}-\value{sujet#1intlitteraireETAHV}-\value{sujet#1intlitteraireETAHL}-\value{sujet#1intlitteraireMMHV}-\value{sujet#1intlitteraireMMHL}-\value{sujet#1intlitteraireEXSHV}-\value{sujet#1intlitteraireEXSHL}-\value{sujet#1intphilosophiqueETACC}-\value{sujet#1intphilosophiqueEXSCC}-\value{sujet#1intphilosophiqueMMCC}-\value{sujet#1intphilosophiqueETAHV}-\value{sujet#1intphilosophiqueETAHL}-\value{sujet#1intphilosophiqueMMHV}-\value{sujet#1intphilosophiqueMMHL}-\value{sujet#1intphilosophiqueEXSHV}-\value{sujet#1intphilosophiqueEXSHL}+
\value{sujet#1RSreflitteraire}+\value{sujet#1HQreflitteraire}+\value{sujet#1HQrefphilosophique}+\value{sujet#1RSrefphilosophique}-\value{sujet#1reflitteraireETACC}-\value{sujet#1reflitteraireEXSCC}-\value{sujet#1reflitteraireMMCC}-\value{sujet#1reflitteraireETAHV}-\value{sujet#1reflitteraireETAHL}-\value{sujet#1reflitteraireMMHV}-\value{sujet#1reflitteraireMMHL}-\value{sujet#1reflitteraireEXSHV}-\value{sujet#1reflitteraireEXSHL}-\value{sujet#1refphilosophiqueETACC}-\value{sujet#1refphilosophiqueEXSCC}-\value{sujet#1refphilosophiqueMMCC}-\value{sujet#1refphilosophiqueETAHV}-\value{sujet#1refphilosophiqueETAHL}-\value{sujet#1refphilosophiqueMMHV}-\value{sujet#1refphilosophiqueMMHL}-\value{sujet#1refphilosophiqueEXSHV}-\value{sujet#1refphilosophiqueEXSHL}
\relax}

%%%Somme par discipline
\setcounter{sujet#1litteraireETA}
{\numexpr%
\value{sujet#1intlitteraireETA}+\value{sujet#1reflitteraireETA}+\value{sujet#1intlitteraireEXSTA}+\value{sujet#1reflitteraireEXSTA}+\value{sujet#1intlitteraireETAMM}+\value{sujet#1reflitteraireETAMM}+
\value{sujet#1intlitteraireETACC}+\value{sujet#1reflitteraireETACC}+\value{sujet#1intlitteraireETAHV}+\value{sujet#1reflitteraireETAHV}+\value{sujet#1intlitteraireETAHL}+\value{sujet#1reflitteraireETAHL}+\value{sujet#1intlitteraireETAPdP}+\value{sujet#1reflitteraireETAPdP}+\value{sujet#1intlitteraireETADFI}+\value{sujet#1reflitteraireETADFI}+\value{sujet#1intlitteraireETADMRC}+\value{sujet#1reflitteraireETADMRC}+\value{sujet#1intlitteraireETAHA}+\value{sujet#1reflitteraireETAHA}
\relax}

\setcounter{sujet#1litteraireEXS}{\numexpr%
\value{sujet#1intlitteraireEXS}+\value{sujet#1intlitteraireEXSTA}+\value{sujet#1intlitteraireEXSMM}+\value{sujet#1intlitteraireEXSHV}+\value{sujet#1intlitteraireEXSHL}+\value{sujet#1intlitteraireEXSCC}+\value{sujet#1intlitteraireEXSPdP}+\value{sujet#1intlitteraireEXSDFI}+\value{sujet#1intlitteraireEXSDMRC}+\value{sujet#1intlitteraireEXSHA}+\value{sujet#1reflitteraireEXS}+\value{sujet#1reflitteraireEXSTA}+\value{sujet#1reflitteraireEXSMM}+\value{sujet#1reflitteraireEXSHV}+\value{sujet#1reflitteraireEXSHL}+\value{sujet#1reflitteraireEXSCC}+\value{sujet#1reflitteraireEXSPdP}+\value{sujet#1reflitteraireEXSDFI}+\value{sujet#1reflitteraireEXSDMRC}+\value{sujet#1reflitteraireEXSHA}
\relax}

\setcounter{sujet#1litteraireMM}{\numexpr%
\value{sujet#1intlitteraireMM}+\value{sujet#1intlitteraireEXSMM}+\value{sujet#1intlitteraireETAMM}+\value{sujet#1intlitteraireMMCC}+\value{sujet#1intlitteraireMMHV}+\value{sujet#1intlitteraireMMHL}+\value{sujet#1intlitteraireMMPdP}+\value{sujet#1intlitteraireMMDFI}+\value{sujet#1intlitteraireMMDMRC}+\value{sujet#1intlitteraireMMHA}%
+\value{sujet#1reflitteraireMM}+\value{sujet#1reflitteraireEXSMM}+\value{sujet#1reflitteraireETAMM}+\value{sujet#1reflitteraireMMCC}+\value{sujet#1reflitteraireMMHV}+\value{sujet#1reflitteraireMMHL}+\value{sujet#1reflitteraireMMPdP}+\value{sujet#1reflitteraireMMDFI}+\value{sujet#1reflitteraireMMDMRC}+\value{sujet#1reflitteraireMMHA}%
\relax}

\setcounter{sujet#1litteraireCC}{\numexpr%
\value{sujet#1intlitteraireCC}+\value{sujet#1intlitteraireHVCC}+\value{sujet#1intlitteraireHLCC}+\value{sujet#1intlitteraireETACC}+\value{sujet#1intlitteraireEXSCC}+\value{sujet#1intlitteraireMMCC}+\value{sujet#1intlitteraireCCPdP}+\value{sujet#1intlitteraireCCDFI}+\value{sujet#1intlitteraireCCDMRC}+\value{sujet#1intlitteraireCCHA}%
+\value{sujet#1reflitteraireCC}+\value{sujet#1reflitteraireHVCC}+\value{sujet#1reflitteraireHLCC}+\value{sujet#1reflitteraireETACC}+\value{sujet#1reflitteraireEXSCC}+\value{sujet#1reflitteraireMMCC}+\value{sujet#1reflitteraireCCPdP}+\value{sujet#1reflitteraireCCDFI}+\value{sujet#1reflitteraireCCDMRC}+\value{sujet#1reflitteraireCCHA}
\relax}

\setcounter{sujet#1litteraireHV}{\numexpr%
\value{sujet#1intlitteraireHV}+\value{sujet#1intlitteraireHVL}+\value{sujet#1intlitteraireHVCC}+\value{sujet#1intlitteraireETAHV}+\value{sujet#1intlitteraireEXSHV}+\value{sujet#1intlitteraireMMHV}+\value{sujet#1intlitteraireHVPdP}+\value{sujet#1intlitteraireHVDFI}+\value{sujet#1intlitteraireHVDMRC}+\value{sujet#1intlitteraireHVHA}%
+\value{sujet#1reflitteraireHV}+\value{sujet#1reflitteraireHVL}+\value{sujet#1reflitteraireHVCC}+\value{sujet#1reflitteraireETAHV}+\value{sujet#1reflitteraireEXSHV}+\value{sujet#1reflitteraireMMHV}+\value{sujet#1reflitteraireHVPdP}+\value{sujet#1reflitteraireHVDFI}+\value{sujet#1reflitteraireHVDMRC}+\value{sujet#1reflitteraireHVHA}%
\relax}

\setcounter{sujet#1litteraireHL}{\numexpr%
\value{sujet#1intlitteraireHL}+\value{sujet#1intlitteraireHVL}+\value{sujet#1intlitteraireHLCC}+\value{sujet#1intlitteraireETAHL}+\value{sujet#1intlitteraireEXSHL}+\value{sujet#1intlitteraireMMHL}+\value{sujet#1intlitteraireHLPdP}+\value{sujet#1intlitteraireHLDFI}+\value{sujet#1intlitteraireHLDMRC}+\value{sujet#1intlitteraireHLHA}%
+\value{sujet#1reflitteraireHL}+\value{sujet#1reflitteraireHVL}+\value{sujet#1reflitteraireHLCC}+\value{sujet#1reflitteraireETAHL}+\value{sujet#1reflitteraireEXSHL}+\value{sujet#1reflitteraireMMHL}+\value{sujet#1reflitteraireHLPdP}+\value{sujet#1reflitteraireHLDFI}+\value{sujet#1reflitteraireHLDMRC}+\value{sujet#1reflitteraireHLHA}%
\relax}

\setcounter{sujet#1litteraireaut}
{\numexpr%
\value{sujet#1intlitteraireOdS}+\value{sujet#1intlitteraireOdH}+\value{sujet#1reflitteraireOdS}+\value{sujet#1reflitteraireOdH}
\relax}

%%%Phil
\setcounter{sujet#1philosophiqueETA}
{\numexpr%
\value{sujet#1intphilosophiqueETA}+\value{sujet#1refphilosophiqueETA}+\value{sujet#1intphilosophiqueEXSTA}+\value{sujet#1refphilosophiqueEXSTA}+\value{sujet#1intphilosophiqueETAMM}+\value{sujet#1refphilosophiqueETAMM}+
\value{sujet#1intphilosophiqueETACC}+\value{sujet#1refphilosophiqueETACC}+\value{sujet#1intphilosophiqueETAHV}+\value{sujet#1refphilosophiqueETAHV}+\value{sujet#1intphilosophiqueETAHL}+\value{sujet#1refphilosophiqueETAHL}+\value{sujet#1intphilosophiqueETAPdP}+\value{sujet#1refphilosophiqueETAPdP}+\value{sujet#1intphilosophiqueETADFI}+\value{sujet#1refphilosophiqueETADFI}+\value{sujet#1intphilosophiqueETADMRC}+\value{sujet#1refphilosophiqueETADMRC}+\value{sujet#1intphilosophiqueETAHA}+\value{sujet#1refphilosophiqueETAHA}
\relax}

\setcounter{sujet#1philosophiqueEXS}{\numexpr%
\value{sujet#1intphilosophiqueEXS}+\value{sujet#1intphilosophiqueEXSTA}+\value{sujet#1intphilosophiqueEXSMM}+\value{sujet#1intphilosophiqueEXSHV}+\value{sujet#1intphilosophiqueEXSHL}+\value{sujet#1intphilosophiqueEXSCC}+\value{sujet#1intphilosophiqueEXSPdP}+\value{sujet#1intphilosophiqueEXSDFI}+\value{sujet#1intphilosophiqueEXSDMRC}+\value{sujet#1intphilosophiqueEXSHA}+\value{sujet#1refphilosophiqueEXS}+\value{sujet#1refphilosophiqueEXSTA}+\value{sujet#1refphilosophiqueEXSMM}+\value{sujet#1refphilosophiqueEXSHV}+\value{sujet#1refphilosophiqueEXSHL}+\value{sujet#1refphilosophiqueEXSCC}+\value{sujet#1refphilosophiqueEXSPdP}+\value{sujet#1refphilosophiqueEXSDFI}+\value{sujet#1refphilosophiqueEXSDMRC}+\value{sujet#1refphilosophiqueEXSHA}
\relax}

\setcounter{sujet#1philosophiqueMM}{\numexpr%
\value{sujet#1intphilosophiqueMM}+\value{sujet#1intphilosophiqueEXSMM}+\value{sujet#1intphilosophiqueETAMM}+\value{sujet#1intphilosophiqueMMCC}+\value{sujet#1intphilosophiqueMMHV}+\value{sujet#1intphilosophiqueMMHL}+\value{sujet#1intphilosophiqueMMPdP}+\value{sujet#1intphilosophiqueMMDFI}+\value{sujet#1intphilosophiqueMMDMRC}+\value{sujet#1intphilosophiqueMMHA}%
+\value{sujet#1refphilosophiqueMM}+\value{sujet#1refphilosophiqueEXSMM}+\value{sujet#1refphilosophiqueETAMM}+\value{sujet#1refphilosophiqueMMCC}+\value{sujet#1refphilosophiqueMMHV}+\value{sujet#1refphilosophiqueMMHL}+\value{sujet#1refphilosophiqueMMPdP}+\value{sujet#1refphilosophiqueMMDFI}+\value{sujet#1refphilosophiqueMMDMRC}+\value{sujet#1refphilosophiqueMMHA}%
\relax}

\setcounter{sujet#1philosophiqueCC}{\numexpr%
\value{sujet#1intphilosophiqueCC}+\value{sujet#1intphilosophiqueHVCC}+\value{sujet#1intphilosophiqueHLCC}+\value{sujet#1intphilosophiqueETACC}+\value{sujet#1intphilosophiqueEXSCC}+\value{sujet#1intphilosophiqueMMCC}+\value{sujet#1intphilosophiqueCCPdP}+\value{sujet#1intphilosophiqueCCDFI}+\value{sujet#1intphilosophiqueCCDMRC}+\value{sujet#1intphilosophiqueCCHA}%
+\value{sujet#1refphilosophiqueCC}+\value{sujet#1refphilosophiqueHVCC}+\value{sujet#1refphilosophiqueHLCC}+\value{sujet#1refphilosophiqueETACC}+\value{sujet#1refphilosophiqueEXSCC}+\value{sujet#1refphilosophiqueMMCC}+\value{sujet#1refphilosophiqueCCPdP}+\value{sujet#1refphilosophiqueCCDFI}+\value{sujet#1refphilosophiqueCCDMRC}+\value{sujet#1refphilosophiqueCCHA}
\relax}

\setcounter{sujet#1philosophiqueHV}{\numexpr%
\value{sujet#1intphilosophiqueHV}+\value{sujet#1intphilosophiqueHVL}+\value{sujet#1intphilosophiqueHVCC}+\value{sujet#1intphilosophiqueETAHV}+\value{sujet#1intphilosophiqueEXSHV}+\value{sujet#1intphilosophiqueMMHV}+\value{sujet#1intphilosophiqueHVPdP}+\value{sujet#1intphilosophiqueHVDFI}+\value{sujet#1intphilosophiqueHVDMRC}+\value{sujet#1intphilosophiqueHVHA}%
+\value{sujet#1refphilosophiqueHV}+\value{sujet#1refphilosophiqueHVL}+\value{sujet#1refphilosophiqueHVCC}+\value{sujet#1refphilosophiqueETAHV}+\value{sujet#1refphilosophiqueEXSHV}+\value{sujet#1refphilosophiqueMMHV}+\value{sujet#1refphilosophiqueHVPdP}+\value{sujet#1refphilosophiqueHVDFI}+\value{sujet#1refphilosophiqueHVDMRC}+\value{sujet#1refphilosophiqueHVHA}%
\relax}

\setcounter{sujet#1philosophiqueHL}{\numexpr%
\value{sujet#1intphilosophiqueHL}+\value{sujet#1intphilosophiqueHVL}+\value{sujet#1intphilosophiqueHLCC}+\value{sujet#1intphilosophiqueETAHL}+\value{sujet#1intphilosophiqueEXSHL}+\value{sujet#1intphilosophiqueMMHL}+\value{sujet#1intphilosophiqueHLPdP}+\value{sujet#1intphilosophiqueHLDFI}+\value{sujet#1intphilosophiqueHLDMRC}+\value{sujet#1intphilosophiqueHLHA}%
+\value{sujet#1refphilosophiqueHL}+\value{sujet#1refphilosophiqueHVL}+\value{sujet#1refphilosophiqueHLCC}+\value{sujet#1refphilosophiqueETAHL}+\value{sujet#1refphilosophiqueEXSHL}+\value{sujet#1refphilosophiqueMMHL}+\value{sujet#1refphilosophiqueHLPdP}+\value{sujet#1refphilosophiqueHLDFI}+\value{sujet#1refphilosophiqueHLDMRC}+\value{sujet#1refphilosophiqueHLHA}%
\relax}

\setcounter{sujet#1philosophiqueaut}
{\numexpr%
\value{sujet#1intphilosophiqueOdS}+\value{sujet#1intphilosophiqueOdH}+\value{sujet#1refphilosophiqueOdS}+\value{sujet#1refphilosophiqueOdH}
\relax}

\setcounter{sujet#1litterairetotal}{\numexpr%
\value{sujet#1litteraireETA}+\value{sujet#1litteraireEXS}+\value{sujet#1litteraireMM}+\value{sujet#1litteraireHV}+\value{sujet#1litteraireHL}+\value{sujet#1litteraireCC}+\value{sujet#1litteraireaut}\relax}

\setcounter{sujet#1philosophiquetotal}{\numexpr%
\value{sujet#1philosophiqueETA}+\value{sujet#1philosophiqueEXS}+\value{sujet#1philosophiqueMM}+\value{sujet#1philosophiqueHV}+\value{sujet#1philosophiqueHL}+\value{sujet#1philosophiqueCC}+\value{sujet#1philosophiqueaut}\relax}

\setcounter{sujet#1allaut}{\numexpr%
\value{sujet#1litteraireaut}+\value{sujet#1philosophiqueaut}\relax}

%%%Pour gérer la production des graphiques
\setcounter{sujet#1exclRS}{\numexpr%
\value{sujet#1ETA}+\value{sujet#1EXS}+\value{sujet#1MM}
\relax}

\setcounter{sujet#1mixtRS}{\numexpr%
\value{sujet#1ETAMM}+\value{sujet#1EXSMM}+\value{sujet#1EXSTA}+\value{sujet#1OdS}
\relax}

\setcounter{sujet#1exclHQ}{\numexpr%
\value{sujet#1CC}+\value{sujet#1HV}+\value{sujet#1HL}
\relax}

\setcounter{sujet#1mixtHQ}{\numexpr%
\value{sujet#1HVCC}+\value{sujet#1HLCC}+\value{sujet#1HVL}+\value{sujet#1OdH}
\relax}

\setcounter{sujet#1mixtHQRS}{\numexpr%
\value{sujet#1ETACC}+\value{sujet#1ETAHV}+\value{sujet#1ETAHL}+\value{sujet#1EXSCC}+\value{sujet#1EXSHV}+\value{sujet#1EXSHL}+\value{sujet#1MMCC}+\value{sujet#1MMHV}+\value{sujet#1MMHL}+\value{sujet#1premterm}
\relax}

\setcounter{sujettrans#1intlitteraire}{\numexpr%
\value{sujet#1intlitterairepremterm}+\value{sujet#1intlitteraireETACC}+\value{sujet#1intlitteraireETAHV}+\value{sujet#1intlitteraireETAHL}+\value{sujet#1intlitteraireEXSCC}+\value{sujet#1intlitteraireEXSHV}+\value{sujet#1intlitteraireEXSHL}+\value{sujet#1intlitteraireMMCC}+\value{sujet#1intlitteraireMMHV}+\value{sujet#1intlitteraireMMHL}+\value{sujet#1intlitteraireOdS}+\value{sujet#1intlitteraireOdH}
\relax}

\setcounter{sujettrans#1reflitteraire}{\numexpr%
\value{sujet#1reflitterairepremterm}+\value{sujet#1reflitteraireETACC}+\value{sujet#1reflitteraireETAHV}+\value{sujet#1reflitteraireETAHL}+\value{sujet#1reflitteraireEXSCC}+\value{sujet#1reflitteraireEXSHV}+\value{sujet#1reflitteraireEXSHL}+\value{sujet#1reflitteraireMMCC}+\value{sujet#1reflitteraireMMHV}+\value{sujet#1reflitteraireMMHL}+
\value{sujet#1reflitteraireOdS}+\value{sujet#1reflitteraireOdH}
\relax}

\setcounter{sujettrans#1intphilosophique}{\numexpr%
\value{sujet#1intphilosophiquepremterm}+\value{sujet#1intphilosophiqueETACC}+\value{sujet#1intphilosophiqueETAHV}+\value{sujet#1intphilosophiqueETAHL}+\value{sujet#1intphilosophiqueEXSCC}+\value{sujet#1intphilosophiqueEXSHV}+\value{sujet#1intphilosophiqueEXSHL}+\value{sujet#1intphilosophiqueMMCC}+\value{sujet#1intphilosophiqueMMHV}+\value{sujet#1intphilosophiqueMMHL}+
\value{sujet#1intphilosophiqueOdS}+\value{sujet#1refphilosophiqueOdH}
\relax}

\setcounter{sujettrans#1refphilosophique}{\numexpr%
\value{sujet#1refphilosophiquepremterm}+\value{sujet#1refphilosophiqueETACC}+\value{sujet#1refphilosophiqueETAHV}+\value{sujet#1refphilosophiqueETAHL}+\value{sujet#1refphilosophiqueEXSCC}+\value{sujet#1refphilosophiqueEXSHV}+\value{sujet#1refphilosophiqueEXSHL}+\value{sujet#1refphilosophiqueMMCC}+\value{sujet#1refphilosophiqueMMHV}+\value{sujet#1refphilosophiqueMMHL}+
\value{sujet#1refphilosophiqueOdS}+\value{sujet#1refphilosophiqueOdH}
\relax}

\setcounter{sujet#1intlitteraire}{\numexpr%
\value{sujet#1intlitteraireEXS}+\value{sujet#1intlitteraireMM}+\value{sujet#1intlitteraireEXSTA}+\value{sujet#1intlitteraireETAMM}+\value{sujet#1intlitteraireETACC}+\value{sujet#1intlitteraireEXSMM}+\value{sujet#1intlitteraireETAHV}+\value{sujet#1intlitteraireETAHL}+\value{sujet#1intlitteraireETAPdP}+\value{sujet#1intlitteraireETADFI}+\value{sujet#1intlitteraireETADMRC}+\value{sujet#1intlitteraireETAHA}+\value{sujet#1intlitteraireEXSHV}+\value{sujet#1intlitteraireEXSHL}+\value{sujet#1intlitteraireEXSCC}+\value{sujet#1intlitteraireEXSPdP}+\value{sujet#1intlitteraireEXSDFI}+\value{sujet#1intlitteraireEXSDMRC}+\value{sujet#1intlitteraireEXSHA}+\value{sujet#1intlitteraireMMCC}+\value{sujet#1intlitteraireMMHV}+\value{sujet#1intlitteraireMMHL}+\value{sujet#1intlitteraireMMPdP}+\value{sujet#1intlitteraireMMDFI}+\value{sujet#1intlitteraireMMDMRC}+\value{sujet#1intlitteraireMMHA}+\value{sujet#1intlitteraireOdS}+\value{sujet#1intlitteraireCC}+\value{sujet#1intlitteraireHV}+\value{sujet#1intlitteraireHL}+\value{sujet#1intlitteraireHVCC}+\value{sujet#1intlitteraireHLCC}+\value{sujet#1intlitteraireHVL}+\value{sujet#1intlitteraireCCPdP}+\value{sujet#1intlitteraireCCDFI}+\value{sujet#1intlitteraireCCDMRC}+\value{sujet#1intlitteraireCCHA}+\value{sujet#1intlitteraireHVPdP}+\value{sujet#1intlitteraireHVDFI}+\value{sujet#1intlitteraireHVDMRC}+\value{sujet#1intlitteraireHVHA}+\value{sujet#1intlitteraireHLPdP}+\value{sujet#1intlitteraireHLDFI}+\value{sujet#1intlitteraireHLDMRC}+\value{sujet#1intlitteraireHLHA}+\value{sujet#1intlitteraireOdH}\relax}

\setcounter{sujet#1reflitteraire}{\numexpr%
\value{sujet#1reflitteraireEXS}+\value{sujet#1reflitteraireMM}+\value{sujet#1reflitteraireEXSTA}+\value{sujet#1reflitteraireETAMM}+\value{sujet#1reflitteraireETACC}+\value{sujet#1reflitteraireEXSMM}+\value{sujet#1reflitteraireETAHV}+\value{sujet#1reflitteraireETAHL}+\value{sujet#1reflitteraireETAPdP}+\value{sujet#1reflitteraireETADFI}+\value{sujet#1reflitteraireETADMRC}+\value{sujet#1reflitteraireETAHA}+\value{sujet#1reflitteraireEXSHV}+\value{sujet#1reflitteraireEXSHL}+\value{sujet#1reflitteraireEXSCC}+\value{sujet#1reflitteraireEXSPdP}+\value{sujet#1reflitteraireEXSDFI}+\value{sujet#1reflitteraireEXSDMRC}+\value{sujet#1reflitteraireEXSHA}+\value{sujet#1reflitteraireMMCC}+\value{sujet#1reflitteraireMMHV}+\value{sujet#1reflitteraireMMHL}+\value{sujet#1reflitteraireMMPdP}+\value{sujet#1reflitteraireMMDFI}+\value{sujet#1reflitteraireMMDMRC}+\value{sujet#1reflitteraireMMHA}+\value{sujet#1reflitteraireOdS}+\value{sujet#1reflitteraireCC}+\value{sujet#1reflitteraireHV}+\value{sujet#1reflitteraireHL}+\value{sujet#1reflitteraireHVCC}+\value{sujet#1reflitteraireHLCC}+\value{sujet#1reflitteraireHVL}+\value{sujet#1reflitteraireCCPdP}+\value{sujet#1reflitteraireCCDFI}+\value{sujet#1reflitteraireCCDMRC}+\value{sujet#1reflitteraireCCHA}+\value{sujet#1reflitteraireHVPdP}+\value{sujet#1reflitteraireHVDFI}+\value{sujet#1reflitteraireHVDMRC}+\value{sujet#1reflitteraireHVHA}+\value{sujet#1reflitteraireHLPdP}+\value{sujet#1reflitteraireHLDFI}+\value{sujet#1reflitteraireHLDMRC}+\value{sujet#1reflitteraireHLHA}+\value{sujet#1reflitteraireOdH}\relax}

\setcounter{sujet#1intphilosophique}{\numexpr%
\value{sujet#1intphilosophiqueEXS}+\value{sujet#1intphilosophiqueMM}+\value{sujet#1intphilosophiqueEXSTA}+\value{sujet#1intphilosophiqueETAMM}+\value{sujet#1intphilosophiqueETACC}+\value{sujet#1intphilosophiqueEXSMM}+\value{sujet#1intphilosophiqueETAHV}+\value{sujet#1intphilosophiqueETAHL}+\value{sujet#1intphilosophiqueETAPdP}+\value{sujet#1intphilosophiqueETADFI}+\value{sujet#1intphilosophiqueETADMRC}+\value{sujet#1intphilosophiqueETAHA}+\value{sujet#1intphilosophiqueEXSHV}+\value{sujet#1intphilosophiqueEXSHL}+\value{sujet#1intphilosophiqueEXSCC}+\value{sujet#1intphilosophiqueEXSPdP}+\value{sujet#1intphilosophiqueEXSDFI}+\value{sujet#1intphilosophiqueEXSDMRC}+\value{sujet#1intphilosophiqueEXSHA}+\value{sujet#1intphilosophiqueMMCC}+\value{sujet#1intphilosophiqueMMHV}+\value{sujet#1intphilosophiqueMMHL}+\value{sujet#1intphilosophiqueMMPdP}+\value{sujet#1intphilosophiqueMMDFI}+\value{sujet#1intphilosophiqueMMDMRC}+\value{sujet#1intphilosophiqueMMHA}+\value{sujet#1intphilosophiqueOdS}+\value{sujet#1intphilosophiqueCC}+\value{sujet#1intphilosophiqueHV}+\value{sujet#1intphilosophiqueHL}+\value{sujet#1intphilosophiqueHVCC}+\value{sujet#1intphilosophiqueHLCC}+\value{sujet#1intphilosophiqueHVL}+\value{sujet#1intphilosophiqueCCPdP}+\value{sujet#1intphilosophiqueCCDFI}+\value{sujet#1intphilosophiqueCCDMRC}+\value{sujet#1intphilosophiqueCCHA}+\value{sujet#1intphilosophiqueHVPdP}+\value{sujet#1intphilosophiqueHVDFI}+\value{sujet#1intphilosophiqueHVDMRC}+\value{sujet#1intphilosophiqueHVHA}+\value{sujet#1intphilosophiqueHLPdP}+\value{sujet#1intphilosophiqueHLDFI}+\value{sujet#1intphilosophiqueHLDMRC}+\value{sujet#1intphilosophiqueHLHA}+\value{sujet#1intphilosophiqueOdH}\relax}

\setcounter{sujet#1refphilosophique}{\numexpr%
\value{sujet#1refphilosophiqueEXS}+\value{sujet#1refphilosophiqueMM}+\value{sujet#1refphilosophiqueEXSTA}+\value{sujet#1refphilosophiqueETAMM}+\value{sujet#1refphilosophiqueETACC}+\value{sujet#1refphilosophiqueEXSMM}+\value{sujet#1refphilosophiqueETAHV}+\value{sujet#1refphilosophiqueETAHL}+\value{sujet#1refphilosophiqueETAPdP}+\value{sujet#1refphilosophiqueETADFI}+\value{sujet#1refphilosophiqueETADMRC}+\value{sujet#1refphilosophiqueETAHA}+\value{sujet#1refphilosophiqueEXSHV}+\value{sujet#1refphilosophiqueEXSHL}+\value{sujet#1refphilosophiqueEXSCC}+\value{sujet#1refphilosophiqueEXSPdP}+\value{sujet#1refphilosophiqueEXSDFI}+\value{sujet#1refphilosophiqueEXSDMRC}+\value{sujet#1refphilosophiqueEXSHA}+\value{sujet#1refphilosophiqueMMCC}+\value{sujet#1refphilosophiqueMMHV}+\value{sujet#1refphilosophiqueMMHL}+\value{sujet#1refphilosophiqueMMPdP}+\value{sujet#1refphilosophiqueMMDFI}+\value{sujet#1refphilosophiqueMMDMRC}+\value{sujet#1refphilosophiqueMMHA}+\value{sujet#1refphilosophiqueOdS}+\value{sujet#1refphilosophiqueCC}+\value{sujet#1refphilosophiqueHV}+\value{sujet#1refphilosophiqueHL}+\value{sujet#1refphilosophiqueHVCC}+\value{sujet#1refphilosophiqueHLCC}+\value{sujet#1refphilosophiqueHVL}+\value{sujet#1refphilosophiqueCCPdP}+\value{sujet#1refphilosophiqueCCDFI}+\value{sujet#1refphilosophiqueCCDMRC}+\value{sujet#1refphilosophiqueCCHA}+\value{sujet#1refphilosophiqueHVPdP}+\value{sujet#1refphilosophiqueHVDFI}+\value{sujet#1refphilosophiqueHVDMRC}+\value{sujet#1refphilosophiqueHVHA}+\value{sujet#1refphilosophiqueHLPdP}+\value{sujet#1refphilosophiqueHLDFI}+\value{sujet#1refphilosophiqueHLDMRC}+\value{sujet#1refphilosophiqueHLHA}+\value{sujet#1refphilosophiqueOdH}\relax}

\setcounter{sujet#1total}{\numexpr%
\value{sujet#1allETA}+\value{sujet#1allEXS}+\value{sujet#1allMM}+\value{sujet#1allCC}+\value{sujet#1allHV}+\value{sujet#1allHL}+\value{sujet#1allaut}
  \relax}%
  }

%%%%%%%%%%%%%%%%%%%%%%Pour l'alignement du texte%%%%%%%%%%%%%%%%%%%%%%%
\newenvironment{envtexte}{\begin{adjustwidth}{15mm}{15mm}\hspace{\parindent}}
{\end{adjustwidth}}


\newcommand{\auteur}[2]{
\ifthenelse{\isempty{#1}}%%%Si l'auteur n'a pas de prénom
{\textsc{#2}\index[a]{\textsc{#2}}}%%Ne mettre que son nom en majuscules
{#1~\textsc{#2}\index[a]{\textsc{#2}, #1}}%%Sinon insérer/indexer tout.
}
\newcommand{\deuxauteurs}[4]{\auteur{#1}{#2}~et~\auteur{#3}{#4}}
%\newcommand{\particule}[3]{\index[a]{\textsc{#1}, #3|see{\textsc{#2}}}}
\newcommand{\chapeau}[1]{\textit{#1} \medskip}
\newcommand{\ouvrage}[1]{\textit{#1}}
\newcommand{\datep}[1]{\hspace{-3pt}(#1)}
\newcommand{\reference}[3]{%%%Inséré à la fin des textes|À améliorer -- Notamment pour le flushright qu'on clôt dans* envtexte%%%
\vspace{-0.3cm}\begin{flushright}\nopagebreak\vspace{3pt}\nopagebreak{{\auteur{#1}{#2}, \ouvrage{#3}}
}\end{flushright}
}%
\newcommand{\referenced}[5]{\vspace*{3pt}{\deuxauteurs{#1}{#2}{#3}{#4}, \ouvrage{#5}}}

\newcommand{\question}[1]{\textbf{#1}}
\newcommand{\spcsujet}{\par \medskip}

\newenvironment{envquestion}
{\begin{adjustwidth}{15mm+\parindent}{15mm}}
{\end{adjustwidth}}

\newenvironment{poesiestyle}{\begin{singlespace}}{\end{singlespace}}

%%%%%%%%%%%%%%%%%%%%%%%%%%%%%%%%%%%%%%%%%%%%%%%%%%%%%%%%%%%%%%%%%%%%%%%
%%%%%%%%%%%%%Commandes pour la mise en page des questions%%%%%%%%%%%%%%
%%%%%%%%%%%%%%%%%%%%%%%%%%%%%%%%%%%%%%%%%%%%%%%%%%%%%%%%%%%%%%%%%%%%%%%
\newcommand{\QI}[2]{
\noindent \question{Question d'interprétation #1~:}\par\nopagebreak
\noindent #2\index[q]{Question d'interprétation #1!#2}
\spcsujet
}
\newcommand{\QR}[2]{
\noindent \question{Essai #1~:}\par\nopagebreak
\noindent #2\index[q]{Essai #1!#2}
\par \bigskip%
}

\setlength{\parskip}{6pt}
\setlength{\headheight}{15.4pt}

%%%%%%%%%%%%%%%%%%%%%%%%%%%%%%%%%%%%%%%%%%%%%%%%%%%%%%%%%%%%%%%%%%%%%%%
%%%%%%%%%%%%%%%%%%%%%%Commandes de mise en page%%%%%%%%%%%%%%%%%%%%%%%%
%%%%%%%%%%%%%%%%%%%%%%%%%%%%%%%%%%%%%%%%%%%%%%%%%%%%%%%%%%%%%%%%%%%%%%%

\newcommand{\chapitre}[4]{
  % nom, chemin de l'image, largeur de l'image, légende
  \pagestyle{plain}
  \cleardoublepage
\noindent  \begin{minipage}[l]{\textwidth}
    \centering
    \chapter{#1}
    \vskip1cm
    \includegraphics[width=#4]{#2}\\\pagebreak
    \medskip
    #3
  \end{minipage}
}

\newcommand{\pageindex}{%
\clearpage
\pagestyle{plain}
\cleardoublepage
\pagestyle{index}
}

\newcommand{\theme}[1]{\index[t]{#1}}

\newcommand{\errata}[3]{{\small “#1”}&\ding{213}&{\small “#2”}&{\footnotesize (#3)}\\}

%%%%%%%%%%%%%%%%%%%%%%%%%%%%%%%%%%%%%%%%%%%%%%%%%%%%%%%
%%%%%%%%%%%%%%Mise en page des en-têtes%%%%%%%%%%%%%%%%
%%%%%%%%%%%%%%%%%%%%%%%%%%%%%%%%%%%%%%%%%%%%%%%%%%%%%%%

%%%%%%%%%En-tête sur la majorité des pages%%%%%%%%%%%%%
\fancypagestyle{normal}
\fancyhead{} % clear all header fields 
\fancyhead[RO]{\leftmark}
\fancyhead[LE]{\rightmark}
\fancyfoot{} %clear all footer fields
\fancyfoot[RO]{\textbf{\thepage ~sur~\pageref*{TotPages}}\hspace*{1.5cm}}
\fancyfoot[LE]{\hspace*{1.5cm}\textbf{\thepage ~sur~\pageref*{TotPages}}}
\renewcommand{\headrulewidth}{0.4pt}
\renewcommand{\footrulewidth}{0pt}

%%%%%%%%%%%%Page blanche et titre%%%%%%%%%%%%%%%%%%%%%%
\fancypagestyle{plain}{%
  \fancyhf{}%
  \fancyfoot{}%
  \renewcommand{\headrulewidth}{0pt}% Line at the header invisible
  \renewcommand{\footrulewidth}{0pt}% Line at the footer invisible
%  \renewcommand{\footrule{\hrule{height 0pt width 0pt}}}%
}
%%%%%%%%%%%%%%%%%%%%%%%%%%%%%%%%%%%%%%%%%%%%%%%%%%%%%%%%%%%%%%%%%%%%%%
%%%%%%%%%%%%%%%%Pour la table des matières%%%%%%%%%%%%%%%%%%%%%%%%%%%%
%%%%%%%%%%%%%%%%%%%%%%%%%%%%%%%%%%%%%%%%%%%%%%%%%%%%%%%%%%%%%%%%%%%%%%
\fancypagestyle{ToC}{%
	\fancyfoot{} % clear all header fields 
	\fancyhead[LE,RO]{\textsc{\rightmark}}
	\fancyfoot{} %clear all footer fields
	\fancyfoot[LE,RO]{\thepage}
	\renewcommand{\headrulewidth}{0.4pt}
	\renewcommand{\footrulewidth}{0pt}
	}
	
\fancypagestyle{version}{%
	\fancyhead{} % clear all header fields 
	\fancyhead[LE,RO]{Notes de version}
	\fancyfoot{} %clear all footer fields
	\fancyfoot[RO]{\textbf{\thepage ~sur~\pageref*{TotPages}}\hspace*{1.5cm}}
	\fancyfoot[LE]{\hspace*{1.5cm}\textbf{\thepage ~sur~\pageref*{TotPages}}}
	\renewcommand{\headrulewidth}{0.4pt}
	\renewcommand{\footrulewidth}{0pt}
	}
	
\fancypagestyle{recap}{%
	\fancyhead{} % clear all header fields 
	\fancyhead[LE,RO]{Récapitulatif des indexations}
	\fancyfoot{} %clear all footer fields
	\fancyfoot[RO]{\textbf{\thepage ~sur~\pageref*{TotPages}}}
	\fancyfoot[LE]{\textbf{\thepage ~sur~\pageref*{TotPages}}}
	\fancyfoot[LO,RE]{\scriptsize{Humanités, Littérature, Philosophie -- Banque des sujets indexés -- Terminale}}
	\renewcommand{\headrulewidth}{0.4pt}
	\renewcommand{\footrulewidth}{0pt}
	}

%\titlecontents*{chapter}% <section-type>
%  [0pt]% <left>
%  {\addvspace{1em}}% <above-code>
%  {\bfseries\chaptername\ \thecontentslabel.\quad}% <numbered-entry-format>
%  {\bfseries}% <numberless-entry-format>
%  {\bfseries\hfill\contentspage}% <filler-page-format>
  
%	\titlecontents*{section}[3.8em] % ie, 1.5em (chapter) + 2.3em
%	{\begin{footnotesize}}
%	{\contentslabel{2.3em}}
%	{\hspace*{-2.3em}}
%	{\titlerule*[1pc]{.}\contentspage}[\end{footnotesize}]
	
\titlecontents*{subsection}[1.5em]
{\filright\scriptsize\hspace*{1em}}
{~\thecontentslabel~}
{}
{,~p.~\thecontentspage}
[][ ---][.]

\titlecontents*{subsubsection}[1.5em]
{~:~\scriptsize}
{\makebox[0pt][r]{~\thecontentslabel~}}
{\makebox[0pt][r]{}}
{,~p.~\thecontentspage}
[ \textbullet~][]

\makeatletter
\renewcommand{\subsubsection}{\@startsection{subsubsection}{3}{\z@}%
  {-3.25ex\@plus -1ex \@minus -.2ex}%
  {1.5ex \@plus .2ex}%
  {\centering\normalfont\normalsize\bfseries}}
\makeatother
%%%%%%%%%%%%%%%%%%%%%%%%%%%%%%%%%%%%%%%%%%%%%%%%%%%%%%%%%%%%%%%%%%%%%%%
%%%%%%%%%%%%%%%%%%%%%%%Pour les index%%%%%%%%%%%%%%%%%%%%%%%%%%%%%%%%%%
%%%%%%%%%%%%%%%%%%%%%%%%%%%%%%%%%%%%%%%%%%%%%%%%%%%%%%%%%%%%%%%%%%%%%%%

\indexsetup{
  %level=\chapter,% <-- already default
  othercode={%
    \thispagestyle{indextitle}%
%    \setlength{\columnseprule}{0.6pt}
    \small\raggedright
  }
}

\fancypagestyle{indextitle}{%
\fancypagestyle{index}{%
  \renewcommand{\headrulewidth}{0pt}%
  \renewcommand{\footrulewidth}{0pt}%
  \fancyhf{}%
  \fancyhead[LE,RO]{{{\itshape\nouppercase{\rightmark}}}}%
  \fancyfoot[LO,RE]{\scriptsize{Humanités, Littérature, Philosophie -- Banque des sujets indexés -- Terminale}}
  \fancyfoot[LE,RO]{\small{\thepage~sur~\pageref*{TotPages}}}
  }}

\fancypagestyle{index}{%
  \renewcommand{\headrulewidth}{0.0pt}%
  \renewcommand{\footrulewidth}{0pt}%
  \fancyhf{}%
  \fancyhead[LE,RO]{{\itshape\nouppercase{\rightmark}}}%
  \fancyfoot[LO,RE]{\scriptsize{Humanités, Littérature, Philosophie -- Banque des sujets indexés -- Terminale}}
  \fancyfoot[LE,RO]{\small{\thepage~sur~\pageref*{TotPages}}}
} 

%%%%%%%%%%%%%%%%%%%%%%%%%%%%%%%%%%%%%%%%%%%%%%%%%%%%%%%%%%%%%%%%%%%%%%%
%%%%%%%%%%%%%%%%%%%%Pour les statistiques%%%%%%%%%%%%%%%%%%%%%%%%%%%%%%
%%%%%%%%%%%%%%%%%%%%%%%%%%%%%%%%%%%%%%%%%%%%%%%%%%%%%%%%%%%%%%%%%%%%%%%

\fancypagestyle{stats}{%
	\renewcommand{\headrulewidth}{0.4pt}
	  \renewcommand{\footrulewidth}{0pt}%
  \fancyhf{}%
  \fancyhead[CE,CO]{\textbf{Statistiques --- Répartitions des sujets selon les thèmes}}%
  \fancyfoot[LO,RE]{\scriptsize{Humanités, Littérature, Philosophie -- Banque des sujets indexés -- Terminale}}
  \fancyfoot[LE,RO]{\small{\thepage~sur~\pageref*{TotPages}}}
}

\newcommand{\centreexamen}{}
\fancypagestyle{stat-centreexam}{%
	\fancyhf{}
	\fancyhead[CE,CO]{\textbf{Statistiques -- Répartition des sujets %
	\IfBeginWith{\centreexamen}{C}{dans les}{en} %
	\centreexamen}}
	\fancyfoot[LE,RO]{\small{\thepage~sur~\pageref*{TotPages}}}
	\fancyfoot[LO,RE]{\scriptsize{Humanités, Littérature, Philosophie -- Banque des sujets indexés -- Terminale}}
}

\fancypagestyle{corr}{%
	\renewcommand{\headrulewidth}{0.4pt}
	  \renewcommand{\footrulewidth}{0pt}%
  \fancyhf{}%
  \fancyhead[CE,CO]{Éléments d'évaluation \textbf{\subsectionmark}}%
  \fancyfoot[LO,RE]{\scriptsize{Humanités, Littérature, Philosophie -- Banque des sujets indexés -- Terminale}}
  \fancyfoot[LE,RO]{\small{\thepage~sur~\pageref*{TotPages}}}
}



%%%%%%%%%%%%%%%%%%%%%%%%%%%%%%%%%%%%%%%%%%%%%%%%%%%%%%%%%%%%%%%%%%%%%%
%%%%%%%%%%%%%%Mise en page du tableau de statistiques%%%%%%%%%%%%%%%%%
%%%%%%%%%%%%%%%%%%%%%%%%%%%%%%%%%%%%%%%%%%%%%%%%%%%%%%%%%%%%%%%%%%%%%%
\newcolumntype{M}[1]{>{\centering\arraybackslash}m{#1}}
\newcolumntype{S}{@{} p{1em} @{}}

\newcommand{\statparlieux}[1]{
  \IfStrEqCase{#1}{%
    {Métropole}{\def\lieuprefix{met}}%
    {Amérique du Nord}{\def\lieuprefix{amnord}}%
    {Centres étrangers}{\def\lieuprefix{cea}}%
    {Polynésie}{\def\lieuprefix{pol}}%
    {Asie}{\def\lieuprefix{asi}}%
    {Nouvelle-Calédonie}{\def\lieuprefix{nca}}%
}[\def\lieuprefix{unk}]

\begin{table}[h]
\begin{tabularx}{\linewidth}{XXXXXXXXXXXXXXXXX}
\hline
\multicolumn{17}{c}{\makecell[cc]{{\large \textbf{#1}}}}\\
\hline
~&~&~&~&~&~&~&~&~&~&~&~&~&~&~&~&~\\
~&~&~&~&~&~&~&~&~&~&~&~&~&~&~&~&~\\\hline
~&\multicolumn{8}{c}{\makecell[cc]{La Recherche de Soi}}&\multicolumn{8}{c}{\makecell[cc]{L'Humanité en question}}\\\hline
\end{tabularx}
\end{table}
}

\newcommand{\tvert}[1]{\parbox[t]{2mm}{\rotatebox[origin=c]{52}{{\footnotesize ~#1~}}}}

%%%%%%%%%%%%%%%%%%%%%%%%%%%%%%%%%%%%%%%%%%%%%%%%%%%%%%%%%%%%%%%%%%%%%%
%%%%%%%%%%%%%%%Mise en forme de la table des matières%%%%%%%%%%%%%%%%%
%%%%%%%%%%%%%%%%%%%%%%%%%%%%%%%%%%%%%%%%%%%%%%%%%%%%%%%%%%%%%%%%%%%%%%
\dottedcontents{section}%
[3.5em]% retrait gauche
{\addvspace{3pt} \footnotesize}% matériel avant
{3.2em}% espacement de contentslabel
{0.5em}% espace entre les . . . .
[]% matériel après

%%%%%%%%%%%%%%%%%%%%%%%%%%%%%%%%%%%%%%%%%%%%%%%%%%%%%%%%%%%%%%%%%%%%%%
%%%%%%%%%%%%%%%%%%%%%%%%%Alignement des ndbp%%%%%%%%%%%%%%%%%%%%%%%%%%
%%%%%%%%%%%%%%%%%%%%%%%%%%avec le texte%%%%%%%%%%%%%%%%%%%%%%%%%%%%%%%
%%%%%%%%%%%%%%%%%%%%%%%%%%%%%%%%%%%%%%%%%%%%%%%%%%%%%%%%%%%%%%%%%%%%%%

\makeatletter
\xpatchcmd{\@footnotetext}{\@parboxrestore}%
{\@parboxrestore\leftskip+12mm\rightskip+15mm}%<--Ce qui change l'alignement
{}{}
\makeatother

\renewcommand{\footnoterule}{%
  \kern -3pt
  \hspace*{15mm}{~} \rule{0.5\textwidth}{1pt} \vspace*{3pt}{~}
  \kern 2pt
}

\definecolor{darkblue}{rgb}{0.0, 0.0, 0.55}
\definecolor{darkviolet}{rgb}{0.58, 0.0, 0.83} 

%\frenchsetup{
%AutoSpaceFootnotes=false
%}

\usepackage{hyperref}

\hypersetup{%
pdffitwindow=true,
pdftitle={Sujets indexés d'HLP Terminale},
pdfauthor={Mikaël Quesseveur},
pdfcreator=pdflatex,
pdfkeywords={Sujets, HLP, Humanités, Littérature, Philosophie, Banque, Recherche, soi, recherche de soi, Humanité en question, Histoire, Violence, limites, création, continuité, rupture},
pdfpagelayout=TwoPageRight,
pdfstartpage=1,
pdfstartview=fitV,
pdfpagemode=UseOutlines,
bookmarksopenlevel=2,
colorlinks,
    linkcolor={darkblue!80!black},
    urlcolor={darkviolet!80!black},
}

\usepackage[type={CC},
modifier={by-nc-sa},
lang={french},
imagemodifier={-eu},
version={4.0},]{doclicense}%%Gestion des droits

\usepackage[numbered,open,openlevel=0]{bookmark}%%%Gestion des signets ; à placer après hyperref : sinon cela désactive les lien des numéros de page dans les index

%\usepackage[french]{minitoc}
%\renewcommand{\mtcSfont}{\footnotesize}

%%%%%%%%%%%%%%%%%%%%%%%%%%%%%%%%%%%%%%%%%%%%%%%%%%%%%%%%%%%%%%%%%%%%%%%
%%%%%%%%%%%%%%%%%%%Liste de toutes les abbrévations%%%%%%%%%%%%%%%%%%%%
%%%%%%%%%%%%%%%%%%%%%%%%%%%%%%%%%%%%%%%%%%%%%%%%%%%%%%%%%%%%%%%%%%%%%%%

%%%%%%%%%%%%%%%%%%%%%%%%%%%%%%%%%%%%%%%%%%%%%%%%%%%%%%%%%%%%%%%%%%%%%%%
%%%%%%%%%%%%%%%%%%3e semestre : La recherche de soi%%%%%%%%%%%%%%%%%%%%
%%%%%%%%%%%%%%%%%%%%%%%%%%%Sujets simples%%%%%%%%%%%%%%%%%%%%%%%%%%%%%%
\newcommand{\RS}{La recherche de soi\stepcounter{sujetRS}}
\newcommand{\ETA}{Education, transmission et émancipation}
\newcommand{\EXS}{Les expressions de la sensibilité}
\newcommand{\MM}{Les métamorphoses du moi}
%%%%%%%%%%%%%%%%%%%%%%%Sujets croisés%%%%%%%%%%%%%%%%%%%%%%%%%%%%%%%%%%
\newcommand{\EXSTA}{\EXS~et \ETA}
\newcommand{\EXSMM}{\EXS~et \MM}
\newcommand{\ETAMM}{\ETA~et \MM}
\newcommand{\EXSTAMM}{Tous les axes sur la recherche de soi}
\newcommand{\OdS}{Sujets originaux sur la recherche de soi}
%%%%%%%%%%%%%%%4e semestre : L'humanité en question%%%%%%%%%%%%%%%%%%%%
%%%%%%%%%%%%%%%%%%%%Sujets simples%%%%%%%%%%%%%%%%%%%%%%%%%%%%%%%%%%%%%
\newcommand{\HQ}{L'humanité en question\stepcounter{sujetHQ}}
\newcommand{\CC}{Création, continuités et ruptures}
\newcommand{\HV}{Histoire et violence}
\newcommand{\HL}{L'humain et ses limites}
%%%%%%%%%%%%%%%%%%%%Sujets mixtes%%%%%%%%%%%%%%%%%%%%%%%%%%%%%%%%%%%%%%
\newcommand{\HVL}{\HV~et \HL}
\newcommand{\HVCC}{\HV~et \CC}
\newcommand{\HLCC}{\HL~et \CC}
\newcommand{\HVLCC}{Tous les axes sur l'humanité en question}
%%%%%%%%%%%%%%%Sujets mixtes semestre%%%%%%%%%%%%%%%%%%%%%%%%%%%%%%%%%%
\newcommand{\ETACC}{\ETA~et \CC}
\newcommand{\ETAHV}{\ETA~et \HV}
\newcommand{\ETAHL}{\ETA~et \HL}
\newcommand{\EXSHV}{\EXS~et \HV}
\newcommand{\EXSHL}{\EXS~et \HL}
\newcommand{\EXSCC}{\EXS~et \CC}
\newcommand{\MMCC}{\MM~et \CC}
\newcommand{\MMHV}{\MM~et \HV}
\newcommand{\MMHL}{\MM~et \HL}
\newcommand{\OdH}{Sujets originaux sur l'humanité en question}
%%% Sujets mixtes sur les années
\newcommand{\ETAPdP}{\ETA~et \PdP}
\newcommand{\ETADFI}{\ETA~et \DFI}
\newcommand{\ETADMRC}{\ETA~et \DMRC}
\newcommand{\ETAHA}{\ETA~et \HA}
\newcommand{\EXSPdP}{\EXS~et \PdP}
\newcommand{\EXSDFI}{\EXS~et \DFI}
\newcommand{\EXSDMRC}{\EXS~et \DMRC}
\newcommand{\EXSHA}{\EXS~et \HA}
\newcommand{\MMPdP}{\MM~et \PdP}
\newcommand{\MMDFI}{\MM~et \DFI}
\newcommand{\MMDMRC}{\MM~et \DMRC}
\newcommand{\MMHA}{\MM~et \HA}
\newcommand{\CCPdP}{\CC~et \PdP}
\newcommand{\CCDFI}{\CC~et \DFI}
\newcommand{\CCDMRC}{\CC~et \DMRC}
\newcommand{\CCHA}{\CC~et \HA}
\newcommand{\HVPdP}{\HV~et \PdP}
\newcommand{\HVDFI}{\HV~et \DFI}
\newcommand{\HVDMRC}{\HV~et \DMRC}
\newcommand{\HVHA}{\HV~et \HA}
\newcommand{\HLPdP}{\HL~et \PdP}
\newcommand{\HLDFI}{\HL~et \DFI}
\newcommand{\HLDMRC}{\HL~et \DMRC}
\newcommand{\HLHA}{\HL~et \HA}
\newcommand{\RSPdP}{\RS~et \PdP}
\newcommand{\RSRdM}{\RS~et \RdM}
\newcommand{\RSDFI}{\RS~et \DFI}
\newcommand{\RSDMRC}{\RS~et \DMRC}
\newcommand{\RSHA}{\RS~et \HA}
\newcommand{\HQPdP}{\HQ~et \PdP}
\newcommand{\HQRdM}{\HQ~et \RdM}
\newcommand{\HQDFI}{\HQ~et \DFI}
\newcommand{\HQDMRC}{\HQ~et \DMRC}
\newcommand{\HQHA}{\HQ~et \HA}

\newcommand{\PdP}{Les pouvoirs de la parole}
%\newcommand{\ArdP}{L'art de la parole}
%\newcommand{\AudP}{L'autorité de la parole}
%\newcommand{\SdP}{Les séductions de la parole}
%%%%%%%%%%%%%%%%%%%%%%%Sujets simples%%%%%%%%%%%%%%%%%%%%%%%%%%%%%%%%%%
%\newcommand{\AAdP}{\ArdP{}~et \AudP{}}
%\newcommand{\ArSdP}{\ArdP{}~et \SdP}
%\newcommand{\AuSdP}{\AudP{}~et \SdP{}}
%\newcommand{\AASdP}{Tous les axes sur la parole}
%\newcommand{\OdP}{Sujets originaux sur la parole}
%%%%%%%%%%%1e semestre : Les représentations du monde%%%%%%%%%%%%%%%%%%
%%%%%%%%%%%%%%%%%%%%Sujets simples%%%%%%%%%%%%%%%%%%%%%%%%%%%%%%%%%%%%%
\newcommand{\RdM}{Les représentations du monde}
\newcommand{\DMRC}{Découverte du monde et rencontre des cultures}
\newcommand{\DFI}{Décrire, figurer, imaginer}
\newcommand{\HA}{L'homme et l'animal}
%%%%%%%%%%%%%%%%%%%%Sujets mixtes%%%%%%%%%%%%%%%%%%%%%%%%%%%%%%%%%%%%%%
\newcommand{\DMFI}{\DMRC~et \DFI}
\newcommand{\DMHA}{\DMRC~et \HA}
\newcommand{\FIHA}{\DFI~et \HA}
\newcommand{\DDHA}{Tous les axes sur les représentations du monde}
\newcommand{\OdR}{Sujets originaux sur les représentations du monde}

\newcommand{\litt}{littéraire}
\newcommand{\phil}{philosophique}

\newcommand{\horsgenres}[1]{Zautres@Autres!#1}
%%%On enregistre à "Z" pour le faire apparaître en dernier
%%%en supprimant le Z à l'affichage

%%%%%%%%%%%Classées par ordre alphabétique%%%%%%%%%%%%%

\newcommand{\aphorisme}{Aphorisme%
\index[g]{Philosophie!Aphorisme}}

\newcommand{\articlel}{Article%
\index[g]{Littérature!Article de journal}}

\newcommand{\autobiographie}{Autobiographie%
\index[g]{Littérature!Autobiographie}}

\newcommand{\biologie}{Biologie%
\index[g]{\horsgenres{Biologie}}}

\newcommand{\carnet}{Carnets%
\index[g]{Philosophie!Carnet}}

\newcommand{\chanson}{Chanson%
\index[g]{Littérature!Chanson}}

\newcommand{\conte}{Conte%
\index[g]{Littérature!Conte}}

\newcommand{\coursp}{Cours~de~philosophie%
\index[g]{Philosophie!Cours}}

\newcommand{\dialoguep}{Dialogue~philosophique%
\index[g]{Philosophie!Dialogue~philosophique}}

\newcommand{\dialoguel}{Dialogue~littéraire%
\index[g]{Littérature!Dialogue~littéraire}}

\newcommand{\discours}{Discours littéraire%
\index[g]{Littérature!Discours~Littéraire}}

\newcommand{\discoursp}{Discours philosophique%
\index[g]{Philosophie!Discours~philosophique}}

\newcommand{\dictionnaire}{Article~de~dictionnaire%
\index[g]{\horsgenres{Article~de~dictionnaire~ou~d'Encyclopédie}}}

\newcommand{\essail}{Essai~littéraire%
\index[g]{Littérature!Essai~littéraire}}

\newcommand{\essaip}{Essai~philosophique%
\index[g]{Philosophie!Essai~philosophique}}

\newcommand{\histoire}{Histoire%
\index[g]{\horsgenres{Histoire}}}

\newcommand{\histpol}{Histoire%
\index[g]{\horsgenres{Histoire politique}}}

\newcommand{\lettrel}{Épistolaire%
\index[g]{Littérature!Epistolaire@Épistolaire}}

\newcommand{\lettrep}{Correspondance%
\index[g]{Philosophie!Correspondance philosophique}}

\newcommand{\memoire}{Mémoire%
\index[g]{Littérature!Mémoire}}

\newcommand{\romanl}{Roman%
\index[g]{Littérature!Roman}}

\newcommand{\romanp}{Roman%
\index[g]{Philosophie!Roman philosophique}}

\newcommand{\poesie}{Poésie%
\index[g]{Littérature!Poésie}}

\newcommand{\prose}{Poésie en prose%
\index[g]{Littérature!Poésie en prose}}

\newcommand{\theatre}{Théâtre%
\index[g]{Littérature!Théâtre}}

\newcommand{\traitep}{Traité%
\index[g]{Philosophie!Traité}}

\newcommand{\voyage}{Récit de voyage%
\index[g]{Littérature!Récit de voyage}}

\newcommand{\genrenondefini}{Genre non-défini%
\index[g]{Genre non-défini}}

%%%%Mise en page spécifique pour la littérature%%%%%
%%%Théâtre%%%
\newcommand{\personnage}[1]{\vspace*{6pt}\begin{center}#1\end{center}\vspace*{-6pt}}
\newcommand{\didascalie}[1]{\begin{flushright}
\textit{#1}
\end{flushright}\vspace*{-6pt}}

%%%%Abréviations%%%
\newcommand{\cad}{c'est-à-dire\xspace}

%%% Gestion d'index %%%
\newcommand{\espsousitem}{}
\newcommand{\espsousousitem}{}